% MODULE ONE: LISTEN, READ AND RESPOND TO (ORALLY AND IN WRITING), SELECT, ORDER, EVALUATE AND/OR PRESENT FACTS, IDEAS AND OPINIONS FROM A VARIETY OF TEXTS, SOURCES AND DEMONSTRATE A POSITIVE ATTITUDE TOWARDS LISTENING AND READING.: Diagnostic Assessment — English Language Upper Sixth — Cours signé OnBuch+
% Official MINESEC syllabus. Compiler avec : tectonic ENL01-module-one-listen-read-and-respond-to-orally-and-in-writing.tex
\newcommand{\DOCMATIERE}{English Language}
\newcommand{\DOCNIVEAU}{Upper Sixth}
\newcommand{\DOCMODULE}{Upper Sixth — English language}
\newcommand{\DOCLECON}{1}
\newcommand{\DOCTITRE}{MODULE ONE: LISTEN, READ AND RESPOND TO (ORALLY AND IN WRITING), SELECT, ORDER, EVALUATE AND/OR PRESENT FACTS, IDEAS AND OPINIONS FROM A VARIETY OF TEXTS, SOURCES AND DEMONSTRATE A POSITIVE ATTITUDE TOWARDS LISTENING AND READING.: Diagnostic Assessment}
% OnBuch+ — préambule « cours signés » ENRICHI (compilé avec tectonic / XeLaTeX)
% Système visuel « Pop » d'OnBuch+ : fond crème (jamais blanc), bordures encre
% épaisses, ombres « dures » décalées, titres Archivo Black en capitales.
% Version MATHÉMATIQUES : graphes (pgfplots/tikz), boîtes pédagogiques
% (définition, propriété, méthode, attention, le savais-tu, exercice résolu,
% exercice, TP, bilan), page de garde et signature OnBuch+ ancrée en bas.
%
% Macros attendues dans main.tex AVANT \input{preamble} :
%   \DOCMATIERE  \DOCNIVEAU  \DOCMODULE  \DOCLECON (numéro)  \DOCTITRE
\documentclass[11pt]{article}

\usepackage{fontspec}
\usepackage[english]{babel}
\usepackage[a4paper,margin=2cm,top=3.4cm,bottom=2.8cm,headheight=1.4cm,footskip=1.3cm]{geometry}
\usepackage{amsmath,amssymb,amsfonts}
\usepackage{mathtools} % \xmapsto, \coloneqq... souvent utilisés par les modèles
\usepackage{cancel} % simplifications barrées (bilans de forces)
% \abs{z} (module, valeur absolue) et \norm{u} (norme), utilisés par les modèles
\providecommand{\abs}{}\let\abs\relax\DeclarePairedDelimiter{\abs}{\lvert}{\rvert}
\providecommand{\norm}{}\let\norm\relax\DeclarePairedDelimiter{\norm}{\lVert}{\rVert}
\usepackage[table]{xcolor} % [table] : fournit aussi \rowcolors (alternance de lignes)
\usepackage{pagecolor}
\usepackage{tikz}
\usetikzlibrary{arrows.meta,positioning,calc,shapes.geometric,shapes.misc,shapes.arrows,decorations.pathmorphing,decorations.pathreplacing,decorations.markings,patterns,shadows,mindmap,trees,fit,backgrounds,matrix,intersections,angles,quotes}
\usepackage{pgfplots}
\usepackage[version=4]{mhchem} % équations nucléaires \ce{^{235}_{92}U}
\usepackage[european,siunitx]{circuitikz} % schémas électriques
\pgfplotsset{compat=1.18}
\usepgfplotslibrary{fillbetween,groupplots} % aires entre courbes (calcul intégral)
\usepackage{enumitem}
\usepackage{fancyhdr}
% [most] charge toutes les bibliothèques tcolorbox (listings, minted, raster,
% documentation...) et gonfle inutilement la mémoire TeX sur un document long
% (~300 boîtes) : on ne charge que ce qui est réellement utilisé.
\usepackage[skins,breakable]{tcolorbox}
\usepackage{graphicx}
\usepackage{colortbl}
\usepackage{booktabs}
\usepackage{tabularx}
\usepackage{diagbox} % case d'en-tête diagonale des tableaux à double entrée
% Colonnes extensibles centrée / gauche / droite (C, L, R), souvent utilisées par les modèles
\newcolumntype{C}{>{\centering\arraybackslash}X}
\newcolumntype{L}{>{\raggedright\arraybackslash}X}
\newcolumntype{R}{>{\raggedleft\arraybackslash}X}
\usepackage{array}
\usepackage{multicol}
\usepackage{siunitx}
% parse-numbers=false : accepte aussi \SI{2,5 \times 10^{-3}}{...}
\sisetup{locale=UK,output-decimal-marker={.},group-minimum-digits=4,parse-numbers=false}
\DeclareSIUnit\ohm{\ensuremath{\symup{\Omega}}} % U+2126 absent de la police
\DeclareSIUnit\division{div} % réglages d'oscilloscope (ms/div, V/div)
\DeclareSIUnit\tr{tr} % tours (vitesse de rotation, ex. \SI{1500}{\tr\per\minute})
\usepackage{qrcode}
% \degree : commande fréquemment utilisée par les modèles pour un angle ou une
% température en dehors de \SI{}{} (non fournie par les paquets chargés).
\providecommand{\degree}{^{\circ}}
% \qed (fin de démonstration), utilisé par les modèles sans amsthm
\providecommand{\qed}{\hfill$\square$}
% \barre : double barre des valeurs interdites dans un tableau de variations (tkz-tab)
\providecommand{\barre}{\ensuremath{\|}}
% environnement proof (amsthm non chargé) utilisé par les modèles
\ifdefined\proof\else\newenvironment{proof}[1][Proof]{\par\noindent\textit{#1.}\ }{\qed\par}\fi
\usepackage{lastpage}
\usepackage{hyperref}
\hypersetup{hidelinks}

% ============================================================
% Palette « Pop » OnBuch+ — mêmes hex que la classe P (plus_ui.dart)
% ============================================================
\definecolor{poporange}{HTML}{F59321}
\definecolor{popdark}{HTML}{E07A0C}
\definecolor{popcream}{HTML}{FFF6E8}
\definecolor{popink}{HTML}{1C1714}
\definecolor{poppurple}{HTML}{7A5AE0}
\definecolor{popgreen}{HTML}{2E9E6A}
\definecolor{poppink}{HTML}{E0457B}
\definecolor{popblue}{HTML}{2F7DE1}
\definecolor{popgold}{HTML}{C98A12}
\definecolor{popmuted}{HTML}{8A7F76}
\definecolor{popline}{HTML}{EADFD2}
\definecolor{popgray}{HTML}{8A7F76}
\colorlet{popgrey}{popgray}
% Alias tolérants (le modèle nomme parfois les couleurs différemment)
\colorlet{popwhite}{white}
\colorlet{popblack}{popink}
\colorlet{popred}{poppink}
\colorlet{popyellow}{popgold}
% Teintes claires pour fonds de tableaux / zones
\colorlet{popblueL}{popblue!10!white}
\colorlet{poporangeL}{poporange!14!white}
\colorlet{popgreenL}{popgreen!12!white}
\colorlet{poppurpleL}{poppurple!12!white}
\colorlet{poppinkL}{poppink!12!white}
\colorlet{popgoldL}{popgold!14!white}

\pagecolor{popcream}

% ============================================================
% Typographie
% ============================================================
\defaultfontfeatures{Ligatures=TeX}
\setmainfont{PlusJakartaSans-Medium.ttf}[Path=../../../fonts/,BoldFont=PlusJakartaSans-Bold.ttf]
% Maths en sans-serif (Fira Math) avec chiffres et lettres droites de Plus
% Jakarta, pour que les formules se fondent dans le texte.
\usepackage[math-style=ISO,bold-style=ISO]{unicode-math}
\setmathfont{FiraMath-Regular.otf}[Path=../../../fonts/]
\setmathfont{PlusJakartaSans-Medium.ttf}[Path=../../../fonts/,range={up/num,up/{Latin,latin},\mathup}]
% (icomma retiré : décimales avec point en anglais)
% Caractères mathématiques Unicode absents des polices de texte (Plus Jakarta,
% Archivo Black) : rendus par leur équivalent mathématique, en texte comme en
% formule — sinon ils s'affichent en carré vide (ex. « Arithmétique dans ℤ »).
\usepackage{newunicodechar}
\newunicodechar{ℝ}{\ensuremath{\mathbb{R}}}
\newunicodechar{ℤ}{\ensuremath{\mathbb{Z}}}
\newunicodechar{ℂ}{\ensuremath{\mathbb{C}}}
\newunicodechar{ℕ}{\ensuremath{\mathbb{N}}}
\newunicodechar{ℚ}{\ensuremath{\mathbb{Q}}}
\newunicodechar{𝔻}{\ensuremath{\mathbb{D}}}
\newunicodechar{α}{\ensuremath{\alpha}}
\newunicodechar{β}{\ensuremath{\beta}}
\newunicodechar{γ}{\ensuremath{\gamma}}
\newunicodechar{δ}{\ensuremath{\delta}}
\newunicodechar{ε}{\ensuremath{\varepsilon}}
\newunicodechar{θ}{\ensuremath{\theta}}
\newunicodechar{λ}{\ensuremath{\lambda}}
\newunicodechar{μ}{\ensuremath{\mu}}
\newunicodechar{σ}{\ensuremath{\sigma}}
\newunicodechar{τ}{\ensuremath{\tau}}
\newunicodechar{φ}{\ensuremath{\varphi}}
\newunicodechar{ψ}{\ensuremath{\psi}}
\newunicodechar{ω}{\ensuremath{\omega}}
\newunicodechar{Σ}{\ensuremath{\Sigma}}
% majuscules produites par \MakeUppercase dans les titres (α-aminés -> Α)
\newunicodechar{Α}{\ensuremath{\alpha}}
\newunicodechar{Β}{\ensuremath{\beta}}
\newunicodechar{Γ}{\ensuremath{\Gamma}}
\newunicodechar{Δ}{\ensuremath{\Delta}}
\newunicodechar{Ε}{\ensuremath{\varepsilon}}
\newunicodechar{Θ}{\ensuremath{\Theta}}
\newunicodechar{Λ}{\ensuremath{\Lambda}}
\newunicodechar{Μ}{\ensuremath{\mu}}
\newunicodechar{Τ}{\ensuremath{\tau}}
\newunicodechar{Φ}{\ensuremath{\Phi}}
\newunicodechar{Ψ}{\ensuremath{\Psi}}
\newunicodechar{Ω}{\ensuremath{\Omega}}
\newunicodechar{↦}{\ensuremath{\mapsto}}
\newunicodechar{∈}{\ensuremath{\in}}
\newunicodechar{∉}{\ensuremath{\notin}}
\newunicodechar{∧}{\ensuremath{\wedge}}
\newunicodechar{⇔}{\ensuremath{\Leftrightarrow}}
\newunicodechar{⟺}{\ensuremath{\Longleftrightarrow}}
\newunicodechar{⇒}{\ensuremath{\Rightarrow}}
\newunicodechar{⟹}{\ensuremath{\Longrightarrow}}
\newunicodechar{∘}{\ensuremath{\circ}}
\newunicodechar{∪}{\ensuremath{\cup}}
\newunicodechar{∩}{\ensuremath{\cap}}
\newunicodechar{≡}{\ensuremath{\equiv}}
\newunicodechar{⊂}{\ensuremath{\subset}}
\newunicodechar{⊥}{\ensuremath{\perp}}
\newunicodechar{∃}{\ensuremath{\exists}}
\newunicodechar{∀}{\ensuremath{\forall}}
\newunicodechar{∛}{\ensuremath{\sqrt[3]{\,}}}
\newunicodechar{∜}{\ensuremath{\sqrt[4]{\,}}}
\newunicodechar{⃗}{\ensuremath{{}^{\to}}}
\newunicodechar{ⁿ}{\ifmmode^{n}\else\textsuperscript{\ensuremath{n}}\fi}
\newunicodechar{ˣ}{\ifmmode^{x}\else\textsuperscript{\ensuremath{x}}\fi}
\newunicodechar{ᵇ}{\ifmmode^{b}\else\textsuperscript{\ensuremath{b}}\fi}
\newunicodechar{ʳ}{\ifmmode^{r}\else\textsuperscript{\ensuremath{r}}\fi}
\newunicodechar{ᵘ}{\ifmmode^{u}\else\textsuperscript{\ensuremath{u}}\fi}
\newunicodechar{ʸ}{\ifmmode^{y}\else\textsuperscript{\ensuremath{y}}\fi}
\newunicodechar{ᵃ}{\ifmmode^{a}\else\textsuperscript{\ensuremath{a}}\fi}
\newunicodechar{ᵉ}{\ifmmode^{e}\else\textsuperscript{\ensuremath{e}}\fi}
\newunicodechar{ᵏ}{\ifmmode^{k}\else\textsuperscript{\ensuremath{k}}\fi}
\newunicodechar{ᵖ}{\ifmmode^{p}\else\textsuperscript{\ensuremath{p}}\fi}
\newunicodechar{ᴺ}{\ifmmode^{N}\else\textsuperscript{\ensuremath{N}}\fi}
\newunicodechar{ᶜ}{\ifmmode^{c}\else\textsuperscript{\ensuremath{c}}\fi}
\newunicodechar{ʰ}{\ifmmode^{h}\else\textsuperscript{\ensuremath{h}}\fi}
\newunicodechar{ᵠ}{\ifmmode^{q}\else\textsuperscript{\ensuremath{q}}\fi}
\newunicodechar{ₙ}{\ifmmode_{n}\else\textsubscript{\ensuremath{n}}\fi}
\newunicodechar{ₐ}{\ifmmode_{a}\else\textsubscript{\ensuremath{a}}\fi}
\newunicodechar{ᵢ}{\ifmmode_{i}\else\textsubscript{\ensuremath{i}}\fi}
\newunicodechar{ᵣ}{\ifmmode_{r}\else\textsubscript{\ensuremath{r}}\fi}
\newunicodechar{ₖ}{\ifmmode_{k}\else\textsubscript{\ensuremath{k}}\fi}
\newunicodechar{ᵦ}{\ifmmode_{\beta}\else\textsubscript{\ensuremath{\beta}}\fi}
\newunicodechar{ₓ}{\ifmmode_{x}\else\textsubscript{\ensuremath{x}}\fi}
\newfontfamily\popdisplay{ArchivoBlack-Regular.ttf}[Path=../../../fonts/]
\newfontfamily\poplogo{SpaceGrotesk-Bold.ttf}[Path=../../../fonts/]
\newfontfamily\popsemi{PlusJakartaSans-SemiBold.ttf}[Path=../../../fonts/]
\newfontfamily\popxbold{PlusJakartaSans-ExtraBold.ttf}[Path=../../../fonts/]
\newfontfamily\popxboldls{PlusJakartaSans-ExtraBold.ttf}[Path=../../../fonts/,LetterSpace=3.0]

\setlist[enumerate]{leftmargin=1.7em,itemsep=5pt,topsep=4pt}
\setlist[itemize]{leftmargin=1.7em,itemsep=4pt,topsep=4pt,label={\color{poporange}\textbullet}}
\setlength{\parindent}{0pt}
\setlength{\parskip}{5pt}
\renewcommand{\arraystretch}{1.25}

% pgfplots : style Pop par défaut
\pgfplotsset{
  every axis/.append style={axis line style={popink,line width=1pt},tick style={popink},
    grid style={popline,line width=0.5pt},label style={font=\small},tick label style={font=\footnotesize}},
  popcurve/.style={poporange,line width=1.8pt,no markers,smooth},
}
% Même style disponible dans un \draw TikZ simple (hors pgfplots)
\tikzset{popcurve/.style={poporange,line width=1.8pt}}
\tikzset{popgrid/.style={popline,very thin}}
\colorlet{popcurve}{poporange} % le nom du style est parfois utilisé comme couleur

% ============================================================
% Pill / squiggle
% ============================================================
\newcommand{\poppill}[3]{%
  \tikz[baseline=-0.55ex]\node[draw=popink,line width=1.2pt,rounded corners=2.6pt,%
    fill=#2,inner xsep=6.5pt,inner ysep=3pt]{%
    \popxboldls\fontsize{7}{7}\selectfont\color{#3}\MakeUppercase{#1}};%
}
\newcommand{\popsquiggle}[1]{%
  \tikz[baseline=-0.4ex]{
    \draw[#1,line width=2.6pt,line cap=round]
      plot [smooth] coordinates {(0,0.09) (0.34,0.01) (0.68,0.21) (1.02,0.01) (1.36,0.10)};
  }%
}
% Mot-clé surligné (marqueur orange sous le texte)
\newcommand{\cle}[1]{{\popxbold\color{popdark}#1}}

% ============================================================
% En-tête / pied
% ============================================================
\pagestyle{fancy}
\fancyhf{}
\renewcommand{\headrulewidth}{0pt}
\renewcommand{\footrulewidth}{0pt}
\lhead{\raisebox{-0.15ex}{\poplogo\fontsize{13}{13}\selectfont\color{popink}OnBuch\color{poporange}+}}
\rhead{\poppill{\DOCMATIERE\ \textbullet\ \DOCNIVEAU\ \textbullet\ Chapter \DOCLECON}{white}{popink}}
\lfoot{\popxbold\fontsize{7.6}{7.6}\selectfont\color{popmuted}\MakeUppercase{Signed course OnBuch+}\ \textbullet\ {\popsemi\DOCTITRE}}
\rfoot{\tikz[baseline=-0.6ex]\node[rounded corners=8.5pt,fill=popink,minimum height=17pt,minimum width=17pt,inner xsep=5pt,inner sep=0]{\popxbold\fontsize{8}{8}\selectfont\color{popcream}\thepage\,/\,\pageref*{LastPage}};}

% ============================================================
% Titres
% ============================================================
\newcommand{\chaptitre}[2]{%
  \begin{tcolorbox}[enhanced,colback=poporange,colframe=popink,boxrule=2.6pt,arc=11pt,
      drop shadow=popink,left=15pt,right=15pt,top=12pt,bottom=11pt,before skip=3pt,after skip=18pt]
    \poppill{#2}{popink}{popcream}\par\vspace{9pt}
    {\raggedright\hyphenpenalty=10000\exhyphenpenalty=10000\popdisplay\fontsize{22}{24}\selectfont\color{popcream}\MakeUppercase{#1}\par}
  \end{tcolorbox}
}
% Section numérotée : pastille orange avec le numéro + titre Archivo
\newcounter{popsec}
\newcommand{\coursec}[1]{%
  \stepcounter{popsec}\setcounter{popsub}{0}%
  \par\vspace{16pt}\noindent
  \begin{minipage}{\linewidth}
  \tikz[baseline=-0.7ex]\node[draw=popink,line width=1.6pt,fill=poporange,rounded corners=4pt,minimum size=22pt,inner sep=0]{\popdisplay\fontsize{12}{12}\selectfont\color{popcream}\thepopsec};%
  \hspace{8pt}{\popdisplay\fontsize{15}{17}\selectfont\color{popink}\MakeUppercase{#1}}\par
  \vspace{4pt}\noindent{\color{poporange}\rule{36pt}{3.6pt}}
  \end{minipage}\par\nopagebreak\vspace{8pt}
}
\newcounter{popsub}
\newcommand{\courssub}[1]{%
  \stepcounter{popsub}%
  \par\vspace{9pt}\noindent{\popxbold\fontsize{11.5}{13}\selectfont\color{popink}{\color{poporange}\thepopsec.\thepopsub}\hspace{6pt}#1}\par\nopagebreak\vspace{4pt}
}

% ============================================================
% Boîtes pédagogiques — même recette : fond blanc, bordure épaisse colorée,
% ombre dure encre, badge plein. Titre optionnel [#1].
% ============================================================
\tcbset{popbase/.style={breakable,enhanced,colback=white,boxrule=2.3pt,arc=9pt,
  shadow={3pt}{-3pt}{0pt}{popink},left=14pt,right=14pt,top=11pt,bottom=11pt,before skip=11pt,after skip=15pt,
  fontupper=\popsemi\fontsize{10.6}{14.5}\selectfont\color{popink},
  fontlower=\popsemi\fontsize{10.6}{14.5}\selectfont\color{popink}}}
% Coche dessinée (le glyphe ✓ n'existe pas dans les polices du document)
\newcommand{\popcheck}[1]{\tikz[baseline=-0.5ex]\draw[#1,line width=1.6pt,line cap=round,line join=round](-0.13,0.0)--(-0.04,-0.09)--(0.13,0.1);}
\providecommand{\coche}{\popcheck{popgreen}}
\providecommand{\resultat}[1]{\par\smallskip\noindent\fcolorbox{popgreen}{popgreenL}{\parbox{\dimexpr\linewidth-2\fboxsep-2\fboxrule\relax}{\textbf{Result.} #1}}\par\smallskip}
\newcommand{\popboxhead}[3]{\poppill{#1}{#2}{white}\ifx\relax#3\relax\else\hspace{6pt}{\popxbold\fontsize{10.6}{12}\selectfont\color{popink}#3}\fi\par\vspace{7pt}}

\newtcolorbox{aretenir}[1][]{popbase,colframe=popgreen,before upper={\popboxhead{Key points}{popgreen}{#1}}}
\newtcolorbox{exemplebox}[1][]{popbase,colframe=poporange,before upper={\popboxhead{Exemple}{poporange}{#1}}}
\newtcolorbox{definition}[1][]{popbase,colframe=popblue,before upper={\popboxhead{Definition}{popblue}{#1}}}
\newtcolorbox{propriete}[1][]{popbase,colframe=poppurple,before upper={\popboxhead{Property}{poppurple}{#1}}}
\newtcolorbox{methode}[1][]{popbase,colframe=popgold,before upper={\popboxhead{Method}{popgold}{#1}}}
\newtcolorbox{attention}[1][]{popbase,colframe=poppink,before upper={\popboxhead{Warning}{poppink}{#1}}}
\newtcolorbox{experience}[1][]{popbase,colframe=popblue,colback=popblueL,before upper={\popboxhead{Experiment}{popblue}{#1}}}
% « Le savais-tu ? » : carte violette pleine (contexte, culture, Cameroun)
\newtcolorbox{savaistu}[1][]{popbase,colback=poppurple,colframe=popink,
  fontupper=\popsemi\fontsize{10.4}{14.2}\selectfont\color{white},
  before upper={\poppill{Did you know?}{white}{poppurple}\ifx\relax#1\relax\else\hspace{6pt}{\popxbold\color{white}#1}\fi\par\vspace{7pt}}}
% Exercice résolu : énoncé en haut, \tcblower puis la solution sur fond crème
\newcounter{popexo}\newcounter{popexr}
\newtcolorbox{exoresolu}[1][]{popbase,colframe=poporange,colbacklower=popcream,
  segmentation style={popink,line width=1.2pt,dashed},
  before upper={\stepcounter{popexr}\popboxhead{Worked example \thepopexr}{poporange}{#1}},
  before lower={\poppill{Solution}{popink}{popcream}\par\vspace{6pt}}}
% Exercice (non résolu en place) — la correction est dans la partie Corrigés
\newtcolorbox{exercice}[1][]{popbase,colframe=popink,
  before upper={\stepcounter{popexo}\popboxhead{Exercise \thepopexo}{popink}{#1}}}
% Corrigé
\newtcolorbox{corrige}[1][]{popbase,colframe=popgreen,colback=popgreenL,
  before upper={\popboxhead{Solution}{popgreen}{#1}}}
% Objectifs / prérequis / bilan
\newtcolorbox{objectifs}{popbase,colframe=popink,colback=poporangeL,
  before upper={\popboxhead{Chapter objectives}{popink}{}}}
\newtcolorbox{prerequis}{popbase,colframe=popink,colback=popblueL,
  before upper={\popboxhead{Prerequisites}{popblue}{}}}
\newtcolorbox{bilan}{popbase,colframe=popink,colback=white,boxrule=2.8pt,
  before upper={\popboxhead{Fiche bilan}{poporange}{L'essentiel à connaître pour l'examen}}}
% Figure encadrée avec légende
\newtcolorbox{popfigure}[1][]{enhanced,breakable=false,colback=white,colframe=popline,boxrule=1.4pt,arc=8pt,
  left=10pt,right=10pt,top=10pt,bottom=8pt,before skip=10pt,after skip=12pt,halign=center,
  lower separated=false,halign lower=center,
  fontlower=\popsemi\fontsize{9}{11.5}\selectfont\color{popmuted},
  #1}
\newcommand{\legende}[1]{\par\vspace{6pt}{\popsemi\fontsize{9}{11.5}\selectfont\color{popmuted}#1}}

% ============================================================
% Signature OnBuch+
% ============================================================
\newcommand{\poplogoinline}[1]{{\poplogo\fontsize{#1}{#1}\selectfont\color{popink}OnBuch\color{poporange}+}}
\newcommand{\signatureonbuch}{%
  \begin{tcolorbox}[enhanced,colback=popink,colframe=popink,boxrule=2.2pt,arc=10pt,
      drop shadow=poporange,left=14pt,right=14pt,top=10pt,bottom=10pt,before skip=0pt,after skip=0pt]
    \tikz[baseline=-0.7ex]\node[circle,fill=popgreen,draw=popcream,line width=1.3pt,minimum size=22pt,inner sep=0]{\popcheck{white}};%
    \hspace{9pt}\begin{minipage}[c]{0.62\linewidth}
      {\popxbold\fontsize{12}{13}\selectfont\color{popcream}Signed\ \textbullet\ {\poplogo OnBuch\color{poporange}+}}\par\vspace{2pt}
      {\raggedright\popsemi\fontsize{8}{10.5}\selectfont\color{popline}\MakeUppercase{OnBuch+ teaching team}\\\MakeUppercase{Follows the official MINESEC syllabus}\par}
    \end{minipage}\hfill
    {\poplogo\fontsize{22}{22}\selectfont\color{popcream}OnBuch\color{poporange}+}
  \end{tcolorbox}%
}

% ============================================================
% Page de garde
% #1 titre · #2 matière · #3 niveau · #4 module · #5 n° leçon · #6 durée ·
% #7 description / objectifs (texte ou itemize)
% ============================================================
\newcommand{\pagegarde}[7]{%
  \thispagestyle{empty}
  \newgeometry{a4paper,margin=1.7cm,top=1.6cm,bottom=1.4cm}%
  \begin{tikzpicture}[remember picture,overlay]
    \fill[poporange!18] ([xshift=-3.2cm,yshift=-3.6cm]current page.north east) circle (4.6cm);
    \fill[poppurple!14] ([xshift=2.4cm,yshift=7.6cm]current page.south west) circle (3.8cm);
  \end{tikzpicture}%
  {\poplogo\fontsize{30}{30}\selectfont\color{popink}OnBuch\color{poporange}+}\hfill
  \raisebox{0.6ex}{\poppill{Signed course \textbullet\ Premium}{popink}{popcream}}\par
  \vspace{1pt}\popsquiggle{poppurple}\par
  \vspace{8pt}
  \begin{tcolorbox}[enhanced,colback=poporange,colframe=popink,boxrule=2.8pt,arc=13pt,
      drop shadow=popink,left=18pt,right=18pt,top=12pt,bottom=13pt,after skip=10pt]
    \poppill{#2 \textbullet\ #3}{popink}{popcream}\hspace{5pt}\poppill{#4}{white}{popink}\par\vspace{8pt}
    {\popdisplay\fontsize{14}{14}\selectfont\color{popink}CHAPTER #5}\par\vspace{4pt}
    {\raggedright\hyphenpenalty=10000\exhyphenpenalty=10000\popdisplay\fontsize{25}{27}\selectfont\color{popcream}\MakeUppercase{#1}\par}
    \vspace{8pt}
    {\popxbold\fontsize{9.5}{9.5}\selectfont\color{popink}\MakeUppercase{Suggested time: #6}}
  \end{tcolorbox}
  \begin{tcolorbox}[enhanced,colback=white,colframe=popink,boxrule=2.2pt,arc=10pt,
      drop shadow=popink,left=16pt,right=16pt,top=10pt,bottom=10pt,after skip=8pt]
    \poppill{In this chapter}{poppurple}{white}\par\vspace{5pt}
    \popsemi\fontsize{9.3}{12.6}\selectfont\color{popink}#7
  \end{tcolorbox}
  \noindent\begin{minipage}[t]{0.487\linewidth}
    \begin{tcolorbox}[enhanced,colback=popgreen,colframe=popink,boxrule=2.2pt,arc=10pt,
        drop shadow=popink,left=11pt,right=11pt,top=10pt,bottom=10pt,height=3.1cm,valign=center]
      \tcbox[colback=white,colframe=popink,boxrule=1.3pt,arc=5pt,boxsep=0pt,nobeforeafter,
        left=3pt,right=3pt,top=3pt,bottom=3pt,valign=center]{\qrcode[height=1.9cm]{https://play.google.com/store/apps/details?id=com.luvvix.onbuch&hl=en}}%
      \hspace{8pt}\begin{minipage}[c]{0.5\linewidth}
        {\popdisplay\fontsize{11}{12.5}\selectfont\color{white}\MakeUppercase{Download OnBuch}}\par\vspace{4pt}
        {\raggedright\popsemi\fontsize{8.4}{11}\selectfont\color{white}Free \textbullet\ Google Play\\AI tutor, past papers, results\par}
      \end{minipage}
    \end{tcolorbox}
  \end{minipage}\hfill
  \begin{minipage}[t]{0.487\linewidth}
    \begin{tcolorbox}[enhanced,colback=poppurple,colframe=popink,boxrule=2.2pt,arc=10pt,
        drop shadow=popink,left=11pt,right=11pt,top=10pt,bottom=10pt,height=3.1cm,valign=center]
      \tikz[baseline=-0.7ex]\node[circle,fill=white,draw=popink,line width=1.3pt,minimum size=1.6cm,inner sep=0]{\popdisplay\fontsize{24}{24}\selectfont\color{poppurple}+};%
      \hspace{8pt}\begin{minipage}[c]{0.6\linewidth}
        {\popdisplay\fontsize{11}{12.5}\selectfont\color{white}\MakeUppercase{Go OnBuch+}}\par\vspace{4pt}
        {\raggedright\popsemi\fontsize{8.4}{11}\selectfont\color{white}All signed courses, unlimited AI tutor, PDF sheets — OnBuch+ tab in the app\par}
      \end{minipage}
    \end{tcolorbox}
  \end{minipage}
  \vspace{8pt}
  \popcontenu
  \vfill
  \signatureonbuch
  \restoregeometry
  \newpage
}

% Rangée « Ce cours contient » (page de garde)
\newcommand{\poptile}[3]{% #1 couleur #2 gros texte #3 légende
  \begin{tcolorbox}[enhanced,colback=#1,colframe=popink,boxrule=2pt,arc=9pt,shadow={2.5pt}{-2.5pt}{0pt}{popink},
      left=6pt,right=6pt,top=9pt,bottom=9pt,halign=center,valign=center,height=1.75cm,width=\linewidth,nobeforeafter]
    {\popdisplay\fontsize{12.5}{13}\selectfont\color{white}\MakeUppercase{#2}}\par\vspace{3pt}
    {\centering\popsemi\fontsize{7.8}{9.5}\selectfont\color{white}#3\par}
  \end{tcolorbox}}
\newcommand{\popcontenu}{%
  \par\noindent{\popxboldls\fontsize{7.5}{7.5}\selectfont\color{popmuted}THIS SIGNED COURSE CONTAINS}\par\vspace{7pt}
  \noindent\begin{minipage}[t]{0.235\linewidth}\poptile{popblue}{Course}{complete, illustrated, step by step}\end{minipage}\hfill
  \begin{minipage}[t]{0.235\linewidth}\poptile{poporange}{Methods}{and worked examples}\end{minipage}\hfill
  \begin{minipage}[t]{0.235\linewidth}\poptile{poppink}{Exercises}{exam style + solutions}\end{minipage}\hfill
  \begin{minipage}[t]{0.235\linewidth}\poptile{popgreen}{Summary sheet}{the essentials to remember}\end{minipage}\par}

% Sommaire
\newtcolorbox{sommaire}{enhanced,colback=white,colframe=popink,boxrule=2.2pt,arc=10pt,
  drop shadow=popink,left=18pt,right=18pt,top=13pt,bottom=13pt,before skip=6pt,after skip=20pt,
  before upper={\poppill{Contents}{poppurple}{white}\par\vspace{9pt}\popsemi\fontsize{10.6}{14.5}\selectfont\color{popink}}}

% Signature de fin de cours — poussée tout en bas de la dernière page
\newcommand{\findecours}{%
  \par\vfill\nopagebreak
  \begin{center}\popsquiggle{poporange}\end{center}\vspace{4pt}
  \signatureonbuch
}
\AtBeginDocument{\def\llbracket{\mathopen{[\mkern-3.5mu[}}\def\rrbracket{\mathclose{]\mkern-3.5mu]}}} % intervalles d'entiers (glyphe ⟦ absent de la police)
% Glyphes absents de Fira Math : redéfinis à partir de glyphes disponibles
\AtBeginDocument{%
  \def\setminus{\mathbin{\backslash}}\let\smallsetminus\setminus
  \def\vdots{\mathord{\vbox{\baselineskip3.2pt\lineskiplimit0pt\kern4pt\hbox{.}\hbox{.}\hbox{.}}}}%
  \def\ddots{\mathinner{\mkern1mu\raise6.5pt\hbox{.}\mkern2mu\raise3.6pt\hbox{.}\mkern2mu\raise0.7pt\hbox{.}\mkern1mu}}%
  \def\triangle{\mathord{\scalebox{0.9}{$\symup{\Delta}$}}}%
  \def\nabla{\mathord{\raisebox{\depth}{\scalebox{1}[-1]{$\symup{\Delta}$}}}}%
  \def\checkmark{\popcheck{popdark}}%
  \def\star{\mathbin{\ast}}\def\bigstar{\mathord{\ast}}%
  \def\diamond{\mathbin{\scalebox{0.6}{$\Diamond$}}}%
  \def\bigcup{\mathop{\mathchoice{\raisebox{-0.3em}{\scalebox{1.7}{$\cup$}}}{\scalebox{1.25}{$\cup$}}{\cup}{\cup}}\displaylimits}%
  \def\bigcap{\mathop{\mathchoice{\raisebox{-0.3em}{\scalebox{1.7}{$\cap$}}}{\scalebox{1.25}{$\cap$}}{\cap}{\cap}}\displaylimits}%
  \def\lesseqqgtr{\mathrel{\lessgtr}}\def\bigoplus{\mathop{\oplus}\displaylimits}%
}
\providecommand{\ssi}{if and only if} % abréviation fréquente
\AtBeginDocument{\providecommand{\wideparen}{\overparen}} % arc de cercle

% Fonctions usuelles que les modèles utilisent sans que amsmath les fournisse
\providecommand{\arsinh}{\operatorname{arsinh}}\providecommand{\arcosh}{\operatorname{arcosh}}
\providecommand{\artanh}{\operatorname{artanh}}\providecommand{\arsech}{\operatorname{arsech}}
\providecommand{\arcsch}{\operatorname{arcsch}}\providecommand{\arcoth}{\operatorname{arcoth}}
\providecommand{\arcsinh}{\operatorname{arsinh}}\providecommand{\arccosh}{\operatorname{arcosh}}
\providecommand{\arctanh}{\operatorname{artanh}}\providecommand{\sech}{\operatorname{sech}}
\providecommand{\csch}{\operatorname{csch}}\providecommand{\arcsec}{\operatorname{arcsec}}
\providecommand{\arccsc}{\operatorname{arccsc}}\providecommand{\arccot}{\operatorname{arccot}}
\providecommand{\sgn}{\operatorname{sgn}}\providecommand{\lcm}{\operatorname{lcm}}
\providecommand{\Var}{\operatorname{Var}}\providecommand{\Cov}{\operatorname{Cov}}

%% FIGFIX-BEGIN v3 — correctifs globaux des figures (docs/onbuch-generation/tools/figfix)
\usepackage{adjustbox}\usepackage{etoolbox}
% toute figure TikZ/pgfplots plus large que la ligne est réduite à la largeur disponible
\BeforeBeginEnvironment{tikzpicture}{\begin{adjustbox}{max width=\linewidth}}
\AfterEndEnvironment{tikzpicture}{\end{adjustbox}}
% légendes pgfplots lisibles par-dessus les courbes
\pgfplotsset{every axis/.append style={legend style={fill=white,fill opacity=0.92,text opacity=1,draw=black!15,inner sep=3pt}}}
% étiquettes d'arêtes (syntaxe edge["..."]) avec fond blanc
\tikzset{every edge quotes/.append style={fill=white,inner sep=1.2pt}}
% bulles de cartes mentales : texte à la ligne dans la bulle, bulle un peu plus grande
\tikzset{every concept/.append style={text width=2.0cm,align=center,minimum size=2.3cm,inner sep=2pt}}
%% FIGFIX-END
\begin{document}

\pagegarde{MODULE ONE: LISTEN, READ AND RESPOND TO (ORALLY AND IN WRITING), SELECT, ORDER, EVALUATE AND/OR PRESENT FACTS, IDEAS AND OPINIONS FROM A VARIETY OF TEXTS, SOURCES AND DEMONSTRATE A POSITIVE ATTITUDE TOWARDS LISTENING AND READING.: Diagnostic Assessment}{English Language}{Upper Sixth}{Upper Sixth — English language}{1}{42 periods}{This opening module of Upper Sixth English Language equips learners with the advanced listening, reading, critical reasoning and writing skills required for GCE Advanced Level Papers 1, 2 and 3. Through complex passages, prescribed literary texts, phonetics, grammar and guided composition, learners learn to select, order, evaluate and present facts, ideas and opinions with accuracy and a positive, critical attitude.
\vspace{4pt}
\begin{multicols}{2}\small
\begin{itemize}
\item By the end of this chapter I can describe the structure and expectations of the three GCE A-Level English papers and assess my own level through a diagnostic test.
\item By the end of this chapter I can analyse complex passages of 900–1000 words to identify main ideas, implicit meaning, tone, attitude and point of view.
\item By the end of this chapter I can evaluate arguments and evidence, distinguishing facts from opinions in persuasive texts.
\item By the end of this chapter I can summarise complex passages and synthesise information from multiple sources into coherent writing.
\item By the end of this chapter I can analyse characters, themes, symbols and style in the prescribed texts (Twilight of Misty Foliage, Unanswered Cries, The Old Man and the Sea, The Swamp Dwellers) and selected poems.
\item By the end of this chapter I can transcribe speech using the full IPA chart and describe connected-speech processes and intonation functions.
\end{itemize}\end{multicols}}

\begin{sommaire}
\begin{multicols}{2}
\begin{enumerate}
\item Section 1: Orientation and Diagnostic Assessment
\item Section 2: Advanced Comprehension of Complex Passages
\item Section 3: Summary, Synthesis and Report Writing
\item Section 4: Grammar Mastery I — Clauses, Tenses and Conditionals
\item Section 5: Grammar Mastery II — Modifiers, Idioms and Word Nuance
\item Section 6: Phonetics and Spoken English
\item Section 7: Argumentative and Synthesis Essay Writing
\item Section 8: Creative and Dramatic Writing
\item Section 9: Prescribed Text Study I — Twilight of Misty Foliage by Ben Jama
\item Section 10: Prescribed Text Study II — Unanswered Cries by Osman Conteh
\item Section 11: Prescribed Text Study III — The Old Man and the Sea by Ernest Hemingway
\item Section 12: Prescribed Text Study IV — The Swamp Dwellers by Wole Soyinka
\item Section 13: Poetry Study — Cosmic Anthology and Selected Poems
\item Section 14: Mid-Year Review and Examination Consolidation
\item Integration activity
\item Methods and worked examples
\item Exercises
\item Solutions to the exercises
\item Summary sheet
\end{enumerate}
\end{multicols}
\end{sommaire}

\begin{objectifs}
\begin{itemize}
\item By the end of this chapter I can describe the structure and expectations of the three GCE A-Level English papers and assess my own level through a diagnostic test.
\item By the end of this chapter I can analyse complex passages of 900–1000 words to identify main ideas, implicit meaning, tone, attitude and point of view.
\item By the end of this chapter I can evaluate arguments and evidence, distinguishing facts from opinions in persuasive texts.
\item By the end of this chapter I can summarise complex passages and synthesise information from multiple sources into coherent writing.
\item By the end of this chapter I can analyse characters, themes, symbols and style in the prescribed texts (Twilight of Misty Foliage, Unanswered Cries, The Old Man and the Sea, The Swamp Dwellers) and selected poems.
\item By the end of this chapter I can transcribe speech using the full IPA chart and describe connected-speech processes and intonation functions.
\item By the end of this chapter I can use advanced grammar accurately: subordinate clauses, conditionals, perfect tenses, modifiers, parts of speech and idioms.
\item By the end of this chapter I can write balanced argumentative essays, formal reports, short plays and synthesis essays using appropriate discourse markers.
\item By the end of this chapter I can edit texts across genres and recognise advanced figures of speech, connotation and implicit meaning such as irony and understatement.
\end{itemize}
\end{objectifs}

\begin{prerequis}
\begin{itemize}
\item Lower Sixth comprehension skills: skimming, scanning and identifying topic sentences
\item Basic literary analysis: plot, character, setting and theme in prose and drama
\item Core grammar: tenses, sentence types, phrases and clauses, basic parts of speech
\item Introduction to phonetics: 44 English phonemes and elementary transcription
\item Essay writing basics: introduction, body paragraphs, conclusion and simple linking words
\end{itemize}
\end{prerequis}

\newpage

\coursec{Section 1: Orientation and Diagnostic Assessment}

\courssub{Why This Section Comes First: From Passive Reader to Strategic Respondent}

Imagine a young graduate in Yaound\'e sitting for a competitive entrance examination: she must read a dense 950-word passage on land tenure in the North West Region, summarise it in 80 words, detect the writer's hidden irony, and then write a balanced argumentative essay --- all in under three hours. This is exactly what the GCE Advanced Level English Language papers demand, and it is also what university, the civil service, and professional life in Cameroon demand. Too many candidates fail not because they cannot read, but because they cannot \cle{listen}, \cle{select}, \cle{order}, \cle{evaluate} and \cle{present} ideas strategically.

Here is the problem this section solves: most Upper Sixth students begin the year without a clear map of the examination they are preparing for, and without an honest picture of their own strengths and weaknesses. They revise ``everything'' and therefore master nothing. The remedy is twofold: first, a precise \cle{orientation} to the syllabus and the GCE papers; second, a \cle{diagnostic assessment} that tells you exactly where you stand before the real work begins.

\begin{definition}[Diagnostic assessment]
A \cle{diagnostic assessment} is a test given \emph{before} instruction begins, designed to reveal a learner's prior knowledge, skills and misconceptions. Unlike a graded examination, its purpose is not to judge but to \emph{guide}: it tells the teacher what to teach and tells the student what to revise. A good diagnostic test samples every major skill area of the syllabus (comprehension, summary, grammar, essay writing, phonetics, literature) so that no weakness remains hidden.
\end{definition}

\begin{definition}[Orientation]
\cle{Orientation} is the process of familiarising yourself with the destination before the journey: the structure of the examination, the types of questions set, the command words used, the time allowed, and the qualities examiners reward. An oriented candidate never meets a surprise in the examination hall.
\end{definition}

\courssub{The Upper Sixth Syllabus and the Three GCE Papers}

The Upper Sixth English Language syllabus is not a random list of topics; it is a carefully engineered preparation for the three papers of the GCE Advanced Level. Every lesson you will study this year --- from subordinate clauses to Hemingway's symbolism, from phonemic transcription to argumentative essays --- feeds directly into at least one paper.

\begin{aretenir}[How the syllabus maps onto the papers]
\begin{itemize}
\item \textbf{Paper 1 (Essay and Comprehension):} argumentative and expository essays, comprehension of complex passages, summary writing, synthesis of sources.
\item \textbf{Paper 2 (Literature and Language use):} the prescribed texts (\emph{Twilight of Misty Foliage}, \emph{Unanswered Cries}, \emph{The Old Man and the Sea}, \emph{The Swamp Dwellers}, selected poems), grammar in context, figures of speech, editing.
\item \textbf{Paper 3 (Oral and Phonetics components, as applicable):} phonemic transcription, connected speech, intonation, stress, and spoken-English skills.
\end{itemize}
\end{aretenir}

\begin{popfigure}
\begin{tikzpicture}[
font=\small,
paper/.style={draw=popblue, fill=popblueL, rounded corners=2pt, text width=3.4cm, align=center, minimum height=1.1cm},
comp/.style={draw=popline, fill=white, rounded corners=2pt, text width=3.4cm, align=center, minimum height=0.7cm},
arr/.style={->, thick, popmuted}]
\node[paper, align=center] (p1) at (0,0){\textbf{Paper 1}\\ Essay \& Comprehension};
\node[paper, align=center] (p2) at (5.2,0){\textbf{Paper 2}\\ Literature \& Language};
\node[paper, align=center] (p3) at (10.4,0){\textbf{Paper 3}\\ Oral \& Phonetics};
\node[comp, align=center] (c1) at (0,-1.6){Argumentative essay\\ Comprehension\\ Summary \& synthesis};
\node[comp, align=center] (c2) at (5.2,-1.6){Prescribed texts\\ Grammar in context\\ Editing, figures of speech};
\node[comp, align=center] (c3) at (10.4,-1.6){IPA transcription\\ Connected speech\\ Intonation \& stress};
\node[comp, fill=popgreenL, draw=popgreen, align=center] (t1) at (0,-3.1){Longer writing time;\\ heavy weighting};
\node[comp, fill=popgreenL, draw=popgreen, align=center] (t2) at (5.2,-3.1){Text-based essays;\\ language items};
\node[comp, fill=popgreenL, draw=popgreen, align=center] (t3) at (10.4,-3.1){Shorter paper;\\ accuracy rewarded};
\draw[arr] (p1) -- (c1); \draw[arr] (c1) -- (t1);
\draw[arr] (p2) -- (c2); \draw[arr] (c2) -- (t2);
\draw[arr] (p3) -- (c3); \draw[arr] (c3) -- (t3);
\end{tikzpicture}
\legende{The three GCE Advanced Level English papers: components and emphasis. Always confirm the current year's exact durations and weightings with your teacher and the official GCE Board syllabus.}
\end{popfigure}

\begin{attention}[Do not memorise invented numbers]
Exact durations and mark weightings are set by the Cameroon GCE Board and may be adjusted. Learn the \emph{structure} and \emph{question types} from past papers, and verify current timings from official sources --- never from rumour.
\end{attention}

\courssub{Examination Structure and Examiner Expectations}

Examiners are not looking for brilliance; they are looking for \emph{evidence of skill}, presented clearly, within the time allowed. Across all three papers, four expectations recur:

\begin{enumerate}
\item \textbf{Relevance:} every sentence must answer the question actually asked.
\item \textbf{Accuracy:} grammar, spelling, punctuation and transcription must be controlled.
\item \textbf{Organisation:} ideas must be selected, ordered and signposted with discourse markers.
\item \textbf{Time discipline:} marks must be earned in proportion to the time spent.
\end{enumerate}

\begin{methode}[Reading a rubric and budgeting time per mark]
\begin{enumerate}
\item \textbf{Read the whole paper first} (2--3 minutes): note the number of questions, which are compulsory, and the mark for each.
\item \textbf{Convert marks into minutes.} If a paper lasts 150 minutes and carries 100 marks, each mark ``costs'' about 1.5 minutes. Reserve 10 minutes for checking, then divide the rest proportionally.
\item \textbf{Underline the command word} in each question (analyse, evaluate, summarise\ldots) before writing anything.
\item \textbf{Plan briefly} for essays and summaries (5 minutes of planning saves 15 minutes of confusion).
\item \textbf{Move on when time is up} for a question; a half-finished answer to every question beats a perfect answer to half of them.
\item \textbf{Check against the rubric} in your final minutes: Have I answered every compulsory question? Every part of each question?
\end{enumerate}
\end{methode}

\begin{exemplebox}[Time budgeting worked out]
A paper lasts 150 minutes and carries 100 marks. Reserve 10 minutes for final checking, leaving 140 working minutes. Time per mark: $140 \div 100 = 1.4$ minutes. A 30-mark essay therefore deserves about $30 \times 1.4 = 42$ minutes (roughly 5 planning, 32 writing, 5 revising); a 10-mark comprehension question deserves about 14 minutes. Write these mini-deadlines in the margin of your question paper and respect them.
\end{exemplebox}

\courssub{Command Words: The Examiner's Instructions}

A command word is the verb in an exam question that tells you \emph{what intellectual operation} to perform. Misreading it is the single most common cause of lost marks at Advanced Level.

\begin{center}
\begin{tabularx}{\linewidth}{|l|X|X|}
\hline
\rowcolor{popblueL} \textbf{Command word} & \textbf{What it demands} & \textbf{Typical trap} \\
\hline
Analyse & Break the text/idea into parts and show how they relate (e.g.\ how tone is created). & Retelling the story instead of examining technique. \\
\hline
Evaluate & Judge the strength of arguments/evidence, giving a reasoned verdict. & Listing points without judging them. \\
\hline
Summarise & Select only essential ideas, in your own words, within the word limit. & Copying sentences; exceeding the limit. \\
\hline
Discuss & Present both sides of an issue before a balanced conclusion. & Arguing one side only. \\
\hline
Compare & Identify similarities \emph{and} differences between texts or characters. & Describing each text separately, never linking them. \\
\hline
Explain & Make clear \emph{how} or \emph{why}, using evidence from the text. & Simply quoting without interpretation. \\
\hline
Justify / To what extent & Defend a position with reasons and textual evidence; measure how far a claim holds. & Asserting without proof; ignoring the counter-case. \\
\hline
\end{tabularx}
\end{center}

\begin{attention}[Common mistake: answering the question you wished for]
Many candidates see the word ``discuss'' and write a one-sided opinion essay, or see ``summarise'' and copy whole sentences from the passage. The examiner marks \emph{the question asked}, not the answer written. Train yourself to underline the command word and to re-read it before you finish each answer.
\end{attention}

\begin{exoresolu}[Command words in action]
\textbf{Statement.} Two questions are set on the same passage about mobile money in Cameroon:
\begin{enumerate}
\item[(a)] \emph{Summarise} in about 60 words the advantages of mobile money mentioned by the writer.
\item[(b)] \emph{Evaluate} the writer's claim that mobile money will soon replace banks entirely.
\end{enumerate}
Explain what each question requires and sketch an appropriate opening sentence for each.
\tcblower
\textbf{Solution.}
\begin{itemize}
\item Question (a) demands \emph{selection and compression}: only the advantages \emph{stated by the writer}, rephrased, no personal opinion, within 60 words. Opening: \emph{``According to the writer, mobile money offers speed, low transaction costs and access for the unbanked rural population\ldots''}
\item Question (b) demands \emph{judgement}: weigh the evidence for and against the claim (network failures, cash culture, regulation) and reach a reasoned verdict. Opening: \emph{``While the writer is right that mobile money has transformed payments, the claim that it will entirely replace banks overlooks\ldots''}
\end{itemize}
The two answers must look completely different: (a) reports the writer; (b) judges the writer.
\end{exoresolu}

\courssub{Designing the Diagnostic Test: Lower Sixth Review}

A sound Upper Sixth diagnostic test samples the four skill domains mastered in Lower Sixth, because these are the foundations on which every Upper Sixth lesson builds.

\begin{enumerate}
\item \textbf{Comprehension (25\%):} a 400--500-word passage with literal, inferential and vocabulary-in-context questions.
\item \textbf{Grammar (25\%):} items on tenses, concord, clauses, prepositions and sentence transformation.
\item \textbf{Summary (25\%):} a short passage to be reduced to a fixed number of words in own words.
\item \textbf{Essay (25\%):} one guided composition (narrative, descriptive or expository) of about 300 words.
\end{enumerate}

Each domain is marked separately, so the result is not one global score but a \emph{profile} of four scores --- and it is the profile, not the total, that matters. Two students can both score 50\% overall and need completely different remedies: one may be weak in grammar but strong in comprehension, the other the reverse.

\begin{popfigure}
\begin{tikzpicture}
\begin{axis}[
ybar, bar width=0.55cm, width=11cm, height=6.5cm,
ymin=0, ymax=100,
ylabel={Average score (\%)},
symbolic x coords={Comprehension, Grammar, Summary, Essay},
xtick=data,
x tick label style={font=\small},
nodes near coords, nodes near coords style={font=\small},
ymajorgrids, grid style={popline!40},
]
\addplot[fill=popblue, draw=popdark] coordinates {(Comprehension,62) (Grammar,48) (Summary,41) (Essay,55)};
\end{axis}
\end{tikzpicture}
\legende{Template for plotting class diagnostic results by skill area. The sample values shown are illustrative only; replace them with your own class data. Here, summary emerges as the weakest domain and should drive the first weeks of revision.}
\end{popfigure}

\courssub{Interpreting Your Results: Strengths, Gaps and Targets}

Read your profile diagnostically, not emotionally. A score below 50\% in a domain signals a \cle{gap} that needs systematic re-teaching; a score between 50\% and 70\% signals a skill that needs \emph{practice and feedback}; above 70\% signals a \cle{strength} to maintain and to exploit in the exam (for example, by answering your strongest question type first).

\begin{methode}[From results to a revision portfolio]
\begin{enumerate}
\item \textbf{Record} each domain score in a table and highlight the weakest two.
\item \textbf{Set SMART targets} --- Specific, Measurable, Achievable, Relevant, Time-bound. Not ``improve grammar'' but ``score at least 14/20 on clause-identification exercises by the end of October''.
\item \textbf{Open an error logbook}: copy each diagnostic error, write the correction, and state the rule in your own words.
\item \textbf{Build a portfolio} containing your diagnostic test, error logbook, monthly targets, and one improved piece of writing per month.
\item \textbf{Review monthly}: re-attempt one diagnostic-style task per domain and update the bar chart of your progress.
\end{enumerate}
\end{methode}

\begin{exemplebox}[A model error logbook entry]
\begin{center}
\begin{tabularx}{\linewidth}{|X|X|X|}
\hline
\rowcolor{popblueL} \textbf{Error (as written)} & \textbf{Correction} & \textbf{Rule in my own words} \\
\hline
The writer say that corruption is bad. & The writer \textbf{says} that corruption is bad. & Third person singular present simple takes -s; report what a text states in the present tense. \\
\hline
\end{tabularx}
\end{center}
\end{exemplebox}

\begin{savaistu}[The error logbook: the cheapest private tutor]
Top candidates at the GCE are rarely those who do the most exercises; they are those who never make the same mistake twice. A simple exercise book in which you record every recurring error --- a misused tense, a copied summary sentence, a misread command word --- and revisit it monthly can be worth more than any extra textbook.
\end{savaistu}

\begin{exercice}[Your personal diagnostic profile]
Using your own diagnostic test results:
\begin{enumerate}
\item Complete a table with your four domain scores (Comprehension, Grammar, Summary, Essay) and rank them from strongest to weakest.
\item Write one SMART target for each of your two weakest domains.
\item Copy three errors from your script into your error logbook, with corrections and rules.
\item State which GCE paper each weak domain affects most, and why.
\end{enumerate}
\end{exercice}

\begin{corrige}[Exercise: model response]
\begin{enumerate}
\item Sample profile: Comprehension 65\%, Grammar 50\%, Summary 40\%, Essay 58\%. Ranking: Comprehension $>$ Essay $>$ Grammar $>$ Summary.
\item Sample targets: (i) ``By the end of November, I will reduce a 300-word passage to 60 words in my own words, keeping all main ideas, in three consecutive practice tests.'' (ii) ``By mid-October, I will correctly identify noun, adjective and adverb clauses in 9 out of 10 sentences.''
\item Sample logbook entry: \emph{Error:} ``The writer say that\ldots'' \emph{Correction:} ``The writer \textbf{says} that\ldots'' \emph{Rule:} third person singular present simple takes -s; use the present tense to report what a text states.
\item Summary weakness affects \textbf{Paper 1} directly (summary section) and indirectly the essay (selection of ideas); clause grammar affects \textbf{Paper 2} language items and the accuracy mark in \textbf{Paper 1} essays.
\end{enumerate}
\end{corrige}

\begin{aretenir}[Key points of Section 1]
\begin{itemize}
\item The Upper Sixth syllabus maps directly onto GCE Papers 1, 2 and 3; know what each paper tests.
\item Command words (analyse, evaluate, summarise, discuss, compare, justify) define the operation required --- underline them.
\item Budget time in proportion to marks (time per mark $=$ working minutes $\div$ total marks), and answer the question asked.
\item A diagnostic assessment produces a \emph{profile}, not a verdict: use it to set SMART targets.
\item Keep an error logbook and a revision portfolio, and revisit them monthly.
\end{itemize}
\end{aretenir}

\coursec{Section 2: Advanced Comprehension of Complex Passages}

In Section 1, you sat a diagnostic assessment and discovered exactly where you stand: which question types you handle confidently and which ones expose gaps in your technique. Many of you found that the greatest challenge was not vocabulary or grammar, but \emph{reading}: long, dense passages of 900 to 1000 words, followed by questions that ask not merely \emph{what} the writer says, but \emph{what the writer means}, \emph{how the writer feels}, and \emph{whether the argument holds together}. That is precisely the territory of this section. By the end of it, you will read complex passages strategically, separate main ideas from details, infer unstated meaning, identify tone and attitude, map the discourse structure, and evaluate arguments like a critical examiner.

\courssub{Strategies for Reading 900--1000 Word Passages}

\textbf{Intuition first.} Imagine you are sent to Douala to buy spare parts in a market you have never visited. If you wander in without a plan, you will waste the whole morning. A smart trader first walks around the market to get its \emph{layout}, then decides which stalls matter, then negotiates. Reading a long passage works the same way: you do not dive into sentence one and crawl to the end; you first survey the terrain.

Three strategies transform how you handle long passages:

\begin{enumerate}
\item \textbf{Previewing.} Before reading properly, spend 60--90 seconds on: the title, the source, the first sentence of each paragraph, and the final paragraph. This gives you the \emph{skeleton} of the text. Previewing tells you the topic, the likely text type (expository, argumentative, narrative), and the writer's probable purpose.
\item \textbf{Chunking.} Do not read word by word. Read in meaningful \emph{chunks} --- groups of words that form one idea: \emph{The government's new agricultural policy / has been welcomed by farmers / across the North West Region.} Chunking doubles reading speed and improves comprehension because the brain processes ideas, not isolated words.
\item \textbf{Annotation.} Mark the text as you read: underline topic sentences, circle signal words (\emph{however, therefore, admittedly}), write two- or three-word margin notes for each paragraph, and put a question mark beside anything surprising or doubtful. Annotation converts passive reading into an active dialogue with the writer.
\end{enumerate}

\begin{methode}[The SQ3R Approach Adapted for Exam Passages]
\begin{enumerate}
\item \textbf{Survey} (about 1 minute): read the title, first and last paragraphs, and the first sentence of each body paragraph. Predict the writer's purpose.
\item \textbf{Question} (30 seconds): turn the title and headings into questions. If the passage is about ``the decline of reading culture among Cameroonian youth'', ask: \emph{Why is it declining? What evidence is given? What solutions are proposed?}
\item \textbf{Read} (4--5 minutes): read the whole passage actively, chunking and annotating. Do not stop for every difficult word; infer meaning from context and move on.
\item \textbf{Recite} (mentally, 30 seconds): after each paragraph, silently state its main idea in your own words. If you cannot, reread that paragraph.
\item \textbf{Review} (before answering questions): glance back at your annotations to rebuild the passage's overall argument. Then answer the questions, returning to the text for evidence --- never answering from memory alone.
\end{enumerate}
\end{methode}

\begin{attention}[Time Trap]
Candidates who read the passage slowly from the first word to the last, without previewing, often run out of time and guess the final questions. The Survey and Question steps feel like a delay, but they save time because you read with a purpose instead of reading blindly.
\end{attention}

\begin{savaistu}[GCE Paper 1 Timing]
In the GCE Advanced Level English Language Paper 1 (Comprehension), you have 2 hours 30 minutes for the whole paper. The comprehension passage (900--1000 words) and its questions typically deserve 50--60 minutes. Allocate roughly: Survey/Question 2 min, Reading/Annotation 8 min, Answering questions 35--40 min, Review 5 min. Practise this split under timed conditions at least once a week.
\end{savaistu}

\courssub{Main Ideas versus Supporting Details}

In dense expository or argumentative prose, every paragraph is built around one \cle{main idea} --- the single point the paragraph exists to make. Everything else --- examples, statistics, quotations, anecdotes, comparisons --- is \cle{supporting detail}.

\begin{exemplebox}[Main Idea or Detail?]
Consider this paragraph:\\
\emph{``Urban farming is quietly transforming Cameroonian cities. In Yaound\'e, families cultivate vegetables in sacks on balconies; in Bafoussam, a cooperative of retired teachers supplies two markets with tomatoes grown on a reclaimed dumpsite; and studies elsewhere in Africa suggest that city gardens can supply a meaningful share of a household's fresh vegetables.''}\\
The \textbf{main idea} is the first sentence: urban farming is transforming cities. The sack gardens, the Bafoussam cooperative and the studies are \textbf{supporting details} --- evidence that makes the claim believable.
\end{exemplebox}

A reliable test: if you delete a sentence and the paragraph still makes its essential point, that sentence was a detail. If deleting it destroys the paragraph, it was the main idea.

\begin{methode}[Locating the Main Idea in Exam Passages]
\begin{enumerate}
\item Check the \textbf{first sentence} (topic sentence) --- most common position.
\item Check the \textbf{last sentence} (concluding sentence) --- common in inductive paragraphs.
\item Look for a \textbf{repeated keyword or synonym} running through the paragraph.
\item Ask: \emph{``What is the one thing the writer wants me to remember from this paragraph?''}
\item Verify with the \textbf{deletion test}: mentally remove each sentence; the one whose removal collapses the paragraph is the main idea.
\end{enumerate}
\end{methode}

\begin{exercice}[Main Idea Identification]
Read the following paragraph and: (a) Underline the main idea sentence. (b) List three supporting details. (c) Write a one-sentence summary of the paragraph.\\
\emph{``The Anglophone crisis in Cameroon has disrupted education for over half a decade. In the North West and South West Regions, thousands of schools have remained closed since 2017. UNICEF reports that more than 700,000 children have been deprived of regular schooling. Many families have fled to Francophone regions or neighbouring Nigeria, where children struggle to integrate into unfamiliar systems. The long-term human capital cost of this interruption is incalculable.''}
\end{exercice}

\begin{corrige}[Main Idea Identification]
(a) Main idea: \emph{``The Anglophone crisis in Cameroon has disrupted education for over half a decade.''} (First sentence; deletion test confirms it.)\\
(b) Supporting details: (i) Thousands of schools closed since 2017; (ii) UNICEF figure of 700,000+ children affected; (iii) Displacement to Francophone regions/Nigeria; (iv) Incalculable long-term cost.\\
(c) Summary: The ongoing Anglophone crisis has caused massive, prolonged educational disruption in Cameroon's North West and South West Regions, affecting hundreds of thousands of children and creating incalculable future costs.
\end{corrige}

\courssub{Inferring Implicit Meaning}

\begin{definition}[Implicit Meaning]
\cle{Implicit meaning} is meaning that is suggested but not directly stated by the writer. The reader must infer it by combining clues in the text (word choice, examples, contrasts, what is emphasised or omitted) with reasonable background knowledge. An inference is only valid if the text genuinely supports it.
\end{definition}

Writers imply rather than state for many reasons: politeness, persuasion, irony, or to make the reader feel clever for drawing the conclusion. Exam questions love implicit meaning because it separates true comprehension from surface reading.

\begin{exemplebox}[Reading Between the Lines]
\emph{``The new headmaster arrived with impressive certificates and a driver who washed his car every morning. The staff room, meanwhile, still had no chalk.''}\\
Nothing is stated directly, yet the writer implies a sharp criticism: the headmaster cares more about appearances and personal comfort than about the school's real needs. The clues are the contrast (\emph{meanwhile}), the ironic detail of the daily car-washing, and the pointed mention of missing chalk.
\end{exemplebox}

\begin{methode}[Making Valid Inferences --- The ICE Method]
\begin{enumerate}
\item \textbf{I} --- \textbf{Identify} the textual clues (specific words, contrasts, omissions, tone markers).
\item \textbf{C} --- \textbf{Connect} the clues to each other and to relevant context (cultural, situational, logical).
\item \textbf{E} --- \textbf{Evaluate} whether the inference is the \emph{most likely} reading, not just a possible one. Ask: ``Does the text \emph{require} this inference, or merely \emph{allow} it?''
\end{enumerate}
\end{methode}

\begin{attention}[Inference vs. Assumption]
An \emph{inference} is grounded in textual evidence. An \emph{assumption} imports outside beliefs without textual support. In the headmaster example, inferring ``the headmaster prioritises image over substance'' is valid (supported by contrast and detail). Assuming ``the headmaster is corrupt'' is not (no evidence of financial misconduct). GCE questions reward inferences, penalise assumptions.
\end{attention}

\courssub{Writer's Attitude, Tone and Point of View}

The \cle{tone} of a text is the writer's attitude towards the subject, conveyed through word choice, sentence rhythm, imagery and emphasis. The \cle{point of view} is the position from which the writer speaks --- an insider, a victim, a detached observer, a satirist. Tone questions at Advanced Level typically offer options such as \emph{ironic, nostalgic, critical, celebratory, detached, indignant, admiring, sceptical, resigned, urgent}.

\begin{center}
\begin{tabularx}{\linewidth}{|l|X|X|}
\hline
\rowcolor{popblueL} \textbf{Tone} & \textbf{How it sounds} & \textbf{Contextual example} \\
\hline
Ironic & Says the opposite of what is meant, often mockingly & ``Brilliant --- another power cut during my oral examination.'' \\
\hline
Nostalgic & Tender longing for the past & ``Those evenings in the village square, when the moon was our electricity, shaped us more than any classroom.'' \\
\hline
Critical & Disapproving, pointing out faults & ``The report praises itself loudly while the hospitals it describes remain without medicine.'' \\
\hline
Celebratory & Joyful, praising & ``When the Indomitable Lions scored, the whole of Bamenda became one voice.'' \\
\hline
Detached & Neutral, clinical, unemotional & ``The data indicate a decline in enrolment over the period under review.'' \\
\hline
Indignant & Angry at injustice & ``It is a scandal that children still sit on bare floors in a country of such wealth.'' \\
\hline
Sceptical & Doubtful, questioning & ``We are told the new curriculum will solve everything; experience suggests otherwise.'' \\
\hline
Resigned & Accepting something unpleasant as inevitable & ``Another year, another strike; we have learnt to wait.'' \\
\hline
\end{tabularx}
\end{center}

\begin{attention}[Reported Opinion Is Not the Writer's Opinion]
The most common tone error is treating an opinion the writer \emph{reports} as the writer's \emph{own} opinion. In ``Supporters of the project call it a miracle; the villagers who lost their land use a different word'', the writer reports two views but personally leans towards sympathy with the villagers --- signalled by the dry closing phrase. Always ask: \emph{whose voice is this sentence in, and what does the writer do with it?}
\end{attention}

\begin{exemplebox}[Tone Analysis Step by Step]
\textbf{Extract:} \emph{``The minister's convoy swept through the potholed market road, sirens blaring, while the women selling tomatoes scrambled to save their produce from the dust. Later, on television, he spoke of his 'deep connection to the grassroots'.''}\\
\textbf{Step 1:} Identify contrasting images --- convoy/sirens vs. women scrambling; TV speech vs. market reality.\\
\textbf{Step 2:} Note loaded words: \emph{swept, blaring, scrambled, 'deep connection' (scare quotes)}.\\
\textbf{Step 3:} Determine attitude: the writer exposes hypocrisy through juxtaposition.\\
\textbf{Step 4:} Name the tone: \textbf{ironic / satirical / critical}. The scare quotes around 'deep connection' are the smoking gun.
\end{exemplebox}

\courssub{Analysing Discourse Structure: Cohesion and Coherence}

\begin{definition}[Cohesion versus Coherence]
\cle{Cohesion} refers to the surface linking devices that stitch sentences together: reference words (\emph{he, this, such}), substitution (\emph{one, do so}), ellipsis, and conjunctions (\emph{however, therefore, moreover}). \cle{Coherence} is the deeper logical unity of ideas --- whether the text makes sense as a whole and each idea follows naturally from the previous one. A text can be cohesive (full of linkers) yet incoherent (the ideas do not actually connect).
\end{definition}

When you analyse a passage, trace the reference chains: who or what does \emph{this}, \emph{they}, \emph{such measures} point back to? GCE comprehension questions frequently ask ``What does the word \emph{it} in line 14 refer to?'' --- a pure cohesion question.

\begin{exemplebox}[Cohesion Devices in Action]
\emph{``The government announced a new scholarship scheme. \textbf{This} was welcomed by student unions. \textbf{However}, \textbf{the details} remained vague. \textbf{Many} feared \textbf{it} would benefit only the well-connected.''}\\
\textbf{This} (line 2) = the scholarship scheme (reference).\\
\textbf{However} (line 3) = adversative conjunction signalling contrast.\\
\textbf{The details} (line 3) = lexical cohesion (definite article + noun linked to 'scheme').\\
\textbf{Many} (line 4) = substitution for 'many student union members / many people'.\\
\textbf{It} (line 4) = reference back to 'the scheme'.
\end{exemplebox}

\begin{methode}[Answering Reference Questions]
\begin{enumerate}
\item Locate the reference word in the text.
\item Read \textbf{backwards} (usually) to find the nearest plausible antecedent.
\item Check \textbf{number} (singular/plural) and \textbf{gender} agreement.
\item Substitute your answer into the sentence: does it make grammatical and logical sense?
\item If the reference is forward (\emph{``This is what I mean: ...''}), read \textbf{forwards}.
\end{enumerate}
\end{methode}

\courssub{Facts, Opinions, Bias and Loaded Language}

\begin{propriete}[Fact versus Opinion]
A \cle{fact} can be verified --- checked against evidence, measurement or records. An \cle{opinion} expresses a judgement, feeling or belief and cannot be proved true or false. (This distinction is not proved in an exam; it is applied.)
\end{propriete}

Persuasive writers blur the line deliberately. Watch for \cle{loaded language} --- words chosen for emotional charge rather than accuracy: \emph{freedom fighter} versus \emph{terrorist}; \emph{thrifty} versus \emph{stingy}; \emph{the regime} versus \emph{the government}. Bias also shows in \emph{selection} (which facts are included or omitted) and \emph{placement} (what gets the headline position).

\begin{exemplebox}[Fact or Opinion?]
\begin{itemize}
\item ``Cameroon hosted the Africa Cup of Nations in 2022.'' --- \textbf{Fact} (verifiable).
\item ``It was the finest tournament ever organised on African soil.'' --- \textbf{Opinion} (a judgement; note the evaluative word \emph{finest}).
\item ``The reckless minister squandered the budget.'' --- \textbf{Opinion with loaded language}: \emph{reckless} and \emph{squandered} carry the writer's hostility.
\end{itemize}
\end{exemplebox}

\begin{exemplebox}[Spotting Bias through Selection and Placement]
\textbf{Headline:} \emph{``Youths Riot in Buea; Police Restore Order''}\\
\textbf{Lead paragraph:} \emph{``Violent clashes erupted yesterday as groups of unemployed youths attacked businesses in Buea. Police intervened swiftly, arresting forty-two suspects. The Senior Divisional Officer praised the forces' restraint.''}\\
\textbf{Buried in paragraph 7:} \emph{``Eyewitnesses, however, claim the violence began when plain-clothes officers seized goods from street vendors without receipts. Three vendors were hospitalised.''}\\
\textbf{Analysis:} The headline and lead frame the youths as aggressors and police as restorers of order. The vendors' perspective is delayed and downplayed. Selection (omitting the trigger) and placement (burying the counter-narrative) create a pro-authority bias.
\end{exemplebox}

\courssub{Critical Reasoning: Evaluating Arguments}

An argument is not a quarrel; it is a structured movement from reasons to a conclusion. Every argument can be mapped using the Toulmin model:

\begin{popfigure}
\begin{center}
\begin{tikzpicture}[
box/.style={draw=popblue, thick, rounded corners, fill=popblueL, text width=3.1cm, align=center, minimum height=1.1cm, font=\small},
arr/.style={-{Stealth[length=3mm]}, thick, popdark}]
\node[box, align=center] (claim) at (0,0){\textbf{Claim}\\ The position the writer defends};
\node[box, align=center] (evid) at (5,0){\textbf{Evidence}\\ Facts, examples, data, testimony};
\node[box, align=center] (warr) at (10,0){\textbf{Warrant}\\ The reasoning that links evidence to claim};
\node[box, fill=popgreenL, draw=popgreen, align=center] (conc) at (5,-2.6){\textbf{Conclusion}\\ What the reader is asked to accept};
\draw[arr] (evid) -- (claim);
\draw[arr] (warr) -- (evid);
\draw[arr] (evid) -- (conc);
\draw[arr, dashed, poppurple] (claim) to[bend right=18] (conc);
\node[font=\scriptsize\itshape, popmuted] at (8.4,-1.9) {claim restated};
\end{tikzpicture}
\end{center}
\legende{The anatomy of an argument: evidence supports the claim through a warrant (often unstated), leading to the conclusion.}
\end{popfigure}

To \emph{evaluate} an argument, interrogate each part. Is the evidence accurate, sufficient, recent and relevant? Is the warrant sound --- does the evidence really lead to the claim? Three fallacies appear constantly in persuasive texts:

\begin{itemize}
\item \textbf{Hasty generalisation}: drawing a sweeping conclusion from too few cases. ``My two neighbours who voted for that party are dishonest, so all its supporters are corrupt.''
\item \textbf{False cause} (\emph{post hoc ergo propter hoc}): assuming that because B followed A, A caused B. ``Since the new mayor took office, rainfall has been poor; his policies have brought drought.''
\item \textbf{Ad hominem}: attacking the person instead of the argument. ``Do not listen to her views on road safety; she failed her driving test twice.''
\item \textbf{Straw man}: misrepresenting an opponent's position to make it easier to attack. ``My opponent wants to abolish all exams; that would destroy education.'' (Opponent actually proposed reform, not abolition.)
\item \textbf{False dilemma}: presenting only two extreme options. ``Either we ban all social media or our children will fail.'' (Ignores regulated use, digital literacy, etc.)
\end{itemize}

\begin{aretenir}[Checklist for Evaluating an Argument]
\begin{itemize}
\item \textbf{Claim:} Is it clear, specific, and debatable?
\item \textbf{Evidence:} Is it accurate? Sufficient? Representative? Recent? From credible sources?
\item \textbf{Warrant:} Is the reasoning valid? Are there hidden assumptions? Does the evidence actually support the claim?
\item \textbf{Counter-arguments:} Does the writer acknowledge and refute them, or ignore them?
\item \textbf{Language:} Is it fair, precise, free of loaded terms and fallacies?
\end{itemize}
\end{aretenir}

\courssub{Implicit Meaning through Irony, Sarcasm and Understatement}

Three figures let writers mean more --- or other --- than they say:

\begin{itemize}
\item \textbf{Irony}: the intended meaning is the opposite of the literal words, or events contradict expectations. ``What a punctual service,'' said about a train three hours late.
\item \textbf{Sarcasm}: a sharper, mocking form of irony intended to wound. ``Oh yes, your plan to study the night before the GCE is truly inspired.''
\item \textbf{Understatement}: deliberately presenting something as less significant than it is, often for ironic or polite effect. ``The famine caused some inconvenience in the region.''
\end{itemize}

Detecting these depends on \emph{context}: a mismatch between words and situation, an exaggerated politeness, or a suspiciously calm description of a disaster.

\begin{exemplebox}[Understatement in Cameroonian Context]
\emph{``The road between Bamenda and Kumbo has \textbf{room for improvement}.''}\\
Context: the road is notorious for deep potholes, landslides, and multi-hour delays. The phrase ``room for improvement'' is a massive understatement, conveying frustration through restraint. The tone is \textbf{ironic / sardonic}.
\end{exemplebox}

\begin{exoresolu}[Analysing a Persuasive Extract]
Read the extract and answer: (a) State the writer's main claim. (b) Identify one fact and one opinion. (c) What is the writer's tone, and how is it conveyed? (d) Name the fallacy in the final sentence.\\
\emph{``Everyone in this country agrees that boarding schools build character. My own years at a boarding school in Kumbo taught me discipline, and the friends I made there remain my brothers today. Any parent who hesitates to send a child to boarding school is simply afraid of excellence. Day students, who go home to soft beds and mother's cooking, never learn to stand alone.''}
\tcblower
\textbf{(a) Main claim:} Boarding schools are superior because they build character and independence.\\
\textbf{(b) Fact:} ``My own years at a boarding school in Kumbo'' --- the writer's attendance is verifiable. \textbf{Opinion:} ``boarding schools build character'' --- a judgement that cannot be proved.\\
\textbf{(c) Tone:} dismissive and slightly contemptuous towards opponents. It is conveyed through loaded language (``soft beds'', ``simply afraid of excellence'') and the belittling picture of day students.\\
\textbf{(d) Fallacy:} \emph{Ad hominem} --- instead of answering parents' concerns, the writer attacks their character (``afraid of excellence''). There is also a hasty generalisation from one person's experience (``my own years'') to ``everyone'' and ``never''.
\end{exoresolu}

\begin{exoresolu}[Analysing Discourse Structure and Implicit Meaning]
\textbf{Extract:} \emph{``The delegation arrived in four-wheel drives, their tinted windows reflecting the thatched roofs of the village. They spoke of 'partnership' and 'capacity building' in air-conditioned comfort. The villagers, who had walked hours to attend, sat on the dust. When the water project was mentioned, the team leader checked his watch and said, 'We must be realistic about timelines.' The meeting ended with photographs and promises. The well is still not dug.''}\\
\textbf{Questions:} (a) What is the writer's attitude towards the delegation? (b) Identify two cohesive devices linking the first three sentences. (c) What is implied by the final sentence? (d) Explain the effect of the contrast between ``tinted windows'' and ``thatched roofs''.\\
\tcblower
\textbf{(a) Attitude:} Critical, cynical, sympathetic to the villagers. The writer exposes the delegation's detachment and empty rhetoric.\\
\textbf{(b) Cohesive devices:} (i) \emph{Their} (line 1) refers back to \emph{The delegation} (anaphoric reference). (ii) \emph{They} (line 2) refers to \emph{The delegation} (reference). (iii) \emph{The villagers} (line 3) links lexically to \emph{the village} (line 1) and contrasts with \emph{The delegation}.\\
\textbf{(c) Implication:} The promises were empty; the project was abandoned or never intended to happen. The final sentence is an understated indictment.\\
\textbf{(d) Contrast effect:} The image juxtaposes wealth/power (4WDs, tinted windows, AC) with poverty/vulnerability (thatched roofs, walking hours, sitting on dust). It visually encodes the power imbalance and the delegation's physical and moral distance from the people they claim to serve.
\end{exoresolu}

\begin{exercice}[Practice: Reading Critically]
Write a short persuasive paragraph (8--10 lines) arguing that mobile phones should be banned in all Cameroonian secondary schools. Deliberately include: one verifiable fact, two opinions, one example of loaded language, and one logical fallacy. Then exchange with a partner and identify each feature in their paragraph.
\end{exercice}

\begin{corrige}[Practice: Reading Critically]
A strong answer might read: ``Mobile phones are destroying education in our schools. In 2019, the Ministry of Secondary Education issued guidelines restricting phone use in schools (fact). Phones are a plague on concentration (opinion, loaded language: \emph{plague}). My cousin failed the GCE because he owned a smartphone (false cause fallacy). Clearly, every serious student agrees that these devices must be banned (hasty generalisation, opinion).'' Marks are awarded for correctly labelling each planted feature and justifying the label.
\end{corrige}

\begin{exercice}[GCE-Style Comprehension Questions]
Read the following passage and answer the questions that follow.\\
\emph{``The proposal to make agriculture compulsory in all secondary schools has sparked heated debate. Proponents argue that it would equip students with practical skills, reduce youth unemployment, and revive the rural economy. 'Agriculture is the backbone of our nation,' declared the Minister of Agriculture at the launch ceremony. 'Every child should know how to feed themselves.'\\
Critics, however, point out that most urban schools lack land, qualified agriculture teachers, and equipment. 'We have laboratories without reagents,' said a principal in Douala. 'Now we are told to start farms?' The teachers' union warns that adding a compulsory subject will overload an already crowded curriculum. Meanwhile, students in rural areas, where farming is a daily reality, wonder why they should be examined on what they already live.\\
The policy document itself is vague on implementation. It speaks of 'school gardens' and 'industrial attachments' but allocates no specific budget. History suggests caution: the 2008 'Green Schools' initiative was announced with similar fanfare; five years later, not a single school garden from that programme survives.''}
\begin{enumerate}
\item State the writer's main purpose in this passage. (2 marks)
\item Identify two arguments \emph{for} and two arguments \emph{against} compulsory agriculture. (4 marks)
\item What is the writer's attitude towards the policy document? Quote a phrase to support your answer. (3 marks)
\item Explain the function of the final sentence. (2 marks)
\item The word ``fanfare'' (last line) suggests: A) careful planning B) publicity without substance C) musical celebration D) government funding. Choose the correct option and explain why. (3 marks)
\item Is the statement ``The 2008 Green Schools initiative failed'' a fact or an opinion in this passage? Justify. (2 marks)
\end{enumerate}
\end{exercice}

\begin{corrige}[GCE-Style Comprehension Questions]
\begin{enumerate}
\item \textbf{Purpose:} To present a balanced overview of the debate on compulsory agriculture in secondary schools, while implicitly criticising the policy's lack of concrete planning and historical precedent. (2 marks: 1 for 'balanced overview/debate', 1 for 'implicit criticism'.)
\item \textbf{For:} (i) Equips students with practical skills; (ii) Reduces youth unemployment; (iii) Revives rural economy. \textbf{Against:} (i) Urban schools lack land/teachers/equipment; (ii) Curriculum already overloaded; (iii) Rural students already live farming, making examination redundant. (4 marks: 1 per distinct argument, max 2 per side.)
\item \textbf{Attitude:} Sceptical / critical / dismissive. \textbf{Support:} ``vague on implementation'', ``allocates no specific budget'', or ``History suggests caution''. (3 marks: 1 for attitude, 2 for apt quotation.)
\item \textbf{Function:} It provides historical evidence (precedent) to undermine confidence in the current proposal; it implies the new policy will suffer the same fate. (2 marks.)
\item \textbf{Answer:} B) publicity without substance. \textbf{Explanation:} ``Fanfare'' implies a loud, showy announcement; the context (``similar fanfare; five years later, not a single school garden... survives'') shows the announcement was not matched by action. (3 marks: 1 for option, 2 for explanation.)
\item \textbf{Fact.} It is presented as a verifiable historical claim (``not a single school garden from that programme survives''), not as the writer's judgement. (2 marks: 1 for 'fact', 1 for justification.)
\end{enumerate}
\end{corrige}

\begin{aretenir}[Key Points for the Examination]
\begin{itemize}
\item Preview, chunk and annotate; use SQ3R to read with purpose, not panic.
\item The main idea survives deletion; supporting details do not.
\item Implicit meaning must be \emph{inferred from clues}, never invented. Use the ICE method.
\item Tone lives in word choice; always separate the writer's voice from reported voices.
\item Cohesion is surface linking; coherence is logical unity. Trace reference chains for cohesion questions.
\item Facts are verifiable; opinions are judgements. Bias hides in loaded words, selection, and placement.
\item Evaluate arguments by testing claim, evidence and warrant; name fallacies precisely (hasty generalisation, false cause, ad hominem, straw man, false dilemma).
\item Irony, sarcasm, and understatement signal implicit meaning through context mismatch.
\item In GCE answers: quote precisely, name the device/fallacy/tone, and explain the effect.
\end{itemize}
\end{aretenir}

\coursec{Section 3: Summary, Synthesis and Report Writing}

In Section 2, you learned to read complex passages closely: to trace main ideas, infer implicit meanings and evaluate a writer's attitude. But comprehension alone is not enough at the GCE Advanced Level. The examiner also wants to know whether you can \emph{do something} with what you have read: compress it into a summary, combine it with other sources into a synthesis, or reorganise it into a formal report. These three skills --- \cle{summary}, \cle{synthesis} and \cle{report writing} --- are the natural continuation of comprehension, and they appear every year in Papers 1 and 2. This section trains you in all three, step by step.

\courssub{Summarising Complex Passages: Selecting Essential Information}

\begin{definition}[Summary]
A \cle{summary} is a shortened restatement of a passage, in your own words, which keeps all the \emph{essential} ideas and removes everything that is not essential: examples, illustrations, repetitions, figures of speech, quotations and minor details.
\end{definition}

Imagine a trader at Mokolo Market who has two hours of conversation with a supplier but must report the outcome to her association in two minutes. She will not repeat the jokes, the greetings or the stories about the journey; she will report only the price agreed, the quantity and the delivery date. That is exactly what a summary does to a text.

What must be \textbf{deleted} when summarising:
\begin{itemize}
\item \textbf{Examples and illustrations} introduced by \emph{for example, for instance, such as}.
\item \textbf{Figures of speech}: metaphors, similes and hyperbole must be replaced by plain statements (``he drowned in a sea of debt'' becomes ``he was heavily indebted'').
\item \textbf{Repetitions and restatements}: writers often say the same thing twice for emphasis; keep it once.
\item \textbf{Quotations, statistics and anecdotes}, unless the question specifically requires them.
\item \textbf{Minor details}: dates, names and places that do not carry the main argument.
\end{itemize}

What must be \textbf{kept}: the topic sentence of each paragraph, the writer's main claims, important contrasts and causes, and the conclusion.

\begin{methode}[The 5-Step Summary Method]
\begin{enumerate}
\item \textbf{Identify topic sentences.} Read the passage once for gist, then locate the sentence in each paragraph that carries its main idea (usually, but not always, the first sentence).
\item \textbf{Underline key points.} Mark only the words and phrases that carry essential meaning. Ignore examples and decorative language.
\item \textbf{Paraphrase.} Rewrite each key point in your own words (see the paraphrasing techniques below).
\item \textbf{Link.} Join your points with suitable discourse markers (\emph{first, moreover, however, consequently, finally}) so the summary reads as one flowing piece, not a list of fragments.
\item \textbf{Count words.} Check that you are within the required limit; if not, compress further (see word-limit discipline below).
\end{enumerate}
\end{methode}

\begin{popfigure}
\begin{center}
\begin{tikzpicture}[
box/.style={draw=popblue, fill=popblueL, rounded corners=2pt, minimum width=2.5cm, minimum height=0.9cm, align=center, font=\small, text=popdark},
arr/.style={->, thick, poporange}
]
\node[box, align=center] (read) at (0,0){\textbf{1. READ}\\gist + topic\\sentences};
\node[box, align=center] (sel) at (3.6,0){\textbf{2. SELECT}\\underline key\\points only};
\node[box, align=center] (para) at (7.2,0){\textbf{3. PARAPHRASE}\\own words,\\plain style};
\node[box, align=center] (comp) at (10.8,0){\textbf{4. COMPRESS}\\link + cut to\\word limit};
\node[box, align=center] (check) at (14.4,0){\textbf{5. CHECK}\\all points?\\own words?};
\draw[arr] (read) -- (sel);
\draw[arr] (sel) -- (para);
\draw[arr] (para) -- (comp);
\draw[arr] (comp) -- (check);
\end{tikzpicture}
\end{center}
\legende{The summary-writing process: a five-stage pipeline from reading to final check.}
\end{popfigure}

\courssub{Paraphrasing Techniques}

A summary must be in \emph{your own words}. Four reliable techniques allow you to re-express an idea without changing its meaning:

\begin{center}
\begin{tabularx}{\linewidth}{|l|X|X|}
\hline
\rowcolor{popblueL} \textbf{Technique} & \textbf{Original} & \textbf{Paraphrase} \\
\hline
Synonym substitution & The harvest was \emph{abundant}. & The harvest was \emph{plentiful}. \\
\hline
Restructuring (word order) & \emph{Because the rains failed}, the farmers suffered. & The farmers suffered \emph{as a result of the failed rains}. \\
\hline
Changing voice & Locusts \emph{destroyed} the crop. & The crop \emph{was destroyed} by locusts. \\
\hline
Changing word class & The government \emph{decided} to act. & The government took a \emph{decision} to act. \\
\hline
\end{tabularx}
\end{center}

The strongest paraphrases combine several techniques at once: ``The authorities \emph{reluctantly acknowledged} that the project had \emph{failed}'' can become ``The authorities \emph{admitted unwillingly} that the project had \emph{not succeeded}'' (synonym + word class).

\begin{attention}[Do not lift whole sentences]
The most common and most heavily penalised mistake in GCE summary questions is \cle{lifting} --- copying whole phrases or sentences from the passage. Examiners call this ``wholesale lifting'' and award little or no credit for it, even if the idea is correct. Change the wording, change the structure, or both. If a technical term cannot be replaced (e.g. \emph{photosynthesis}, \emph{inflation}), keep the term but rebuild the sentence around it.
\end{attention}

\courssub{Word-Limit Discipline: Drafting, Counting and Compressing}

GCE summary questions typically impose a limit such as ``in about 100 words'' or ``in not more than 120 words.'' Discipline with the limit is itself a tested skill.

\begin{methode}[Compressing to the Limit]
\begin{enumerate}
\item \textbf{Draft freely first.} Write all your points without worrying about length.
\item \textbf{Count honestly.} Count every word (contractions like \emph{don't} count as one word; hyphenated words count as one).
\item \textbf{Cut in layers.} First remove adjectives and adverbs that add no information (\emph{very, really, extremely}); then merge two short sentences into one; finally, replace phrases with single words (\emph{in spite of the fact that} $\rightarrow$ \emph{although}).
\item \textbf{Never cut a content point to save words.} Cut \emph{expression}, not \emph{ideas}. One relevant point lost is one mark lost.
\end{enumerate}
\end{methode}

\begin{exemplebox}[Compression in action]
Draft (34 words): \emph{``The old man, who was very tired and completely exhausted after his long struggle, finally managed to bring the great fish back towards the shore.''}

Compressed (14 words): \emph{``Exhausted by his struggle, the old man finally brought the great fish towards shore.''}

All three content points (exhaustion, struggle, return with the fish) survive; only padding has gone.
\end{exemplebox}

\begin{savaistu}[GCE Word-Count Rules]
In the Cameroon GCE, words are counted as follows: each separate token separated by a space counts as one word. Contractions (\emph{can't, won't}) count as one word. Hyphenated compounds (\emph{well-known, up-to-date}) count as one word. Numbers written in digits (\emph{2025}) count as one word; written in full (\emph{two thousand and twenty-five}) count as separate words. Proper nouns (\emph{Yaound\'e}) count as one word. Always count your final draft before submitting.
\end{savaistu}

\courssub{Synthesising Information from Multiple Passages}

\begin{definition}[Synthesis]
A \cle{synthesis} is the combination of ideas from several sources into a new, unified whole. Unlike a summary, which shortens \emph{one} text, a synthesis \emph{compares and merges} two or more texts, noting where they agree, where they contradict each other, and what overall picture emerges.
\end{definition}

\begin{methode}[From Two Passages to One Synthesis]
\begin{enumerate}
\item \textbf{Read both passages} and summarise each one's position in a single sentence.
\item \textbf{Draw a comparison grid}: list the points of \emph{agreement} on one side and the points of \emph{contradiction} on the other.
\item \textbf{Group the material by theme}, not by passage. A synthesis organised ``Passage A says\ldots Passage B says\ldots'' is two summaries, not a synthesis.
\item \textbf{Attribute ideas} to their sources as you write (\emph{According to the first writer\ldots; the second passage, however, argues\ldots}).
\item \textbf{Conclude with the combined picture}: what do the sources, taken together, establish?
\end{enumerate}
\end{methode}

\begin{exemplebox}[Comparing viewpoints]
\textbf{Passage A} (a ministry report): boarding schools improve discipline and academic results in rural areas. \textbf{Passage B} (a parents' association letter): boarding schools separate young children from their families and weaken cultural transmission.

\textbf{Agreement}: both accept that rural education needs urgent improvement. \textbf{Contradiction}: A sees separation from home as beneficial (discipline); B sees it as harmful (loss of family and culture).

\textbf{Synthesis sentence}: \emph{``While both sources agree that rural schooling must improve, they disagree sharply over boarding schools: the ministry report praises the discipline they instil, whereas the parents' association condemns the separation of children from their families and culture.''}
\end{exemplebox}

\courssub{Synthesis Writing: Integrating Multiple Perspectives}

In the examination, synthesis writing means producing \emph{one coherent essay} that weaves several perspectives together, with clear \cle{attribution} (saying whose idea each point is) and with your own controlling line of argument.

Useful attribution and contrast language:
\begin{itemize}
\item \textbf{Attribution}: \emph{According to Source 1\ldots; The second writer maintains that\ldots; Both passages acknowledge\ldots}
\item \textbf{Convergence}: \emph{Similarly; In the same way; Both sources agree that\ldots}
\item \textbf{Divergence}: \emph{However; By contrast; Whereas Source 1 claims\ldots, Source 2 insists\ldots}
\item \textbf{Resolution}: \emph{Taken together, the sources suggest that\ldots; On balance\ldots}
\end{itemize}

\begin{exoresolu}[Synthesis task]
\textbf{Statement.} Source 1 (editorial): ``Social media has given young Cameroonians a voice in national affairs for the first time.'' Source 2 (lecturer's article): ``Social media spreads rumours faster than facts and distracts youth from serious study.'' Write one paragraph (about 80 words) synthesising the two views.
\tcblower
\textbf{Model answer.} Both sources recognise that social media now shapes the lives of young Cameroonians, but they judge its value differently. According to the editorial, these platforms are beneficial because they give the youth a voice in national affairs for the first time. The lecturer, by contrast, warns that the same platforms spread rumours faster than facts and draw students away from serious study. Taken together, the sources suggest that social media is a powerful tool whose value depends on how responsibly the youth use it. (79 words)
\end{exoresolu}

\courssub{Reconstructing Formal Reports: Structure and Language}

A \cle{formal report} presents the findings of an investigation to an official audience (a principal, a ministry, a committee). At the GCE, you may be given jumbled notes or a passage and asked to \emph{reconstruct} it as a proper report.

\begin{propriete}[Language of the formal report]
A formal report uses an \cle{impersonal tone}, \cle{passive constructions} and \cle{clear headings}. Write \emph{``It was observed that the laboratory lacks equipment''} rather than \emph{``I saw that the lab has no equipment.''} Avoid contractions, slang, rhetorical questions and emotional vocabulary. (This property is directly examinable: candidates are routinely asked to rewrite informal material in report style.)
\end{propriete}

The standard structure has five parts:
\begin{enumerate}
\item \textbf{Title} --- names the subject of the report.
\item \textbf{Terms of reference} --- who asked for the report, why, and when.
\item \textbf{Findings} --- the facts discovered, presented under numbered headings, in impersonal language.
\item \textbf{Conclusion} --- what the findings mean, drawn logically from them.
\item \textbf{Recommendations} --- specific, practical proposals, often introduced by \emph{It is recommended that\ldots}
\end{enumerate}

\begin{popfigure}
\begin{center}
\begin{tikzpicture}[
head/.style={draw=popblue, fill=popblue, text=white, rounded corners=2pt, minimum width=11.5cm, minimum height=0.7cm, font=\small\bfseries, align=center},
body/.style={draw=popline, fill=white, rounded corners=2pt, minimum width=11.5cm, minimum height=0.85cm, font=\small, align=center, text=popdark},
note/.style={font=\footnotesize\itshape, text=popmuted, align=left}
]
\node[head] (t) at (0,0) {TITLE: Report on\ldots};
\node[body] (tor) at (0,-1.0) {TERMS OF REFERENCE --- who commissioned it, purpose, date};
\node[body] (f1) at (0,-2.0) {FINDINGS --- 1. \quad 2. \quad 3. \quad (numbered, impersonal, factual)};
\node[body] (c) at (0,-3.0) {CONCLUSION --- interpretation of the findings};
\node[body] (r) at (0,-4.0) {RECOMMENDATIONS --- It is recommended that\ldots};
\node[body, fill=popblueL] (s) at (0,-5.0) {Signature, name, position, date};
\node[note] at (7.6,-1.0) {past tense, passive voice};
\node[note] at (7.6,-3.0) {no new facts here};
\node[note] at (7.6,-4.0) {must follow from findings};
\end{tikzpicture}
\end{center}
\legende{Annotated layout of a formal report, with the language expectations for each part.}
\end{popfigure}

\begin{exemplebox}[From informal to report style]
\textbf{Informal note}: ``Honestly, the students just don't have enough textbooks and the library is a mess. Something must be done!''

\textbf{Report style}: \emph{``It was observed that the supply of textbooks is insufficient and that the library is poorly maintained. It is therefore recommended that additional textbooks be purchased and that the library be reorganised.''}
\end{exemplebox}

\begin{experience}[Report Reconstruction Practice]
You are the Senior Prefect of GBHS Mfoundi. The Principal has asked you to investigate the causes of late coming among Form Five students and write a report. You have gathered the following notes. Reorganise them into a formal report with the five standard headings.

\textbf{Notes:}
\begin{itemize}
\item Many students live far from school and public transport is unreliable.
\item Some students stay up late watching films and oversleep.
\item The school gate closes at 7:30 a.m. but assembly starts at 7:15 a.m.
\item Teachers have complained that latecomers disrupt first-period lessons.
\item The PTA could organise a car-pool system for distant students.
\item The school could introduce a supervised study hall from 6:30 a.m. for early arrivals.
\item A survey of 120 Form Five students showed 65\% cite transport, 25\% cite poor time management, 10\% cite domestic chores.
\item The Principal requested this report on 10 September 2025.
\end{itemize}
\end{experience}

\courssub{Discourse Markers for Staging, Sequencing and Concluding}

Summaries, syntheses and reports all depend on \cle{discourse markers} --- the signposts that tell the reader where the text is going.

\begin{center}
\begin{tabularx}{\linewidth}{|l|X|}
\hline
\rowcolor{popblueL} \textbf{Function} & \textbf{Markers} \\
\hline
Staging (announcing the plan) & \emph{This report examines\ldots; The passage discusses\ldots; Three main points emerge.} \\
\hline
Sequencing & \emph{First(ly); Next; Subsequently; In addition; Moreover; Finally.} \\
\hline
Contrast & \emph{However; Nevertheless; By contrast; On the other hand.} \\
\hline
Cause and result & \emph{Consequently; As a result; Therefore; This explains why\ldots} \\
\hline
Concluding & \emph{In conclusion; On balance; Taken together; It is clear that\ldots} \\
\hline
\end{tabularx}
\end{center}

\begin{attention}[Marker overload]
Do not begin every sentence with a marker. Two or three well-placed markers per paragraph are enough; a text stuffed with \emph{moreover, furthermore, additionally} in every line sounds mechanical and loses marks for style.
\end{attention}

\begin{aretenir}[Key Points of Section 3]
\begin{itemize}
\item A summary keeps essential ideas and deletes examples, figures of speech, repetitions and minor details --- always in your own words.
\item Paraphrase using synonyms, restructuring, change of voice and change of word class.
\item Respect the word limit: draft, count, then compress expression, never content.
\item A synthesis merges several sources by theme, noting agreement and contradiction, with clear attribution.
\item A formal report follows Title -- Terms of reference -- Findings -- Conclusion -- Recommendations, in impersonal, often passive, language.
\item Discourse markers stage, sequence and conclude your writing.
\end{itemize}
\end{aretenir}

\begin{exercice}[Summary practice]
Read the following extract, then summarise the causes of examination failure it mentions in \textbf{not more than 50 words}, in your own words.

\emph{``Many candidates fail the examination, not because they are dull, but because they prepare badly. For instance, some students, like a certain boy I knew in Bamenda, begin serious revision only a week before the papers. Others drown in a sea of past questions without ever reading their textbooks, mistaking quantity for understanding. Still others burn the midnight oil until they are exhausted, and then sit the paper with foggy minds. In short, poor planning, shallow study and fatigue --- not lack of intelligence --- are the real enemies.''}
\end{exercice}

\begin{corrige}[Exercise 1]
\textbf{Model summary} (44 words): \emph{According to the writer, candidates fail mainly because of poor preparation rather than low intelligence: they start revising too late, study past questions superficially without mastering their textbooks, and overwork themselves until exhaustion, which leaves them mentally unfit for the examination.}

\textbf{Marking guide}: 1 mark per relevant point (late revision; shallow use of past questions; fatigue from overwork; failure not due to low intelligence) + 1 mark for own words and correct limit. Note how the Bamenda anecdote and the metaphors (``sea of past questions'', ``burn the midnight oil'') were deleted or translated into plain language.
\end{corrige}

\begin{exercice}[Report reconstruction]
The following sentences are the jumbled parts of a report on a school water shortage. Rewrite them as a proper report with the five standard headings, improving the language where necessary: (a) ``We therefore advise that a second borehole be drilled.'' (b) ``Report on the Water Shortage at GBHS Example College.'' (c) ``It was found that the only borehole breaks down frequently and that students spend two hours daily fetching water.'' (d) ``Following a request by the Principal in January, the committee investigated the water problem.'' (e) ``The shortage clearly disrupts classes and threatens hygiene.''
\end{exercice}

\begin{corrige}[Exercise 2]
\textbf{Title}: Report on the Water Shortage at GBHS Example College. \textbf{Terms of reference}: Following a request by the Principal in January, the committee investigated the water problem. \textbf{Findings}: It was found that the only borehole breaks down frequently and that students spend two hours daily fetching water. \textbf{Conclusion}: The shortage clearly disrupts classes and threatens hygiene. \textbf{Recommendations}: It is recommended that a second borehole be drilled. (Note the impersonal reconstruction of (a): ``We advise'' $\rightarrow$ ``It is recommended that\ldots'')
\end{corrige}

\begin{exercice}[Synthesis practice]
Read the two short texts below and write a synthesis of about 80 words.

\textbf{Text A (Ministry of Public Health communiqu\'e)}: ``The free malaria treatment policy for children under five has reduced infant mortality by 18\% in the last three years. Health centres report higher attendance and earlier presentation of cases.''

\textbf{Text B (NGO field report)}: ``While the policy is welcome, drug stock-outs last an average of three weeks per quarter in rural clinics. Mothers travel long distances only to be turned away. The policy exists on paper but not in the pharmacy.''
\end{exercice}

\begin{corrige}[Exercise 3]
\textbf{Model synthesis} (78 words): Both texts address the free malaria treatment policy for under-fives. The Ministry communiqu\'e highlights its success: an 18\% drop in infant mortality and increased early attendance at health centres. The NGO report, however, reveals a critical implementation gap: rural clinics suffer drug stock-outs averaging three weeks per quarter, forcing mothers to travel in vain. While the policy's intent is proven beneficial, its impact is undermined by supply-chain failures. The sources together suggest that policy success depends on reliable drug distribution, not merely declaration.
\end{corrige}

You now possess the full toolkit for transforming reading into writing: compressing single texts, merging several texts and reorganising information into official formats. These skills will be exercised throughout the year --- in the prescribed-text sections, where you will summarise chapters and synthesise critical opinions, and in Section 7, where synthesis feeds directly into argumentative essay writing.

\coursec{Section 4: Grammar Mastery I — Clauses, Tenses and Conditionals}

In Section 3 you learnt to compress long texts into summaries and to rebuild formal reports. Both tasks depend on one hidden skill: the ability to \emph{control sentence structure}. A summary of a 900-word passage forces you to pack several ideas into single sentences using subordinate clauses; a report demands precise tense sequences (\emph{The committee \textbf{has been meeting} since January}) and conditional reasoning (\emph{Were the funds released, the project would proceed}). Grammar is therefore not a separate subject from comprehension and writing — it is the machinery underneath them. In this section we overhaul that machinery: first the parts of speech that carry meaning, then the three families of subordinate clauses, then the perfect tenses, and finally the full system of conditionals. We close with sentence-transformation practice in the exact format the GCE examiners use.

\courssub{Parts of Speech Review: Nouns, Pronouns and Verbs at Advanced Level}

You have known the parts of speech since primary school. At Advanced Level, the question is no longer \emph{What is a noun?} but \emph{What work can a noun, pronoun or verb do inside a complex sentence?}

\textbf{Nouns beyond the basics.}\par
A noun is a naming word, but Upper Sixth candidates must handle its advanced functions:
\begin{itemize}
\item \textbf{Abstract nouns} built from verbs and adjectives: \emph{govern} $\rightarrow$ \emph{governance}; \emph{corrupt} $\rightarrow$ \emph{corruption}; \emph{free} $\rightarrow$ \emph{freedom}; \emph{decide} $\rightarrow$ \emph{decision}; \emph{visible} $\rightarrow$ \emph{visibility}. These are essential in summary writing, where \emph{``the minister stole public money and the people protested''} becomes \emph{``the minister's corruption provoked public protest''}.
\item \textbf{Nominalisation}: turning a whole clause into a noun phrase. \emph{Because the rains failed, famine spread} $\rightarrow$ \emph{The failure of the rains caused famine to spread}. \emph{The students arrived late because the bus broke down} $\rightarrow$ \emph{The students' late arrival resulted from the bus breakdown}. Nominalisation is the engine of formal, academic style.
\item \textbf{Apposition}: a noun phrase renaming another: \emph{Ben Jama, a Cameroonian novelist, explores rural change.} \emph{The principal, Mr. Nfor, addressed the assembly.}
\item \textbf{Collective nouns and agreement}: \emph{The committee has decided} (singular focus) vs. \emph{The committee are arguing among themselves} (plural focus). GCE examiners test this distinction.
\end{itemize}

\textbf{Pronouns and reference.}\par
A pronoun replaces a noun or noun phrase. The examiner tests whether you can track \cle{reference} — what a pronoun points back to — in comprehension passages, and whether your own pronouns are unambiguous in essays.
\begin{itemize}
\item \textbf{Relative pronouns}: \emph{who, whom, whose, which, that} — they introduce adjective clauses (below).
\item \textbf{Demonstrative and indefinite pronouns}: \emph{this, that, each, none, either, neither, one, all, some}.
\item \textbf{Reflexive and emphatic pronouns}: \emph{The chief himself signed the decree.} \emph{She prepared herself for the interview.}
\item \textbf{Reciprocal pronouns}: \emph{each other, one another} — \emph{The candidates helped one another.}
\end{itemize}

\begin{attention}[Ambiguous pronoun reference]
In \emph{``The teacher told the student that he had failed''}, who failed? In your essays, every \emph{he, it, this, they} must point to one clear noun. Rewrite as \emph{``The teacher told the student, `You have failed'\,''} or \emph{``The teacher admitted that he himself had failed the student.''} Ambiguous reference costs marks in Paper 2 and confuses readers in summary questions. \textbf{Watch especially for \emph{this} and \emph{which} referring to whole clauses.}
\end{attention}

\textbf{Verbs: the heart of the clause.}\par
Every clause contains a finite verb, and the verb determines the clause's tense, voice and mood. At this level you must distinguish:
\begin{itemize}
\item \textbf{Transitive} verbs (take an object): \emph{The drought destroyed the harvest.}
\item \textbf{Intransitive} verbs (no object): \emph{The old man slept.}
\item \textbf{Linking} verbs (connect subject to complement): \emph{The evidence seems convincing.}
\item \textbf{Ditransitive} verbs (two objects): \emph{The examiner gave the candidates a break.}
\item \textbf{Auxiliaries}: primary (\emph{be, have, do}) and modal (\emph{can, could, may, might, must, shall, should, will, would}). Modals carry attitude — certainty, obligation, possibility — which is why they matter when you analyse a writer's tone.
\item \textbf{Ergative verbs} (same form, transitive/intransitive): \emph{The sun melted the ice} / \emph{The ice melted.}
\end{itemize}

\begin{aretenir}[Why this review matters]
Comprehension questions such as \emph{``What does `it' in line 14 refer to?''} test pronoun reference. Summary questions reward nominalisation. Error-correction questions test verb agreement. The parts of speech are not theory; they are marks.
\end{aretenir}

\courssub{Subordinate Clauses: Noun, Adjective and Adverb Clauses}

Consider these two sentences:
\begin{enumerate}
\item \emph{The rains failed. The village starved.}
\item \emph{Because the rains failed, the village starved.}
\end{enumerate}
Sentence (2) is superior for academic writing because it shows the \emph{relationship} between the ideas. The words \emph{because the rains failed} form a \cle{subordinate clause}: it has a subject and a verb, but it cannot stand alone.

\begin{definition}[Subordinate clause]
A \cle{subordinate clause} is a clause that cannot stand alone as a complete sentence and that functions as a \textbf{noun}, an \textbf{adjective} or an \textbf{adverb} within a larger sentence. It is introduced by a subordinating conjunction (\emph{because, although, when, so that, if\ldots}) or by a relative pronoun (\emph{who, which, that, whose\ldots}).
\end{definition}

\begin{methode}[Identifying the clause type]
To classify any subordinate clause, ask what \emph{question} it answers in the sentence:
\begin{enumerate}
\item Does the clause act as \textbf{subject, object or complement} of the main verb? (Ask: \emph{who? what?}) $\rightarrow$ \textbf{Noun clause}. \emph{What the witness said} shocked the court.
\item Does the clause \textbf{describe a noun or pronoun}? (Ask: \emph{which one? what kind?}) $\rightarrow$ \textbf{Adjective (relative) clause}. The girl \emph{who won the scholarship} comes from Bamenda.
\item Does the clause modify the \textbf{verb}, telling \emph{when, why, although, for what purpose, under what condition}? $\rightarrow$ \textbf{Adverb clause}. \emph{Although the road was flooded}, the lorry reached Wum.
\end{enumerate}
\end{methode}

\begin{popfigure}
\begin{tikzpicture}[
grow=down,
level 1/.style={sibling distance=4.6cm, level distance=1.6cm},
level 2/.style={sibling distance=4.4cm, level distance=2.2cm},
every node/.style={draw=popline, rounded corners=2pt, fill=white, align=center, font=\small, text width=3.9cm},
edge from parent/.style={draw=popmuted, thick}
]
\node[fill=popblueL, font=\small\bfseries] {Subordinate Clauses}
  child { node[fill=poppurple!12, align=center]{\textbf{Noun clause}\\ \emph{That he lied} was obvious.}
    child { node {subject: \emph{What she said} matters} }
    child { node {object: I believe \emph{that God is just}} }
    child { node {complement: The fact is \emph{that he left}} }
  }
  child { node[fill=popgreenL, align=center]{\textbf{Adjective clause}\\ the man \emph{whose farm burnt}}
    child { node {defining: students \emph{who cheat} fail} }
    child { node {non-defining: My uncle, \emph{who lives in Buea}, is a teacher} }
    child { node {reduced: The girl \emph{sitting there} is my sister} }
  }
  child { node[fill=poporange!15, align=center]{\textbf{Adverb clause}\\ \emph{Although tired}, she read on}
    child { node {time/reason: \emph{When the bell rang}; \emph{because it rained}} }
    child { node {concession/purpose: \emph{though poor}; \emph{so that all may pass}} }
    child { node {condition/place: \emph{if ready}; \emph{where the river bends}} }
  };
\end{tikzpicture}
\legende{Classification of subordinate clauses with examples of each function.}
\end{popfigure}

\textbf{Noun clauses.}\par
A noun clause does the job of a noun. It can be:
\begin{itemize}
\item \textbf{Subject}: \emph{That the earth is warming} is no longer disputed. \emph{Whether he will come} remains uncertain.
\item \textbf{Object}: The report stated \emph{that enrolment had doubled}. I wonder \emph{why she left}.
\item \textbf{Complement}: The truth is \emph{that nobody was prepared}. The question is \emph{whether we can afford it}.
\item \textbf{Object of a preposition}: They argued about \emph{who should lead}.
\end{itemize}
Noun clauses are introduced by \emph{that}, \emph{whether/if}, or question words (\emph{what, who, how, why, where, when}).

\textbf{Adjective (relative) clauses.}\par
These describe a noun and answer \emph{which one?} Distinguish:
\begin{itemize}
\item \textbf{Defining} (no commas; the clause identifies the noun): \emph{Candidates who arrive late will not be admitted.} \emph{The book that you lent me is fascinating.}
\item \textbf{Non-defining} (commas; extra information): \emph{Mount Cameroon, which rises over 4\,000 metres, dominates the coast.} \emph{My father, who is a doctor, works in Yaoundé.}
\item \textbf{Reduced relative clauses} (participle phrases): \emph{The girl \textbf{sitting} by the window is my cousin.} (= \emph{who is sitting}) \emph{The bridge \textbf{built} last year has collapsed.} (= \emph{which was built})
\end{itemize}

\begin{attention}[Relative pronoun choice]
\begin{itemize}
\item \emph{Who/whom} for people; \emph{which} for things; \emph{that} for both (defining only).
\item \emph{Whose} for possession (people and things): \emph{The house whose roof leaked.}
\item \emph{Where, when, why} as relative adverbs: \emph{The village where I was born.} \emph{The year when independence came.}
\item Never use \emph{what} as a relative pronoun in standard English: \emph{*The thing what he said} $\rightarrow$ \emph{The thing that/which he said}.
\end{itemize}
\end{attention}

\textbf{Adverb clauses.}\par
These modify the verb of the main clause:
\begin{itemize}
\item \textbf{Time}: \emph{When the harmattan came}, the wells dried up. (\emph{when, while, before, after, until, as soon as, once, since})
\item \textbf{Reason}: \emph{Because the bridge collapsed}, traders used canoes. (\emph{because, since, as, now that})
\item \textbf{Concession}: \emph{Although he studied hard}, he found Paper 1 demanding. (\emph{although, though, even though, whereas, while})
\item \textbf{Purpose}: She revised at night \emph{so that she could pass}. (\emph{so that, in order that, lest})
\item \textbf{Condition}: \emph{If you work hard}, you will pass. (\emph{if, unless, provided that, as long as, on condition that})
\item \textbf{Place}: \emph{Where the road ends}, the forest begins. (\emph{where, wherever})
\item \textbf{Manner}: \emph{He spoke as if he knew everything.} (\emph{as, as if, as though})
\item \textbf{Result}: \emph{It rained so heavily that the match was cancelled.} (\emph{so...that, such...that})
\end{itemize}

\begin{attention}[Fragments and comma splices]
A subordinate clause alone is a \emph{fragment}: \emph{Because the examination was difficult.} is not a sentence. Conversely, joining two main clauses with only a comma is a \emph{comma splice}: \emph{The examination was difficult, many candidates failed.} Correct it with a subordinator (\emph{Because the examination was difficult, many candidates failed}), a semicolon (\emph{The examination was difficult; many candidates failed}), or a full stop.
\end{attention}

\courssub{Advanced Tenses: Perfect and Perfect Continuous Forms}

The perfect tenses connect two moments in time. The \cle{perfect} aspect (\emph{have + past participle}) presents an event as \emph{completed with relevance} to a reference point; the \cle{perfect continuous} (\emph{have been + -ing}) presents it as \emph{ongoing up to} that reference point, often emphasising duration.

\begin{popfigure}
\begin{tikzpicture}[scale=1.0]
% axis
\draw[->, thick, popdark] (0,0) -- (10.5,0) node[below left, font=\small] {time};
\draw[thick, popdark] (5,0.15) -- (5,-0.15) node[below, font=\small] {now};
% present perfect continuous: ongoing up to now
\draw[very thick, popblue] (1,0.6) -- (5,0.6);
\draw[->, very thick, popblue] (5,0.6) -- (5.6,0.6);
\fill[popblue] (1,0.6) circle (1.5pt);
\node[above, font=\small, text=popblue, align=center] at (3,0.65) {has been raining (since 6 a.m.)};
% present perfect: completed before now, result now
\draw[very thick, poporange] (2,1.4) -- (4.5,1.4);
\fill[poporange] (4.5,1.4) circle (1.5pt);
\node[above, font=\small, text=poporange, align=center] at (3.25,1.45) {has built three colleges};
% past perfect: completed before past ref
\draw[very thick, popgreen] (6,2.2) -- (8,2.2);
\fill[popgreen] (8,2.2) circle (1.5pt);
\draw[dashed, popmuted] (8,2.2) -- (8,0) node[below, font=\small] {past ref.};
\node[above, font=\small, text=popgreen!60!black, align=center] at (7,2.25) {had assembled (before examiners arrived)};
% future perfect: completed before future ref
\draw[very thick, poppurple] (5,3.0) -- (8.5,3.0);
\fill[poppurple] (8.5,3.0) circle (1.5pt);
\draw[dashed, popmuted] (8.5,3.0) -- (8.5,0) node[below, font=\small] {future ref.};
\node[above, font=\small, text=poppurple, align=center] at (6.75,3.05) {will have trained (by 2030)};
% past perfect continuous
\draw[very thick, poppink] (6.5,3.8) -- (8.5,3.8);
\draw[->, very thick, poppink] (8.5,3.8) -- (9,3.8);
\fill[poppink] (6.5,3.8) circle (1.5pt);
\draw[dashed, popmuted] (8.5,3.8) -- (8.5,0);
\node[above, font=\small, text=poppink, align=center] at (7.5,3.85) {had been teaching 20 yrs (before retirement)};
\end{tikzpicture}
\legende{Perfect forms anchor an event to a reference point: duration up to now (perfect continuous), completion before a past moment (past perfect), completion before a future moment (future perfect).}
\end{popfigure}

\begin{center}
\begin{tabularx}{\linewidth}{|l|X|X|}
\hline
\rowcolor{popblueL} \textbf{Tense} & \textbf{Form} & \textbf{Authentic example} \\
\hline
Present perfect & have/has + past participle & The government \emph{has built} three new colleges. (result visible now) \\
\hline
Present perfect continuous & have/has been + -ing & Farmers \emph{have been waiting} for rain since March. (duration up to now) \\
\hline
Past perfect & had + past participle & When the examiners arrived, the candidates \emph{had assembled}. (earlier of two past events) \\
\hline
Past perfect continuous & had been + -ing & She \emph{had been teaching} for twenty years before she retired. (duration before a past event) \\
\hline
Future perfect & will have + past participle & By 2030, Cameroon \emph{will have trained} thousands of engineers. (completed before a future point) \\
\hline
Future perfect continuous & will have been + -ing & In June, we \emph{will have been studying} this text for six months. \\
\hline
\end{tabularx}
\end{center}

\textbf{Perfect versus perfect continuous: the meaning contrast.}\par
\begin{itemize}
\item \emph{I \textbf{have read} the novel.} $\rightarrow$ focus on \textbf{completion/result} (I finished it).
\item \emph{I \textbf{have been reading} the novel.} $\rightarrow$ focus on \textbf{activity/duration} (perhaps unfinished; explains why I am tired).
\item \emph{She \textbf{has written} three letters.} (completed actions, countable)
\item \emph{She \textbf{has been writing} letters all morning.} (ongoing activity, duration)
\end{itemize}
With verbs that do not take the continuous form (\emph{know, believe, own, understand, love, hate, seem, contain, consist}), only the simple perfect is possible: \emph{I have known her for years} (never \emph{*have been knowing}). \emph{This book has belonged to me since 2010.}

\begin{exoresolu}[Choosing the correct perfect form]
Complete with the correct perfect or perfect continuous form: \emph{By the time the invigilator collected the scripts, most candidates (write) for three hours, and everyone (finish) the comprehension section.}
\tcblower
\textbf{Solution.} Two events are anchored to a past reference point (\emph{collected}). The writing is presented as a \emph{duration leading up to} that point, so we use the \textbf{past perfect continuous}: \emph{had been writing}. The finishing is presented as a \emph{completed result} before that point, so we use the \textbf{past perfect simple}: \emph{had finished}.
\emph{By the time the invigilator collected the scripts, most candidates \textbf{had been writing} for three hours, and everyone \textbf{had finished} the comprehension section.}
\end{exoresolu}

\begin{exoresolu}[Perfect tense in reported speech and conditionals]
Transform into reported speech: \emph{``I have finished the essay,'' she said.} \quad \emph{``We have been waiting since 8 a.m.,'' they complained.}
\tcblower
\textbf{Solution.} Present perfect $\rightarrow$ past perfect; present perfect continuous $\rightarrow$ past perfect continuous.
\emph{She said (that) she \textbf{had finished} the essay.} \quad \emph{They complained (that) they \textbf{had been waiting} since 8 a.m.}
\end{exoresolu}

\courssub{Conditional Sentences: Zero to Mixed}

Conditionals express the link between a condition (the \emph{if}-clause, or protasis) and a result (the main clause, or apodosis). The tense combination signals \emph{how real} the speaker considers the condition.

\begin{center}
\begin{tabularx}{\linewidth}{|l|X|X|}
\hline
\rowcolor{popblueL} \textbf{Type} & \textbf{Structure} & \textbf{Meaning and example} \\
\hline
Zero & If + present, present & General truth: \emph{If you heat water to 100\textdegree C, it boils.} \\
\hline
First & If + present, will + verb & Real future possibility: \emph{If it rains, the match will be postponed.} \\
\hline
Second & If + past, would + verb & Unreal present/future: \emph{If I were the principal, I would abolish corporal punishment.} \\
\hline
Third & If + past perfect, would have + past participle & Unreal past: \emph{If the driver had been careful, the accident would not have happened.} \\
\hline
Mixed (Type 3/2) & If + past perfect, would + verb & Unreal past condition, present result: \emph{If I had taken science, I would be a doctor now.} \\
\hline
Mixed (Type 2/3) & If + past, would have + past participle & Unreal present condition, past result: \emph{If he were honest, he would have returned the money.} \\
\hline
\end{tabularx}
\end{center}

\begin{propriete}[The third conditional]
The \cle{third conditional} expresses an \textbf{unreal past}: a condition that was \emph{not} fulfilled and a result that therefore did \emph{not} occur.
\[ \text{If} + \text{past perfect},\ \ldots\ \text{would have} + \text{past participle} \]
\emph{If the students \textbf{had revised} the conditionals, they \textbf{would have passed} the grammar section.} (They did not revise; they did not pass.) The structure is examinable in sentence transformation and in error correction; memorise the formula exactly.
\end{propriete}

\textbf{Meaning contrasts to master.}\par
\begin{itemize}
\item \emph{If she \textbf{works} hard, she \textbf{will succeed}.} — I think success is genuinely possible (first).
\item \emph{If she \textbf{worked} hard, she \textbf{would succeed}.} — She does not work hard; the situation is hypothetical (second).
\item \emph{If she \textbf{had worked} hard, she \textbf{would have succeeded}.} — She did not work hard, and the chance is gone (third).
\item \emph{If she \textbf{had worked} hard, she \textbf{would be} confident now.} — past cause, present consequence (mixed 3/2).
\item \emph{If she \textbf{were} honest, she \textbf{would have told} the truth.} — present trait, past consequence (mixed 2/3).
\end{itemize}

\begin{attention}[Never ``would'' in the if-clause]
The most frequent GCE error: \emph{*If I \textbf{would have} known, I \textbf{would have} come.} The if-clause takes the \textbf{past perfect}, never \emph{would have}. Correct: \emph{If I \textbf{had known}, I would have come.} Similarly: \emph{*If it \textbf{would rain}\ldots} must be \emph{If it \textbf{rained}\ldots} (second conditional). Note also the formal inversion: \emph{\textbf{Had I known}, I would have come} (= \emph{If I had known}) and \emph{\textbf{Were I} the minister\ldots} (= \emph{If I were}). \textbf{Should} can replace \emph{if} in first conditionals: \emph{Should you need help, call me.}
\end{attention}

\begin{exoresolu}[Conditional identification and correction]
Identify the conditional type and correct any error:
\begin{enumerate}
\item \emph{If the government will increase salaries, teachers will call off the strike.}
\item \emph{If I was you, I would apply for that scholarship.}
\item \emph{Had they not intervened, the situation would be worse now.}
\end{enumerate}
\tcblower
\textbf{Solutions.}
\begin{enumerate}
\item \textbf{Error}: \emph{will} in if-clause. \textbf{Type}: First conditional. \textbf{Correction}: \emph{If the government \textbf{increases} salaries, teachers will call off the strike.}
\item \textbf{Type}: Second conditional. \textbf{Correction}: \emph{If I \textbf{were} you, I would apply\ldots} (Formal English prefers \emph{were} for all persons; \emph{was} is colloquial.)
\item \textbf{Type}: Mixed (3/2) — inverted third conditional with present result. \textbf{Correct as written}: \emph{Had they not intervened, the situation would be worse now.} (= \emph{If they had not intervened\ldots})
\end{enumerate}
\end{exoresolu}

\courssub{Sentence Transformation: Combining Clauses, Tenses and Conditionals}

GCE Paper 2 regularly asks you to rewrite a sentence beginning with a given word, keeping the meaning. This tests clauses, tenses and conditionals together.

\begin{methode}[Sentence transformation]
\begin{enumerate}
\item \textbf{Identify} the grammatical relationship in the original (cause? condition? concession? time?).
\item \textbf{Choose} the structure the new opening word demands (e.g. \emph{Had\ldots} $\rightarrow$ inverted third conditional; \emph{Despite\ldots} $\rightarrow$ noun phrase or -ing clause; \emph{Not only\ldots} $\rightarrow$ inversion).
\item \textbf{Adjust} tenses consistently (a past condition keeps the past perfect).
\item \textbf{Check} that no meaning has been added or lost, and that the sentence is grammatically complete.
\end{enumerate}
\end{methode}

\begin{exoresolu}[Exam-style transformations]
Rewrite each sentence beginning with the words given, without changing the meaning.
\begin{enumerate}
\item \emph{Because the bridge was flooded, the traders could not cross.} \\ Begin: \emph{The flooding of the bridge\ldots}
\item \emph{She did not revise the tenses, so she failed the grammar section.} \\ Begin: \emph{Had\ldots}
\item \emph{Although he is poor, he is honest.} \\ Begin: \emph{Poor\ldots}
\item \emph{The students worked hard. They passed the examination.} \\ Begin: \emph{Had the students not\ldots}
\item \emph{``Don't touch the wire,'' the electrician warned us.} \\ Begin: \emph{The electrician warned us\ldots}
\end{enumerate}
\tcblower
\textbf{Solutions.}
\begin{enumerate}
\item \emph{The flooding of the bridge \textbf{prevented the traders from crossing}.} (Nominalisation turns the adverb clause of reason into a noun phrase as subject.)
\item \emph{\textbf{Had she revised} the tenses, \textbf{she would not have failed} the grammar section.} (Inverted third conditional: \emph{Had + subject + past participle}; the result clause keeps \emph{would have + past participle}, negated to preserve meaning.)
\item \emph{\textbf{Poor as/though he is}, he is honest.} (Fronted concessive structure; the adjective moves to the front with \emph{as} or \emph{though}.)
\item \emph{\textbf{Had the students not worked hard}, they \textbf{would not have passed} the examination.} (Inverted third conditional, negative condition.)
\item \emph{\textbf{The electrician warned us not to touch the wire}.} (Reported command: \emph{warn + object + not to-infinitive}.)
\end{enumerate}
\end{exoresolu}

\begin{exoresolu}[Advanced transformation: cleft sentences and passive]
Rewrite beginning with the words given:
\begin{enumerate}
\item \emph{The principal announced the results yesterday.} \\ Begin: \emph{It was\ldots}
\item \emph{They are building a new library.} \\ Begin: \emph{A new library\ldots}
\item \emph{He not only passed the exam but also won a prize.} \\ Begin: \emph{Not only\ldots}
\end{enumerate}
\tcblower
\textbf{Solutions.}
\begin{enumerate}
\item \emph{\textbf{It was yesterday that the principal announced the results}.} (Cleft sentence for focus on time.)
\item \emph{\textbf{A new library is being built}.} (Passive voice, present continuous.)
\item \emph{\textbf{Not only did he pass the exam, but he also won a prize}.} (Inversion after fronted negative adverbial.)
\end{enumerate}
\end{exoresolu}

\begin{exercice}[Consolidation practice]
\begin{enumerate}
\item Classify the underlined clause as noun, adjective or adverb: \emph{The minister announced \textbf{that the exams would begin in June}.}
\item Combine into one sentence using a relative clause: \emph{The novel is set in a village. We studied the novel last term.}
\item Correct the error: \emph{If I would have had money, I would have bought the textbook.}
\item Rewrite beginning with the given words: \emph{He worked hard, yet he failed.} Begin: \emph{Despite\ldots}
\item Fill the gap: \emph{By next April, the class \underline{\hspace{2cm}} (study) four prescribed texts.}
\item Choose the correct form: \emph{If the bus (come / comes / came / had come) on time, we would not have missed the lecture.}
\item Transform: \emph{The rain was so heavy that the match was cancelled.} Begin: \emph{So heavy\ldots}
\item Identify the tense: \emph{By the time you arrive, I \textbf{will have been waiting} for two hours.}
\item Correct: \emph{The teacher asked that the students \textbf{should to} bring their textbooks.}
\item Rewrite: \emph{``I have never seen such a beautiful sunset,'' she said.} Begin: \emph{She exclaimed\ldots}
\end{enumerate}
\end{exercice}

\begin{corrige}[Exercise n° 1]
\begin{enumerate}
\item \textbf{Noun clause} — it is the object of \emph{announced} (answers: announced \emph{what?}).
\item \emph{The novel \textbf{which/that we studied last term} is set in a village.} (Defining relative clause.)
\item \emph{If I \textbf{had had} money, I would have bought the textbook.} (No \emph{would} in the if-clause; third conditional requires the past perfect.)
\item \emph{\textbf{Despite} working hard, he failed.} (or \emph{Despite his hard work, he failed.})
\item \emph{\textbf{will have studied}} (future perfect: completion before a future reference point).
\item \emph{\textbf{had come}} (third conditional: unreal past condition).
\item \emph{\textbf{So heavy was the rain that the match was cancelled}.} (Fronted \emph{so + adjective} with inversion.)
\item \emph{\textbf{Future perfect continuous}} (duration up to a future reference point).
\item \emph{The teacher asked that the students \textbf{bring} (or \textbf{should bring}) their textbooks.} (Subjunctive after \emph{ask/demand/insist}; no \emph{to} after \emph{should}.)
\item \emph{\textbf{She exclaimed that she had never seen such a beautiful sunset}.} (Reported exclamation; present perfect $\rightarrow$ past perfect.)
\end{enumerate}
\end{corrige}

\begin{aretenir}[Key points of Section 4]
\begin{itemize}
\item A subordinate clause functions as a noun, adjective or adverb; identify it by the question it answers.
\item Perfect = completion with relevance to a reference point; perfect continuous = duration up to that point.
\item Stative verbs (\emph{know, believe, own, love, hate, seem, contain}) do not take continuous forms.
\item Third conditional: \emph{If + past perfect, would have + past participle}; never \emph{would} in the if-clause.
\item Mixed conditionals combine different time references (past condition / present result, or present condition / past result).
\item Formal inversion: \emph{Had I known\ldots}, \emph{Were I\ldots}, \emph{Should you\ldots}.
\item Transformations test meaning preservation: identify the relationship, apply the required structure, keep tenses consistent.
\end{itemize}
\end{aretenir}

With this grammatical machinery in place, you are ready for Section 5, where we refine the sentence further: dangling and misplaced modifiers, idiomatic expressions, and the fine nuances of connotation and denotation.

\coursec{Section 5: Grammar Mastery II — Modifiers, Idioms and Word Nuance}

In Section 4 you mastered the \emph{architecture} of the sentence: clauses, tenses and conditionals. You can now build grammatically correct sentences. But correctness alone does not earn top marks at the GCE Advanced Level. Examiners also reward \emph{precision} — placing modifiers where they belong, choosing the word with exactly the right shade of meaning, deploying an idiom or a figure of speech with control, and editing a faulty text into clean, formal English. That precision is the business of this section. We move from \emph{writing correct sentences} to \emph{writing exact sentences}.

\courssub{Dangling and Misplaced Modifiers}

\textbf{The intuition.}\par
Read this sentence aloud: \emph{Walking to school, the rain started.} Something feels wrong, doesn't it? Your ear tells you that the rain was walking to school — which is absurd. English has a powerful convention: a modifying phrase placed at the start of a sentence is understood to describe the \textbf{subject of the main clause}. When that subject is not the doer of the action in the phrase, the modifier is left \emph{dangling} — attached to nothing, or attached to the wrong thing.

\begin{definition}[Dangling Modifier]
A \cle{dangling modifier} is a modifying word or phrase (usually a participial, infinitive or elliptical phrase) whose intended subject is \textbf{missing from the sentence}, so that the phrase appears to modify the wrong noun — or nothing at all.
\end{definition}

\begin{definition}[Misplaced Modifier]
A \cle{misplaced modifier} is a word or phrase placed too far from the word it is meant to modify, so that it seems to modify a different word and creates ambiguity or absurdity. Example: \emph{She served tea to the guests in plastic cups.} (Were the guests in plastic cups?)
\end{definition}

\begin{attention}[Common Mistake — The Walking Rain]
\emph{Walking to school, the rain started.} This claims the rain was walking. Correct it either by supplying the true subject — \emph{Walking to school, \textbf{I was caught by the rain}} — or by converting the phrase into a full clause — \emph{\textbf{As I was walking to school}, the rain started}. Never leave the introductory participle guessing at its subject.
\end{attention}

\begin{methode}[Correcting Dangling and Misplaced Modifiers]
\begin{enumerate}
\item \textbf{Identify} the modifier and ask: \emph{who or what} is doing/being described?
\item \textbf{Locate} that doer in the sentence. If it is absent, the modifier dangles; if it is present but distant, the modifier is misplaced.
\item \textbf{Repair} by one of three strategies: (a) name the true subject immediately after the phrase; (b) expand the phrase into a full subordinate clause with its own subject; (c) move the modifier directly beside the word it modifies.
\item \textbf{Re-read} the sentence and check that the nearest noun to the modifier is genuinely the one described.
\end{enumerate}
\end{methode}

\begin{exemplebox}[Diagnosis and Repair]
\begin{center}
\begin{tabularx}{\linewidth}{|l|X|X|}
\hline
\rowcolor{popblueL} \textbf{Faulty sentence} & \textbf{Problem} & \textbf{Correction} \\
\hline
Flying over Mount Cameroon, the plains looked endless. & The plains were not flying (dangling). & Flying over Mount Cameroon, \textbf{we saw} the endless plains. \\
\hline
To pass Paper 1, regular practice is needed. & Who must practise? (dangling infinitive) & To pass Paper 1, \textbf{a candidate must practise} regularly. \\
\hline
The principal punished the boy with a stern warning. & Did the boy carry a warning? (misplaced) & With a stern warning, the principal punished the boy. \\
\hline
While eating fufu, the phone rang. & The phone was not eating (dangling). & While \textbf{I was} eating fufu, the phone rang. \\
\hline
\end{tabularx}
\end{center}
\end{exemplebox}

\begin{attention}[The Adverb Trap: \emph{only}, \emph{almost}, \emph{just}, \emph{nearly}]
Single-word modifiers are the most frequently misplaced. Compare:
\begin{itemize}
\item \emph{\textbf{Only} she lent me 500 francs} (nobody else lent me money).
\item \emph{She \textbf{only} lent me 500 francs} (she did not give me the money; she merely lent it).
\item \emph{She lent me \textbf{only} 500 francs} (the amount was no more than 500 francs).
\end{itemize}
Place \emph{only}, \emph{almost}, \emph{just}, \emph{nearly} immediately before the word they limit.
\end{attention}

\courssub{Connotation and Denotation — Advanced Nuance in Word Choice}

\textbf{The intuition.}\par
Suppose you describe your grandmother as \emph{frugal}. She smiles. Now call her \emph{stingy}. The smile disappears. Both words point to the same core idea — carefulness with money — yet one praises and the other insults. The core idea is the \cle{denotation}: the literal, dictionary meaning of a word. The emotional and cultural associations that cling to it are its \cle{connotation}. Advanced writing, and the comprehension questions that ask you to comment on the writer's \emph{choice of words}, depend entirely on this distinction.

\begin{definition}[Denotation and Connotation]
The \cle{denotation} of a word is its explicit, objective, dictionary meaning. The \cle{connotation} of a word is the set of feelings, attitudes and associations it evokes beyond its literal meaning — positive, negative or neutral.
\end{definition}

\begin{popfigure}
\begin{center}
\begin{tabularx}{\linewidth}{|l|X|X|X|}
\hline
\rowcolor{popblueL} \textbf{Core idea (denotation)} & \textbf{Positive connotation} & \textbf{Neutral connotation} & \textbf{Negative connotation} \\
\hline
Not fat & slim, slender & thin & skinny, scrawny \\
\hline
Careful with money & frugal, thrifty & economical & stingy, miserly \\
\hline
Not young & mature, elderly & old & ancient, decrepit \\
\hline
Firm in opinion & determined, resolute & firm & stubborn, pig-headed \\
\hline
Interested in others' affairs & curious, inquisitive & questioning & nosy, prying \\
\hline
Inexpensive & affordable, low-cost & cheap & shoddy, cut-rate \\
\hline
\end{tabularx}
\end{center}
\legende{Graded word sets: same denotation, different connotations. In essays and comprehension answers, always justify a word choice by its connotative effect.}
\end{popfigure}

\begin{exemplebox}[Nuance in Action]
Compare: \emph{The minister gave a \textbf{frugal} reply} (admiring — he wasted no words) and \emph{The minister gave a \textbf{stingy} reply} (critical — he withheld what he should have shared). The denotation — a small reply — is identical; the connotation reverses the judgement. In summary writing, replacing the author's \emph{slim} with \emph{skinny} silently changes the author's attitude, and you will lose marks for distortion.
\end{exemplebox}

\begin{attention}[Common Mistake in Summaries and Essays]
When summarising, do not ``improve'' the author's vocabulary with a near-synonym of a different connotation. When writing essays, do not reach for a grand word whose connotation you have not checked: calling a freedom fighter a \emph{notorious} activist (notorious = famous \emph{for something bad}) undermines your own argument.
\end{attention}

\courssub{Advanced Idiomatic Expressions}

\textbf{The intuition.}\par
If a friend says the HOD \emph{let the cat out of the bag} about the strike action, you know no animal was involved: a secret was revealed. An \cle{idiom} is a fixed expression whose meaning cannot be deduced from the literal meanings of its parts. Idioms are learned whole, like vocabulary items, and they carry a register: many belong to informal speech, some are neutral, and a few are perfectly at home in formal prose.

\begin{definition}[Idiom]
An \cle{idiom} is a fixed, conventional phrase whose overall meaning is figurative and cannot be predicted from the meanings of its individual words (e.g. \emph{to burn the midnight oil} = to study or work late into the night).
\end{definition}

\begin{exemplebox}[Idioms by Register]
\begin{center}
\begin{tabularx}{\linewidth}{|l|X|X|}
\hline
\rowcolor{popblueL} \textbf{Idiom} & \textbf{Meaning} & \textbf{Register} \\
\hline
to burn the midnight oil & to work/study late & neutral — fine in essays \\
\hline
to turn a blind eye (to) & to deliberately ignore & neutral \\
\hline
to break the ice & to ease initial social tension & informal/neutral \\
\hline
to hit the nail on the head & to state exactly the truth & informal/neutral \\
\hline
to kick the bucket & to die & very informal — avoid in essays \\
\hline
to bite the dust & to fail, to be defeated & informal \\
\hline
a blessing in disguise & a hidden benefit & neutral \\
\hline
\end{tabularx}
\end{center}
\end{exemplebox}

\begin{aretenir}[Exam Tip — Idioms in Essays]
Idioms impress examiners \textbf{only when used accurately and sparingly}. One or two well-placed idioms (\emph{The government turned a blind eye to the corruption eating into the council's funds}) signal maturity. Five idioms in one paragraph signal showing off, and a single mangled idiom (\emph{he burnt the midnight candle}) signals carelessness. If you are not sure of the exact wording, use plain language instead.
\end{aretenir}

\begin{attention}[Common Mistakes with Idioms]
\begin{itemize}
\item \textbf{Altering the fixed form}: idioms resist change — \emph{let the cat out of the bag}, never \emph{release the cat from the bag}.
\item \textbf{Literal translation}: do not translate idioms from French or your mother tongue word for word; the result is usually not English.
\item \textbf{Register clash}: \emph{The deceased kicked the bucket} in a formal report or funeral speech is offensive flippancy.
\end{itemize}
\end{attention}

\courssub{Advanced Figures of Speech — Irony, Paradox, Hyperbole, Synecdoche, Metonymy}

You already know simile, metaphor and personification from Lower Sixth. The GCE A-Level expects you to recognise and comment on five more demanding figures.

\begin{definition}[Irony]
\cle{Irony} is a figure in which the intended meaning is the opposite of, or sharply different from, the literal meaning. \emph{Verbal irony}: saying the reverse of what you mean (\emph{``Lovely weather!''} during a downpour). \emph{Situational irony}: an outcome contrary to expectation (a fire station burning down). \emph{Dramatic irony}: the audience knows what a character does not.
\end{definition}

\begin{definition}[Paradox]
A \cle{paradox} is a statement that seems self-contradictory yet contains a truth: \emph{``The child is father of the man''}; \emph{``Less is more.''} Its effect is to force the reader to think beyond surface logic.
\end{definition}

\begin{definition}[Hyperbole]
\cle{Hyperbole} is deliberate exaggeration for emphasis or effect, not meant literally: \emph{``I have told you a million times''}; \emph{``The whole of Yaoundé was at the stadium.''}
\end{definition}

\begin{definition}[Synecdoche and Metonymy]
\cle{Synecdoche} names a thing by a \textbf{part of it} (or the whole by a part): \emph{``All hands on deck''} (hands = sailors); \emph{``fifty head of cattle''}. \cle{Metonymy} names a thing by something \textbf{closely associated} with it: \emph{``The Crown announced new taxes''} (the monarchy); \emph{``Etoudi has responded''} (the Presidency of Cameroon); \emph{``the pen is mightier than the sword''} (writing vs.\ warfare). Test: in synecdoche the word used is \emph{physically part of} the thing meant; in metonymy it is merely \emph{linked to} it.
\end{definition}

\begin{popfigure}
\begin{center}
\begin{tikzpicture}[
box/.style={draw=popblue, thick, rounded corners, fill=popblueL, text width=4.6cm, align=center, inner sep=4pt, font=\small},
ttl/.style={font=\small\bfseries, text=popdark}]
\node[box, align=center] (irony) at (0,0){\textcolor{popdark}{\textbf{IRONY}}\\ meaning = opposite of words\\ \emph{``What a genius!''} (of a blunder)};
\node[box, align=center] (paradox) at (5.6,0){\textcolor{popdark}{\textbf{PARADOX}}\\ contradiction that is true\\ \emph{``Less is more.''}};
\node[box, align=center] (hyper) at (11.2,0){\textcolor{popdark}{\textbf{HYPERBOLE}}\\ deliberate exaggeration\\ \emph{``I waited for ages.''}};
\node[box, align=center] (syn) at (2.8,-3.4){\textcolor{popdark}{\textbf{SYNECDOCHE}}\\ part stands for whole\\ \emph{``All hands on deck.''}};
\node[box, align=center] (met) at (8.4,-3.4){\textcolor{popdark}{\textbf{METONYMY}}\\ associated term for thing\\ \emph{``The Crown has spoken.''}};
\draw[->, thick, poporange] (syn) -- (met) node[midway, above, font=\scriptsize\itshape, text=popdark]{part vs.\ association};
\end{tikzpicture}
\end{center}
\legende{Five advanced figures of speech: definition plus one-line example. Synecdoche and metonymy are cousins — distinguish them by the part/association test.}
\end{popfigure}

\begin{exemplebox}[Recognition and Effect]
Sentence: \emph{``The bench laughed at the witness's trembling story — a thousand tongues could not have saved him.''} \emph{The bench} = metonymy (the judges, associated with where they sit); \emph{a thousand tongues} = hyperbole plus synecdoche (tongues = speakers). Effect: the courtroom feels hostile and the witness doomed. In the exam, always name the figure \textbf{and} state its effect; naming alone earns only half the marks.
\end{exemplebox}

\courssub{Editing Across Text Types — Newspapers, Journals and Academic Papers}

Editing tasks at A-Level present a flawed passage — a news report, an editorial, an abstract — and ask you to correct grammar, register and style. Each text type has its own expectations: newspapers want clarity and concision; academic papers want formality, precision and impersonal constructions; journals sit between the two. What never changes is the underlying grammar.

\begin{methode}[The Five-Point Editing Checklist]
\begin{enumerate}
\item \textbf{Agreement}: does every verb agree with its true subject? (Beware long subjects: \emph{The quality of the roads \textbf{has} — not \emph{have} — declined.})
\item \textbf{Modifiers}: hunt for dangling participles and misplaced \emph{only}/\emph{almost}/\emph{nearly}; reposition or supply the missing subject.
\item \textbf{Tense consistency}: a narrative begun in the past stays in the past unless there is a reason to shift; reports of research use the present for established findings.
\item \textbf{Register}: replace contractions (\emph{don't} $\rightarrow$ \emph{do not}), slang and vague words (\emph{stuff}, \emph{a lot of}) with formal equivalents suited to the text type.
\item \textbf{Punctuation}: check comma splices, apostrophes (\emph{its} vs.\ \emph{it's}), and capitalisation of titles and proper nouns.
\end{enumerate}
\end{methode}

\begin{exoresolu}[Editing a Newspaper Extract]
\textbf{Statement.} The following extract from a school magazine report contains five errors. Identify and correct them.\par
\emph{``Rushing to the scene, the fire was brought under control by the brigade. The amount of damages have not yet been estimated. The mayor, who only arrived at midnight, said the council don't intend to demolish the market. Its a great relief that no lives was lost.''}
\tcblower
\textbf{Detailed solution.}
\begin{enumerate}
\item \emph{Rushing to the scene, the fire\ldots} — dangling modifier (the fire was not rushing). Correction: \emph{Rushing to the scene, \textbf{the brigade brought} the fire under control.}
\item \emph{The amount of damages \textbf{have}} — agreement error; the subject is \emph{amount} (singular). Correction: \emph{The amount of damage \textbf{has} not yet been estimated.} (Also \emph{damages} $\rightarrow$ \emph{damage}: \emph{damages} means legal compensation.)
\item \emph{who \textbf{only} arrived at midnight} — misplaced modifier suggesting he did nothing but arrive. Correction: \emph{who arrived \textbf{only} at midnight.}
\item \emph{the council \textbf{don't} intend} — register/agreement error in formal reporting. Correction: \emph{the council \textbf{does} not intend.}
\item \emph{\textbf{Its} a great relief\ldots no lives \textbf{was} lost} — apostrophe error and agreement error. Correction: \emph{\textbf{It is} a great relief that no lives \textbf{were} lost.}
\end{enumerate}
Corrected passage: \emph{``Rushing to the scene, the brigade brought the fire under control. The amount of damage has not yet been estimated. The mayor, who arrived only at midnight, said the council does not intend to demolish the market. It is a great relief that no lives were lost.''}
\end{exoresolu}

\begin{exercice}[Practice — Modifiers, Idioms and Figures]
\begin{enumerate}
\item Correct: \emph{Having finished the exam, the bell rang.}
\item Correct: \emph{The teacher nearly marked fifty scripts in one evening.} (Say what the sentence wrongly claims, then fix it.)
\item Choose the word with the appropriate connotation for a formal tribute: the late chief was (stingy / frugal / miserly) and (stubborn / resolute / pig-headed).
\item Identify the figure of speech: (a) \emph{``Etoudi remained silent on the matter.''} (b) \emph{``I could sleep for a hundred years.''} (c) \emph{``New wheels were parked outside the mansion.''} (d) \emph{``The surgeon died of a simple infection he could have treated blindfolded.''}
\item Rewrite for a formal academic report: \emph{``Loads of farmers can't be bothered with the new fertiliser 'cause it's way too dear.''}
\end{enumerate}
\end{exercice}

\begin{corrige}[Exercise — Modifiers, Idioms and Figures]
\begin{enumerate}
\item Dangling modifier: \emph{Having finished the exam, \textbf{the candidates heard} the bell ring} (or \emph{\textbf{When the exam was finished}, the bell rang}).
\item The sentence \emph{The teacher nearly marked fifty scripts} places \emph{nearly} before the verb, claiming the teacher almost performed the action of marking fifty scripts but did not. If the intended meaning is ``approximately fifty'', write: \emph{The teacher marked \textbf{nearly} fifty scripts in one evening}. If the meaning is ``almost all fifty'', write: \emph{The teacher marked \textbf{almost all} fifty scripts}. The lesson: \emph{nearly} must sit immediately before the word it limits.
\item \emph{frugal} (positive) and \emph{resolute} (positive) — a tribute requires favourable connotations.
\item (a) Metonymy (the Presidency named by its associated place). (b) Hyperbole. (c) Synecdoche (wheels, a part, for cars, the whole). (d) Situational irony (the healer dying of what he could cure).
\item \emph{``A considerable number of farmers are reluctant to adopt the new fertiliser because its cost is prohibitively high.''} (Formal register: no contractions, no slang, precise quantifiers.)
\end{enumerate}
\end{corrige}

\begin{aretenir}[Key Points of Section 5]
\begin{itemize}
\item A modifier must stand next to the word it describes; if the doer is missing, the modifier dangles — supply the subject or expand the phrase into a clause.
\item Denotation is the dictionary meaning; connotation is the emotional colour. Choose words for both.
\item Idioms are fixed, figurative and register-sensitive: use them rarely, exactly and appropriately.
\item Irony, paradox, hyperbole, synecdoche and metonymy: name the figure \emph{and} explain its effect.
\item Edit any text with the five-point checklist: agreement, modifiers, tense consistency, register, punctuation.
\end{itemize}
\end{aretenir}

\coursec{Section 6: Phonetics and Spoken English}

In Section 5 we polished the written word: modifiers placed with precision, idioms used naturally, and vocabulary chosen for its exact nuance. But the GCE Advanced Level does not test only what you write. Paper 3 and the oral components of the examination demand that you \emph{hear} English accurately and \emph{produce} it intelligibly. A candidate who writes a flawless essay but pronounces \emph{ship} and \emph{sheep} identically, or who reads a question with the falling intonation of a statement, loses marks that grammar alone cannot save. This section therefore turns from the page to the ear and the mouth: phonemic transcription, connected speech, intonation, and the art of monitoring your own speech.

\courssub{Why Phonemes? From Spelling to Sound}

English spelling is famously treacherous. Consider the letters \textbf{ough}: \textbf{though}, \textbf{through}, \textbf{rough}, \textbf{cough}, \textbf{bough}. Five different pronunciations for one spelling. Conversely, the single vowel sound in \textbf{see}, \textbf{sea}, \textbf{scene}, \textbf{machine} and \textbf{people} is spelt five different ways. If we relied on spelling to discuss pronunciation, we would drown in confusion.

The solution is the \cle{International Phonetic Alphabet} (IPA): a system in which \textbf{one symbol always represents exactly one sound}, and one sound is always written with the same symbol. When we transcribe, we place symbols between slashes: \emph{cat} is transcribed /kæt/.

\begin{definition}[Phoneme]
A \cle{phoneme} is the smallest unit of sound that can distinguish one word from another in a language. Changing one phoneme changes the word: /pɪn/ (\emph{pin}) versus /bɪn/ (\emph{bin}) differ only in the first phoneme, /p/ versus /b/. Such pairs are called \cle{minimal pairs}.
\end{definition}

Standard British English (Received Pronunciation, the reference accent used in GCE phonetics questions) has \textbf{44 phonemes}: 20 vowels (12 monophthongs and 8 diphthongs) and 24 consonants.

\courssub{The Vowels: Monophthongs and Diphthongs}

\courssub{Monophthongs and the Vowel Quadrilateral}

A \cle{monophthong} is a pure vowel: the tongue holds one position throughout. Vowels are classified by three questions: How \textbf{high} is the tongue (close, mid, open)? How far \textbf{forward} is it (front, central, back)? Are the lips \textbf{rounded}?

\begin{popfigure}
\begin{tikzpicture}[scale=1.05]
\draw[thick, popink] (0,0) -- (4.2,0.6) -- (3.4,4.2) -- (0,4.2) -- cycle;
\node[popmuted, font=\small] at (-0.55,4.2) {close};
\node[popmuted, font=\small] at (-0.55,2.1) {mid};
\node[popmuted, font=\small] at (-0.55,0) {open};
\node[popmuted, font=\small] at (0,4.55) {front};
\node[popmuted, font=\small] at (3.6,4.55) {back};
% close vowels
\node[popblue] at (0.35,3.9) {/iː/};
\node[popblue] at (1.15,3.55) {/ɪ/};
\node[popblue] at (2.85,3.55) {/ʊ/};
\node[popblue] at (3.45,3.75) {/uː/};
% mid vowels
\node[poporange] at (0.6,2.35) {/e/};
\node[poporange] at (1.85,2.35) {/ə/};
\node[poporange] at (2.5,2.05) {/ɜː/};
\node[poporange] at (3.15,2.2) {/ɔː/};
% open vowels
\node[poppurple] at (0.85,0.75) {/æ/};
\node[poppurple] at (1.95,0.9) {/ʌ/};
\node[poppurple] at (2.85,0.7) {/ɑː/};
\node[poppurple] at (3.35,1.15) {/ɒ/};
\end{tikzpicture}
\legende{The IPA vowel quadrilateral: the twelve English monophthongs placed by tongue height and tongue position.}
\end{popfigure}

Learn them in contrastive groups, with a key word for each:

\begin{center}
\begin{tabularx}{\linewidth}{|l|X|l|X|}
\hline
\rowcolor{popblueL} \textbf{Symbol} & \textbf{Key word} & \textbf{Symbol} & \textbf{Key word} \\
\hline
/iː/ & \emph{sheep}, \emph{beat} & /ʊ/ & \emph{put}, \emph{good} \\
\hline
/ɪ/ & \emph{ship}, \emph{bit} & /uː/ & \emph{food}, \emph{two} \\
\hline
/e/ & \emph{bed}, \emph{said} & /ɔː/ & \emph{cord}, \emph{law} \\
\hline
/æ/ & \emph{cat}, \emph{bad} & /ɒ/ & \emph{pot}, \emph{hot} \\
\hline
/ɜː/ & \emph{bird}, \emph{work} & /ɑː/ & \emph{card}, \emph{half} \\
\hline
/ə/ & \emph{about}, \emph{teacher} & /ʌ/ & \emph{cup}, \emph{love} \\
\hline
\end{tabularx}
\end{center}

Note in particular /ə/, the \cle{schwa}: the most frequent vowel in English, always \textbf{unstressed}, as in the first syllable of \emph{about} /əˈbaʊt/ and the last of \emph{doctor} /ˈdɒktə/. In connected speech, many grammatical words (articles, prepositions, auxiliary verbs) reduce to a schwa; mastering these \cle{weak forms} is essential for natural rhythm.

\courssub{Diphthongs}

A \cle{diphthong} is a glide: the tongue moves from one vowel position to another within a single syllable. English has eight:

\begin{center}
\begin{tabularx}{\linewidth}{|l|X|l|X|}
\hline
\rowcolor{popblueL} \textbf{Symbol} & \textbf{Key word} & \textbf{Symbol} & \textbf{Key word} \\
\hline
/eɪ/ & \emph{day}, \emph{great} & /əʊ/ & \emph{go}, \emph{home} \\
\hline
/aɪ/ & \emph{my}, \emph{time} & /aʊ/ & \emph{now}, \emph{house} \\
\hline
/ɔɪ/ & \emph{boy}, \emph{coin} & /ɪə/ & \emph{here}, \emph{idea} \\
\hline
/eə/ & \emph{there}, \emph{care} & /ʊə/ & \emph{tour}, \emph{pure} \\
\hline
\end{tabularx}
\end{center}

\courssub{The Consonants}

Consonants are described by \textbf{voicing} (do the vocal cords vibrate?), \textbf{place of articulation} (where is the airflow obstructed?) and \textbf{manner of articulation} (how is it obstructed?).

\begin{center}
\begin{tabularx}{\linewidth}{|l|X|X|X|X|X|}
\hline
\rowcolor{popblueL} \textbf{Manner} & \textbf{Bilabial} & \textbf{Labiodental} & \textbf{Dental} & \textbf{Alveolar} & \textbf{Velar / Glottal} \\
\hline
Plosive & p b &  &  & t d & k g \\
\hline
Fricative &  & f v & θ ð & s z & ʃ ʒ ; h \\
\hline
Affricate &  &  &  & tʃ dʒ &  \\
\hline
Nasal & m &  &  & n & ŋ \\
\hline
Approximant & w &  &  & r l & j \\
\hline
\end{tabularx}
\end{center}

In each voiceless/voiced pair (e.g. /p/–/b/, /s/–/z/, /θ/–/ð/), the articulation is identical; only the vocal cords differ. Place your fingers on your throat and say /sss/ then /zzz/: you will feel the vibration begin.

\begin{attention}[The dental fricatives: /θ/ and /ð/]
The commonest transcription error among Cameroonian candidates is confusing /θ/ (voiceless, as in \emph{think} /θɪŋk/, \emph{bath} /bɑːθ/) and /ð/ (voiced, as in \emph{this} /ðɪs/, \emph{bathe} /beɪð/) with /t/, /d/, /s/ or /z/. In many Cameroonian languages these dental sounds do not exist, so speakers substitute them: \emph{three} pronounced like \emph{tree}, \emph{then} like \emph{den}. In transcription, \textbf{never} write /t/ for \emph{th}: decide whether the sound is voiceless (/θ/) or voiced (/ð/). A useful guide: function words (\emph{the, this, that, them, though}) take /ð/; most content words (\emph{think, thank, mouth, health}) take /θ/.
\end{attention}

Other frequent Cameroonian interference points include: merging /ɪ/ and /iː/ (\emph{ship}/\emph{sheep}), merging /ɒ/ and /ɔː/ (\emph{cot}/\emph{caught}), pronouncing /h/ weakly or dropping it (\emph{house} → \emph{ouse}), and adding a vowel after final consonant clusters (\emph{asked} pronounced ``ask-ed'' with two syllables instead of /ɑːskt/).

\begin{methode}[How to transcribe accurately]
\begin{enumerate}
\item \textbf{Say the word aloud} naturally, at normal speed — do not spell it in your head.
\item \textbf{Transcribe what you hear, not what you expect from spelling.} \emph{Women} is /ˈwɪmɪn/, not ``wo-men''; \emph{said} is /sed/, not /seɪd/.
\item \textbf{Count the sounds}, not the letters: \emph{church} has 6 letters but 3 phonemes: /tʃɜːtʃ/.
\item \textbf{Mark the primary stress} with /ˈ/ before the stressed syllable: \emph{important} /ɪmˈpɔːtnt/.
\item \textbf{Check each symbol} against the chart: is that vowel short or long? Is that \emph{th} voiced or voiceless?
\end{enumerate}
\end{methode}

\begin{exoresolu}[Transcribing words and an utterance]
Transcribe the following into phonemic script: (a) \emph{examination}; (b) \emph{thought}; (c) \emph{The teacher asked three questions.}
\tcblower
\textbf{Solution.}
\begin{itemize}
\item (a) /ɪɡˌzæmɪˈneɪʃn/ — five syllables, primary stress on the fourth; note /ʃ/ for \emph{-ti-} and the diphthong /eɪ/.
\item (b) /θɔːt/ — three phonemes only; the spelling \emph{-ough} gives the long vowel /ɔː/, and \emph{th} here is voiceless /θ/.
\item (c) /ðə ˈtiːtʃə ɑːskt θriː ˈkwestʃənz/ — note /ðə/ (voiced dental in \emph{the}), the schwa ending of \emph{teacher}, the cluster /skt/ in \emph{asked} (one syllable!), /θ/ in \emph{three}, and final voiced /z/ in \emph{questions}.
\end{itemize}
\end{exoresolu}

\courssub{Connected Speech: Elision, Assimilation, Intrusion, Catenation}

Read this sentence aloud at natural speed: \emph{``I must go and see my aunt and uncle.''} You almost certainly did not pronounce every sound of every word. Natural English is not a string of separate words; words flow into one another and sounds change, disappear, or appear. These processes are called \cle{connected speech} phenomena, and understanding them is essential both for listening comprehension and for the oral examination.

\begin{definition}[Elision]
\cle{Elision} is the omission of a sound (usually /t/ or /d/) in connected speech, typically when it is trapped between two consonants. Examples: \emph{next day} → /neks deɪ/ (the /t/ disappears); \emph{last night} → /lɑːs naɪt/; \emph{friendship} → /frenʃɪp/.
\end{definition}

\begin{definition}[Assimilation]
\cle{Assimilation} is a sound changing to become more like a neighbouring sound, so that the organs of speech move less. Examples: \emph{good boy} → /gʊb bɔɪ/ (/d/ becomes /b/ before bilabial /b/); \emph{ten minutes} → /tem ˈmɪnɪts/ (/n/ becomes /m/ before /m/); \emph{this year} → /ðɪʃ jɪə/ (/s/ becomes /ʃ/ before /j/).
\end{definition}

Two further processes complete the picture:

\begin{itemize}
\item \textbf{\cle{Intrusion}}: an extra sound is inserted between two vowels to smooth the glide. \emph{Go on} → /gəʊ w ɒn/ (intrusive /w/); \emph{I agree} → /aɪ j əˈgriː/ (intrusive /j/); \emph{law and order} → /lɔː r ənd ˈɔːdə/ (intrusive /r/, the famous ``linking/intrusive r'').
\item \textbf{\cle{Catenation}} (linking): a final consonant joins directly to an initial vowel of the next word, so the words sound fused. \emph{An apple} → /əˈnæpl/; \emph{turn off} → /tɜːˈnɒf/; \emph{get up} → /geˈtʌp/.
\end{itemize}

\begin{exemplebox}[One sentence, four processes]
\emph{``Would you ask him about it?''} said naturally: /wʊdʒu ɑːsk ɪm əˈbaʊt ɪt/. Here /d/ + /j/ assimilate to /dʒ/ (\emph{would you} → ``wudju''); the /h/ of \emph{him} is elided; and \emph{about it} is linked by catenation. A listener who expects every written sound will simply not find it in the stream of speech.
\end{exemplebox}

\courssub{Intonation: The Music of English}

Intonation is the rise and fall of the voice across an utterance. It carries meaning beyond the words themselves. English has three core nuclear tones:

\begin{itemize}
\item \textbf{Falling tone} (\searrow): the voice falls on the stressed syllable. Used for \textbf{statements}, \textbf{commands}, \textbf{Wh-questions} and \textbf{exclamations}: \emph{She lives in Bamenda.} \emph{Where did you go?}
\item \textbf{Rising tone} (\nearrow): the voice rises. Used for \textbf{yes/no questions}, \textbf{polite requests} and \textbf{listing items} (all but the last): \emph{Are you ready?} \emph{We bought yams \nearrow, rice \nearrow, plantains \nearrow\ and fish \searrow.}
\item \textbf{Fall-rise} (\searrow\nearrow): the voice falls then rises within the stressed syllable.
\end{itemize}

\begin{propriete}[The attitudinal fall-rise]
A fall-rise intonation typically signals \cle{doubt}, \cle{reservation}, \cle{politeness} or an implied ``but''. Compare \emph{It's \searrow fine} (genuine approval) with \emph{It's \searrow\nearrow fine...} (I have reservations). Examiners may ask you to identify the attitude conveyed by a given tone; the fall-rise is the classic marker of uncertainty or polite correction. (Identifying tones in context is examinable; memorise the tone--function correspondences.)
\end{propriete}

\begin{popfigure}
\begin{tikzpicture}[scale=1.0]
% Falling
\begin{scope}
\draw[->, popmuted] (0,0) -- (3.6,0) node[right, font=\small]{time};
\draw[->, popmuted] (0,0) -- (0,1.8) node[above, font=\small]{pitch};
\draw[very thick, popblue] (0.3,1.4) .. controls (1.2,1.45) and (1.8,1.3) .. (3.1,0.25);
\node[popink, font=\small\bfseries] at (1.8,-0.55) {Falling: ``She \emph{passed}.''};
\end{scope}
% Rising
\begin{scope}[xshift=5.2cm]
\draw[->, popmuted] (0,0) -- (3.6,0) node[right, font=\small]{time};
\draw[->, popmuted] (0,0) -- (0,1.8) node[above, font=\small]{pitch};
\draw[very thick, poporange] (0.3,0.35) .. controls (1.4,0.4) and (2.0,0.7) .. (3.1,1.5);
\node[popink, font=\small\bfseries] at (1.8,-0.55) {Rising: ``She \emph{passed}?''};
\end{scope}
% Fall-rise
\begin{scope}[xshift=10.4cm]
\draw[->, popmuted] (0,0) -- (3.6,0) node[right, font=\small]{time};
\draw[->, popmuted] (0,0) -- (0,1.8) node[above, font=\small]{pitch};
\draw[very thick, poppurple] (0.3,1.3) .. controls (1.2,0.15) and (1.9,0.15) .. (3.1,1.1);
\node[popink, font=\small\bfseries] at (1.8,-0.55) {Fall-rise: ``\emph{Maybe}...''};
\end{scope}
\end{tikzpicture}
\legende{The three core intonation contours. The same words, ``she passed'', become a statement, a surprised question, or an expression of doubt depending on the tone.}
\end{popfigure}

Notice how intonation performs a \textbf{grammatical} function (distinguishing statement from question, marking the end of a list with a fall) and an \textbf{attitudinal} function (a high rise can signal surprise: \emph{You paid \nearrow how much?}; a warm fall-rise can soften advice: \emph{You might \searrow\nearrow try again}).

\courssub{Self-Monitoring in Speech}

Even advanced speakers make slips. What distinguishes an excellent oral candidate is the ability to \cle{monitor} their own speech and repair it gracefully, exactly as a good driver constantly corrects the steering.

\begin{methode}[Correction and modification strategies]
\begin{enumerate}
\item \textbf{Rephrasing}: when a word escapes you, say it another way instead of freezing. ``We went to the --- the place where they treat sick people --- the hospital.''
\item \textbf{Self-repair}: correct yourself openly and move on. ``He go --- he \emph{went} --- to Douala last week.'' A quick, calm repair impresses examiners; pretending the error never happened does not.
\item \textbf{Pacing}: slow down at difficult words and at key points. Speaking slightly slower than feels natural almost always sounds \emph{more} fluent, because it eliminates the false starts caused by rushing.
\item \textbf{Pausing}: use short silent pauses at clause boundaries to plan the next idea, rather than filling every gap with ``ehhh''.
\item \textbf{Checking understanding}: in dialogue, confirm with rising intonation: ``You mean the second chapter?'' This is both a listening strategy and a repair strategy.
\end{enumerate}
\end{methode}

\begin{attention}[Over-correction and hypercorrection]
Do not ``repair'' what is not broken. Some candidates, aware that /h/-dropping is stigmatised, begin adding /h/ where it does not belong (``h-is'', ``h-apple''). Monitor yourself against the phonemic chart, not against anxiety.
\end{attention}

\begin{exoresolu}[Identifying connected speech processes]
In natural rapid speech, the sentence \emph{``I can't go and see her this afternoon''} is realised as /aɪ kɑːŋ gəʊ ən siː ə ðɪs ɑːftəˈnuːn/. Identify the connected speech processes involved.
\tcblower
\textbf{Solution.}
\begin{itemize}
\item \emph{can't go}: \textbf{assimilation} — final /t/ of \emph{can't} is elided and /n/ assimilates to /ŋ/ before the velar /g/ (written here as /kɑːŋ/).
\item \emph{go and}: \textbf{intrusion} of /w/ between /əʊ/ and /æ/ (``go-w-and''), plus reduction of \emph{and} to /ən/.
\item \emph{see her}: \textbf{elision} of /h/ (\emph{her} → /ə/), a standard weak form.
\item \emph{and see}: \textbf{catenation} — /n/ links directly to /siː/.
\end{itemize}
\end{exoresolu}

\begin{exercice}[Practice: transcription, tone and connected speech]
\begin{enumerate}
\item Transcribe: \emph{thought}, \emph{measure}, \emph{question}, \emph{comfortable}, \emph{She bought a new umbrella.}
\item State whether the \emph{th} in each word is /θ/ or /ð/: \emph{either}, \emph{method}, \emph{breathe}, \emph{breath}, \emph{clothing}.
\item What tone (fall, rise, fall-rise) would you expect on: (a) \emph{Where is the market?} (b) \emph{Did you lock the door?} (c) \emph{Well, I \emph{suppose} so...}?
\item Rewrite in natural connected speech, naming each process: \emph{last Saturday}, \emph{good morning}, \emph{I saw it}, \emph{stand up}.
\end{enumerate}
\end{exercice}

\begin{corrige}[Exercise: Phonetics and Spoken English]
\begin{enumerate}
\item /θɔːt/; /ˈmeʒə/; /ˈkwestʃən/; /ˈkʌmftəbl/; /ʃi bɔːt ə njuː ʌmˈbrelə/.
\item \emph{either} /ð/; \emph{method} /θ/; \emph{breathe} /ð/; \emph{breath} /θ/; \emph{clothing} /ð/.
\item (a) Falling (Wh-question); (b) Rising (yes/no question); (c) Fall-rise (doubt, reservation).
\item \emph{last Saturday} /lɑːs ˈsætədeɪ/ — elision of /t/; \emph{good morning} /gʊb ˈmɔːnɪŋ/ — assimilation of /d/ to /b/; \emph{I saw it} /aɪ sɔː r ɪt/ — intrusive /r/; \emph{stand up} /stænˈdʌp/ — catenation (and possible elision of /d/: /stæn ˈʌp/).
\end{enumerate}
\end{corrige}

\begin{aretenir}[Key points of Section 6]
\begin{itemize}
\item English has 44 phonemes: 12 monophthongs, 8 diphthongs, 24 consonants; transcribe sounds, never spellings.
\item Master the traps: /θ/ vs /ð/, /ɪ/ vs /iː/, the schwa /ə/, and stress marking /ˈ/.
\item Connected speech: \textbf{elision} (sound dropped), \textbf{assimilation} (sound changes), \textbf{intrusion} (sound added), \textbf{catenation} (words linked).
\item Falling tone = statements, commands, Wh-questions; rising tone = yes/no questions and open lists; fall-rise = doubt, reservation, politeness.
\item In oral work, monitor yourself: rephrase, repair calmly, control your pace, and pause to plan.
\end{itemize}
\end{aretenir}

With the ear trained and the voice under control, we are ready to return to writing at the highest level: Section 7 tackles the argumentative and synthesis essays that dominate Paper 2.

\coursec{Section 7: Argumentative and Synthesis Essay Writing}

In Section 6 you trained your ear and your voice: elision, assimilation, intonation, self-monitoring in speech. You learned that spoken English persuades through rhythm, stress and tone of voice. We now move from the spoken word to the written argument. In the examination room, no intonation can rescue you; only the \emph{structure} of your reasoning, the \emph{quality} of your evidence and the \emph{clarity} of your language will convince the examiner. Paper 2 of the GCE Advanced Level regularly demands an argumentative or discursive essay of 450--600 words, and the synthesis of documents is increasingly tested. This section gives you the complete architecture.

\courssub{The Balanced Argumentative Essay}

\begin{definition}[Argumentative essay]
An \cle{argumentative essay} is a piece of formal writing in which the writer examines a debatable proposition, presents the strongest arguments on \emph{both} sides, weighs them, and arrives at a clear, defended personal position. It is not a list of opinions; it is a reasoned demonstration.
\end{definition}

Consider a concrete situation. Your school debating club in Bamenda is given the motion: \emph{``Boarding schools do more harm than good.''} A weak speaker shouts only the points favouring his side and ignores the rest. A strong speaker says: ``My opponents rightly observe that boarding schools build independence; however, they underestimate the psychological cost of separating a twelve-year-old from the family.'' The second speaker wins because he shows he has \emph{heard} the other side before defeating it. That is exactly what the GCE examiner rewards.

\begin{methode}[The balanced essay plan]
Follow these steps every time you meet an argumentative topic:
\begin{enumerate}
\item \textbf{Analyse the question.} Underline the key terms and identify the two possible positions. Decide your \emph{final} stance before writing a single line.
\item \textbf{Brainstorm both sides.} List at least three arguments \emph{for} and three \emph{against}. Select the two strongest on each side.
\item \textbf{Write the introduction.} Set the context in one or two sentences, then state your \cle{thesis} --- the precise claim your essay will defend.
\item \textbf{Develop two `for' paragraphs} (or the side you disagree with first, so you can dismantle it).
\item \textbf{Develop two `against' paragraphs} (your own side, presented more strongly and in more depth).
\item \textbf{Write a weighed conclusion.} Do not merely repeat; \emph{weigh} the two sides and commit to a clear final position, ideally with a memorable closing sentence.
\end{enumerate}
\end{methode}

\begin{attention}[One-sidedness and fence-sitting]
The two fatal errors are: (1) presenting \emph{only one side} of the question, which reduces your essay to propaganda and costs you content marks; and (2) \emph{sitting on the fence} --- concluding with ``both sides have good points, so it depends'' without committing yourself. The examiner wants a mind at work: acknowledge the opposition, then \textbf{defeat it} and state where you stand.
\end{attention}

\begin{popfigure}
\begin{tikzpicture}[
box/.style={draw=popblue, fill=popblueL, rounded corners=2pt, text width=4.6cm, align=center, minimum height=1.05cm, font=\small},
opp/.style={draw=poporange, fill=white, rounded corners=2pt, text width=4.6cm, align=center, minimum height=1.05cm, font=\small},
own/.style={draw=popgreen, fill=popgreenL, rounded corners=2pt, text width=4.6cm, align=center, minimum height=1.05cm, font=\small},
arr/.style={->, thick, popdark}]
\node[box, align=center] (intro) at (0,0){\textbf{Introduction}\\ context + clear thesis};
\node[opp, align=center] (a1) at (0,-1.7){\textbf{Paragraph 2}\\ opposing argument 1 + concession};
\node[opp, align=center] (a2) at (0,-3.0){\textbf{Paragraph 3}\\ opposing argument 2 + rebuttal};
\node[own, align=center] (b1) at (0,-4.3){\textbf{Paragraph 4}\\ own argument 1 + evidence};
\node[own, align=center] (b2) at (0,-5.6){\textbf{Paragraph 5}\\ own argument 2 + evidence};
\node[box, align=center] (conc) at (0,-7.3){\textbf{Conclusion}\\ weigh both sides, firm stance};
\draw[arr] (intro) -- (a1);
\draw[arr] (a1) -- (a2);
\draw[arr] (a2) -- (b1);
\draw[arr] (b1) -- (b2);
\draw[arr] (b2) -- (conc);
\node[font=\small\itshape, text=popmuted, align=left, text width=4.2cm] at (5.6,-3.0) {Opposition first,\\ then your stronger\\ case: the essay\\ rises towards\\ your position.};
\end{tikzpicture}
\legende{Architecture of the balanced argumentative essay: concede and rebut the opposing view, then build your own case to a weighed conclusion.}
\end{popfigure}

\courssub{Developing Paragraphs: the TEEL Structure}

A paragraph is not a random bundle of sentences; it is a miniature argument. The reliable model is \cle{TEEL}:

\begin{itemize}
\item \textbf{T --- Topic sentence:} one sentence stating the single idea of the paragraph.
\item \textbf{E --- Evidence:} a fact, example, statistic (only if certain), or reference that supports the idea.
\item \textbf{E --- Explanation:} show \emph{how} the evidence proves the topic sentence; this is where marks are earned.
\item \textbf{L --- Link:} a closing sentence that links back to the thesis or forward to the next paragraph.
\end{itemize}

\begin{exemplebox}[A TEEL paragraph in action]
\textbf{(T)} Mobile money has transformed rural commerce in Cameroon. \textbf{(E)} A farmer in Bafoussam can now receive payment for her cassava by phone before the lorry even leaves the farm. \textbf{(E)} Because payment is instant and traceable, she no longer depends on middlemen who once delayed or reduced her earnings; the risk of carrying cash on unsafe roads also disappears. \textbf{(L)} This single technology therefore illustrates how digital tools can redistribute economic power to the countryside --- a point developed further below.
\end{exemplebox}

\begin{attention}[Evidence without explanation]
Stating a fact and moving on is the most common weakness in Upper Sixth scripts. The examiner asks: \emph{so what?} Always add the explanation sentence that interprets the evidence for the reader.
\end{attention}

\courssub{Discourse Markers: Staging, Sequencing and Concluding}

Discourse markers are the signposts that tell the reader where your argument is going. Used accurately, they create \emph{cohesion}; used mechanically, they irritate the examiner. Learn them by \emph{function}, not as a random list.

\begin{center}
\begin{tabularx}{\linewidth}{|l|X|X|}
\hline
\rowcolor{popblueL} \textbf{Function} & \textbf{Markers} & \textbf{Example in use} \\
\hline
Sequencing & firstly, subsequently, finally & \emph{Firstly, the cost of data excludes the poorest households.} \\
\hline
Adding & moreover, furthermore, in addition & \emph{Moreover, the scheme created two hundred jobs.} \\
\hline
Contrasting & however, nevertheless, whereas & \emph{The policy is popular; however, it remains underfunded.} \\
\hline
Conceding & admittedly, granted, it is true that & \emph{Admittedly, boarding schools foster independence.} \\
\hline
Illustrating & for instance, such as, namely & \emph{Many crops, for instance plantain, perish quickly.} \\
\hline
Cause / result & consequently, therefore, as a result & \emph{Roads were cut off; consequently, prices doubled.} \\
\hline
Concluding & in the final analysis, on balance, to conclude & \emph{In the final analysis, the benefits outweigh the risks.} \\
\hline
\end{tabularx}
\end{center}

\begin{attention}[Marker misuse]
Do not open every paragraph with ``Firstly'' and never write ``In a nutshell'' in a formal essay. Note also the comma: \emph{However,} with a comma after it when it begins a sentence; \emph{however} without a comma means ``in whatever way''.
\end{attention}

\courssub{Synthesis Writing: Integrating Multiple Perspectives}

\begin{definition}[Synthesis essay]
A \cle{synthesis essay} is an essay that builds a single, coherent argument out of \emph{several sources} --- articles, reports, speeches, statistics --- rather than out of personal opinion alone. The writer compares the sources, identifies agreements and tensions, and weaves them into his or her own line of argument with accurate attribution.
\end{definition}

The key skill is \cle{attribution}: naming where each idea comes from using precise citation language:

\begin{itemize}
\item \emph{According to the Ministry's report, \ldots}
\item \emph{As Nganang argues, \ldots} / \emph{The editorial contends that \ldots}
\item \emph{Source B challenges this view, claiming that \ldots}
\item \emph{Both the journalist and the researcher agree on one point: \ldots}
\end{itemize}

\begin{methode}[Writing a synthesis paragraph]
\begin{enumerate}
\item Read all sources and annotate each with its \emph{claim} and its \emph{evidence}.
\item Group the sources: which agree? Which clash? Which complement each other?
\item Formulate \emph{your} thesis that takes all sources into account.
\item In each paragraph, combine at least two sources around one idea, using attribution verbs (\emph{argues, claims, concedes, demonstrates, dismisses}).
\item End by drawing the threads together --- the synthesis is \emph{yours}, not a summary of each source in turn.
\end{enumerate}
\end{methode}

\courssub{Avoiding Logical Fallacies and Overgeneralisation}

An argument collapses when its reasoning is faulty. Learn to spot --- and avoid --- these fallacies in your own writing:

\begin{itemize}
\item \textbf{Hasty generalisation:} \emph{``All politicians are corrupt.''} One example does not prove a universal claim. Prefer: \emph{``Some officials have abused public funds.''}
\item \textbf{False cause (post hoc):} \emph{``Since the new headmaster arrived, results have improved; he must be the cause.''} Correlation is not causation.
\item \textbf{Ad hominem:} attacking the person instead of the argument: \emph{``His view on fees is worthless because he is rich.''}
\item \textbf{Straw man:} misrepresenting the opposing view to defeat it easily: \emph{``Those who oppose the dam want Cameroonians to live in darkness.''}
\item \textbf{Sweeping absolute:} words like \emph{always, never, everyone, no one} are almost always false and signal weak thinking. Qualify: \emph{most, often, in many cases}.
\end{itemize}

\begin{aretenir}[Exam tip]
Examiners reward a \textbf{clear, consistent line of argument} far more than long or exotic vocabulary. A simple sentence that advances the argument beats a decorated sentence that says nothing. Choose precision over display: \emph{every paragraph must earn its place}.
\end{aretenir}

\courssub{Worked Example: A Balanced Essay Under GCE Conditions}

\begin{exoresolu}[Model argumentative essay (abridged)]
\textbf{Topic:} \emph{``Social media does more harm than good to the Cameroonian youth.''} Discuss. (40 minutes, 450--600 words.)
\tcblower
\textbf{Introduction.} In Douala and in the remotest villages of the North West, young Cameroonians now live as much online as offline. While critics denounce social media as a poison, I shall argue that, used with discipline, it does more good than harm.

\textbf{Opposing view 1 (concession).} Admittedly, social media consumes study time. Many candidates scroll through videos late into the night and arrive at school exhausted. However, the fault lies not in the tool but in its use: a machete can clear a farm or wound a man.

\textbf{Opposing view 2 (rebuttal).} Moreover, critics point to misinformation and online fraud. Yet these dangers existed before the internet --- rumours and tricksters are as old as markets. What social media changes is speed, and the remedy is education in critical reading, not prohibition.

\textbf{Own argument 1.} Firstly, these platforms democratise knowledge. A student in Nkambe can follow free tutorials, download past GCE papers and join study groups that no village library could offer. Consequently, the gap between urban and rural learners narrows.

\textbf{Own argument 2.} Furthermore, social media creates livelihoods. Young entrepreneurs in Yaound\'e market shoes, cakes and tutoring services online at almost no cost, turning a phone into a shop. In an economy where formal jobs are scarce, this is not trivial.

\textbf{Conclusion.} On balance, the harms of social media are real but avoidable, while its benefits --- education, enterprise, connection --- are structural. In the final analysis, the answer is not to disconnect the youth but to teach them to navigate wisely.
\end{exoresolu}

\courssub{Timed Practice and Assessment}

Under GCE conditions, allocate your 40--45 minutes as follows: 5 minutes planning, 30 minutes writing, 5--10 minutes proofreading (checking verb agreement, article use, and that every paragraph has a TEEL spine). After writing, assess yourself --- and a partner --- with the examiner's eye: Is there a clear thesis? Are both sides treated? Is each paragraph developed with evidence and explanation? Is the conclusion a genuine verdict?

\begin{exercice}[Practice essay]
Write a balanced argumentative essay (450--600 words) on: \emph{``Examinations are the fairest way to select students for university.''} Your essay must contain: a clear thesis, two concession-and-rebuttal paragraphs, two paragraphs for your own side using TEEL, at least five discourse markers from the table, and a weighed conclusion. Then exchange scripts with a partner and award marks out of 20 using the criteria: argument (8), organisation (6), language accuracy (6).
\end{exercice}

\begin{corrige}[Practice essay]
A strong script will: (1) state a definite thesis, e.g. \emph{``Examinations, despite their imperfections, remain the fairest selection tool available''}; (2) concede opposing points --- exam anxiety, coaching advantages for the wealthy --- and rebut them (\emph{``yet interviews and quotas are even easier to manipulate''}); (3) defend fairness through anonymity of scripts and uniformity of standards, with Cameroonian illustrations such as the national marking of the GCE; (4) use markers accurately (\emph{admittedly, however, moreover, consequently, in the final analysis}); (5) end with a firm, weighed verdict, not a summary. Penalise: one-sidedness (max 4/8 for argument), unexplained evidence, fallacies such as \emph{``everyone knows exams are stressful''}, and fence-sitting conclusions. Reward precise, modest language over inflated vocabulary.
\end{corrige}

\coursec{Section 8: Creative and Dramatic Writing}

In Section 7 you learned to argue: to build a balanced, well-evidenced case and to synthesise several voices into one coherent essay. In this section we do something that looks very different but draws on exactly the same discipline of structure and economy: we write for the \cle{stage}. In an argumentative essay, the writer tells the reader what to think through reasoned prose. In a play, the writer must \emph{never} tell; characters must \emph{show} everything through what they say and what they do. The GCE Advanced Level Paper 2 (Composition) regularly offers a creative writing option, and candidates who can produce a tight, well-formatted short play or dramatic scene often distinguish themselves — but only if the script has genuine plot, conflict and theme, not merely pleasant conversation. This section teaches you the craft: natural dialogue, economical stage directions, one-act structure, adaptation of themes from your prescribed texts, and the peer-performance techniques that professional dramatists use to test their work.

\courssub{From Essay to Stage: What Changes?}

Consider how you would express the theme of \emph{tradition versus modernity} — the theme you met in \emph{The Swamp Dwellers} — in an essay. You might write: ``Young Cameroonians increasingly reject the customs of their parents, creating painful generational conflict.'' That is a thesis statement: clear, direct, abstract. A playwright is forbidden such shortcuts. The same idea must emerge from a scene: perhaps a daughter returns from university in Yaoundé and refuses to pour libation at her grandfather's shrine, and the two argue. Nothing is stated; everything is \emph{dramatised}.

\begin{definition}[Drama and the Play Script]
A \cle{play script} is a text written to be \emph{performed}, not merely read. It consists of two elements only: \cle{dialogue} (the words characters speak) and \cle{stage directions} (instructions about movement, tone, setting and delivery, which are not spoken aloud). Everything the audience knows must reach them through these two channels.
\end{definition}

This constraint is precisely what makes playwriting excellent training for the GCE candidate: it forces economy, precision and structural control — the same virtues the examiners reward in summary and essay writing.

\courssub{Stage Directions: Conventions and Economy}

\begin{definition}[Stage Direction]
A \cle{stage direction} is an instruction written by the playwright, conventionally placed in brackets and printed in italics, telling actors and directors about \emph{setting}, \emph{movement} (entrances, exits, gestures), \emph{tone of voice} and \emph{stage business}. Stage directions are \textbf{never spoken aloud} by the actors.
\end{definition}

Observe the conventions carefully, because examiners penalise scripts that ignore them:

\begin{itemize}
\item \textbf{Character names} are written in capitals, followed by a colon or placed centred above the speech; choose one style and remain consistent.
\item \textbf{Stage directions} appear in \emph{italics} within brackets: \emph{(rising angrily)}, \emph{(she crosses to the window)}.
\item Directions are written in the \textbf{present tense}: \emph{(He enters)}, not \emph{(He entered)}.
\item Directions are \textbf{economical}: they give only what cannot be conveyed through dialogue. A direction such as \emph{(He walks slowly and sadly across the room thinking about his dead mother)} is bad — the grief must come through the words.
\item Standard stage positions use fixed vocabulary: \emph{upstage} (away from the audience), \emph{downstage} (towards the audience), \emph{stage left} and \emph{stage right} (from the actor's point of view).
\end{itemize}

\begin{attention}[Common Mistake: Over-Directing]
Weak candidates write a stage direction before every single speech, turning the script into a puppet show. Direct only what matters: entrances, exits, significant gestures, and moments where tone contradicts or sharpens the words (e.g. \emph{(smiling, but coldly)}). Trust your dialogue to carry the rest.
\end{attention}

\courssub{Dialogue That Reveals Character and Advances Conflict}

Read these two versions of the same moment aloud:

\begin{exemplebox}[Essay Dialogue vs. Dramatic Dialogue]
\textbf{Version A (essay-like):}\par
\emph{NGOH:} Father, I must inform you that I have decided not to participate in the annual libation ceremony because my university education has led me to question such traditional practices.\par
\emph{PA TABI:} I am extremely disappointed and angry at your decision, my daughter.\par\medskip
\textbf{Version B (dramatic):}\par
\emph{NGOH:} Papa, about tomorrow's libation\ldots{} I'm not going.\par
\emph{PA TABI:} \emph{(slowly putting down his cup)} Say that again.\par
\emph{NGOH:} You heard me, Papa.\par
\emph{PA TABI:} So. Yaoundé has finally swallowed my daughter.
\end{exemplebox}

Version A is grammatically perfect and dramatically dead: the characters announce their feelings and motives like debaters. Version B works because real people speak in \emph{fragments}, \emph{interruptions}, \emph{short questions} and \emph{loaded silences}; and because the conflict \emph{escalates} with each line. Pa Tabi's final metaphor — Yaoundé ``swallowing'' his daughter — reveals his world-view far more powerfully than any statement of anger.

\begin{methode}[Show Character Through Speech and Action, Not Narration]
\begin{enumerate}
\item \textbf{Decide what each character wants} in the scene (their \cle{objective}). Conflict arises when two objectives collide.
\item \textbf{Give each character a distinct voice}: vocabulary, rhythm, favourite expressions. An elder from the village and a university student must not sound identical.
\item \textbf{Cut all self-explanation.} Delete any line in which a character states what they feel (``I am angry''). Replace it with words or actions that \emph{demonstrate} the feeling.
\item \textbf{Make every speech do work}: it must either reveal character, advance the conflict, or — ideally — both. If a line does neither, delete it.
\item \textbf{Read the scene aloud.} If a sentence cannot be said naturally in one breath, rewrite it.
\end{enumerate}
\end{methode}

\begin{attention}[Common Mistake: Dialogue That Sounds Like an Essay]
The single most frequent fault in GCE play scripts is dialogue written in formal, complete, polished sentences. Real speech is broken, elliptical and repetitive. Compare: ``I strongly disagree with your position on this matter'' (essay) with ``No. No, Papa — you can't mean that'' (stage). When in doubt, speak the line aloud before writing it down.
\end{attention}

\courssub{The Structure of a One-Act Play}

A short play — the only kind feasible in an examination — still obeys the classical arc of dramatic structure. Because time is limited, the arc is compressed: the exposition is brief, the rising action steep, and the climax arrives quickly.

\begin{definition}[The One-Act Dramatic Arc]
\begin{itemize}
\item \cle{Exposition}: the opening moments establish setting, characters and the initial situation — economically, through dialogue, not narration.
\item \cle{Inciting incident}: the event that disturbs the balance and launches the conflict (Ngoh's announcement, in our example).
\item \cle{Rising action}: a series of exchanges in which the conflict intensifies; each speech raises the stakes.
\item \cle{Climax}: the point of maximum tension, where a decision or revelation becomes irreversible.
\item \cle{Resolution} (or \emph{dénouement}): the new balance after the climax. In a one-act play it may be brief — a single line, a gesture, an exit — but it must exist.
\end{itemize}
\end{definition}

\begin{popfigure}
\begin{tikzpicture}[scale=1.0]
% axes
\draw[->, thick, popline] (0,0) -- (10.6,0) node[below left, popmuted, font=\small] {stage time};
\draw[->, thick, popline] (0,0) -- (0,4.6) node[above, popmuted, font=\small, align=center] {tension};
% the arc
\draw[very thick, popblue] (0.4,0.5)
  .. controls (1.6,0.7) and (2.6,1.1) .. (3.4,1.5)
  .. controls (4.6,2.1) and (5.6,2.9) .. (6.8,4.0)
  .. controls (7.3,4.3) and (7.7,3.4) .. (8.4,1.6)
  .. controls (8.8,0.9) and (9.4,0.7) .. (10.0,0.7);
% points
\fill[popdark] (0.4,0.5) circle (2.2pt);
\fill[poporange] (3.4,1.5) circle (2.2pt);
\fill[poppink] (6.8,4.0) circle (2.6pt);
\fill[popgreen] (10.0,0.7) circle (2.2pt);
% labels
\node[below, popdark, font=\small, align=center] at (0.9,0.35) {Exposition};
\node[below, poporange, font=\small, align=center] at (3.4,1.1) {Inciting\\incident};
\node[above, popblue, font=\small, align=center] at (5.0,2.6) {Rising action};
\node[above, poppink, font=\small, align=center] at (6.9,4.25) {Climax};
\node[right, popgreen, font=\small, align=center] at (9.6,1.15) {Resolution};
% dashed guide
\draw[dashed, popmuted] (6.8,4.0) -- (6.8,0);
\end{tikzpicture}
\legende{The dramatic arc of a one-act play: compressed exposition, steep rising action, a single sharp climax, and a brief resolution.}
\end{popfigure}

\begin{attention}[Exam Tip]
In the composition paper, a play script is still a \emph{composition}: it must have a clear plot, a discernible theme, and a satisfying shape. A script that is ``just conversation'' — however lively — will be marked down for lack of structure. Before writing, sketch your arc in four lines: situation, disturbance, climax, outcome.
\end{attention}

\courssub{Formatting the Script Page}

Examiners reward correct presentation. Study the model page below, then imitate it exactly.

\begin{popfigure}
\begin{tikzpicture}
% page
\fill[white] (0,0) rectangle (12.6,9.2);
\draw[popline, thick] (0,0) rectangle (12.6,9.2);
% title
\node[font=\bfseries] at (6.3,8.6) {THE LIBATION};
\node[font=\itshape\small, popmuted] at (6.3,8.2) {A one-act play};
% setting direction
\node[anchor=west, font=\itshape\small, text width=11.4cm] at (0.5,7.5)
{(A village parlour near Bamenda. Early morning. A wooden shrine stands in one corner. PA TABI, about sixty, sits drinking tea. NGOH, twenty-two, enters wearing jeans and carrying a suitcase.)};
% dialogue
\node[anchor=west, font=\small\bfseries] at (0.5,6.3) {NGOH:};
\node[anchor=west, font=\small, text width=11.2cm] at (1.9,6.3) {Papa, I'm home. \emph{(He does not look up. She sets down the suitcase.)} Papa?};
\node[anchor=west, font=\small\bfseries] at (0.5,5.3) {PA TABI:};
\node[anchor=west, font=\small, text width=11.2cm] at (2.3,5.3) {\emph{(slowly)} So the city remembers the way back.};
\node[anchor=west, font=\small\bfseries] at (0.5,4.5) {NGOH:};
\node[anchor=west, font=\small, text width=11.2cm] at (1.9,4.5) {About tomorrow's libation\ldots{} I'm not going.};
\node[anchor=west, font=\small\bfseries] at (0.5,3.7) {PA TABI:};
\node[anchor=west, font=\small, text width=11.2cm] at (2.3,3.7) {\emph{(putting down his cup)} Say that again.};
% annotation labels
\node[popblue, font=\scriptsize\bfseries, anchor=west] at (0.5,2.9) {KEY:};
\node[popblue, font=\scriptsize, anchor=west, text width=11.6cm] at (0.5,2.4)
{CAPITALS + colon = speaker's name \quad $\bullet$ \quad \emph{(italics in brackets)} = stage direction, never spoken \quad $\bullet$ \quad present tense in directions \quad $\bullet$ \quad opening direction fixes time, place, characters};
\node[popmuted, font=\scriptsize, anchor=west, text width=11.6cm] at (0.5,1.5)
{Note the economy: the silence after ``Papa, I'm home'' is a direction of three words, yet it tells us the conflict has already begun.};
\end{tikzpicture}
\legende{Annotated model script page: title, opening stage direction (setting and characters), and dialogue with bracketed, italicised directions.}
\end{popfigure}

\courssub{Adapting Themes from Studied Texts}

The prescribed texts you study are a goldmine of dramatic material. \emph{The Swamp Dwellers} dramatises tradition against modernity; \emph{Unanswered Cries} stages the clash between custom and a young girl's rights; \emph{Twilight of Misty Foliage} explores community, change and identity. Adapting such a theme into an \emph{original} scene is excellent examination preparation — but note the word \emph{original}: you must not simply copy Soyinka's characters.

\begin{methode}[From Studied Theme to Original Scene]
\begin{enumerate}
\item \textbf{Extract the abstract conflict} (e.g. tradition vs. modernity; individual desire vs. communal duty).
\item \textbf{Change the concrete situation}: new characters, new setting, new stakes. Soyinka's swamp becomes, say, a family compound in Buea deciding whether to sell ancestral land to a developer.
\item \textbf{Assign opposing objectives} to two or three characters.
\item \textbf{Build the arc}: brief exposition, inciting announcement, escalating exchanges, one irreversible climax, a short resolution.
\item \textbf{Write, then read aloud and revise.}
\end{enumerate}
\end{methode}

\begin{exoresolu}[Planning and Drafting a Short Dramatic Scene]
\textbf{Task:} Write the opening of a one-act play (about one page) dramatising the conflict between tradition and modernity in a contemporary Cameroonian family.\tcblower
\textbf{Step 1 — Plan the arc.} Situation: the Ekane family in Buea; Grandmother Malle wants the family land kept for the annual harvest rite. Disturbance: her grandson Ako, a young engineer, announces he has signed papers to lease the land for a telecom mast. Climax: Grandmother picks up the ancestral staff and declares she will stand on the land when the machines come. Resolution: Ako kneels — not in surrender, but to ask her to bless a compromise.\par
\textbf{Step 2 — Draft (extract).}\par
\emph{(A compound in Buea. Evening. MALLE, seventy, tends a small fire. AKO, twenty-six, enters with a folder.)}\par
\textbf{AKO:} Ma, it's done. The company signed today.\par
\textbf{MALLE:} \emph{(not turning)} Signed what, my son?\par
\textbf{AKO:} The lease. The east plot. Ma, that money will pay Nadia's fees for three years.\par
\textbf{MALLE:} \emph{(turning slowly)} The east plot. Where your grandfather planted the first kola.\par
\textbf{AKO:} It's just land, Ma.\par
\textbf{MALLE:} \emph{(lifting the staff from beside the door)} Just land. \emph{(beat)} Then you will not mind when I sleep on ``just land'' the night their machines come.\par
\textbf{Step 3 — Commentary.} Notice: no character says ``I am angry'' or ``tradition matters to me''; the staff, the fire and the kola tree carry the theme. Each speech either reveals character or raises the stakes, and the fragment ``Just land'' echoed back becomes the pivot of the scene.
\end{exoresolu}

\courssub{Peer Performance and Feedback: Testing Dialogue Aloud}

Professional dramatists never trust the silent page; they hear the play. You should do the same. Reading a script aloud with classmates is the fastest way to detect false dialogue, flat passages and structural sag.

\begin{methode}[Running a Peer Reading]
\begin{enumerate}
\item \textbf{Cast the scene}: assign one reader per character; the writer stays silent and takes notes.
\item \textbf{Read at performance pace}, observing every stage direction.
\item \textbf{Mark the stumbles}: wherever a reader trips over a sentence, the line is probably too long or too formal.
\item \textbf{Ask three questions} afterwards: Where did you lose interest? Which line sounded least like real speech? What does each character want — could you tell?
\item \textbf{Revise} the two weakest moments, then read again.
\end{enumerate}
\end{methode}

\begin{savaistu}[Drama in Cameroon]
Cameroon has a rich living theatre culture: playwrights such as Bole Butake and Guillaume Oyono-Mbia built their reputations on exactly the skills in this section — sharp dialogue, generational conflict, and themes drawn from village and city life. Oyono-Mbia's comedies of tradition clashing with modernity are close cousins of the scene you drafted above. When you write a short play, you join a national tradition.
\end{savaistu}

\begin{exercice}[Writing Your Own One-Act Opening]
Choose one theme from a prescribed text (e.g. tradition vs. modernity, individual rights vs. custom, man vs. nature). Write the opening page of an original one-act play (10--15 speeches), including: a title; an opening stage direction fixing time, place and characters; at least three meaningful stage directions; and dialogue that escalates a conflict. Then read it aloud with a partner and revise the two weakest lines.
\end{exercice}

\begin{corrige}[Exercise: Writing Your Own One-Act Opening]
A strong answer will show the following features (modelled on the theme \emph{man vs. nature}, adapted from \emph{The Old Man and the Sea}):\par
\textbf{Title:} THE LAST CATCH.\par
\textbf{Opening direction:} \emph{(A fishing beach at Limbe, dawn. Old MOLUA, sixty-five, mends a net. His nephew EKWE, thirty, enters with a newspaper.)}\par
\textbf{Sample dialogue:}\par
\textbf{EKWE:} Uncle, it's in the paper. The trawlers have been licensed for our waters.\par
\textbf{MOLUA:} \emph{(still mending)} Paper. Does paper swim?\par
\textbf{EKWE:} They will empty the sea, Uncle. We must sell the canoe while it is worth something.\par
\textbf{MOLUA:} \emph{(holding up the net)} My father knotted this twine. Sell the canoe\ldots{} \emph{(beat)} Ekwe, a man does not sell his hands.\par
\textbf{EKWE:} And a man does not eat pride!\par
\textbf{Marking guide:} correct format (capitals, colons, bracketed italic directions in present tense) — 4 marks; economical, meaningful directions — 2 marks; dialogue that reveals character and escalates conflict — 6 marks; evidence of a beginning arc (situation established, disturbance introduced) — 3 marks. Typical faults to correct in revision: characters announcing feelings, over-long speeches, and stage directions that merely repeat the dialogue.
\end{corrige}

\begin{aretenir}[Key Points of Section 8]
\begin{itemize}
\item A play script communicates through \cle{dialogue} and \cle{stage directions} only; nothing may be narrated.
\item Stage directions are bracketed, italicised, in the present tense, economical — and never spoken.
\item Good dialogue is fragmented, distinctive to each character, and always advances conflict or reveals character — never an essay aloud.
\item Even a one-act play needs the full arc: exposition, inciting incident, rising action, climax, resolution.
\item Adapt \emph{themes} from prescribed texts into \emph{original} situations; never copy characters or scenes.
\item Always test a script by reading it aloud; revise where readers stumble or lose interest.
\end{itemize}
\end{aretenir}

\coursec{Section 9: Prescribed Text Study I — Twilight of Misty Foliage by Ben Jama}

In Section 8, you learned to \emph{create} literature — to write dialogue, stage directions and dramatic structure yourselves. We now reverse the direction of travel: instead of producing a text, you will learn to \emph{interrogate} one. Paper 3 of the GCE Advanced Level does not ask whether you have read the prescribed novel; it asks whether you can \cle{analyse} it — that is, whether you can move from what happens in the story to what the story \emph{means}, and how the writer makes it mean it. Our first prescribed prose text is \emph{Twilight of Misty Foliage} by the Cameroonian novelist Ben Jama, a work firmly rooted in the social landscape of Anglophone Cameroon.

\begin{attention}[A note on editions and quotations]
Editions of this novel circulate in different printings, and page numbers vary. The short quotations used in this section are illustrative models of the \emph{kind} of evidence examiners reward. When you revise, replace them with the exact wording of your own copy, and always check quotations against the text before the examination.
\end{attention}

\courssub{Historical and Cultural Context}

No novel is written in a vacuum. To understand \emph{Twilight of Misty Foliage}, you must first place it on the map of Cameroonian history. The novel belongs to the tradition of Anglophone Cameroonian writing that emerged after independence and reunification (1960--1961), a tradition preoccupied with three great questions:

\begin{itemize}
\item \textbf{The colonial inheritance.} The territory that became Anglophone Cameroon passed from German to British administration before joining the Republic of Cameroon. Villages like the one evoked in the novel carry this layered history: mission schools, colonial courts, and imported languages sitting alongside older institutions.
\item \textbf{Tradition versus change.} The village community in the novel is governed by elders, custom and the rhythms of the farming season, yet it is increasingly penetrated by money, formal education, urban migration and new religions. The drama of the novel grows out of this friction.
\item \textbf{The individual versus the community.} In the world of the novel, a person is never merely an individual; he or she is a son, a daughter, a member of a lineage. Personal choices — whom to marry, whether to leave for the town, whether to obey an elder — are therefore always public events.
\end{itemize}

\begin{savaistu}[Why context earns marks]
In a context question, examiners award marks for linking the passage to its \emph{moment}: the social pressures, customs and historical forces at work. A candidate who writes ``this shows the conflict between tradition and modernity in post-colonial Cameroon'' and supports it with a detail from the passage will always outscore one who merely paraphrases the extract.
\end{savaistu}

\courssub{Plot Overview and Narrative Structure}

The novel traces the fortunes of a village community — and of the young protagonist at its centre — as the certainties of the old order are tested by new forces. Without reducing the book to a summary (you must read it in full), keep this structural skeleton in mind:

\begin{enumerate}
\item \textbf{Exposition:} the village world is established — its customs, its elders, its land, and the protagonist's place within it.
\item \textbf{Disruption:} an outside force (education, money, a stranger, a death, a contested decision) disturbs the equilibrium.
\item \textbf{Complication:} loyalties split; the protagonist is pulled between obedience to tradition and the pull of change.
\item \textbf{Climax:} the conflict becomes irreversible; a choice or event seals the characters' fates.
\item \textbf{Resolution:} the ``twilight'' of the title — nothing is simply restored; something has ended, and something uncertain has begun.
\end{enumerate}

Notice that the title itself announces the structure: \emph{twilight} is the moment between day and night, and \emph{misty foliage} is a landscape half-hidden, half-revealed. The novel is built as a story of \cle{transition} — of a world seen through mist, where nothing is entirely clear.

\courssub{Character Analysis and Development}

\begin{methode}[The character analysis frame]
For any character, build your answer in four moves:
\begin{enumerate}
\item \textbf{Traits:} name two or three dominant qualities (e.g.\ dutiful, proud, restless).
\item \textbf{Evidence:} support each trait with a short quotation or a precise reference to an incident.
\item \textbf{Development:} show how the character changes (or refuses to change) between the beginning and the end.
\item \textbf{Significance:} state what the character \emph{represents} — which theme or social force he or she embodies.
\end{enumerate}
\end{methode}

\textbf{The protagonist.}\par
The young central figure of the novel is best read as a \cle{liminal} character — one who stands on a threshold. Educated enough to see the limits of the village, rooted enough to feel its claims, the protagonist embodies the novel's central tension. In analysis, trace the arc: early obedience and belonging $\rightarrow$ growing doubt and inner conflict $\rightarrow$ a decisive act $\rightarrow$ a mature, often bittersweet, understanding. The protagonist's development \emph{is} the theme of identity in motion.

\textbf{The elders.}\par
The elder figures (parents, chiefs, custodians of custom) are not villains; Jama draws them with sympathy. They represent continuity, wisdom and communal memory — but also rigidity. The best essays avoid the lazy formula ``tradition is bad'' and instead show how the novel dramatises the \emph{cost} of both obedience and rebellion.

\textbf{The agents of change.}\par
Secondary characters associated with the town, money, schooling or new belief systems function as catalysts. They matter less as individuals than as forces: through them, the outside world enters the village.

\begin{popfigure}
\begin{center}
\begin{tikzpicture}[
 every node/.style={font=\small},
 main/.style={draw=popblue, fill=popblueL, rounded corners, thick, minimum width=3.2cm, minimum height=1cm, align=center},
 sec/.style={draw=poppurple, fill=white, rounded corners, thick, minimum width=2.6cm, minimum height=0.9cm, align=center},
 rel/.style={font=\scriptsize\itshape, text=popmuted}
]
\node[main, align=center] (p) at (0,0){\textbf{The Protagonist}\\(threshold figure)};
\node[sec, align=center] (e) at (-4.2,2.2){\textbf{The Elders}\\tradition, lineage};
\node[sec, align=center] (a) at (4.2,2.2){\textbf{Agents of Change}\\town, school, money};
\node[sec, align=center] (c) at (-4.2,-2.2){\textbf{The Community}\\village opinion};
\node[sec, align=center] (i) at (4.2,-2.2){\textbf{Intimates}\\friend / love interest};
\draw[thick, popline] (p) -- (e) node[rel, midway, above left] {duty vs.\ doubt};
\draw[thick, popline] (p) -- (a) node[rel, midway, above right] {attraction, risk};
\draw[thick, popline] (p) -- (c) node[rel, midway, below left] {judgement, shame};
\draw[thick, popline] (p) -- (i) node[rel, midway, below right] {trust, confession};
\draw[dashed, popmuted] (e) -- (a) node[rel, midway, above] {the central conflict};
\end{tikzpicture}
\end{center}
\legende{Character relationship map: every relationship is a channel through which the tradition--change conflict reaches the protagonist.}
\end{popfigure}

\courssub{Themes, Symbols and Motifs}

\begin{definition}[Motif]
A \cle{motif} is a recurring element — an image, phrase, object or situation — that appears throughout a work and reinforces a \cle{theme}. A theme is the abstract idea (e.g.\ identity); a motif is the concrete, repeated vehicle that carries it (e.g.\ mist that blurs vision whenever a character faces a moral choice).
\end{definition}

The title gives you the novel's master motif: \textbf{mist} and \textbf{foliage}. Mist obscures; it makes the familiar strange and hides the path ahead. Foliage is dense, layered, alive — beautiful but difficult to see through. Together they image a society in which truth is hard to discern and the future is veiled. Track where mist, fog, thick bush, dawn and twilight appear in the narrative, and ask what is uncertain at each of those moments.

\begin{center}
\begin{tabularx}{\linewidth}{|l|X|X|}
\hline
\rowcolor{popblueL} \textbf{Motif / Symbol} & \textbf{Linked theme} & \textbf{How to use it in an essay} \\
\hline
Mist, fog, blurred sight & Uncertainty of identity; a future that cannot yet be seen & ``The recurring mist externalises the protagonist's inner confusion\ldots'' \\
\hline
Twilight (between day and night) & Transition between tradition and modernity; the end of an era & ``Twilight is neither day nor night: the village, like the hero, is suspended between two worlds.'' \\
\hline
Dense foliage / the bush & The weight of tradition — protective yet concealing & ``Foliage shelters but also hides; custom both protects and imprisons.'' \\
\hline
The path / road to town & Change, ambition, departure and its costs & ``Each journey along the road marks a stage in the hero's separation from home.'' \\
\hline
The seasons and farming cycle & Continuity of communal life against individual time & ``The seasons repeat; the characters cannot.'' \\
\hline
The iroko / ancestral tree & Rootedness, memory, and the permanence of the land & ``The iroko stands witness; human disputes are brief beneath its shade.'' \\
\hline
\end{tabularx}
\end{center}

\begin{aretenir}[Key themes to master]
\begin{itemize}
\item \textbf{Identity and belonging:} who am I when my community's answer no longer satisfies me?
\item \textbf{Tradition and change:} neither is simply right; the novel weighs both.
\item \textbf{Generational conflict:} elders and youth read the same world differently.
\item \textbf{The individual and the community:} private desire versus public duty.
\item \textbf{Alienation and return:} the price of leaving and the impossibility of a full return.
\end{itemize}
\end{aretenir}

\courssub{The Writer's Style and Language}

Jama's narrative voice combines three registers, and recognising them is itself an examination skill:

\begin{itemize}
\item \textbf{Lyrical description:} the landscape passages, heavy with mist and vegetation imagery, slow the narrative and create atmosphere. This is \cle{local colour} raised to symbolism.
\item \textbf{Realist dialogue:} characters speak with the rhythms of Cameroonian village English, proverbs included. Proverbs signal communal wisdom — note who uses them and who stops using them.
\item \textbf{Reflective narration:} the third-person narrator moves close to the protagonist's consciousness (free indirect style), letting us feel doubt from the inside.
\end{itemize}

\begin{exemplebox}[Analysing a descriptive sentence]
Suppose the narrator writes that ``the mist came down over the compound until the iroko tree was only a rumour of itself.'' A weak answer says: \emph{this describes the weather.} A strong answer says: \emph{the personified ``rumour'' turns a solid landmark into hearsay, dramatising how the certainties of the village world are dissolving for the protagonist — description has become theme.}
\end{exemplebox}

\begin{exemplebox}[Analysing free indirect style]
Consider: ``He would not go. The elders had spoken; the matter was closed. But the road to the motor park lay open, a pale scar in the bush.'' The narration slides from the communal voice (``The elders had spoken'') to the protagonist's private perception (``a pale scar''). This technique forces the reader to inhabit the conflict rather than observe it from outside.
\end{exemplebox}

\courssub{Answering Context and Essay Questions}

\begin{attention}[The cardinal sin: plot retelling]
The most common failure in Paper 3 is narrating the story instead of analysing it. ``Then the hero refuses the elders and leaves the village'' earns almost nothing. ``The hero's refusal — signalled by the sudden disappearance of proverbs from his speech — marks the moment individual identity defeats communal duty'' earns marks. Rule: \textbf{every point must carry a quotation or a precise textual reference, and every quotation must be followed by analysis.}
\end{attention}

\begin{exoresolu}[Context-style question on a descriptive extract]
\textbf{Question.} Read the passage in which the narrator describes the village at dusk, with mist gathering as the protagonist waits for the elders' decision. (a) What atmosphere does the writer create? (b) How does the description reflect the protagonist's state of mind?
\tcblower
\textbf{Model answer.} (a) The writer creates an atmosphere of suspension and unease. The gathering mist blurs familiar objects, so that the village — normally a place of certainty — becomes strange and unreadable. The falling dusk reinforces this: the scene takes place between day and night, in a ``twilight'' that echoes the title. (b) The landscape mirrors the protagonist's mind. Just as he cannot see through the mist, he cannot foresee the elders' decision or his own future. The veiled path is an objective correlative for his uncertainty; nature is made to think and feel with him. The passage is therefore not decorative description but psychological and thematic writing: the mist motif fuses setting, character and the theme of an obscured future.
\end{exoresolu}

\begin{exoresolu}[Essay question on character]
\textbf{Question.} ``The protagonist of \emph{Twilight of Misty Foliage} changes; the village does not.'' Discuss.
\tcblower
\textbf{Model plan and opening.} \emph{Introduction:} define the terms — the protagonist's arc is one of growing self-knowledge; the village's ``refusal to change'' must be qualified, since the community too is altered by the events. \emph{Paragraph 1:} the protagonist at the start — dutiful, embedded in custom (evidence: early scenes of obedience, proverb-laden speech). \emph{Paragraph 2:} the turning points — encounters with the agents of change; the disappearance of certainty, tracked through the mist motif. \emph{Paragraph 3:} the final state — a character who has chosen, and paid for the choice. \emph{Paragraph 4 (counter-argument):} the village \emph{appears} unchanged — the seasons turn, the elders still sit — but the departure of the young has already transformed it; twilight has fallen on the old order. \emph{Conclusion:} the statement is half true, and the half that is false is the novel's point.
\end{exoresolu}

\begin{exercice}[Practice tasks]
\begin{enumerate}
\item Choose one elder character and one ``agent of change'' character. Apply the four-step character analysis frame to each, with one quotation per trait.
\item Find three occurrences of the mist or foliage motif in the novel. For each, write two sentences linking the image to a theme.
\item Write a 250-word answer to: ``How does Jama use the natural landscape to comment on social change?''
\end{enumerate}
\end{exercice}

\begin{corrige}[Exercise 1 -- guidance]
\textbf{1.} A strong frame entry looks like: \emph{Trait:} the elder is dignified and inflexible. \emph{Evidence:} quote his refusal to reconsider the decision, noting the formal, proverb-heavy register. \emph{Development:} he does not change — and that fixity is itself the point. \emph{Significance:} he embodies tradition's strength and its tragic limit. \textbf{2.} Each occurrence should be tied to a moment of uncertainty or transition; the motif is never random. \textbf{3.} Structure: landscape as atmosphere $\rightarrow$ landscape as symbol (mist, twilight, path) $\rightarrow$ landscape as commentary (nature endures; human orders pass). Conclude that setting in this novel is an argument, not a backdrop.
\end{corrige}

\begin{aretenir}[Exam tip: a quotation bank]
Memorise eight to ten \textbf{short, flexible quotations} — phrases of three to seven words — that can serve several themes at once (a mist image can illustrate identity, uncertainty \emph{and} tradition--change). A short quotation accurately recalled is worth more than a long one half-remembered. Write each on a revision card with the themes it can serve.
\end{aretenir}

\coursec{Section 10: Prescribed Text Study II — Unanswered Cries by Osman Conteh}

\courssub{From Misty Foliage to Unanswered Cries: A Natural Transition}

In Section 9, we studied Ben Jama's \emph{Twilight of Misty Foliage}, a novel deeply rooted in the Cameroonian landscape and concerned with the tensions between communal tradition and individual aspiration. We now cross West Africa to Sierra Leone, where Osman Conteh's \emph{Unanswered Cries} confronts us with one of the most painful social questions a society can face: what happens when a tradition harms the very people it claims to protect? The analytical tools you mastered with Jama --- character analysis, theme tracking, conflict mapping, and the PEE paragraph --- remain exactly the same. What changes is the urgency of the subject matter: Conteh writes not merely to entertain but to \emph{advocate}, and understanding that purpose is the key to unlocking the whole novel.

\begin{savaistu}[Paper 3 alert: the "advocacy novel" in the GCE]
The GCE Advanced Level Paper 3 (Literary Appreciation) frequently sets questions on \cle{writer's purpose} and \cle{intended effect on the reader}. An advocacy novel like \emph{Unanswered Cries} is a gift: the purpose is explicit (to expose and condemn FGM), so you can spend your energy analysing \emph{how} the writer achieves that purpose through character, conflict, irony, narrative perspective and legal realism. Always ask: "What does Conteh want me to feel, think or do at this precise moment?"
\end{savaistu}

\courssub{Context of the Novel: Sierra Leone, Society and the Writer's Purpose}

Osman Conteh is a Sierra Leonean writer and lawyer, and \emph{Unanswered Cries} grows directly out of his country's social reality. The novel is set in Sierra Leone and centres on the practice of \cle{female circumcision} (also called female genital mutilation or cutting), defended within the story by traditionalists as an essential rite of passage into womanhood, and resisted by the young protagonist who experiences it as violence.

To read the novel well, you must hold three contextual facts in mind:

\begin{itemize}
\item \textbf{The practice is communal, not individual.} It is sustained by families, secret societies (the \emph{Bondo} society) and respected elders --- often by women themselves --- which is why a single girl's refusal creates a crisis far bigger than herself.
\item \textbf{The law and modern institutions exist alongside tradition.} Courts, lawyers, schools and city life represent an alternative order, and the novel stages the collision between these two orders.
\item \textbf{The writer is an advocate.} Conteh's legal background shows in the courtroom dimension of the story and in the novel's persuasive energy: he wants readers to hear the ``cries'' of vulnerable girls that society refuses to answer.
\end{itemize}

\begin{savaistu}[Why this matters for Cameroonian candidates]
The GCE examiner does not expect you to be an expert on Sierra Leone. What is rewarded is your ability to connect a text to its \cle{social context}: to show that literature in Africa often carries a reformist mission, speaking for those whom custom silences. You have seen the same impulse in Cameroonian writing that questions harmful customs while respecting culture as a whole. In your essays, you may briefly compare Conteh's approach to that of Cameroonian writers like Bate Besong or Francis Nyamnjoh who also interrogate tradition.
\end{savaistu}

\courssub{Character Introduction and Analysis}

The novel's power comes from a small, sharply drawn cast, each figure embodying a position in the central debate.

\textbf{Olabisi}, the young protagonist, is an intelligent, spirited girl sent from the city to stay with her mother in the village. She embodies \cle{modernity, education and individual rights}. Her refusal to submit to circumcision is the spark that ignites the plot. Analyse her not as a perfect heroine but as a frightened, determined child whose courage grows under pressure. Key moment: her internal monologue when she first hears of the plan --- "I will not let them cut me. My body belongs to me."

\textbf{Makalay}, Olabisi's mother, is caught between love for her daughter and loyalty to the tradition that shaped her own life. She embodies the \emph{internal conflict} of a whole generation: she has suffered the practice, yet has been taught to see it as honourable. Her silence is as telling as her speech. Key moment: the night before the ceremony, when she enters Olabisi's room but cannot speak the protection her daughter needs.

\textbf{Yah Posseh}, the formidable grandmother, represents \cle{tradition at its most uncompromising}. She is not a cartoon villain: within her worldview she is protecting her granddaughter's marriageability and honour. Conteh's skill is to make us understand her even as we resist her. Key moment: her proverbial defence of the rite --- "The palm nut does not fall far from the tree; a woman uncut is like a house without a roof."

\textbf{Ade Jones}, the lawyer who takes up Olabisi's cause, embodies the \emph{modern legal conscience} --- the belief that the state must protect the vulnerable when custom will not. His professional detachment is laced with personal outrage. Around these central figures orbit the circumcisers (the \emph{Soweis}), elders and community members whose pressure makes the conflict societal rather than merely domestic.

\begin{attention}[Common mistake: flat character judgements]
Weak scripts call Yah Posseh ``evil'' and Olabisi ``good'' and stop there. The examiner rewards \emph{complexity}: show that Yah Posseh acts from conviction, that Makalay is torn, and that Olabisi herself feels fear and doubt. Literature essays score highly when characters are treated as human beings, not slogans. Remember: \cle{characterisation serves theme} --- each character's complexity deepens the novel's exploration of how tradition operates from within.
\end{attention}

\begin{exemplebox}[Character analysis: Makalay's divided self]
\textbf{Question:} How does Conteh present Makalay's internal conflict in Chapter 7?
\textbf{Model response:} Conteh uses \cle{free indirect discourse} to blend Makalay's thoughts with the narrator's voice: "She had bled on that same mat, had bitten the same cloth to stifle her screams. And yet, and yet --- the honour of the family, the whispers of the elders, the fear of a daughter unmarriageable." The repetition of "and yet" mimics the oscillation of her mind. The sensory memory ("bled", "bitten", "screams") clashes with abstract nouns ("honour", "whispers", "fear"), dramatising the war between lived trauma and internalised ideology. This technique makes her complicity tragic rather than villainous.
\end{exemplebox}

\courssub{Themes of the Novel}

\begin{definition}[Theme]
A \cle{theme} is the central idea or message that a literary work explores through its characters, plot, setting and language. A theme is best expressed as a full statement (``the novel shows that tradition can become oppression when it denies individuals the right to choose'') rather than a single word (``tradition'').
\end{definition}

The major themes of \emph{Unanswered Cries} are:

\begin{itemize}
\item \textbf{Tradition versus modernity.} The village's ancient rite collides with education, law and city values. Conteh does not reject tradition wholesale; he rejects the \emph{coercion} that hides behind it. The novel asks: can a culture survive without harming its daughters?
\item \textbf{Gender and the control of women's bodies.} The practice is done \emph{to} girls, largely \emph{by} women, \emph{for} a marriage market controlled by men --- a chain of oppression the novel exposes link by link. The \emph{Bondo} society, though female-run, ultimately serves patriarchal interests.
\item \textbf{Justice and the law.} The courtroom strand asks whether written law can protect a child against the weight of communal custom. Ade Jones's legal strategy --- invoking the Constitution, the Children's Act, international conventions --- shows law as a tool, not a magic wand.
\item \textbf{The unanswered cries of the vulnerable.} The title itself is a theme: the cries of girls go ``unanswered'' because adults choose not to hear. The novel is Conteh's answer --- literature as a hearing finally granted. The motif of \cle{silence vs voice} runs throughout: Olabisi's scream, the drumming that drowns it, the courtroom where she finally speaks.
\item \textbf{Courage and individual choice.} Olabisi's stand affirms that one voice, however young, can challenge an entire system. Her courage is not absence of fear but action despite it.
\item \textbf{Education as liberation.} Olabisi's city schooling gives her the vocabulary of rights ("my body", "the law", "choice") that the village lacks. The novel valorises education not as Westernisation but as empowerment.
\end{itemize}

\begin{attention}[Common mistake: theme-spotting without anchoring]
The single most frequent weakness in GCE literature essays is discussing themes in the abstract --- ``the novel is about tradition and modernity'' --- without anchoring every claim in a \cle{specific episode}. Rule: no theme statement without a scene, a decision, a quotation or a consequence from the text to support it. Use the \textbf{PEE} method (Point, Evidence, Explanation) for every thematic claim.
\end{attention}

\begin{exemplebox}[Thematic analysis with textual anchoring]
\textbf{Theme:} Gender and the control of women's bodies.
\textbf{Episode:} The scene where the \emph{Soweis} inspect Olabisi's body before the ceremony.
\textbf{Analysis:} The women's clinical handling of the girl --- "fingers probing, eyes assessing, tongues clucking" --- reduces a human being to a commodity. The verb "probing" suggests violation; "assessing" suggests valuation for the marriage market. That the violators are women deepens the horror: the patriarchy has colonised their gaze. Conteh writes, "They did not see a child; they saw a vessel." This metaphor encapsulates the theme: the female body as vessel for tradition, not person with rights.
\end{exemplebox}

\courssub{Conflict in the Novel}

\begin{propriete}[Conflict drives plot]
Conflict is the engine of narrative: \cle{identifying the central conflict clarifies the whole novel}, because every episode either intensifies it, complicates it or moves it towards resolution. In \emph{Unanswered Cries} the central conflict is Olabisi's refusal of circumcision versus the community's determination to enforce it; all other conflicts radiate from this core. (Stating the central conflict precisely is itself an examinable skill in Paper 3.)
\end{propriete}

The novel operates on three interlocking levels of conflict:

\begin{enumerate}
\item \textbf{Internal conflict} --- within one character: Makalay torn between maternal love and customary duty; Olabisi torn between fear and resolve; even Yah Posseh, in quieter moments, questions whether the old ways still serve.
\item \textbf{Interpersonal conflict} --- between characters: Olabisi against Yah Posseh (granddaughter vs grandmother, future vs past); Makalay against Yah Posseh (daughter vs mother, doubt vs certainty); Ade Jones against the elders (lawyer vs custom, individual rights vs communal authority).
\item \textbf{Societal conflict} --- the individual against the community and its institutions: one girl against the secret society, the elders and public opinion; modern law against customary law; the city against the village.
\end{enumerate}

\begin{popfigure}
\begin{center}
\begin{tikzpicture}[
  every node/.style={font=\small},
  chara/.style={draw=popblue, fill=popblueL, rounded corners=2pt, align=center, text width=2.6cm, inner sep=3pt},
  conf/.style={draw=poporange, fill=white, rounded corners=2pt, align=center, text width=2.4cm, inner sep=3pt, font=\footnotesize},
  lk/.style={thick, popmuted}
]
\node[chara, fill=popgreenL, draw=popgreen, align=center] (olabisi) at (0,0){\textbf{Olabisi}\\ (refusal, courage)};
\node[chara, align=center] (yah) at (-4.5,2.2){\textbf{Yah Posseh}\\ (tradition)};
\node[chara, align=center] (makalay) at (4.5,2.2){\textbf{Makalay}\\ (divided mother)};
\node[chara, align=center] (ade) at (4.5,-2.2){\textbf{Ade Jones}\\ (law, justice)};
\node[chara, align=center] (comm) at (-4.5,-2.2){\textbf{Community \&}\\ \textbf{secret society}};
\node[conf, align=center] (c1) at (-2.6,1.35){interpersonal:\\ refusal vs enforcement};
\node[conf, align=center] (c2) at (2.7,1.35){internal:\\ love vs duty};
\node[conf, align=center] (c3) at (2.6,-1.35){societal:\\ modern law vs custom};
\node[conf, align=center] (c4) at (-2.7,-1.35){societal:\\ individual vs community};
\draw[lk] (olabisi) -- (yah);
\draw[lk] (olabisi) -- (makalay);
\draw[lk] (olabisi) -- (ade);
\draw[lk] (olabisi) -- (comm);
\draw[lk, dashed] (yah) -- (comm);
\draw[lk, dashed] (makalay) -- (ade);
\end{tikzpicture}
\end{center}
\legende{Conflict web: every character relationship in \emph{Unanswered Cries} embodies one of the three conflict types, all radiating from Olabisi's central refusal. Dashed lines indicate indirect but powerful influence.}
\end{popfigure}

\begin{methode}[Mapping conflict for exam essays]
\begin{enumerate}
\item Identify the \textbf{central conflict} in one sentence.
\item List the \textbf{three levels} (internal, interpersonal, societal) with one example each.
\item Trace how a key episode \textbf{escalates} the conflict (e.g., the night Olabisi flees: internal fear $\rightarrow$ interpersonal chase $\rightarrow$ societal manhunt).
\item Note the \textbf{resolution} of each level: does internal conflict resolve? (Makalay's final choice); interpersonal? (Olabisi/Yah Posseh); societal? (court verdict).
\end{enumerate}
\end{methode}

\courssub{Plot Overview: Timeline of Key Events}

\begin{popfigure}
\begin{center}
\begin{tikzpicture}[font=\small]
\draw[thick, popblue, ->] (0,0) -- (13.2,0);
\foreach \x/\lab in {0.4/{Olabisi sent to\\ the village}, 3.0/{Bond with\\ mother; village life}, 5.6/{Plans for the\\ rite revealed}, 8.2/{Olabisi refuses;\\ pressure mounts}, 10.8/{Flight and legal\\ battle}, 12.8/{Cries finally\\ heard}}{
  \draw[popblue, thick] (\x,0.12) -- (\x,-0.12);
  \node[align=center, text width=2.1cm, below, font=\footnotesize] at (\x,-0.15) {\lab};
}
\node[above, font=\footnotesize\itshape, popmuted] at (6.6,0.25) {rising tension $\longrightarrow$ climax $\longrightarrow$ resolution};
\end{tikzpicture}
\end{center}
\legende{Timeline of key plot movements in \emph{Unanswered Cries}: from Olabisi's arrival in the village to the legal confrontation that gives her cries an answer.}
\end{popfigure}

\begin{aretenir}[Detailed plot points for revision]
\begin{itemize}
\item \textbf{Exposition (Ch. 1--3):} Olabisi arrives; warmth with Makalay; village routines; subtle hints of the coming "ceremony".
\item \textbf{Inciting incident (Ch. 4):} Yah Posseh announces the date; Olabisi overhears and realises what "becoming a woman" means.
\item \textbf{Rising action (Ch. 5--9):} Olabisi's refusal; family conferences; Yah Posseh's proverbs vs Olabisi's school knowledge; Makalay's silence; the \emph{Soweis} arrive; physical preparation begins.
\item \textbf{Climax (Ch. 10--11):} Night of the ceremony; Olabisi flees into the bush; pursuit; Ade Jones intervenes; police protection granted.
\item \textbf{Falling action (Ch. 12--13):} Court case: customary law vs constitutional rights; expert testimony; Yah Posseh's dignified defence; Olabisi's testimony.
\item \textbf{Resolution (Ch. 14):} Judgment in Olabisi's favour; Makalay finally embraces her daughter openly; Yah Posseh's ambiguous silence; Olabisi returns to city, changed but whole.
\end{itemize}
\end{aretenir}

\courssub{Literary Style: Narrative Technique, Language and Emotional Appeal}

Conteh's style serves his advocacy. Note five features:

\begin{itemize}
\item \textbf{Accessible, direct prose.} The language is clear enough for young readers --- the very audience most affected by the issue --- yet capable of great emotional intensity. Short sentences dominate moments of crisis: "She ran. She did not look back."
\item \textbf{Emotional appeal (pathos).} Scenes of fear, pain and loneliness are rendered so that the reader \emph{feels} what society refuses to hear. The title's ``cries'' become a recurring emotional motif: literal screams, metaphorical silences, the courtroom as amplification.
\item \textbf{Dialogue as battleground.} Confrontations between Olabisi, Yah Posseh and the elders are carried in dialogue, where proverbs and customary arguments clash with the plain language of rights. Conteh uses \cle{code-switching} between Krio-inflected English and standard English to mark identity and power.
\item \textbf{Legal realism.} The courtroom episodes draw on Conteh's legal experience, giving the novel a documentary credibility that strengthens its persuasive force. Legal terminology ("affidavit", "injunction", "precedent", "best interests of the child") is woven naturally into narrative.
\item \textbf{Symbolism and motif.} The \cle{mat} (where cutting happens, where Makalay bled, where Olabisi refuses to lie); the \cle{drums} (drowning cries, summoning community, rhythm of tradition); the \cle{mirror} (Olabisi seeing herself whole, then fragmented); the \cle{pen} (Ade Jones's tool, Olabisi's school tool, instrument of rights).
\end{itemize}

\begin{exemplebox}[Analysing a stylistic choice: contrast of registers]
When Conteh lets the circumcisers speak in the confident idiom of tradition --- "The knife cleanses; the blood seals the covenant" --- while Olabisi answers in short, frightened but firm sentences --- "No. It is my body. The law says no." --- the \emph{contrast of registers} itself dramatises the theme: an old, elaborate, communal rhetoric against a single, simple human voice. Quoting this contrast is far stronger than merely saying ``the novel is emotional''.
\end{exemplebox}

\begin{exemplebox}[Narrative perspective: third-person limited with strategic shifts]
The novel primarily uses \cle{third-person limited perspective} focalised through Olabisi, giving us direct access to her fear, reasoning and growing resolve. At key moments, the perspective shifts to Makalay (her divided thoughts) or Ade Jones (legal analysis). This technique:
\begin{itemize}
\item Aligns reader sympathy with the victim (we \emph{are} Olabisi).
\item Humanises the "oppressors" by showing their internal logic (Makalay's chapters).
\item Provides authoritative legal context without authorial intrusion (Ade Jones's chapters).
\end{itemize}
\textbf{Exam tip:} If asked about narrative technique, identify the \textbf{dominant mode} (third-person limited) and analyse the \textbf{effect of shifts}.
\end{exemplebox}

\courssub{Answering Examination Questions: Method and Model}

\begin{methode}[The PEE paragraph for literature essays]
\begin{enumerate}
\item \textbf{Point.} Open with one clear claim answering the question (``Conteh presents tradition as oppressive when it silences individual choice'').
\item \textbf{Evidence.} Anchor it in a specific episode, decision or brief quotation from the novel. Use \emph{ellipsis} for long quotes: "The knife cleanses... the blood seals the covenant."
\item \textbf{Explanation.} Show \emph{how} the evidence proves the point --- comment on technique (imagery, syntax, perspective), effect on reader, and link to theme --- then connect back to the question.
\end{enumerate}
A full essay is a chain of PEE paragraphs framed by a focused introduction and a genuine conclusion.
\end{methode}

\begin{methode}[Passage-based questions (Paper 3, Section A)]
\begin{enumerate}
\item \textbf{Read actively:} annotate for \textbf{conflict}, \textbf{tone}, \textbf{imagery}, \textbf{character voice}, \textbf{narrative perspective}.
\item \textbf{Contextualise:} where in the novel? what precedes/follows?
\item \textbf{Answer the specific sub-questions:} (a) usually identification; (b) usually analysis of technique/effect; (c) usually wider significance.
\item \textbf{Use the passage:} every claim must be supported by a word, phrase or detail \emph{from the extract}. Do not import outside knowledge unless asked.
\end{enumerate}
\end{methode}

\begin{exoresolu}[Passage-based question: identifying conflict and tone]
\textbf{Question.} Read a passage in which Yah Posseh insists that Olabisi must undergo the rite ``for her own good and the honour of the family'', while Olabisi weeps and refuses. (a) Identify two types of conflict in the passage. (b) Comment on the writer's attitude towards the tradition as revealed here.
\tcblower
\textbf{Model answer.} (a) The passage presents an \emph{interpersonal} conflict between Olabisi and her grandmother, who demands obedience, and, beneath it, a \emph{societal} conflict between the individual girl and the communal institution of the rite. (b) Conteh's attitude is critical though not contemptuous. The phrase ``for her own good'' is deeply ironic: the claimed benevolence is contradicted by Olabisi's tears, so the reader is guided to see how tradition disguises coercion as care. By giving the grandmother dignified language --- "the honour of the family", "the way of our mothers" --- yet centring the child's suffering through focalisation ("Olabisi felt the words like blows"), Conteh invites sympathy for Olabisi while showing that the oppressor sincerely believes in her own values --- which makes the critique of the \emph{system}, rather than of one woman, all the more powerful.
\end{exoresolu}

\begin{exoresolu}[Essay question: character as vehicle for theme]
\textbf{Question.} ``Makalay is the most tragic figure in the novel.'' How far do you agree? Refer to at least three episodes.
\tcblower
\textbf{Model answer structure.}
\textbf{Introduction:} Define "tragic" in this context: a character who suffers due to a flaw or circumstance beyond their control, evoking pity and fear. Thesis: Makalay's tragedy lies in her consciousness --- she \emph{knows} the harm yet cannot break free, making her more tragic than the ignorant Yah Posseh or the victorious Olabisi.
\textbf{PEE 1 (Episode: the night before the ceremony).} Point: Makalay's love is paralysed by indoctrination. Evidence: she enters Olabisi's room, "her hand hovering over the child's head", but leaves without a word. Explanation: the hovering hand is a visual metaphor for arrested agency; her silence contrasts with Olabisi's screams, showing how tradition colonises even maternal instinct.
\textbf{PEE 2 (Episode: the courtroom).} Point: Makalay's testimony reveals the depth of her fracture. Evidence: she tells the court "I did not protect her" rather than defending the custom. Explanation: this admission breaks the chain of silence; it is her moment of moral clarity, yet it comes too late to prevent the trauma, heightening the tragedy.
\textbf{PEE 3 (Episode: the final embrace).} Point: The resolution offers redemption but not erasure. Evidence: Makalay holds Olabisi "as if to stitch the torn places with her arms". Explanation: the simile acknowledges permanent damage; tragedy is not undone by the verdict. Makalay remains the woman who almost failed her daughter.
\textbf{Conclusion:} Makalay's tragedy is the novel's moral centre: she proves that tradition's victims include its enforcers.
\end{exoresolu}

\begin{exercice}[Essay practice --- timed condition: 45 minutes]
``In \emph{Unanswered Cries}, the real villain is not any character but custom itself.'' Discuss this view, referring closely to at least three episodes from the novel. Write a full essay with introduction, at least three PEE paragraphs, and a conclusion.
\end{exercice}

\begin{corrige}[Essay practice --- full model answer]
\textbf{Introduction.} Osman Conteh's novel dramatises the suffering of a girl who refuses female circumcision, yet its characters are rarely monsters: Yah Posseh acts from conviction and Makalay from divided love. This essay argues that Conteh deliberately directs our judgement away from individuals and towards the custom that programmes them, using the courtroom conflict, Makalay's paralysis, and Olabisi's isolation to expose tradition itself as the antagonist.

\textbf{PEE paragraph 1.} \emph{Point:} Conteh first establishes that even loving characters become agents of harm when custom speaks through them. \emph{Evidence:} Makalay, who has herself endured the rite and who cherishes her daughter, still cannot bring herself to shield Olabisi from it. In Chapter 7, she thinks: "She had bled on that same mat... And yet --- the honour of the family." \emph{Explanation:} A mother who fails to protect her child is normally condemned; Conteh instead makes us pity her, because her failure is shown to be the product of lifelong indoctrination. The free indirect discourse ("And yet") mimics the internalised voice of tradition. The effect is to shift the reader's anger from the woman to the system that formed her --- precisely the claim in the essay title.

\textbf{PEE paragraph 2.} \emph{Point:} Secondly, Yah Posseh's villainy is explicitly framed as the custom's voice, not her own invention. \emph{Evidence:} When challenged by Ade Jones, she replies: "I did not make the custom. The custom made me." \emph{Explanation:} This declarative sentence, stripped of metaphor, positions her as vessel rather than author. Her proverbs ("The palm nut does not fall far from the tree") are communal property, not personal malice. Conteh thus depersonalises the oppression: the true antagonist is the \emph{ideology} that turns grandmothers into guardians of the knife.

\textbf{PEE paragraph 3.} \emph{Point:} Thirdly, the legal strand of the novel frames custom, not persons, as what must be defeated. \emph{Evidence:} Ade Jones's closing argument: "We do not prosecute a grandmother; we challenge a practice that maims children." The judgment grants an injunction against \emph{the rite}, not against Yah Posseh. \emph{Explanation:} By making the climax a contest between written law and customary law, Conteh universalises the conflict: one girl's case stands for all the ``unanswered cries''. The resolution therefore reads as a verdict on a tradition, confirming that custom itself is the true villain of the novel.

\textbf{Conclusion.} Conteh's achievement is to make us mourn the enforcers even as we condemn the enforcement. Yah Posseh is not evil; she is \emph{used}. Makalay is not weak; she is \emph{colonised}. Only by naming custom as the villain can the novel offer hope: laws can change, minds can open, but a "villain" must be defeated. Olabisi's final walk away from the village, "her shadow long in the evening light", suggests that the custom's power, though ancient, is not eternal.
\end{corrige}

\begin{exercice}[Passage-based practice]
\textbf{Read the following extract (invented but representative of Conteh's style) and answer the questions.}

The drums had stopped. In the sudden silence, Olabisi heard her own heart hammering like a trapped bird. Yah Posseh stood before her, the ceremonial knife catching the morning light. "Child," the old woman said, her voice surprisingly gentle, "this is not punishment. This is belonging. Every woman in this village has walked this path. Your mother. Her mother. Me. Do you think we did not cry? We cried. And we became women."

Olabisi's throat closed. The words she had rehearsed --- \emph{my body, my right, the law} --- dissolved in the face of that ancient certainty. She saw, with a flash of terrible clarity, that Yah Posseh was not lying. The old woman believed every syllable. And that was what made it impossible to hate her, and impossible to yield.

\textbf{Questions:}
(a) Identify the narrative perspective in this extract and explain its effect. (3 marks)
(b) Analyse how Conteh uses contrast in this passage. (4 marks)
(c) What does the extract reveal about the nature of tradition as portrayed in the novel? (3 marks)
\end{exercice}

\begin{corrige}[Passage-based practice]
\textbf{(a)} Third-person limited focalised through Olabisi. We access her sensory experience ("heard her own heart", "saw... that Yah Posseh was not lying") and internal processes ("words she had rehearsed... dissolved", "flash of terrible clarity"). Effect: intimate alignment with the victim; we feel her fear and moral complexity.
\textbf{(b)} Contrasts: (i) Sound/silence --- drums stop, then "sudden silence" amplifies internal terror. (ii) Violence/gentleness --- "ceremonial knife" vs "voice surprisingly gentle". (iii) Past/present --- ancestral line ("your mother, her mother, me") vs Olabisi's modern vocabulary ("my body, my right, the law"). (iv) Knowledge/innocence --- Yah Posseh's "ancient certainty" vs Olabisi's "dissolved" rehearsed arguments. Each contrast deepens the tragedy: tradition is not brute force but seductive certainty.
\textbf{(c)} Tradition is portrayed as \emph{self-authenticating} (it validates itself through lineage), \emph{embodied in women} (matrilineal transmission), and \emph{impervious to external logic} (legal rights "dissolve" before it). The extract supports the novel's central insight: tradition persists not because people are evil but because it provides meaning and belonging that abstract rights cannot immediately replace.
\end{corrige}

\begin{aretenir}[Key points to remember for \emph{Unanswered Cries}]
\begin{itemize}
\item \textbf{Context:} Sierra Leone; advocacy novel against FGM by lawyer-novelist Osman Conteh; \emph{Bondo} society as institutional backdrop.
\item \textbf{Characters as positions:} Olabisi (choice, modernity, rights), Yah Posseh (tradition, communal continuity), Makalay (divided loyalty, internalised oppression), Ade Jones (law, justice, external intervention).
\item \textbf{Themes (as full statements):} Tradition becomes oppression when it denies choice; women's bodies are controlled by a chain linking female enforcers to male beneficiaries; law can protect the vulnerable but requires courage to activate; silence enables violence; education empowers resistance.
\item \textbf{Three conflict levels} --- internal (Makalay, Olabisi), interpersonal (Olabisi/Yah Posseh, Makalay/Yah Posseh), societal (individual vs community, statute law vs customary law) --- all radiating from Olabisi's refusal.
\item \textbf{Style:} direct prose, pathos, dialogic confrontation (proverbs vs rights language), legal realism, symbolism (mat, drums, mirror, pen), third-person limited with strategic focalisation shifts.
\item \textbf{Examination technique:} Anchor every theme in a specific episode; build essays from PEE paragraphs; for passage questions, analyse technique (perspective, contrast, irony, imagery) and link to whole-novel concerns.
\item \textbf{Title significance:} "Unanswered Cries" = the girls' screams drowned by drums; the community's wilful deafness; the legal system's previous silence; the novel itself as the answer.
\end{itemize}
\end{aretenir}

\begin{savaistu}[Connecting to the wider syllabus]
The skills you use here --- close reading, PEE paragraphs, conflict mapping, thematic anchoring --- are transferable to \emph{The Old Man and the Sea} (Section 11), \emph{The Swamp Dwellers} (Section 12) and the poetry in Section 13. In each case, ask: \textbf{What is the writer's purpose? How do character, conflict and style serve that purpose? What is the social context?} The GCE rewards consistent, rigorous application of these tools across all set texts.
\end{savaistu}

\coursec{Section 11: Prescribed Text Study III — The Old Man and the Sea by Ernest Hemingway}

In Section 10 we studied Osman Conteh's \emph{Unanswered Cries}, a novel in which a young girl's suffering exposes the conflict between tradition and human rights in West Africa. We now cross the Atlantic — geographically and stylistically — to a very different kind of text: a short, deceptively simple \cle{novella} about an old Cuban fisherman and a great fish. Ernest Hemingway's \emph{The Old Man and the Sea} (1952) is the shortest of your prescribed texts, yet it is arguably the densest in symbolic meaning. Where Conteh's power lies in social realism and emotional appeal, Hemingway's power lies in \emph{restraint}: almost everything important is left unsaid. Learning to read what is \emph{not} on the page is the central skill of this section — and it is precisely the skill the GCE Advanced Level literature paper rewards most highly.

\courssub{Hemingway, the Novella and the Cuban Context}

Ernest Hemingway (1899--1961) was an American novelist and journalist whose experiences as an ambulance driver in the First World War, as a war correspondent, and as a deep-sea fisherman shaped both his subjects and his style. \emph{The Old Man and the Sea}, published in 1952, was his last major work of fiction published during his lifetime and contributed decisively to his receiving the Nobel Prize in Literature in 1954.

The story is set in a small fishing village near Havana, Cuba, and in the Gulf Stream. Santiago, an old fisherman, has gone eighty-four days without catching a fish. His young apprentice, Manolin, has been forced by his parents to work on a luckier boat, though the boy still loves and cares for the old man. On the eighty-fifth day Santiago sails far out alone, hooks a gigantic marlin, and endures a three-day struggle to kill it. He lashes the fish to his skiff, but sharks attack on the journey home and strip the marlin to its skeleton. Santiago returns exhausted, carrying only the great bones.

\begin{savaistu}[Why a novella?]
A \cle{novella} is longer than a short story but shorter than a novel — typically focused on a single character, a single conflict and a single stretch of time. Hemingway called the form ideal for this story because nothing could be cut and nothing needed to be added. In the examination, you may be asked why the brevity matters: the concentrated form mirrors the concentrated struggle — one man, one fish, one sea. Note also that the novella form allows Hemingway to sustain the \cle{iceberg theory} (see below) without the digressions a novel might require.
\end{savaistu}

\begin{methode}[Approaching a prescribed novella for GCE Advanced Level]
\begin{enumerate}
\item \textbf{Read actively:} annotate for symbols, repeated motifs, shifts in tone, and moments of interior monologue.
\item \textbf{Master the plot timeline:} know the sequence of the 84 days, the three days at sea, and the return.
\item \textbf{Build a quotation bank:} select 10--15 short, versatile quotations covering character, theme, style and symbolism.
\item \textbf{Practise the "feature + evidence + effect" chain:} every analytical paragraph must follow this structure.
\item \textbf{Link to the syllabus themes:} struggle, dignity, pride, isolation/interdependence, man and nature.
\end{enumerate}
\end{methode}

\courssub{Symbolism in the Novella}

A \cle{symbol} is a concrete object, person or event that carries a meaning beyond its literal role in the story. Hemingway publicly insisted that the book was simply about "an old man, a boy, a sea and a fish" — yet generations of readers have found rich symbolic layers. Your task in the examination is not to choose between these views but to \emph{interpret the symbols with evidence}.

\begin{exemplebox}[Reading the key symbols: the marlin]
Consider the \cle{marlin}. Literally, it is the fish Santiago catches. Symbolically, it is his worthy opponent and even his brother: Santiago calls it beautiful and noble, and says he has "no luck" but must kill it out of love and respect. The marlin can therefore represent the ideal achievement — the great thing a person devotes a life to reaching — which is why its destruction by the sharks is so painful. \textbf{Key quotations:} "He is beautiful and noble and knows no panic" ; "I love you and respect you very much. But I will kill you dead before this day ends."
\end{exemplebox}

\begin{popfigure}
\begin{center}
\begin{tabularx}{\linewidth}{|l|X|X|}
\hline
\rowcolor{popblueL} \textbf{Symbol} & \textbf{Literal role} & \textbf{Symbolic meaning (with textual support)} \\
\hline
The sea (\emph{la mar}) & The setting; Santiago's workplace & Life itself — beautiful, generous and cruel at once; Santiago calls it \emph{la mar} (feminine, loved), unlike younger fishermen who treat it as an enemy (\emph{el mar}). Quotation: "He always thought of the sea as \emph{la mar} which is what people call her in Spanish when they love her." \\
\hline
The marlin & The great fish Santiago catches & The noble goal or worthy adversary; dignity shared between hunter and hunted. Quotation: "Never have I seen a greater, or more beautiful, or a calmer or more noble thing than you, brother." \\
\hline
The sharks & Scavengers that destroy the catch & Destructive forces — time, decay, critics, fate — that can destroy achievement without defeating the achiever. Quotation: "They beat me, Manolin. They truly beat me." \\
\hline
The lions on the beach & Santiago's recurring dream of African lions seen in his youth & Youth, strength, hope and the harmony of power with peace; the dream returns at the end, signalling an undefeated spirit. Quotation: "He was asleep in a short time and he dreamed of Africa when he was a boy and the long golden beaches and the white beaches, so white they hurt your eyes, and the high capes and the great brown mountains." \\
\hline
The mast and harpoon & Objects Santiago carries up the hill & Cross-like imagery: Santiago shoulders the mast as a burden of suffering, echoing endurance and quiet martyrdom. Quotation: "He started to climb again and at the top he fell and lay for some time with the mast across his shoulder." \\
\hline
DiMaggio / baseball & Santiago's hero, a baseball player with a bone spur & The ideal of professional excellence despite physical pain; a parallel to Santiago's own struggle. Quotation: "I would like to take the great DiMaggio fishing... They say his father was a fisherman. Maybe he was as poor as we are and would understand." \\
\hline
\end{tabularx}
\end{center}
\legende{Symbol interpretation chart: from literal role to symbolic meaning with textual evidence}
\end{popfigure}

\begin{attention}[Common mistake: fixing one meaning to a symbol]
Candidates often write "the sea \emph{means} life" as if symbols were codes with single answers. Symbols in literature are \emph{suggestive}, not mathematical. The examiners expect you to (1) propose an interpretation, (2) support it with a quotation or event, and (3) acknowledge that the symbol works on the literal level too. Never state a symbolic meaning without evidence from the text. Avoid phrases like "the sea represents life" — prefer "the sea \emph{suggests} the dual nature of life, as seen when Santiago calls it \emph{la mar}..."
\end{attention}

\courssub{Hemingway's Style: The Iceberg Theory}

Read any page of the novella aloud. The sentences are short. The words are simple. Nothing seems hidden. And yet readers feel enormous depth. Hemingway explained this with his famous \cle{iceberg theory} (also called the theory of omission).

\begin{definition}[The iceberg theory]
The \cle{iceberg theory} states that the deeper meaning of a story lies beneath the surface of the simple prose. Like an iceberg, only about one-eighth is visible above the water; the remaining seven-eighths — emotion, philosophy, symbolic weight — is submerged. The writer who truly knows his subject can \emph{omit} things, and the reader will feel their presence more strongly than if they were stated.
\end{definition}

\begin{popfigure}
\begin{center}
\begin{tikzpicture}[scale=1.0]
% water line
\draw[popblue, thick] (-4,0) -- (4,0);
\node[popblue, font=\small] at (4.6,0) {sea level};
% visible tip
\fill[popblueL, draw=popblue, thick] (-0.7,0) -- (0,1.3) -- (0.7,0) -- cycle;
% submerged mass
\fill[popblueL!60, draw=popblue, thick] (-0.7,0) -- (-1.9,-2.6) -- (0,-3.4) -- (1.9,-2.6) -- (0.7,0) -- cycle;
\node[align=center, font=\small, text=popdark] at (0,0.55) {surface text};
\node[align=center, font=\scriptsize, text=popdark] at (0,0.15) {short sentences};
\node[align=center, font=\scriptsize, text=popdark] at (0,-0.25) {simple diction, facts,};
\node[align=center, font=\scriptsize, text=popdark] at (0,-0.55) {actions, dialogue};
\node[align=center, font=\small, text=popdark] at (0,-2.6) {submerged meaning};
\node[align=center, font=\scriptsize, text=popdark] at (0,-3.05) {emotion, symbolism, philosophy};
% arrows
\draw[->, poporange, thick] (2.2,0.6) -- (0.6,0.75);
\node[poporange, font=\scriptsize, align=center] at (3.1,0.6) {what the\\ reader sees};
\draw[->, poppurple, thick] (2.2,-2.2) -- (1.2,-2.3);
\node[poppurple, font=\scriptsize, align=center] at (3.3,-2.2) {what the\\ reader feels};
\end{tikzpicture}
\end{center}
\legende{The iceberg theory: visible prose above, implied meaning below}
\end{popfigure}

Three stylistic features create this effect:

\begin{itemize}
\item \textbf{Simple diction.} Hemingway uses short, concrete, mostly Anglo-Saxon words: "He was an old man who fished alone in a skiff in the Gulf Stream." No adjective is decorative; every detail works.
\item \textbf{Understatement.} Extreme emotion is expressed minimally. After losing the marlin to the sharks, Santiago says only: "They beat me, Manolin. They truly beat me." The flatness of the sentence makes the grief \emph{more} powerful, not less.
\item \textbf{Repetition.} Key words and phrases ("I wish the boy were here", "I will show him what a man can do") recur like a refrain, charting the old man's loneliness and resolve.
\item \textbf{Polysyndeton and asyndeton.} Hemingway often strings clauses with repeated "and" (polysyndeton) to create a rhythmic, almost biblical cadence, or omits conjunctions (asyndeton) for speed and tension.
\item \textbf{Free indirect style.} The narrator slips into Santiago's consciousness without quotation marks, blending third-person objectivity with first-person intimacy.
\end{itemize}

\begin{propriete}[Style-effect link (examinable skill)]
In every literature answer, you must connect a stylistic feature to its effect. The formula is: \textbf{feature + quotation + effect on meaning or reader}. Example: "Hemingway's short declarative sentences ('He was an old man who fished alone') create a tone of calm acceptance, which makes Santiago's suffering appear dignified rather than pitiful."
\end{propriete}

\begin{exemplebox}[Analysing style in an exam paragraph]
\textbf{Question:} How does Hemingway's style convey Santiago's endurance?
\textbf{Model response:} Hemingway uses \textbf{repetition} of the phrase "I wish the boy were here" (four times during the struggle) to \textbf{effect} a structural emphasis on Santiago's isolation and his need for human connection, which paradoxically highlights his solitary endurance. The \textbf{simple diction} — "The old man hit him on the head for kindness" — strips violence of glamour, presenting killing as a necessary craft, reinforcing the theme of \textbf{dignity in labour}. The \textbf{polysyndeton} in "He was an old man who fished alone in a skiff in the Gulf Stream and he had gone eighty-four days now without taking a fish" creates a cumulative, inevitable rhythm that mirrors the relentless passage of time.
\end{exemplebox}

\courssub{Character Study: Santiago, the Code Hero}

\begin{definition}[The code hero]
Hemingway's \cle{code hero} is a protagonist who lives by a personal code of honour, courage and endurance, and who demonstrates \emph{grace under pressure} — remaining disciplined, skilful and dignified in the face of suffering and certain defeat. The code hero does not speak of his code; he \emph{lives} it.
\end{definition}

Santiago is the purest example of the code hero. Observe how he embodies the code:

\begin{itemize}
\item \textbf{Endurance:} his hands are cut by the line, his back is cramped, he is old and alone — yet he holds on for three days: "A man can be destroyed but not defeated."
\item \textbf{Skill and pride in craft:} he fishes with perfect technique, keeps his lines straighter than anyone, and respects the sea and the fish. "I keep my lines straighter than anyone."
\item \textbf{Grace in defeat:} he does not complain or curse the sharks; he accepts the loss, blames only himself for going "too far out", and plans to fish again.
\item \textbf{Humility before nature:} he asks the fish's forgiveness, calls it brother, and thanks the sea for the struggle.
\end{itemize}

\textbf{Manolin} is the necessary complement. The boy represents love, loyalty and continuity: he brings food, bait and conversation, and his faith in the old man restores what society (the other fishermen, the boy's parents) has withdrawn. Their relationship shows that the code hero's stoic independence is \emph{not} total isolation — Santiago's repeated wish, "I wish the boy were here," admits his need for human connection. This tension between \cle{isolation} and \cle{interdependence} is one of the novella's deepest themes.

\begin{aretenir}[Santiago vs. the other fishermen]
The younger fishermen use motorboats, treat the sea as \emph{el mar} (masculine, an enemy or a place of work), and mock or pity Santiago. Santiago uses a skiff, calls the sea \emph{la mar} (feminine, a partner), and fishes with traditional skill. This contrast establishes Santiago as the bearer of an older, more respectful relationship with nature — a key point for "man and nature" theme questions.
\end{aretenir}

\courssub{Narrative Technique}

The novella is told in the \cle{third person} by an omniscient narrator who stays close to Santiago's consciousness (limited omniscience). Within this frame, Hemingway uses two special devices:

\begin{itemize}
\item \textbf{Interior monologue.} Alone at sea, Santiago speaks his thoughts aloud ("'Fish,' he said, 'I love you and respect you very much. But I will kill you dead before this day ends.'"). This device lets us hear the old man's mind directly while remaining technically outside it — a bridge between third-person narration and first-person intimacy. It also characterises Santiago as a man who talks to himself because he is alone, yet his self-talk is purposeful, not mad.
\item \textbf{Dreams.} Santiago no longer dreams of storms, women or fights; he dreams of the \emph{lions on the African beach} of his youth. The dreams open a window into his deepest self — his longing for youthful strength and peaceful power — and frame the novella: it begins with the boy noticing the old man's dream and ends with Santiago dreaming of the lions again. The circular structure suggests renewal.
\item \textbf{Time manipulation.} The three days at sea are narrated in near real-time (scene), while the 84 days of bad luck are summarised in a single sentence (summary). This pacing focuses the reader's attention on the central trial.
\end{itemize}

\begin{methode}[Analysing narrative technique in an exam]
\begin{enumerate}
\item Identify the narrative voice (third-person limited omniscient).
\item Locate shifts into interior monologue or free indirect style.
\item Note the function of dreams: character revelation, thematic framing, structural symmetry.
\item Comment on pacing: scene vs. summary, and why the struggle gets "real time" treatment.
\item Link technique to theme: e.g., interior monologue reinforces isolation; dreams reinforce hope.
\end{enumerate}
\end{methode}

\courssub{Themes of the Novella}

\begin{aretenir}[Key themes to master — with thesis statements for essays]
\begin{itemize}
\item \textbf{Struggle and endurance:} Life is presented as a test of how much a person can bear while remaining true to a personal code. \emph{Thesis:} The novella redefines heroism not as victory but as the quality of the struggle itself.
\item \textbf{Dignity in defeat:} Santiago loses the fish but wins moral victory; "a man can be destroyed but not defeated." \emph{Thesis:} Destruction is physical and external; defeat is spiritual and internal — Santiago avoids the latter.
\item \textbf{Pride:} Pride drives Santiago too far out to sea, causing his loss — pride is both his strength and his flaw. \emph{Thesis:} Santiago's tragic flaw (hamartia) is the same quality that makes him heroic: his pride in his craft.
\item \textbf{Isolation and interdependence:} The old man is physically alone yet spiritually connected — to the boy, to the marlin, to the sea, even to DiMaggio. \emph{Thesis:} True independence requires the foundation of human connection.
\item \textbf{Man and nature:} Nature is neither enemy nor servant; the marlin is a brother, and killing it is both necessary and tragic. \emph{Thesis:} The novella presents a non-anthropocentric view: humans participate in nature's cycles rather than dominate them.
\item \textbf{Mentorship and legacy:} The relationship between Santiago and Manolin ensures the transmission of the code. \emph{Thesis:} The novella ends not with death but with apprenticeship — the code survives.
\end{itemize}
\end{aretenir}

\begin{exoresolu}[Analysing a passage for style and symbolism]
\textbf{Statement:} Read this passage: "He was an old man who fished alone in a skiff in the Gulf Stream and he had gone eighty-four days now without taking a fish." Identify two stylistic features and explain their effect; then state one theme the sentence introduces.
\tcblower
\textbf{Solution:}
\textbf{Feature 1 — simple diction and a single long declarative sentence.} Hemingway uses plain, concrete words ("old man", "skiff", "fish") with no ornament. \emph{Effect:} the tone is factual and calm, presenting hardship without self-pity, which establishes Santiago's dignity from the first line.
\textbf{Feature 2 — precise factual detail ("eighty-four days").} The exact number gives the struggle a measured, almost epic weight. \emph{Effect:} the reader feels the length of the old man's bad luck and anticipates the test of endurance to come.
\textbf{Feature 3 — polysyndeton ("and... and").} The repeated conjunction creates a cumulative rhythm that mirrors the relentless, day-after-day nature of Santiago's trial.
\textbf{Theme introduced:} struggle and endurance — the sentence immediately frames the novella as the story of a man defined by how he responds to prolonged failure.
\end{exoresolu}

\begin{exoresolu}[Comparative passage analysis: Santiago and the marlin]
\textbf{Statement:} Analyse the following passage for its presentation of the relationship between Santiago and the marlin: "Then the fish came alive, with his death in him, and rose high out of the water showing his great length and width and all his power and his beauty. He seemed to hang in the air above the old man in the skiff. Then he fell into the water with a crash that sent spray over the old man and over all of the skiff."
\tcblower
\textbf{Solution:}
\textbf{Symbolic equality:} The marlin "came alive, with his death in him" — a paradox that fuses vitality and mortality, suggesting the fish transcends the moment of killing. The vertical rise ("rose high out of the water") and the old man looking up create a visual parity: neither is above the other.
\textbf{Shared beauty and power:} The catalogue "great length and width and all his power and his beauty" uses polysyndeton to give each quality equal weight. The spray covering both man and fish ("over the old man and over all of the skiff") is a physical baptism, uniting them.
\textbf{Style-effect link:} The short, declarative sentences ("Then he fell into the water with a crash") mimic the sudden violence, while the absence of emotive adjectives forces the reader to supply the awe. This is the iceberg theory in action: the surface reports facts; the depth conveys reverence.
\textbf{Theme:} Man and nature as partners in a shared ordeal; the hunter honours the hunted.
\end{exoresolu}

\begin{exercice}[Exam-style essay practice]
Answer in about 350--400 words: "Santiago is defeated by the sharks but victorious as a man." Discuss this statement with close reference to Hemingway's style, symbolism and characterisation in \emph{The Old Man and the Sea}.
\end{exercice}

\begin{corrige}[Exercise 1]
\textbf{Model plan and answer outline:}
\textbf{Introduction:} Accept the paradox as the novella's central idea; define the code hero and the theme of dignity in defeat. Thesis: Hemingway separates \emph{destruction} (physical, external) from \emph{defeat} (spiritual, internal).
\textbf{Paragraph 1 (material defeat):} The sharks strip the marlin; Santiago returns with a skeleton. Symbolically, destructive forces (time, decay, fate) can ruin achievement. Quotation: "They beat me, Manolin. They truly beat me." Style: understatement magnifies loss.
\textbf{Paragraph 2 (moral victory — endurance):} Santiago's endurance over three days, his refusal to give up, his famous line "a man can be destroyed but not defeated"; grace under pressure defines the code hero. Quotation: "But man is not made for defeat... A man can be destroyed but not defeated."
\textbf{Paragraph 3 (moral victory — skill and respect):} He fishes with perfect technique, honours the marlin ("brother"), and accepts responsibility ("I went too far out"). This is the code in action.
\textbf{Paragraph 4 (style as evidence of victory):} Understatement ("They beat me") and simple diction make his acceptance dignified; the final dream of the lions shows an unbroken spirit; Manolin's renewed devotion ("Now we fish together again") confirms the victory is recognised by the only witness that matters.
\textbf{Conclusion:} Santiago is destroyed in body and catch, but never defeated in spirit — the essence of the code hero. The novella ends not with the skeleton on the beach but with the dream of lions: the future belongs to those who endure.
\end{corrige}

\begin{exercice}[Contextual linking exercise — Paper 2 style]
The GCE Paper 2 often asks you to compare across texts. Write a paragraph (150 words) comparing the presentation of \textbf{the individual vs. society} in \emph{The Old Man and the Sea} and \emph{Unanswered Cries} (Section 10).
\end{exercice}

\begin{corrige}[Exercise 2]
\textbf{Model answer:} In \emph{Unanswered Cries}, society (tradition, family, patriarchy) actively destroys the individual (Olabisi) through female genital mutilation and forced marriage; the collective is the antagonist. In \emph{The Old Man and the Sea}, society (the other fishermen, the boy's parents) marginalises Santiago — they call him \emph{salao}, the worst form of unlucky — yet the individual transcends society through a personal code. Olabisi's tragedy is that she has no code to shield her; Santiago's triumph is that his code (skill, endurance, respect) operates independently of social approval. Both texts, however, show that the individual needs an ally: Olabisi has none (her cries are unanswered), while Santiago has Manolin, whose loyalty bridges the gap between the old man's isolation and the community. Hemingway suggests dignity is forged in solitude but sustained by love; Conteh suggests dignity is denied when society silences the vulnerable.
\end{corrige}

\begin{attention}[Examination traps for \emph{The Old Man and the Sea}]
\begin{itemize}
\item \textbf{Don't retell the plot.} The examiner knows the story. Use plot details only as evidence for analytical points.
\item \textbf{Don't treat symbols as allegory.} The marlin is not \emph{only} Christ or \emph{only} a goal. It functions on multiple levels simultaneously.
\item \textbf{Don't ignore the boy.} Manolin is not a minor character; he is the emotional anchor and the future of the code. Essays that omit him lose marks.
\item \textbf{Don't confuse "simple style" with "simple mind".} The iceberg theory means the simplicity is a deliberate craft choice, not a limitation. Always analyse \emph{why} Hemingway chooses plain words.
\item \textbf{Watch the wording of questions.} "How does Hemingway present...?" requires style analysis. "Discuss the theme of...?" requires thematic tracking with evidence. "What is the significance of...?" requires symbolic and structural analysis.
\end{itemize}
\end{attention}

\begin{savaistu}[Cameroon connection: the fisherman's code]
Although Hemingway's Cuba is far from Cameroon, the figure of the skilled, respected elder who teaches a younger generation by example resonates strongly in many Cameroonian coastal and riverine communities (e.g., the Wouri estuary, the Sanaga river, Lake Chad). The relationship between Santiago and Manolin mirrors the \cle{apprenticeship} model still alive in traditional fishing, carving, and weaving crafts across the country. When you write about "mentorship and legacy", you can legitimately note this universal human pattern — but keep your textual evidence firmly in Hemingway's novella.
\end{savaistu}

\coursec{Section 12: Prescribed Text Study IV — The Swamp Dwellers by Wole Soyinka}

In Section 11, we followed an old fisherman alone on the sea and discovered how Hemingway's \emph{novella} uses a single, spare narrative voice to build a universal symbol of human endurance. We now move from prose fiction to \cle{drama}, and from the open sea to the enclosed, waterlogged world of a Nigerian village. Wole Soyinka's \emph{The Swamp Dwellers} asks many of the same questions as \emph{The Old Man and the Sea} — what does a human being owe to the land, to tradition, to hope? — but it asks them through \textbf{dialogue, stage action and ritual}, not through narration. This changes everything about how we read, analyse and write about the text. A play is not a novel that happens to be written in speeches; it is a blueprint for \emph{performance}. Keep that truth in mind throughout this section.

\courssub{Context: Soyinka, the Village and the Play's Social Vision}

Wole Soyinka, the Nigerian playwright, poet and essayist, was awarded the Nobel Prize in Literature in 1986 — the first African writer to receive it. \emph{The Swamp Dwellers} is one of his early plays, written in 1958, at the moment when Nigeria was approaching independence (1960) and debating what kind of society it wanted to become. The play is a compact, one-act tragedy set in a village in the Niger Delta swamps, a region of creeks, seasonal floods and waterlogged farmland.

\begin{experience}[Entering the world of the play]
Before reading any analysis, picture the scene: a hut on stilts in the swamps; an old couple, Makuri and Alu, waiting for the flood season to pass; their son Igwezu gone to the city to seek his fortune; the land itself hostile, drowning the crops year after year. Into this waiting world comes a blind Beggar from the parched North, looking for a lost land of promise. Everything in the play grows out of this single situation: \emph{people waiting for a salvation that may never come}.
\end{experience}

Soyinka's social vision is double-edged. On one side, he exposes the \textbf{corruption of traditional and religious authority}: the Kadiye, the priest of the Serpent of the Swamp, grows fat on the sacrifices of poor villagers while their fields drown. On the other side, he refuses to romanticise \textbf{modernity}: the city to which the young men drift — where Igwezu's twin brother Awuchike has already disappeared — is shown as a place of broken promises, where Igwezu loses his wife (Desala) and his money. Neither the old world nor the new one offers easy hope. Hope, if it exists at all, lies in \emph{honest work on the land}, symbolised by the Beggar's determination to drain and farm the swamp.

\begin{savaistu}[Why this matters to a Cameroonian candidate]
The questions Soyinka raises are immediately recognisable in Cameroon: the drift of young people from villages to Douala, Yaoundé or abroad; the tension between traditional authority and modern life; the exploitation of the rural poor by those who claim spiritual power. In the examination, a brief, well-chosen parallel with your own society can strengthen an essay — but it must \emph{support} textual analysis, never replace it.
\end{savaistu}

\courssub{Character Introduction and the Character–Function Map}

\begin{propriete}[Character in drama is revealed through dialogue and action]
In a novel, a narrator can tell us directly what a character thinks and feels. In a play, there is (normally) no narrator: character is revealed almost entirely through \textbf{what the character says}, \textbf{what others say about him or her}, and \textbf{what the character does on stage} — including silences, gestures and movements indicated in stage directions. Therefore, every claim you make about a character in an essay must be anchored in a \emph{quotation of dialogue} or a \emph{stage direction}.
\end{propriete}

The principal characters, and the thematic function each one embodies, are mapped below.

\begin{center}
\begin{tabularx}{\linewidth}{|l|X|X|}
\hline
\rowcolor{popblueL} \textbf{Character} & \textbf{Role in the plot} & \textbf{Theme embodied} \\
\hline
The Old Man & Father figure of the village; interrogates the Beggar; embodies communal memory & Tradition, communal justice, the weight of the past \\
\hline
Alu & Mother of Igwezu and Awuchike; clings to hope that her lost son lives & Maternal faith, endurance, the cost of waiting \\
\hline
Makuri & Husband of Alu; sharp-tongued, suspicious, cynical about the Kadiye & Scepticism towards religious authority; rural poverty \\
\hline
Igwezu & Son returned ruined from the city; confronts the Kadiye; chooses the land & The prodigal son; integrity; honest labour versus corruption \\
\hline
The Kadiye & Fat priest of the Serpent; lives off sacrifices; never speaks on stage & Corruption of religious authority; exploitation \\
\hline
The Beggar & Blind stranger from the drought-stricken North; asks to farm the swamp & Faith in work; hope; the land as redemption \\
\hline
Awuchike & Igwezu's twin; never appears; lost to city/cult & Vanished youth; the cost of urban drift \\
\hline
Desala & Igwezu's wife; mentioned only; leaves him in the city & Betrayal; the city's destructive power \\
\hline
\end{tabularx}
\end{center}

Notice two structural choices by Soyinka. First, \textbf{Awuchike}, the lost twin, never appears: he exists only in the speeches of others, which makes him a pure symbol — the village's vanished youth. Second, the \textbf{Kadiye never speaks}: his silent, bloated presence is more damning than any speech, a masterstroke of stagecraft. Third, \textbf{Desala} is absent, her betrayal reported only through Igwezu's bitter account.

\courssub{Themes: Tradition versus Modernity, Corruption, Rural Poverty and Urban Drift}

\textbf{Tradition versus modernity.} The play refuses a simple opposition. Tradition is represented both by the Old Man's dignity and communal values \emph{and} by the Kadiye's parasitic cult; modernity is represented both by the city's promise of wealth \emph{and} by its betrayal of Igwezu. Soyinka's point is that neither world is automatically good: what matters is whether a practice — old or new — serves human life.

\textbf{Corruption of religious authority.} The Serpent cult demands sacrifices — money, produce, even (it is implied) human victims like Awuchike — while giving nothing back. Igwezu's confrontation with the Kadiye, in which he seizes the priest and exposes his obesity as proof of theft, is the dramatic climax of the play: the sacred is unmasked as a business.

\textbf{Rural poverty and urban drift.} The floods destroy the villagers' crops; the city destroys the villagers' sons. The Beggar's arrival closes the circle: a man fleeing drought in the North chooses to stay and drain the swamp, suggesting that the answer lies neither in flight nor in fatalism, but in \emph{transforming the land itself}.

\begin{definition}[Dramatic irony]
\cle{Dramatic irony} occurs when the audience (or reader) knows more than a character on stage, so that the character's words or hopes acquire a second, poignant meaning. In \emph{The Swamp Dwellers}, Alu's persistent hope that Awuchike is alive and prospering in the city is deeply ironic: the audience, piecing together the hints about the Serpent cult's victims, suspects what she refuses to know. Her hope, spoken aloud, becomes a measure of her tragedy.
\end{definition}

\begin{definition}[Symbol]
A \cle{symbol} in drama is an object, action, or character that carries a meaning beyond its literal existence. The \textbf{swamp} is the play's central symbol: it is at once a real, hostile environment (it floods, it drowns crops), a symbol of stagnation (a society stuck between a corrupt past and a false future), and — through the Beggar's plan to drain it — a symbol of what human effort could redeem. The \textbf{Serpent}, by contrast, symbolises tradition turned predatory: a god who eats his worshippers.
\end{definition}

\courssub{Language, Style and Dramatic Techniques}

Soyinka's dialogue moves between registers: the proverb-rich, earthy speech of Makuri and Alu; the restrained dignity of the Old Man; the bitter precision of Igwezu; the formal, almost incantatory language of the ritual moments. This variety is itself a theme: each way of speaking encodes a way of seeing the world.

\begin{center}
\begin{tabularx}{\linewidth}{|l|X|X|}
\hline
\rowcolor{popblueL} \textbf{Technique} & \textbf{Example from the play} & \textbf{Effect / Meaning} \\
\hline
Proverbs / idioms & Makuri: ``The toad does not run in the daytime for nothing.'' & Encodes communal wisdom; signals oral tradition; often ironic. \\
\hline
Ritual language & Drumming, chants, processional movements & Creates spectacle; juxtaposes cult splendour with villagers' misery. \\
\hline
Silence & The Kadiye never speaks; pauses before Igwezu's attack & Silence as power (Kadiye) or tension (Igwezu); visual evidence speaks louder. \\
\hline
Contrast of registers & Alu's prayerful hope vs. Makuri's cynical prose & Embodies the play's central conflict: faith vs. scepticism. \\
\hline
Stage directions as characterisation & ``The Kadiye enters, a huge man, dripping with sweat.'' & Physicality replaces dialogue; the audience \emph{sees} corruption. \\
\hline
Sound & Off-stage drumming, distant cries, the sound of rain & Extends the world beyond the hut; builds atmosphere of threat. \\
\hline
Lighting (implied) & The play moves from day to night; the hut interior vs. swamp exterior & Marks temporal passage; isolates characters; focuses attention. \\
\hline
\end{tabularx}
\end{center}

\textbf{Ritual} punctuates the play: the sacrifice procession, the drumming, the Kadiye's appearances. Ritual on stage is powerful because it is \emph{seen and heard}, not described — and Soyinka uses it ironically, juxtaposing the splendour of the cult with the misery of its worshippers.

\begin{attention}[Common mistake: analysing a play like a novel]
The most frequent error in drama essays is to treat the play as a story and to ignore \textbf{stage directions, setting, entrances, exits, sound and visual effects}. Writing ``Igwezu is angry'' is weak; writing ``Igwezu's anger is staged physically — he seizes the Kadiye, and the audience sees the priest's fat body as visible proof of the exploitation the dialogue has described'' is strong. Always ask: \emph{what would the audience see and hear at this moment?}
\end{attention}

\courssub{Stagecraft: Setting, Entrances, Exits and Tension in a One-Act Structure}

The play is a single act in a single setting — the area around Makuri and Alu's hut in the swamps. This unity of place concentrates tension: nothing can be escaped; everyone must come to this one spot. Entrances are used as dramatic events: the Beggar's arrival interrupts a domestic quarrel; the Kadiye's procession interrupts Igwezu's return; each entrance redirects the conflict.

\begin{popfigure}
\begin{tikzpicture}[scale=0.95, every node/.style={font=\small}]
\draw[thick, popink] (0,0) rectangle (10,6);
\node[popmuted] at (5,6.35) {THE STAGE (single setting: the swamp village)};
% Hut on stilts
\draw[thick, popdark] (0.6,2.6) rectangle (3.2,4.6);
\draw[thick, popdark] (0.4,4.6) -- (1.9,5.5) -- (3.4,4.6);
\draw[popdark] (1.0,2.6) -- (1.0,1.9);
\draw[popdark] (2.8,2.6) -- (2.8,1.9);
\node[align=center] at (1.9,3.6) {Makuri and\\Alu's hut};
\node[popmuted, font=\footnotesize] at (1.9,1.6) {on stilts above the water};
% Swamp water
\fill[popblueL] (0,0) rectangle (10,1.2);
\draw[popblue, thick] (0,1.2) -- (10,1.2);
\node[popblue] at (5,0.6) {the swamp — water, mud, flooded fields};
% Path / entrance left
\draw[thick, popgreen, ->] (-0.8,3.4) -- (0.5,3.4);
\node[popgreen, align=center, font=\footnotesize] at (-0.2,4.3) {Beggar's\\entrance (North)};
% Entrance right
\draw[thick, poporange, ->] (10.8,3.4) -- (9.5,3.4);
\node[poporange, align=center, font=\footnotesize] at (10.2,4.3) {Kadiye's\\procession (South)};
% Playing area
\draw[dashed, poppurple] (4.2,1.6) rectangle (8.8,4.8);
\node[poppurple, align=center] at (6.5,3.2) {central playing area:\\waiting, quarrels,\\confrontation};
% Audience
\node[popmuted] at (5,-0.6) {AUDIENCE};
\end{tikzpicture}
\legende{Sketch of the stage layout: one enclosed setting, two opposed entrances (the Beggar from the North, the Kadiye's procession from the South), and the hut on stilts above the ever-present swamp.}
\end{popfigure}

\begin{methode}[Writing a drama essay: always return to the stage]
\begin{enumerate}
\item \textbf{Read the question twice} and identify whether it targets character, theme or technique.
\item \textbf{Select three or four key moments} of the play relevant to the question — moments, not vague summaries.
\item For each moment, \textbf{quote a short piece of dialogue or a stage direction}.
\item \textbf{Comment on performance}: what does the audience see, hear, feel? How does the staging create meaning or tension?
\item \textbf{Link the moment to the theme} and to Soyinka's wider social vision.
\item Conclude by evaluating the effect of the whole: what does the one-act structure, the single setting, or the final image leave with the audience?
\end{enumerate}
\end{methode}

\begin{exoresolu}[Worked essay extract: How does Soyinka expose the corruption of religious authority?]
\textbf{Question:} ``In \emph{The Swamp Dwellers}, Soyinka presents religious authority as corrupt and exploitative.'' Discuss with close reference to the play.
\tcblower
\textbf{Plan and model development.} Introduction: define the cult of the Serpent and its priest, the Kadiye; state the thesis — Soyinka exposes corruption through \emph{dialogue, visual stagecraft and irony}.

Paragraph 1 (dialogue): Makuri's cynical remarks about the sacrifices establish, early, that the villagers themselves perceive the cult as a business; the audience hears the grievance before it sees the priest. \emph{Quotation}: ``We feed the Serpent and the Serpent feeds on us.'' This metaphor turns the religious relationship into an economic one.

Paragraph 2 (stagecraft): when the Kadiye finally appears, he is fat and silent. On stage, his body is the evidence: in a village ruined by floods, only the priest is well fed. His silence is crucial — he never justifies himself because the cult never needs to justify itself to the powerless. \emph{Stage direction}: ``The Kadiye enters, a huge man, dripping with sweat.'' The visual contrast between his bulk and the villagers' emaciation is immediate proof.

Paragraph 3 (climax and irony): Igwezu, ruined by the city, returns to find that the village god has taken his offerings and given nothing. His physical seizure of the Kadiye turns accusation into action; the audience watches the sacred manhandled, and the ritual dignity collapses. The deeper irony is that the villagers' piety has financed their own oppression. \emph{Quotation}: Igwezu: ``You are the Serpent... You eat the offerings and you eat the people.''

Paragraph 4 (the Beggar as counterpoint): The Beggar's blind faith in work — not in the Serpent — offers an alternative spirituality. His refusal to fear the cult (``I have no eyes to see your god'') undermines the priest's authority without violence.

Conclusion: Soyinka does not attack faith itself — the Beggar's quiet faith in work is treated with respect — but faith captured by greedy intermediaries. The one-act structure leaves the exposure unresolved: the Kadiye survives, which is Soyinka's bleakest point.
\end{exoresolu}

\begin{exoresolu}[Worked essay extract: The dramatic function of the Beggar]
\textbf{Question:} ``The Beggar is the most important character in the play.'' How far do you agree?
\tcblower
\textbf{Plan and model development.} Introduction: acknowledge the Beggar's late entrance but argue his structural and symbolic centrality.

Paragraph 1 (structural pivot): The Beggar enters at the midpoint, disrupting the domestic stasis. His request — ``Let me farm this swamp'' — reframes the play's central problem: the land is not cursed, it is neglected.

Paragraph 2 (symbolic reversal): He is blind, yet he ``sees'' the land's potential. The sighted villagers (Makuri, Alu, even Igwezu) have accepted defeat. The irony: physical blindness vs. moral vision. \emph{Quotation}: ``I have hands. I have a back. The land is there.''

Paragraph 3 (thematic synthesis): He embodies the synthesis the play searches for — neither blind tradition (Kadiye) nor blind modernity (city), but \emph{labour}. He connects North (drought) and South (flood), suggesting a pan-African solidarity of the poor.

Paragraph 4 (dramatic contrast): His quiet dignity contrasts with the Kadiye's bloated silence and Igwezu's bitter eloquence. He speaks plainly, without proverbs or rhetoric.

Conclusion: The Beggar does not ``solve'' the play's problems — the floods will return, the Kadiye remains — but he offers the only viable \emph{attitude}. In a GCE essay, evaluate ``importance'' through structural placement, symbolic weight, and thematic resolution.
\end{exoresolu}

\begin{exercice}[Practice: character and technique]
\begin{enumerate}
\item ``Alu and Makuri represent two different responses to suffering.'' Discuss, quoting dialogue and referring to at least one stage direction.
\item Explain the symbolic function of the swamp and show how the Beggar transforms its meaning by the end of the play.
\item Why does Soyinka keep Awuchike off stage? Write one analytical paragraph on this technique.
\item Analyse the dramatic significance of the Kadiye's procession. How do sound, movement and visual spectacle create irony?
\item ``The play's ending offers no hope, only endurance.'' Evaluate this view with reference to the final moments of the play.
\end{enumerate}
\end{exercice}

\begin{corrige}[Exercise — guidance]
\textbf{1.} Alu responds with faith and endurance: she waits, hopes and refuses to believe Awuchike is dead; quote her hopeful speeches (e.g., ``My son will come back. He will bring money. He will build us a house on dry land'') and note the dramatic irony. Makuri responds with suspicion and bitter humour: quote his cynical exchanges about the Kadiye and the flood (e.g., ``The Serpent has a big belly. We fill it.''). Anchor each point in a quoted line, then comment on how the audience perceives the contrast when the two share the stage — the quarrels are staged as a duet of hope and scepticism.

\textbf{2.} Initially the swamp symbolises stagnation and hostility: it drowns crops and traps the villagers. The Beggar, a blind stranger, chooses to stay and drain it; the swamp thus becomes a symbol of redeemable land and of hope through labour. Note the irony that a blind man ``sees'' the land's potential better than those who have lived on it all their lives. The final image — the Beggar staying, the hut remaining — suggests endurance is active, not passive.

\textbf{3.} Keeping Awuchike off stage makes him a pure symbol — the village's vanished youth — and sustains dramatic irony: the audience suspects his fate while Alu hopes. It also universalises him: he could be any young man swallowed by the city or by the cult. On stage, an absent character is created entirely through other people's speeches, which shows how drama builds meaning through dialogue alone.

\textbf{4.} The Kadiye's procession (drumming, chanting, attendants, the priest carried or walking slowly) creates a spectacle of power. The irony: the villagers watch in awe, yet the audience sees the cost — their sacrifices feed this display. Sound (drums) drowns dissent; movement (procession) claims space; visual spectacle (robes, bulk) asserts authority. When Igwezu later disrupts it, the contrast exposes the fragility of that authority.

\textbf{5.} The ending: the Kadiye leaves unharmed; the floods will return; Igwezu stays but is broken; the Beggar stays but is one man. Yet ``endurance'' is not passive: Igwezu rejects the city; the Beggar chooses the swamp; Alu's hope, however deluded, keeps her alive. A strong essay balances bleak realism (the Kadiye survives) with the play's affirmation of labour and choice.
\end{corrige}

\begin{aretenir}[Key points of Section 12]
\begin{itemize}
\item \emph{The Swamp Dwellers} is a one-act play of waiting, set in a flooded Niger Delta village, exposing both corrupt tradition and false modernity.
\item Character in drama is revealed through \cle{dialogue}, \cle{action} and \cle{stage directions} — quote them, and describe what the audience sees.
\item The Kadiye's silent fatness, Awuchike's absence and Alu's ironic hope are Soyinka's finest dramatic devices.
\item The swamp is the central symbol: stagnation that honest labour — the Beggar's choice — might redeem.
\item In every drama essay, follow the method: key moment $\rightarrow$ quotation $\rightarrow$ performance comment $\rightarrow$ theme.
\item Register variation (proverbs, ritual language, silence, plain speech) encodes the play's ideological conflicts.
\item The one-act, single-setting structure creates pressure-cooker tension; entrances are dramatic events.
\end{itemize}
\end{aretenir}

\begin{experience}[GCE Examination Practice: Past Paper Style Questions]
\textbf{Paper 2 (Essay) typical questions:}
\begin{enumerate}
\item ``Soyinka uses the character of the Kadiye to criticise religious hypocrisy.'' Discuss with close reference to the play.
\item How does the playwright use the contrast between the city and the village to explore the theme of betrayal?
\item ``The Swamp Dwellers is a play about waiting.'' How far do you agree?
\item Analyse the dramatic significance of the Beggar's blindness.
\item In what ways does Soyinka use stagecraft (setting, entrances, sound, silence) to create tension?
\end{enumerate}

\textbf{Paper 1 (Comprehension) typical tasks:}
\begin{itemize}
\item An extract from the play (300–400 words) followed by questions on: meaning of words in context; dramatic effect of a stage direction; character revelation through dialogue; irony; thematic significance.
\item A comparison task: an extract from \emph{The Swamp Dwellers} paired with a non-literary text (e.g., a newspaper article on rural-urban migration) testing synthesis and evaluation.
\end{itemize}

\textbf{Paper 3 (Oral) typical tasks:}
\begin{itemize}
\item Reading aloud a passage with attention to tone, pause, and register shift (e.g., Makuri's proverbs vs. Igwezu's bitter prose).
\item Discussion: ``How does the play reflect tensions in your own society?'' — require textual anchoring.
\end{itemize}
\end{experience}

\begin{savaistu}[Examiner's tip: the "So what?" test]
For every point you make in a drama essay, ask: \emph{So what? What does the audience experience?} If you write ``Makuri uses proverbs,'' add: ``This roots his scepticism in communal wisdom, making his critique of the Kadiye feel like the voice of the village itself, not just one man's grievance.'' That extra sentence moves your answer from description to analysis — and from a C to an A.
\end{savaistu}

\coursec{Section 13: Poetry Study — Cosmic Anthology and Selected Poems}

In Section 12 we studied Wole Soyinka's \emph{The Swamp Dwellers}, a play in which language, symbol and dramatic tension carry the clash between tradition and modernity. We now turn from drama to the most concentrated of all literary forms: \cle{poetry}. Where a playwright has two hours of stage time, a poet may have only fourteen lines. Every word, every line break, every sound must therefore work. This section introduces the \emph{Cosmic Anthology} and trains you to analyse two prescribed poems — Claude McKay's \emph{If We Must Die} and Roland Tombekai Dempster's \emph{Africa's Plea} — and to write a full poetry appreciation essay in GCE format.

\courssub{What Makes Poetry Different: Forms and Devices}

Read these two sentences aloud:

\begin{exemplebox}[Prose versus verse]
\emph{Prose:} The rain fell heavily on the zinc roofs of Buea all night and nobody slept.\\[2pt]
\emph{Verse:}\\
All night the rain interrogated the zinc,\\
and Buea, sleepless, answered in drips.
\end{exemplebox}

The facts are identical, but the verse version does more: it \emph{personifies} the rain, it breaks the thought into \cle{lines}, and it creates a music that prose does not attempt. Poetry is language organised for maximum \textbf{density of meaning and sound}. To analyse it, you need two toolkits: \textbf{form} (the container) and \textbf{devices} (the machinery inside).

\textbf{The major poetic forms}

\begin{center}
\begin{tabularx}{\linewidth}{|l|X|X|}
\hline
\rowcolor{popblueL} \textbf{Form} & \textbf{Defining features} & \textbf{Typical effect} \\
\hline
Sonnet & 14 lines, usually iambic pentameter; Shakespearean rhyme \emph{abab cdcd efef gg}; a ``turn'' (\cle{volta}) near line 9 or in the final couplet & Controlled, disciplined argument or meditation \\
\hline
Lyric & Short poem expressing personal emotion; no fixed shape & Intimacy, immediacy of feeling \\
\hline
Ballad & Quatrains, often \emph{abcb}; strong rhythm; tells a story, often tragic or heroic & Song-like narrative; oral tradition \\
\hline
Ode & Elevated address to a person, thing or idea; serious and ceremonious & Praise, celebration, public feeling \\
\hline
Elegy & Poem of mourning for the dead or for something lost & Lament, consolation, reflection \\
\hline
Free verse & No regular metre or rhyme; line breaks and rhythm do the organising & Freedom, natural speech, modernity \\
\hline
\end{tabularx}
\end{center}

\textbf{The core devices}

\begin{definition}[Enjambment]
\cle{Enjambment} occurs when a sentence or clause runs over the end of a line into the next line without a pause. It creates pace, suspense or emphasis, because the reader is pulled forward to complete the thought. The opposite is an \emph{end-stopped} line, where sense and line end together, producing calm or finality.
\end{definition}

Other devices you must command: \cle{imagery} (language appealing to the senses), \cle{alliteration} (repetition of initial consonant sounds), \cle{assonance} (repetition of vowel sounds), \cle{rhyme scheme} (the pattern of end-rhymes, labelled \emph{abab} etc.), \cle{repetition} and \cle{anaphora} (repetition at the start of lines), and the figures of speech from Section 5 — metaphor, simile, personification, irony, hyperbole, synecdoche, metonymy.

\begin{methode}[The poetry analysis frame]
Approach every unseen or prescribed poem in this fixed order:
\begin{enumerate}
\item \textbf{Subject:} what is the poem literally about? (one sentence)
\item \textbf{Theme:} what larger idea does the subject carry? (dignity, colonialism, love, death\ldots)
\item \textbf{Tone:} the poet's attitude — defiant, pleading, mournful, ironic? Find the words that prove it.
\item \textbf{Devices:} identify form, imagery, sound devices, figures of speech.
\item \textbf{Effect:} for \emph{each} device, explain what it does to meaning and feeling.
\item \textbf{Personal response:} a brief, honest evaluation — does the poem succeed, and why?
\end{enumerate}
\end{methode}

\begin{attention}[The ``device-spotting'' trap]
The most common examination error is \textbf{listing devices without explaining their effect}. Writing ``the poet uses alliteration in line 5'' earns almost nothing. Writing ``the harsh alliteration of \emph{k} sounds in `killing' and `kinsmen' (line 5) mimics the brutality of the violence described, making the reader almost hear the blows'' earns full credit. Every device named must be followed by \emph{because\ldots} or \emph{which\ldots}.
\end{attention}

\courssub{Analysing \emph{If We Must Die} by Claude McKay}

\textbf{Context.} Claude McKay (1889--1948) was a Jamaican-born poet of the \cle{Harlem Renaissance}, the flowering of Black American art and literature centred in Harlem, New York, in the 1920s. \emph{If We Must Die} (1919) was written in response to the wave of racial violence against Black communities in the United States known as the ``Red Summer'' of 1919. Crucially, the poem never names race: its address to an oppressed band of ``kinsmen'' facing a murderous ``foe'' gives it a universal reach, which is why it has been quoted in struggles far beyond its original moment.

\textbf{The poem (quoted in full for annotation).}
\begin{quote}
If we must die, let it not be like hogs\\
Hunted and penned in an inglorious spot,\\
While round us bark the mad and hungry dogs,\\
Making their mock at our accursed lot.\\
If we must die, O let us nobly die,\\
So that our precious blood may not be shed\\
In vain; then even the monsters we defy\\
Shall be constrained to honour us though dead!\\
O kinsmen! we must meet the common foe!\\
Though far outnumbered let us show us brave,\\
And for their thousand blows deal one death-blow!\\
What though before us lies the open grave?\\
Like men we'll face the murderous, cowardly pack,\\
Pressed to the wall, dying, but fighting back!
\end{quote}

\begin{popfigure}
\begin{center}
\begin{tikzpicture}[
  box/.style={draw=popblue, rounded corners=2pt, fill=popblueL, text width=4.6cm, align=left, font=\small, inner sep=4pt},
  ln/.style={font=\small, text=popink, align=left}
]
\node[ln, anchor=west] (l1) at (0,0) {If we must die, let it not be like hogs};
\node[ln, anchor=west] (l2) at (0,-0.55) {Hunted and penned in an inglorious spot,};
\node[ln, anchor=west] (l5) at (0,-1.6) {If we must die, O let us nobly die,};
\node[ln, anchor=west] (l9) at (0,-2.65) {O kinsmen! we must meet the common foe!};
\node[ln, anchor=west] (l13) at (0,-3.7) {Like men we'll face the murderous, cowardly pack,};
\node[ln, anchor=west] (l14) at (0,-4.25) {Pressed to the wall, dying, but fighting back!};

\node[box] (a) at (9.5,0.1) {\textbf{Simile:} ``like hogs'' — death without dignity; the animal imagery degrades the victims};
\draw[->, popblue] (a.west) -- (l1.east);

\node[box] (b) at (9.5,-1.5) {\textbf{Volta} (line 5): the poem turns from \emph{how not} to die to \emph{how} to die — ``nobly''};
\draw[->, popblue] (b.west) -- (l5.east);

\node[box] (c) at (9.5,-3.0) {\textbf{Direct address} ``O kinsmen!'': unites the oppressed; rhetorical rallying cry};
\draw[->, popblue] (c.west) -- (l9.east);

\node[box] (d) at (9.5,-4.4) {\textbf{Final couplet:} rhyme ``pack/back'' slams shut like a blow; ``fighting back'' = defiance to the end};
\draw[->, popblue] (d.west) -- (l14.east);
\end{tikzpicture}
\end{center}
\legende{Annotated skeleton of \emph{If We Must Die}: form, devices and the volta.}
\end{popfigure}

\textbf{Form.} The poem is a \cle{Shakespearean sonnet}: fourteen lines of iambic pentameter rhyming \emph{abab cdcd efef gg}. This is a deliberate, meaningful choice. The sonnet is the most disciplined form in English — and McKay fills it with a message of \emph{disciplined} resistance. The tight form \textbf{contains} the rage, transforming raw anger into controlled defiance. The \cle{volta} comes early, at line 5 (``O let us nobly die''), shifting the poem from the horror of an ``inglorious'' death to the programme of a noble one; a second surge of energy arrives at line 9 (``O kinsmen!''), where the meditation becomes a direct call to arms.

\textbf{Themes.} The central theme is \textbf{dignity in the face of death}: if death is unavoidable, its \emph{manner} remains a choice. Linked themes are \textbf{resistance} (``fighting back!''), \textbf{solidarity} (``kinsmen'', ``we'' — the poem never says ``I''), and \textbf{honour} that even the enemy must concede (``constrained to honour us though dead'').

\textbf{Devices and effects.} The \emph{animal imagery} divides the world morally: the victims are threatened with becoming ``hogs'', while the oppressors are ``mad and hungry dogs'', ``monsters'', a ``cowardly pack'' — the poet reverses the expected hierarchy, making the powerful the true beasts. The \emph{rhetorical question} ``What though before us lies the open grave?'' dismisses death itself as a weak argument. The \emph{contrast} ``their thousand blows'' / ``one death-blow'' turns numerical weakness into moral strength. Note also the \emph{conditional frame} ``If we must die\ldots'' (compare Section 4 on conditionals): the poem concedes the worst hypothesis only to master it.

\courssub{Analysing \emph{Africa's Plea} by Roland Tombekai Dempster}

Roland Tombekai Dempster was a Liberian poet and diplomat, one of the pioneers of modern Liberian literature in English. In \emph{Africa's Plea}, the continent itself speaks.

\textbf{The poem's strategy.} The governing device is \cle{personification}: Africa is a suffering, pleading human being — a mother or supplicant — who addresses her children and the world. This single choice shapes everything else:

\begin{itemize}
\item \textbf{Tone:} pleading and wounded, yet dignified; not accusation alone but an \emph{appeal} — the tone of one who still believes her children can change.
\item \textbf{Apostrophe and direct address:} the speaker calls out to those who have wronged or abandoned her, making the poem feel like a speech or prayer.
\item \textbf{Implicit meaning and irony:} the continent that gave humanity so much must \emph{beg} for recognition — the gap between Africa's worth and her treatment is the poem's quiet irony (compare Section 5 on irony and understatement).
\item \textbf{Connotation:} words of motherhood, wounds and supplication carry emotional weight far beyond their dictionary meanings (link back to Section 5: denotation versus connotation).
\end{itemize}

\textbf{Message.} The poem asks Africa's children — including the diaspora — to remember, honour and rebuild the mother continent. Where McKay \emph{commands} resistance, Dempster \emph{appeals} to conscience and love. Both poems are calls to action, but one is a battle cry and the other a plea.

\courssub{Comparative Poetry Analysis}

GCE questions frequently ask you to compare two poems. The method is to compare \emph{category by category}, never poem by poem in two separate blocks.

\begin{popfigure}
\begin{center}
\begin{tabularx}{\linewidth}{|l|X|X|}
\hline
\rowcolor{popblueL} \textbf{Criterion} & \textbf{\emph{If We Must Die} (McKay)} & \textbf{\emph{Africa's Plea} (Dempster)} \\
\hline
Theme & Dignity and honourable resistance against violent oppression & The wounded mother continent appealing to her children \\
\hline
Tone & Defiant, rallying, heroic & Pleading, mournful, yet hopeful \\
\hline
Form & Shakespearean sonnet; \emph{abab cdcd efef gg}; volta at line 5 & Lyric built on sustained personification; looser, speech-like structure \\
\hline
Key devices & Simile (``like hogs''), animal imagery, rhetorical question, contrast & Personification, apostrophe, direct address, connotative diction \\
\hline
Speaker & One of the oppressed (``we'', ``kinsmen'') & Africa herself, personified \\
\hline
Effect on reader & Rouses courage and solidarity & Awakens pity, guilt and filial duty \\
\hline
\end{tabularx}
\end{center}
\legende{Comparison grid: the two prescribed poems across six criteria.}
\end{popfigure}

\begin{exoresolu}[Comparing tone in the two poems]
\textbf{Question:} Compare the tone of McKay's \emph{If We Must Die} with that of Dempster's \emph{Africa's Plea}, showing how each poet's choice of devices creates that tone. (10 marks)
\tcblower
\textbf{Model answer (abridged).} Both poems address oppression, but their tones differ sharply. McKay's tone is \emph{defiant and heroic}. The imperative ``let us nobly die'' and the rallying cry ``O kinsmen!'' speak as a commander to soldiers, while the final couplet — ``Pressed to the wall, dying, but fighting back!'' — uses the hard rhyme ``pack/back'' to end the poem like a struck blow. The strict sonnet form itself disciplines the anger into resolve. Dempster's tone, by contrast, is \emph{pleading and wounded}. By personifying Africa as a suffering mother who must beg her own children, the poet replaces command with appeal; the reader is moved not to fight but to feel pity and responsibility. Thus McKay stirs courage through martial devices and a closed, forceful form, while Dempster stirs conscience through personification and supplication. Both, however, share a deeper similarity: each refuses despair, and each ends by affirming the dignity of the oppressed.
\end{exoresolu}

\begin{attention}[Exam tip: always comment on form]
Whatever the question, find a way to discuss \textbf{form}. For McKay, ask: \emph{why a sonnet?} Because the most controlled form in English perfectly embodies \emph{controlled} defiance — rage held in check, channelled into resolve. For Dempster, ask: \emph{why personification sustained across the whole poem?} Because a continent cannot literally speak; giving it a human voice converts geography into family, and politics into a debt of love. A paragraph on form distinguishes a good script from an excellent one.
\end{attention}

\courssub{Writing the Poetry Appreciation Essay}

A full GCE poetry appreciation essay (about 450--600 words) follows this architecture:

\begin{methode}[Structuring the appreciation essay]
\begin{enumerate}
\item \textbf{Introduction:} name the poem and poet; state the subject and the central theme in one or two sentences; give a one-line context if relevant (e.g. the Harlem Renaissance for McKay).
\item \textbf{Body paragraph 1 — theme and tone:} develop the theme(s) and establish the tone with two or three quoted words as evidence.
\item \textbf{Body paragraph 2 — form and structure:} sonnet or free verse? rhyme scheme? volta? enjambment? Explain \emph{why} the form suits the message.
\item \textbf{Body paragraph 3 — devices and effects:} discuss three or four devices, each with quotation plus effect (never a bare list).
\item \textbf{Conclusion:} your personal response — evaluate the poem's success and, where natural, its relevance (e.g. what McKay's dignity or Dempster's plea says to a young Cameroonian today).
\end{enumerate}
\end{methode}

\begin{exemplebox}[A model opening paragraph]
In \emph{If We Must Die}, the Harlem Renaissance poet Claude McKay transforms the threat of racial violence into a manifesto of human dignity. Written in 1919 against the backdrop of brutal attacks on Black communities in America, the sonnet argues that while death may be forced upon the oppressed, dishonour is not: ``If we must die, O let us nobly die.'' Through disciplined form, savage animal imagery and a rousing collective voice, McKay turns a poem about dying into a poem about how to live.
\end{exemplebox}

\begin{exercice}[Poetry appreciation practice]
\begin{enumerate}
\item Identify the rhyme scheme of McKay's poem from the extract quoted above, and state where the volta occurs.
\item Explain the effect of the simile ``like hogs'' in line 1.
\item Write the introduction and one body paragraph (form and structure) of an appreciation essay on \emph{Africa's Plea}.
\item In one paragraph, compare how McKay and Dempster each use \emph{direct address} (``O kinsmen!'' / Africa's appeal to her children).
\end{enumerate}
\end{exercice}

\begin{corrige}[Exercise 1]
\begin{enumerate}
\item The rhyme scheme is \emph{abab cdcd efef gg} (hogs/spot/dogs/lot; die/shed/defy/dead; foe/brave/death-blow/grave; pack/back). The volta occurs at line 5, ``O let us nobly die'', where the poem turns from rejecting an inglorious death to prescribing a noble one.
\item The simile compares the victims' potential death to that of slaughtered pigs: hunted, penned and helpless. It establishes the poem's central stakes — the issue is not death itself but \emph{degradation} — and so justifies the volta's demand for a ``noble'' death. It also begins the animal imagery that morally reverses the poem's world: the oppressors, not the victims, become the true beasts (``mad and hungry dogs'', ``cowardly pack'').
\item \emph{Introduction:} should name Dempster, identify the personified Africa as speaker, state the theme (the continent's appeal to her children) and the tone (pleading yet dignified). \emph{Body paragraph on form:} should note the lyric, speech-like structure, the sustained personification as the poem's organising principle, and the effect — turning a political subject into an intimate family appeal.
\item Both poets use direct address to create solidarity, but with opposite emotional temperatures. McKay's ``O kinsmen!'' is a battle cry: it binds the oppressed into an army and urges action (``we must meet the common foe''). Dempster's Africa addresses her children as a mother: the address binds through love and guilt rather than courage, and urges remembrance rather than combat. In both cases the device collapses the distance between speaker and audience, making the reader a member of the group addressed.
\end{enumerate}
\end{corrige}

\begin{aretenir}[Key points of Section 13]
\begin{itemize}
\item Poetry analysis follows a fixed frame: \textbf{subject $\to$ theme $\to$ tone $\to$ devices $\to$ effect $\to$ personal response}.
\item Know the forms: \textbf{sonnet, lyric, ballad, ode, elegy, free verse}; know the devices: \textbf{imagery, alliteration, enjambment, rhyme scheme, personification}.
\item \emph{If We Must Die}: Harlem Renaissance sonnet; volta at line 5; themes of dignity, solidarity and resistance; controlled form mirrors controlled defiance.
\item \emph{Africa's Plea}: sustained personification of Africa as a pleading mother; tone of wounded appeal; irony in a rich continent reduced to begging.
\item Compare poems \textbf{criterion by criterion}; never list devices without explaining their \textbf{effect}; always comment on \textbf{form}.
\end{itemize}
\end{aretenir}

With the tools of poetry analysis now in hand, we move to Section 14, the Mid-Year Review, where the comprehension, summary, grammar, phonetics and writing skills of the whole term — including this poetry work — are consolidated against the three GCE papers.

\coursec{Section 14: Mid-Year Review and Examination Consolidation}

In Section 13 we closed our study of poetry, learning how form, device and nuance combine in poems such as ``If We Must Die'' and ``Africa's Plea''. That was the last \emph{new} content of the first half of the year. We now pause, look backwards, and ask the most important question a candidate can ask: \emph{How far have I actually come, and what exactly must I fix before June?} This section turns you from a student who merely works hard into a candidate who works \emph{strategically}. We shall review each GCE paper against your diagnostic baseline, run an integrated timed practice, build an error logbook, and design a revision plan for the remaining terms.

\courssub{Why a Mid-Year Review? From Effort to Evidence}

\begin{experience}[The two candidates]
Consider two Upper Sixth students in Bamenda. Acho studies every evening, re-reading her favourite notes on figures of speech. Beri studies less often, but every fortnight he re-sits a past-paper section under timing, marks it with the official-style mark scheme, records every error in a notebook, and re-tests himself on exactly those weaknesses. In June, Beri outperforms Acho --- not because he is more intelligent, but because his revision is \cle{evidence-driven}. The mid-year review is the moment you become Beri.
\end{experience}

\begin{definition}[Diagnostic baseline]
A \cle{diagnostic baseline} is the profile of scores you obtained in the diagnostic assessment at the start of the year (Section 1), broken down by skill: essay, comprehension, summary, grammar, transcription, literature. It is the reference point against which all later progress is measured.
\end{definition}

\begin{propriete}[Comparing like with like]
Measurable progress requires comparing like with like: a mid-year score is only meaningful if it was obtained under the \emph{same conditions} (same paper type, same timing, same mark scheme) as the diagnostic score. A summary marked generously by a friend cannot be compared with a summary marked strictly against the official grid. \emph{(Not examinable as such, but essential for honest self-assessment.)}
\end{propriete}

\begin{attention}[The favourite-topic trap]
The most common revision mistake is revising only what you already enjoy or master. Re-reading your strong topics feels pleasant and produces \emph{no} progress. The mid-year review must deliberately target your \textbf{weakest} skills, because that is where marks are cheapest to gain.
\end{attention}

\courssub{Reviewing Paper 1: Essay and Comprehension}

Paper 1 tests your ability to write (essays, reports, letters) and to read (comprehension, summary). To review it honestly, proceed in three steps.

\begin{methode}[Auditing your Paper 1 progress]
\begin{enumerate}
\item \textbf{Retrieve the evidence.} Collect your diagnostic essay and comprehension script, plus at least two essays and two comprehension exercises written since (argumentative essays from Section 7, summaries from Section 3).
\item \textbf{Re-mark with the same grid.} Score both old and new work against the same criteria: content/relevance, organisation, expression (grammar, vocabulary, mechanics), and --- for comprehension --- accuracy of answers and use of your own words.
\item \textbf{Classify the gap.} For every mark lost, name the cause precisely: weak thesis statement, unbalanced argument, dangling modifier, wrong conditional, misread implicit meaning, summary exceeding the word limit. Vague labels like ``bad English'' are useless; precise labels like ``third conditional with mixed time frame'' are actionable.
\end{enumerate}
\end{methode}

\begin{exemplebox}[A sample Paper 1 audit]
Ngum's diagnostic argumentative essay scored 11/20: good ideas, but the conclusion introduced a new argument and three dangling modifiers cost expression marks. Her mid-year essay scored 15/20: balanced structure achieved, but she still misuses the second conditional (``If I \emph{would} go\ldots''). Conclusion: structure remediated; conditionals become the next targeted item. This is exactly how an audit converts marks into a to-do list.
\end{exemplebox}

\courssub{Reviewing Paper 2: Literature and Language Skills}

Paper 2 rewards \emph{textual knowledge} plus \emph{analytical skill}. A text-knowledge audit checks whether you can still deploy the prescribed texts studied so far: \emph{Twilight of Misty Foliage} (Ben Jama), \emph{Unanswered Cries} (Osman Conteh), \emph{The Old Man and the Sea} (Hemingway), \emph{The Swamp Dwellers} (Soyinka), and the selected poems.

\begin{methode}[The text-knowledge audit]
For each prescribed text, test yourself from memory --- no notes --- on five questions:
\begin{enumerate}
\item Can I state the context (author, setting, period) in two sentences?
\item Can I trace the development of two major characters?
\item Can I name and illustrate three themes with specific episodes or quotations?
\item Can I identify the writer's key techniques (symbolism, irony, narrative point of view, dramatic devices)?
\item Can I write a one-paragraph critical opinion of the text?
\end{enumerate}
Any ``no'' marks a revision slot in your plan. For language skills, re-test the grammar of Sections 4 and 5 (clauses, tenses, conditionals, modifiers, idioms, connotation) using short self-made quizzes.
\end{methode}

\begin{attention}[Plot summary is not analysis]
In Paper 2, retelling the story of Santiago or of the Swamp Dwellers earns very few marks. Every point must be tied to the \emph{question asked} and supported by a precise reference. When auditing, check whether your practice paragraphs \emph{argue} or merely \emph{narrate}.
\end{attention}

\courssub{Reviewing Paper 3 and the Oral Components}

Progress in phonetics is easy to measure but easy to neglect. Revisit Section 6: the complete IPA chart, connected speech (elision, assimilation, intrusion, catenation), and intonation (attitudinal and grammatical functions).

\begin{exemplebox}[A phonetics self-check]
Transcribe: \emph{``Would you like a cup of tea?''} In careful speech: /wUd ju: laIk @ kVp @v ti:/. In natural connected speech, expect assimilation and elision: /wUdZu laIk @ kVp @ ti:/ --- \emph{would you} fuses into /wUdZu/, and the /v/ of \emph{of} is elided. If your transcription ignores connected speech, schedule a phonetics slot. Similarly, record yourself reading a dialogue and check whether your falling and rising tones match the attitude intended (certainty vs.\ polite surprise).
\end{exemplebox}

\courssub{Measuring Progress Visually}

Numbers become motivating when they are displayed. Plot your diagnostic and mid-year scores side by side, skill by skill.

\begin{popfigure}
\begin{center}
\begin{tikzpicture}
\begin{axis}[
    ybar, bar width=0.32cm,
    width=0.95\linewidth, height=6.2cm,
    ymin=0, ymax=20,
    ylabel={Score (out of 20)},
    symbolic x coords={Essay,Compr.,Summary,Grammar,Phonetics,Literature},
    xtick=data,
    x tick label style={font=\small, rotate=15, anchor=east},
    legend style={at={(0.5,1.02)}, anchor=south, legend columns=2, font=\small},
    ymajorgrids=true, grid style={popline, dashed},
]
\addplot[fill=popmuted] coordinates {(Essay,11) (Compr.,9) (Summary,8) (Grammar,10) (Phonetics,7) (Literature,9)};
\addplot[fill=popblue] coordinates {(Essay,15) (Compr.,13) (Summary,12) (Grammar,13) (Phonetics,11) (Literature,14)};
\legend{Diagnostic (Sept.), Mid-year (Jan.)}
\end{axis}
\end{tikzpicture}
\end{center}
\legende{Sample progress chart: diagnostic versus mid-year scores by skill. The weakest bar (here, Phonetics) defines the priority for the next term.}
\end{popfigure}

The chart teaches a lesson at a glance: progress is real but uneven, and the shortest bar --- not the tallest --- dictates the revision plan.

\courssub{Integrated Timed Practice: Simulating the Real Examination}

Reading about timing is not the same as writing against the clock. From this point onward, simulation becomes a habit.

\begin{savaistu}[Exam tip]
Simulate real timing at least once a fortnight from now until the examination: full comprehension passage (900--1000 words) + summary + essay, in one sitting, without a dictionary, phone or pause. Candidates who have never written for three consecutive hours discover hand fatigue and attention drop \emph{on the day itself}; candidates who have trained treat the real paper as just another rehearsal.
\end{savaistu}

\begin{exoresolu}[Worked exercise: planning and executing a timed sequence]
\textbf{Statement.} You have 2 hours 30 minutes for the following sequence: (a) read a 950-word passage on rural electrification in Cameroon and answer five comprehension questions; (b) summarise the passage in 120 words; (c) write a 500-word argumentative essay: \emph{``Rural electrification is the key to development in Cameroon. Discuss.''} Allocate your time and outline the essay.\par
\tcblower
\textbf{Solution.} \emph{Time allocation (like with like, exam conditions):}
\begin{itemize}
\item 5 min: skim the passage; note the thesis and paragraph functions.
\item 35 min: comprehension answers, in your own words, quoting only when asked.
\item 25 min: summary --- select topic sentences, merge, cut examples, count words.
\item 10 min: essay plan (thesis + two arguments for, one against, rebuttal, conclusion).
\item 65 min: essay writing.
\item 10 min: proofreading (conditionals, modifiers, agreement).
\end{itemize}
\emph{Essay outline:} \textbf{Introduction:} define rural electrification; thesis --- it is a powerful driver of development but not sufficient alone. \textbf{For 1:} economic life --- welding shops in Nkambe, cold storage for tomatoes in Santa, phone charging businesses. \textbf{For 2:} education and health --- evening study, vaccine refrigeration. \textbf{Against:} without roads, water and markets, electricity alone cannot transform a village. \textbf{Rebuttal/conclusion:} electricity is the \emph{spark} that makes other investments profitable; therefore it deserves priority, as part of a package. This plan guarantees balance --- the very skill drilled in Section 7.
\end{exoresolu}

\courssub{The Error Logbook: Classifying and Remediating}

An error repeated twice is a habit; recorded once and targeted, it becomes a mark gained. Keep a dedicated notebook --- your \cle{error logbook} --- with four columns: \emph{Error}, \emph{Category}, \emph{Correction}, \emph{Re-test date}.

\begin{center}
\begin{tabularx}{\linewidth}{|l|l|X|l|}
\hline
\rowcolor{popblueL} \textbf{Error observed} & \textbf{Category} & \textbf{Correction / rule} & \textbf{Re-test} \\
\hline
``If I would go\ldots'' & Conditional & Second conditional: \emph{If I went} & 1 week \\
\hline
``Walking home, the rain soaked me.'' & Dangling modifier & \emph{Walking home, I was soaked.} & 1 week \\
\hline
Summary of 165 words & Summary technique & Count words; cut examples first & 2 weeks \\
\hline
Retelling Santiago's plot & Literature analysis & Point + evidence + link to question & 2 weeks \\
\hline
\end{tabularx}
\end{center}

\begin{methode}[Spaced revision]
Do not cram. Revisit each skill at \emph{increasing intervals}:
\begin{enumerate}
\item Learn or correct the item today.
\item Re-test after \textbf{1 day}, then \textbf{1 week}, then \textbf{2 weeks}, then \textbf{1 month}.
\item If a re-test fails, restart the interval for that item only.
\item Keep items you pass three times in a ``mastered'' list, reviewed monthly.
\end{enumerate}
Spaced revision exploits the way memory consolidates: each near-forgotten recall strengthens the trace far more than massed repetition.
\end{methode}

\courssub{Planning the Remaining Terms}

Convert everything above into a visible plan. The wheel below places every skill family of this chapter on one page; shade each sector as it is mastered, and rotate through the wheel so that no skill is left untouched for more than two weeks.

\begin{popfigure}
\begin{center}
\begin{tikzpicture}[scale=1.05]
\foreach \i/\name in {0/{Essay \& argument},45/{Comprehension},90/{Summary \& report},135/{Grammar},180/{Phonetics \& oral},225/{Literature},270/{Poetry},315/{Editing \& discourse}} {
  \draw[fill=popblueL, draw=popdark] (0,0) -- (\i:2.9) arc (\i:\i+45:2.9) -- cycle;
  \node[font=\small, align=center, text width=2.1cm] at (\i+22.5:2.05) {\name};
}
\draw[fill=white, draw=popdark] (0,0) circle (0.85);
\node[font=\small\bfseries, align=center, text width=1.5cm] at (0,0) {Revision wheel};
\end{tikzpicture}
\end{center}
\legende{The revision-planning wheel: every sector of the chapter revisited in rotation, weakest sectors most often.}
\end{popfigure}

\begin{methode}[Building your past-paper schedule]
\begin{enumerate}
\item List the weeks remaining before the examination.
\item Assign one full past paper (or half paper) per fortnight, always timed, always marked with the same grid.
\item Between papers, fill the gaps with error-logbook items and wheel sectors --- weakest first.
\item Keep a self-assessment grid: date, paper, score per skill, top three errors, action taken.
\item Every month, redraw the progress chart. If a bar refuses to rise, change the \emph{method} for that skill (peer explanation, teacher consultation, model answers), not merely the time spent.
\end{enumerate}
\end{methode}

\begin{aretenir}[Key points of Section 14]
\begin{itemize}
\item Progress is measured against the \cle{diagnostic baseline}, under identical conditions and mark schemes.
\item Audit each paper: Paper 1 (essay, comprehension, summary), Paper 2 (text knowledge + analysis), Paper 3/oral (IPA, connected speech, intonation).
\item Run a full timed sequence at least once a fortnight.
\item Keep an error logbook; classify errors precisely; remediate with spaced revision.
\item Target weaknesses, not favourite topics; let the shortest bar on the chart set the agenda.
\end{itemize}
\end{aretenir}

\begin{exercice}[Your personal consolidation]
\begin{enumerate}
\item Using your own diagnostic and mid-year results, draw a progress chart modelled on the figure above. Identify your two weakest skills.
\item For each weak skill, write three error-logbook entries (real errors from your own scripts) with corrections and re-test dates.
\item Design a four-week revision timetable covering all eight sectors of the wheel, with one timed past-paper session per fortnight.
\item Write a 200-word reflective report: \emph{``My progress since September and my strategy for the remaining terms.''} Use at least three discourse markers of staging and concluding.
\end{enumerate}
\end{exercice}

\begin{corrige}[Exercise n°1 -- guidance]
(1) The chart must use the same scale and mark scheme for both bars per skill; weakest skills are the shortest mid-year bars, not the largest gaps. (2) Each entry must name a precise category (e.g.\ ``mixed conditional'', ``misplaced modifier'', ``word-limit breach''), state the rule, and set a realistic re-test date following the 1-day / 1-week / 2-weeks / 1-month pattern. (3) A valid timetable shows rotation through all eight sectors, extra weight on the two weak ones, and two timed sessions. (4) The report should open with staging markers (\emph{To begin with, Looking back over the term}), develop with sequencing (\emph{Furthermore, As a result}), and close with concluding markers (\emph{In conclusion, All things considered}), with honest reference to the baseline and a concrete strategy.
\end{corrige}

You began the year by listening, reading and responding; you end this first half of the year by \emph{evaluating yourself} --- the highest skill on the syllabus. Carry the logbook, the wheel and the timetable into the next term, and the examination will hold no surprises.

\newpage

\coursec{Integration activity}

\coursec{Integration Activity: The Plastic Sachet Ban Dossier}

\courssub{Aims of the activity}

This integration activity is designed to consolidate the core competences developed in Module One. By the end of the activity, you should be able to:
\begin{itemize}
    \item Read, annotate and critically evaluate complex texts (900--1000 words) presenting conflicting viewpoints on a contemporary Cameroonian issue.
    \item Distinguish facts from opinions, identify implicit meaning, and recognise the writer's tone, attitude and rhetorical strategies.
    \item Synthesise information from multiple sources into a concise, neutral summary (about 100 words).
    \item Demonstrate mastery of connected speech features (elision, assimilation, intrusion, catenation) and intonation patterns in a simulated oral interaction.
    \item Transcribe selected utterances into broad IPA notation.
    \item Write a balanced argumentative essay (350--400 words) that integrates multiple perspectives, uses a full range of subordinate clauses (noun, adjective, adverb), conditional forms, and discourse markers for staging, sequencing and concluding.
    \item Apply an editing checklist to peer-review writing for grammar, cohesion, register and GCE examination requirements.
\end{itemize}

\begin{attention}[A note on the texts]
The three documents below are \cle{model texts} written for this activity. They imitate real Cameroonian text types (a newspaper editorial, an NGO technical report, an open letter) and contain invented figures and events for practice purposes. In the GCE examination, always base your answers \emph{only} on the passage in front of you, and never import outside ``facts'' you cannot verify.
\end{attention}

\courssub{Situation}

The Ministry of Environment, Protection of Nature and Sustainable Development (MINEPDED) has announced a nationwide consultation on the proposed \emph{Plastic Sachet Ban Bill}. The bill seeks to prohibit the production, importation, sale and use of non-biodegradable plastic sachets (commonly called ``\cle{buying bags}'' or ``\cle{pure water sachets}''). Three stakeholders have submitted position documents to the consultation portal. As a class, you constitute the \emph{Youth Parliamentary Committee on Environmental Legislation}. Your task is to analyse the dossier, synthesise the arguments, simulate a press conference, and produce an individual policy position essay.

\textbf{Text A --- Newspaper Editorial: ``The Sachet Scourge: Why Cameroon Must Act Now''} \\
\textit{The Cameroon Tribune, Opinion Desk}

Cameroon stands at an environmental crossroads. Every day, millions of plastic sachets are discarded across our national territory. In Yaoundé alone, municipal services collect tonnes of sachet waste daily, yet a large share escapes formal collection, choking drains, littering streets and eventually feeding the Wouri and Sanaga rivers. The science is unequivocal: low-density polyethylene (LDPE) sachets take centuries to photodegrade, fragmenting into microplastics that enter the food chain. Researchers who have sampled fish from the Wouri estuary report the presence of microplastic particles --- fish that end up on dinner tables from Bonamoussadi to Buea.

Opponents of the ban cite economic hardship for sachet-water vendors and small retailers. This argument, while emotionally resonant, is economically short-sighted. Rwanda's ban on plastic bags did not collapse its informal sector; it catalysed a thriving industry in reusable packaging. Kenya's ban, enforced with heavy fines, achieved high compliance within a few years. Cameroon's own earlier decree banning non-biodegradable packaging below a specified thickness exists on paper but lacks enforcement mechanisms and public awareness campaigns. The proposed Bill corrects this by mandating a transition period, subsidised biodegradable alternatives for registered vendors, and an \cle{Extended Producer Responsibility} (EPR) scheme forcing manufacturers to fund collection and recycling.

Critics also claim the ban disproportionately affects the poor who rely on affordable sachet water. Yet the poor suffer most from the consequences: blocked drains cause seasonal flooding in neighbourhoods like Mvog-Ada and Nkolbisson, breeding malaria and cholera. A ban is not an attack on the poor; it is a public health intervention.

The government must resist lobbying from the plastics manufacturing lobby, which contributes modestly to GDP but externalises cleanup costs that run into billions of CFA francs annually. Political courage is required. The Bill should be passed with strict enforcement timelines, transparent monitoring by civil society, and a massive behavioural-change campaign in schools, markets and motor parks. Cameroon cannot afford to be the dumping ground of Central Africa. The time for half-measures is over.

\textbf{Text B --- NGO Technical Report: ``Assessment of the Socio-Economic Impacts of a Plastic Sachet Ban in Cameroon''} \\
\textit{Green Horizon Cameroon, Executive Summary}

This report evaluates the potential impacts of the proposed Plastic Sachet Ban Bill on vulnerable populations and the informal economy. Using a mixed-methods approach --- household surveys across five regions, key-informant interviews, and focus-group discussions with women's cooperatives --- we present evidence-based recommendations.

\textbf{Key Findings:}
\begin{enumerate}
    \item \textbf{Water Access:} A large majority of low-income urban households rely on sachet water (250 ml at CFA 50) as their primary drinking source. Bottled water and household filtration systems are financially inaccessible. A ban without a simultaneous, subsidised safe-water infrastructure plan risks undermining access to safe drinking water, which the United Nations recognises as a human right.
    \item \textbf{Employment:} The sachet value chain employs tens of thousands of people directly (production, distribution, street vending) and many more indirectly. A large majority are women and youth. Rwanda's transition took several years with donor-funded retraining; Kenya's enforcement led to vendor arrests in its first year. Cameroon's proposed short transition is unrealistic for structural transformation.
    \item \textbf{Alternative Materials:} Locally produced biodegradable sachets currently cost several times the price of LDPE sachets. No domestic manufacturer operates at scale. Import dependency would shift from plastic resins to biopolymers, worsening the trade deficit. Reusable containers require behavioural change campaigns lasting many months to achieve significant adoption (diffusion-of-innovation theory).
    \item \textbf{Waste Management:} Cameroon's recycling rate for plastics is very low. A ban without parallel investment in waste segregation at source, material recovery facilities, and formalisation of waste-picker cooperatives will merely displace waste to other polymer types (bottles, bags).
\end{enumerate}

\textbf{Recommendations:}
\begin{enumerate}
    \item Adopt a \cle{phased approach}: Phase 1 (Years 1--2) --- ban large sachets above 500 ml; mandate recycled content in all sachets; subsidise biodegradable alternatives for registered vendors. Phase 2 (Years 3--5) --- total ban contingent on near-universal urban piped-water coverage and an operational EPR scheme.
    \item Establish a \cle{Just Transition Fund} financed by polymer producer levies, supporting vendor retraining, micro-credit for reusable-container micro-enterprises, and school-based environmental education.
    \item Recognise and integrate waste pickers (\cle{``chiffonniers''}) into municipal waste management contracts with social protection.
    \item Commission an independent regulatory impact assessment (RIA) before parliamentary debate.
\end{enumerate}

A total ban \emph{is} environmentally necessary but socially reckless if implemented without these safeguards. Policy must be both green and just.

\textbf{Text C --- Open Letter: ``Our Livelihoods Are Not Waste''} \\
\textit{By Madeleine Ngo Bitang, President, Yaoundé Sachet Water Vendors' Cooperative (COOVEYA), published in \textit{Mutations}}

Honourable Ministers, Honourable Parliamentarians, Fellow Cameroonians,

I write not as a statistician but as a mother of three who has sold \cle{pure water} at the Carrefour Mvog-Ada junction for seventeen years. My day begins at 4:30 a.m., pushing a wheelbarrow loaded with 200 sachets from the depot in Nsam. By 6:00 a.m., taxi drivers, students, office cleaners, and construction workers buy at CFA 50 per sachet --- the only safe water they can afford. I earn CFA 3,500 profit on a good day. From this, I pay school fees for my daughter at Lycée de Mvog-Ada, rent for our two-room house in Essos, and medical bills for my husband's diabetes.

You call our sachets ``pollution''. We call them survival. Have any of you walked through Mvog-Ada after a heavy rain? The drains are blocked not by our sachets alone but by silt, household refuse, and the plastic bottles your ministries distribute at conferences. Yet you target the poorest link in the chain. The big factories in Douala produce millions of plastic bottles yearly. Where is their ban? Where is their EPR fund? We, the vendors, are visible, voiceless, and voteless --- easy scapegoats.

You speak of Rwanda and Kenya. In Kigali, a vendor gets a government loan to buy a water dispenser and reusable cups. In Nairobi, the county government built water kiosks before enforcing the ban. Here, you give us a few months and a promise of ``subsidised alternatives'' that do not exist in any market in Yaoundé. A biodegradable sachet costs more than three times the price of an ordinary one today. Who will buy water at that price? My customers will drink from doubtful wells and streams. Cholera will return, and you will blame us again.

We are not enemies of the environment. We live in it. Our cooperative has proposed a \cle{deposit-return scheme}: CFA 10 per sachet returned to collection points we manage. This creates jobs, cleans streets, and recycles plastic. We presented it to MINEPDED. No reply. We proposed reusable 5-litre jerrycans branded with our cooperative logo, sold at cost. Silence.

If you ban our sachets tomorrow, thousands of families lose their daily bread. You will not see us in the unemployment statistics because we do not exist in your formal economy. But you will see our children drop out of school, our husbands turn to desperate measures, our communities fracture. Is this the ``emergence'' you promised?

We ask for three things: (1) A realistic transition of at least three years with monthly stipends during retraining. (2) Inclusion in the design of alternatives --- we know our customers. (3) Equal enforcement: if sachets are banned, so must plastic bottles, styrofoam plates, and the packaging of every imported product in our markets.

Do not sign our death warrant with a pen you hold in an air-conditioned office. Come to Carrefour Mvog-Ada at 5:00 a.m. See the faces. Hear the stories. Then legislate with conscience.

\emph{Madeleine Ngo Bitang, for the vendors of Yaoundé.}

\courssub{Tasks}

\begin{enumerate}
    \item \textbf{Critical Reading and Annotation (Group Work --- 45 minutes)} \\
    In groups of four, read Texts A, B and C. For each text, complete the table below. Use a separate sheet.
    \begin{center}
    \begin{tabularx}{\linewidth}{|l|X|X|X|}\hline
    \rowcolor{popblueL}
    \textbf{Analytical Lens} & \textbf{Text A (Editorial)} & \textbf{Text B (NGO Report)} & \textbf{Text C (Open Letter)} \\ \hline
    Main claim / thesis & & & \\ \hline
    Text type and audience & & & \\ \hline
    Writer's tone and attitude & & & \\ \hline
    Three key facts (verifiable) & & & \\ \hline
    Three key opinions / value judgements & & & \\ \hline
    Implicit assumptions & & & \\ \hline
    Rhetorical strategies (logos, pathos, ethos) & & & \\ \hline
    Proposed solution / call to action & & & \\ \hline
    \end{tabularx}
    \end{center}

    \item \textbf{Neutral Synthesis Summary (Individual --- 20 minutes)} \\
    Write a \textbf{100-word} summary (between 90 and 110 words) that neutrally synthesises the three positions. Do not include your own opinion. Use discourse markers of contrast, addition and consequence. Count and indicate the word count at the end.

    \item \textbf{Simulated Press Conference (Whole Class --- 40 minutes)} \\
    \begin{enumerate}
        \item The class divides into four panels: \textbf{Panel 1} (Government/Ministry), \textbf{Panel 2} (NGO/Experts), \textbf{Panel 3} (Vendors/Cooperative), \textbf{Panel 4} (Journalists).
        \item Each panel prepares 3--4 key statements or questions (2 minutes each) defending or probing their assigned position.
        \item \textbf{Oral focus:} Speakers must consciously use connected speech features (elision, assimilation, intrusion, catenation) and appropriate intonation patterns (falling for assertions, rising for genuine questions, fall-rise for reservations).
        \item \textbf{Transcription task:} Panel 4 (Journalists) transcribes \textbf{two} key sentences from each of the other panels into \textbf{broad IPA}, marking connected speech features and tonic syllables. Example: \texttt{/ðə ɡʌvənmənt mʌst rɪˈzɪst ˈlɒbɪɪŋ/} becomes \texttt{/ðə ɡʌvənmən mʌs rɪˈzɪs ˈlɒbɪɪŋ/} (elision of /t/, weakening of unstressed syllables).
    \end{enumerate}

    \item \textbf{Balanced Argumentative Essay (Individual --- 45 minutes)} \\
    Write an essay of \textbf{350--400 words} on the following title:
    \begin{center}
    \textbf{``A total ban on plastic sachets in Cameroon is environmentally necessary but socially unjust unless accompanied by a comprehensive just-transition plan.'' Discuss.}
    \end{center}
    Your essay must:
    \begin{itemize}
        \item Present a clear thesis statement in the introduction.
        \item Develop at least three well-structured body paragraphs, each integrating evidence from \textbf{at least two} of the three texts.
        \item Use a range of \textbf{subordinate clauses} (noun clauses as subject, object or complement; adjective clauses with relative pronouns; adverb clauses of concession, condition, reason and result).
        \item Include at least \textbf{two conditional sentences} (mixed or third conditionals preferred for hypothetical evaluation).
        \item Employ \textbf{discourse markers} for staging (\textit{Firstly, Moreover, Conversely, Consequently, In conclusion}), sequencing and concluding.
        \item End with a nuanced final stance that acknowledges complexity.
    \end{itemize}

    \item \textbf{Peer Editing and Self-Monitoring (Pairs --- 20 minutes)} \\
    Exchange essays with a partner. Use the \textbf{Editing Checklist} below. Annotate the essay directly. Write a brief feedback note (about 50 words) highlighting two strengths and two targeted improvements. Return the essay. The writer then revises and submits a final polished version.
    \begin{center}
    \begin{tabularx}{\linewidth}{|l|X|c|}\hline
    \rowcolor{popblueL}
    \textbf{Criterion} & \textbf{Guiding Questions} & \textbf{Mark} \\ \hline
    Task fulfilment & Word count (350--400)? Title addressed? Clear stance? & \\ \hline
    Structure & Thesis statement? Topic sentences? Logical paragraphing? Conclusion? & \\ \hline
    Synthesis & Evidence from at least two texts per paragraph? Neutral attribution? & \\ \hline
    Grammar: Clauses & Noun, adjective and adverb clauses used correctly? Variety? & \\ \hline
    Grammar: Conditionals & At least two conditionals, correct in form and meaning? & \\ \hline
    Discourse markers & Staging, sequencing, contrast and conclusion used appropriately? & \\ \hline
    Cohesion & Reference, substitution, ellipsis, lexical chains? & \\ \hline
    Register & Formal, objective, persuasive? No colloquialisms? & \\ \hline
    Mechanics & Spelling, punctuation, capitalisation? & \\ \hline
    \end{tabularx}
    \end{center}
\end{enumerate}

\courssub{Model answer}

\textbf{Task 1 --- Completed Analytical Table (Sample Entries)}

\begin{center}
\begin{tabularx}{\linewidth}{|l|X|X|X|}\hline
\rowcolor{popblueL}
\textbf{Analytical Lens} & \textbf{Text A (Editorial)} & \textbf{Text B (NGO Report)} & \textbf{Text C (Open Letter)} \\ \hline
Main claim / thesis & Total ban is urgent, feasible and morally imperative; Rwanda and Kenya prove it works. & Ban is environmentally necessary but socially reckless without a phased approach, a Just Transition Fund and water infrastructure. & Ban as proposed will destroy thousands of livelihoods; vendors propose a deposit-return scheme and demand a 3-year transition with stipends. \\ \hline
Text type and audience & Newspaper editorial; general public, policymakers, opinion leaders. & Technical NGO report; policymakers, donors, civil society. & Open letter in a daily newspaper; ministers, parliamentarians, general public. \\ \hline
Writer's tone and attitude & Urgent, authoritative, moralistic, dismissive of counter-arguments (``short-sighted''). & Evidence-based, cautious, balanced, advocacy-oriented (``green and just''). & Personal, emotional, indignant, pleading, proud (``We are not enemies of the environment''). \\ \hline
Three key facts & Millions of sachets discarded daily; tonnes collected daily in Yaoundé; microplastics reported in Wouri estuary fish. & Most poor households rely on sachet water; tens of thousands of direct jobs, mostly women; biodegradable sachets cost several times more. & Vendor earns CFA 3,500 profit on a good day; a biodegradable sachet costs over three times the ordinary price; the cooperative proposed a deposit-return scheme. \\ \hline
Three key opinions & ``Political courage is required''; ``A ban is not an attack on the poor''; ``The time for half-measures is over''. & ``The short transition is unrealistic''; ``A total ban is socially reckless''; ``Policy must be both green and just''. & ``You target the poorest link''; ``You will not see us in the unemployment statistics''; ``Do not sign our death warrant''. \\ \hline
Implicit assumptions & Enforcement capacity exists; behavioural change is swift; EPR will be transparent; manufacturers will comply. & Donor funding is available for a Just Transition Fund; government can implement phased regulation; waste pickers can be formalised. & Government ignores informal-sector voices; big corporations escape scrutiny; alternatives are unavailable or unaffordable. \\ \hline
Rhetorical strategies & Logos (statistics, Rwanda/Kenya precedents); pathos (cholera, flooding); ethos (``The science is unequivocal''). & Logos (survey data, diffusion theory); ethos (mixed methods, institutional partners); pathos (the human right to water). & Ethos (seventeen years of experience, mother of three); pathos (daily routine, school fees, illness); logos (deposit-return proposal). \\ \hline
Proposed solution & Pass the Bill with strict timelines, EPR, subsidies and massive campaigns. & Phased ban (2 plus 3 years), Just Transition Fund, waste-picker integration, independent RIA. & Three-year transition with stipends, vendor inclusion in design, equal enforcement on all plastics, deposit-return scheme. \\ \hline
\end{tabularx}
\end{center}

\textbf{Task 2 --- Neutral Synthesis Summary (Model, 100 words)}

The proposed plastic sachet ban has elicited three divergent positions. The \textit{Cameroon Tribune} editorial argues that an immediate, strictly enforced ban is environmentally urgent and economically viable, citing Rwanda's and Kenya's successes and the severe health costs of pollution. Conversely, Green Horizon's technical report warns that a short transition is unrealistic: most poor households depend on affordable sachet water, tens of thousands of vendors (mostly women) risk job loss, and biodegradable alternatives cost several times more. It therefore recommends a phased ban, a Just Transition Fund and prior investment in water infrastructure. Meanwhile, vendor leader Madeleine Ngo Bitang condemns the ban as a death warrant for livelihoods, demanding a three-year transition with stipends, vendor inclusion in policy design and equal enforcement on all plastics, while proposing a cooperative-run deposit-return scheme. (100 words)

\textbf{Task 3 --- Sample IPA Transcriptions (Panel 4 Journalist Notes)}

\begin{enumerate}
    \item \textbf{Panel 1 (Government):} ``The ministry will enforce the ban without compromise.'' \\
    \textit{Citation form:} \texttt{/ðə ˈmɪnɪstri wɪl ɪnˈfɔːs ðə bæn wɪˈðaʊt ˈkɒmprəmaɪz/} \\
    \textit{Connected speech:} \texttt{/ðə ˈmɪnɪstri wɪl ɪnˈfɔːs ðə ˈbæm wɪˈðaʊ ˈkɒmprəmaɪz/} --- assimilation of /n/ to /m/ before /b/ in ``the ban'', elision of /t/ in ``without''.
    \item \textbf{Panel 2 (NGO):} ``If the transition were phased, vendors would adapt.'' \\
    \textit{Citation form:} \texttt{/ɪf ðə trænˈzɪʃən wɜː feɪzd ˈvendəz wʊd əˈdæpt/} \\
    \textit{Connected speech:} \texttt{/ɪf ðə trænˈzɪʃən wə feɪzd ˈvendəz wʊd əˈdæpt/} --- weak form /wə/ for ``were'', catenation (linking) of final /d/ into the following vowel in ``phased, vendors''.
    \item \textbf{Panel 3 (Vendors):} ``We have fed this city for seventeen years.'' \\
    \textit{Citation form:} \texttt{/wiː hæv fed ðɪs ˈsɪti fɔː ˌsevənˈtiːn jɪəz/} \\
    \textit{Connected speech:} \texttt{/wi əv fed ðɪs ˈsɪti fə ˌsevənˈtiːn jɪəz/} --- elision of /h/ and weak form /əv/ in ``have'', weak form /fə/ for ``for'', intrusive /j/ glide between ``seventeen'' and ``years''.
\end{enumerate}

\textbf{Task 4 --- Model Argumentative Essay (about 390 words)}

\textbf{A Total Ban on Plastic Sachets in Cameroon Is Environmentally Necessary but Socially Unjust Unless Accompanied by a Comprehensive Just-Transition Plan}

Cameroon's proposed Plastic Sachet Ban Bill has ignited a debate that pits ecological imperatives against the survival strategies of the urban poor. While the environmental case for a ban is scientifically robust, the social consequences of abrupt enforcement --- without a credible just-transition framework --- would deepen inequality and restrict access to safe water. This essay argues that a total ban is necessary but must be embedded in a phased, financed and inclusive plan that protects vendors and guarantees safe water access.

Firstly, the ecological evidence is overwhelming. As the \textit{Cameroon Tribune} editorial highlights, millions of sachets are discarded daily, choking drains and leaching microplastics into the Wouri estuary, where researchers have detected plastic particles in fish. These externalities impose cleanup costs of billions of CFA francs on municipalities every year. Rwanda's and Kenya's bans demonstrate that decisive legislation can achieve high compliance and stimulate green industries. However, the editorial's dismissal of economic concerns as ``short-sighted'' overlooks structural differences: both countries invested heavily in water infrastructure and vendor retraining \emph{before} enforcement. Had Cameroon replicated that sequencing, the transition would be less precarious.

Conversely, the Green Horizon report provides empirical proof that a short transition is a fantasy. Since most low-income households depend on CFA 50 sachets and biodegradable alternatives cost several times more, a sudden ban would price safe water out of reach. Moreover, tens of thousands of direct jobs --- most of them held by women --- would vanish overnight. The report's recommendation of a phased approach, under which a total ban would come only after near-universal piped-water coverage, is pragmatic. If the government had commissioned an independent regulatory impact assessment earlier, these gaps would have been anticipated.

Furthermore, Madeleine Ngo Bitang's open letter humanises the statistics. Her cooperative's deposit-return proposal --- CFA 10 per sachet returned --- illustrates that vendors are not passive polluters but potential agents of a circular economy. Yet their submission received no reply, which reveals a democratic deficit in policy formulation. A just transition would integrate such grassroots innovations, formalise waste pickers, and equalise enforcement across all plastic categories, including the bottles produced by corporate giants who escape scrutiny.

In conclusion, the ban is environmentally non-negotiable but socially indefensible in its current form. A credible pathway requires three pillars: a phased legislative timeline aligned with the rollout of water infrastructure; a Just Transition Fund financed by polymer levies to support vendor retraining and micro-enterprise; and genuine co-design with affected communities. Only then will Cameroon achieve a ban that is both green and just.

\textbf{Task 5 --- Peer Editing Feedback (Sample)}

\textit{Strengths:} (1) Clear thesis and nuanced stance maintained throughout. (2) Effective synthesis: each body paragraph integrates at least two texts with accurate attribution. \textit{Improvements:} (1) Paragraph 2 contains a well-formed mixed conditional (``Had Cameroon replicated\ldots the transition would be less precarious''), but the hypothetical in paragraph 3 (``If the government had commissioned\ldots would have been anticipated'') is a third conditional; varying the conditional types deliberately would show greater grammatical range. (2) The discourse marker ``Furthermore'' in paragraph 4 adds a point but does not signal the shift from NGO data to the vendor's voice; ``However'' or ``By contrast'' would better reflect that movement.

\begin{aretenir}[What this activity tests]
\begin{itemize}
    \item \textbf{Reading:} separating fact from opinion, detecting tone, attitude and implicit assumptions across three contrasting text types.
    \item \textbf{Summary:} neutral synthesis in about 100 words, with contrast and consequence markers, and no personal opinion.
    \item \textbf{Speaking:} connected speech (elision, assimilation, intrusion, catenation), intonation choices, and broad IPA transcription.
    \item \textbf{Writing:} a balanced argumentative essay of 350--400 words using subordinate clauses, conditionals and discourse markers.
    \item \textbf{Editing:} systematic peer review against task fulfilment, structure, grammar, cohesion, register and mechanics.
\end{itemize}
\end{aretenir}

\newpage

\coursec{Methods and worked examples}

\courssub{Skill 1 — Reading a Complex Passage: Main Ideas and Implicit Meaning}

\begin{definition}[Key Terms for Comprehension]
The \cle{main idea} of a passage is the central claim or message that unifies all its paragraphs; everything else is support. An idea is \cle{explicit} when it is stated directly in words; it is \cle{implicit} (or implied) when the reader must infer it from hints such as word choice, examples, contrasts or concessions. An \cle{inference} is a conclusion the reader draws by combining textual evidence with logic — never with mere personal opinion.
\end{definition}

\begin{methode}[Tackling a 900--1000 Word Comprehension Passage]
\begin{enumerate}
\item \textbf{Preview before you read.} Read the title, the source line and the questions first. The questions tell you \emph{what to hunt for}; the title often announces the writer's angle.
\item \textbf{First reading — global.} Read the whole passage once without stopping. At the end of each paragraph, mentally summarise it in five to eight words (e.g. ``cause of the problem'', ``writer's personal example'').
\item \textbf{Locate topic sentences.} In formal expository prose, the main idea of a paragraph usually sits in its first or last sentence. Underline these; they form the skeleton of the passage.
\item \textbf{Second reading — targeted.} Return to the questions. For each one, find the \cle{anchor lines} in the passage and read two lines above and two lines below before answering.
\item \textbf{Distinguish explicit from implicit.} For inference questions, ask: \emph{What must be true for this sentence to make sense?} Then check that your inference is supported by specific words in the text.
\item \textbf{Answer in your own words} unless the question says ``quote''. Lifting whole sentences when not required costs marks, because it does not prove understanding.
\item \textbf{Vocabulary-in-context:} replace the target word with your synonym and check that the sentence still makes sense grammatically and logically.
\end{enumerate}
\textbf{Reflexes:} never answer from general knowledge; always justify an inference with evidence from the text; manage time at roughly one mark per minute.
\end{methode}

\begin{exoresolu}[Comprehension of a Complex Passage]
Read the extract and answer the questions.

\emph{``The Cameroonian student who memorises a textbook the way a parrot learns a song may pass the examination, but he will fail the far sterner test that life sets afterwards. True education is not the accumulation of facts; it is the training of the mind to weigh evidence, to detect the flaw in a sweet-sounding argument, and to hold a conviction only as firmly as the evidence permits. A certificate can open a door, certainly; but only a disciplined mind can keep it open.''}

\textbf{Questions:} (a) State the writer's main idea in one sentence. (b) What does the writer imply about examinations? (c) Explain in your own words: ``the far sterner test that life sets afterwards''. (d) Give the meaning of ``weigh evidence'' as used in the passage.
\tcblower
\textbf{Detailed solution.}

\textbf{(a) Main idea.} The writer argues that genuine education consists in developing critical judgement rather than merely memorising facts for examinations. \emph{Technique:} the topic sentence is the second sentence (``True education is not\ldots it is\ldots''); the first sentence only introduces the problem through the parrot comparison.

\textbf{(b) Implication about examinations.} The writer implies that examinations, as presently passed by rote learning, are an inadequate measure of real ability: a candidate can succeed in them and still be intellectually unprepared for real life. Note the signal ``may pass\ldots but he will fail'' — the concession structure shows the writer grants examinations limited value (they ``open a door'') while denying them ultimate value. This is an \emph{implicit} point: the passage never says ``examinations are inadequate''; we infer it from the contrast.

\textbf{(c) Explanation.} The phrase means the harder, more demanding challenges of real life after school — making decisions, judging people, solving problems — where memorised facts alone are useless. ``Sterner'' means more severe; ``sets'' means presents, as an examiner sets a paper. The metaphor extends the examination image: life itself is an examiner, but a harsher one.

\textbf{(d) Vocabulary in context.} ``Weigh evidence'' means to examine and assess information carefully in order to judge its value or reliability — as one weighs objects on a balance before deciding. \emph{Check:} ``the training of the mind to assess evidence carefully'' — the sentence still reads naturally, so the synonym is correct.

\textbf{Marking reflex:} each answer is anchored to a specific line of the text; none relies on outside knowledge.
\end{exoresolu}

\courssub{Skill 2 — Identifying the Writer's Attitude, Tone and Point of View}

\begin{definition}[Tone, Attitude, Point of View]
\cle{Tone} is the writer's emotional attitude towards the subject as revealed by the language (admiring, bitter, ironic, detached, urgent\ldots). \cle{Point of view} is the position from which the writer speaks: first or third person, insider or outsider, partisan or neutral. \cle{Irony} is a figure in which the intended meaning is the opposite of (or sharply different from) the literal meaning; \cle{sarcasm} is a biting, mocking form of irony intended to wound.
\end{definition}

\begin{methode}[Diagnosing Tone and Attitude]
\begin{enumerate}
\item \textbf{Define your terms} before you write: tone concerns feeling; point of view concerns position. Do not confuse them in your answer.
\item \textbf{Collect the evidence.} Tone is carried by (i) \emph{connotative vocabulary} (``freedom fighter'' versus ``rebel''), (ii) \emph{figurative language}, (iii) \emph{sentence rhythm} (short exclamations signal emotion; long balanced sentences signal calm reasoning), and (iv) \emph{what is omitted or mocked}.
\item \textbf{Watch for irony.} If the literal meaning contradicts the context or the writer's evident purpose, the tone is ironic or sarcastic. Irony praises in order to blame, or blames in order to praise.
\item \textbf{Grade the intensity.} Distinguish neighbouring tones: \emph{critical} (finds faults) is not \emph{hostile} (attacks); \emph{nostalgic} (longs for the past) is not \emph{sentimental} (wallows in emotion); \emph{objective} (reports) is not \emph{indifferent} (does not care).
\item \textbf{Justify with quotations.} An unsupported tone label earns little credit at Advanced Level: always name the tone \emph{and} quote or cite the words that create it.
\end{enumerate}
\textbf{Key idea:} tone = adjective + evidence. Formula for your answer: \emph{``The tone is X, created by Y (quotation), which suggests Z.''}
\end{methode}

\begin{exoresolu}[Tone and Point of View]
Comment on the writer's tone and point of view in this extract, supporting your answer with evidence.

\emph{``Our leaders never tire of reminding us that the nation is a great family. Indeed it is: a family in which the father eats the chicken, the children smell the aroma, and the dog is lectured on patience. How fortunate we are to be governed by such devoted parents!''}
\tcblower
\textbf{Detailed solution.}

\textbf{Step 1 — Literal versus intended meaning.} On the surface the writer praises the leaders (``great family'', ``devoted parents'', ``how fortunate''). But the central image contradicts the praise: the ``father'' monopolises the chicken while the children only ``smell the aroma''. The literal meaning and the described reality clash — the unmistakable signal of \cle{irony}.

\textbf{Step 2 — Name the tone precisely.} The tone is \emph{ironic and bitterly satirical}, not merely critical. The writer mocks the leaders' rhetoric by pretending to accept it and then exposing its falseness. The final exclamation (``How fortunate we are\ldots!'') is sarcasm: it means exactly the opposite of what it says.

\textbf{Step 3 — Identify the point of view.} The writer speaks as an \emph{insider-victim}, using the first person plural ``we'' and ``us''. He places himself among the governed, not above them, which makes the satire a cry of shared grievance rather than a detached lecture. The leaders are distanced as ``father'' — self-serving figures who preach patience to others.

\textbf{Step 4 — Assemble the answer (formula: tone + evidence + effect).}

\emph{Model answer:} The writer adopts an ironic, satirical tone towards the country's leadership. He pretends to endorse the official metaphor of the nation as ``a great family'', then subverts it with an extended image in which ``the father eats the chicken'' while ``the children smell the aroma'' — a bitter picture of leaders who monopolise national wealth while citizens receive only promises. The sarcastic exclamation ``How fortunate we are to be governed by such devoted parents!'' says the opposite of what it means, deepening the mockery. Writing in the first person plural (``us'', ``we''), the author speaks as one of the suffering citizens, which gives the satire the force of a collective complaint rather than an outsider's accusation.

\textbf{Why this scores full marks:} a precise tone label (ironic/satirical, not the vague ``negative''), two quoted pieces of evidence, an explanation of \emph{how} the language works, and a comment on point of view.
\end{exoresolu}

\courssub{Skill 3 — Summarising a Complex Passage}

\begin{methode}[Writing a Summary in the Required Number of Words]
\begin{enumerate}
\item \textbf{Read the instruction twice.} Note the word limit (e.g. ``in about 100 words''), the portion of the passage concerned, and whether one sentence or several are demanded.
\item \textbf{Identify the skeleton.} Using your paragraph-by-paragraph notes, list the \cle{main points} only. A main point is one without which the passage's argument collapses.
\item \textbf{Delete ruthlessly:} examples, illustrations, statistics, repetitions, quotations, anecdotes, figures of speech — all are \emph{supports}, not points.
\item \textbf{Generalise.} Replace lists by a category word (``palm wine, kola nuts and bush meat'' becomes ``food''); replace concrete cases by the general truth they illustrate.
\item \textbf{Draft in your own words.} Combine related points with linkers (\emph{because, although, therefore, moreover}). Do not copy phrases of more than three consecutive words from the passage.
\item \textbf{Count and trim.} If over the limit, cut adjectives, adverbs and relative clauses first — never cut a main point.
\item \textbf{Check grammar and sense.} The summary must read as a coherent, self-contained text understandable by someone who has not read the passage.
\end{enumerate}
\textbf{Warnings:} no personal opinion, no new information, no direct speech, no ``the writer says'' padding.
\end{methode}

\begin{exoresolu}[Summary Writing]
Summarise the following passage in \textbf{about 60 words}.

\emph{``Water scarcity in the Sahel is not merely a matter of insufficient rainfall. Year after year, shallow wells are dug without any hydrological survey, so that many run dry within a single season; pumps imported at great cost break down and rust away because no local technician has been trained to repair them; and the few functioning boreholes are often monopolised by the powerful, who fence them off and sell what should be a common gift. Meanwhile women and girls walk further each year — four, six, sometimes ten kilometres — carrying on their heads the water that engineering and honesty could have brought to their doorsteps. Studies in several villages have shown that when a reliable well is sunk within five hundred metres of homes, girls' school attendance rises sharply, for the hours once spent walking are returned to the classroom.''}
\tcblower
\textbf{Detailed solution.}

\textbf{Step 1 — Extract the main points.} (i) Scarcity is caused not only by low rainfall but by human failures; (ii) wells are dug without surveys and dry up; (iii) pumps fail for lack of trained technicians; (iv) the powerful monopolise working boreholes; (v) women and girls bear the burden through long walks; (vi) nearby wells raise girls' school attendance.

\textbf{Step 2 — Compress and generalise.} Points (ii)--(iv) are all instances of \emph{mismanagement and inequality} — merge them. Point (vi) is the consequence worth keeping because it completes the argument.

\textbf{Step 3 — Draft (own words).}

\emph{Model summary (60 words):} ``Water scarcity in the Sahel stems less from poor rainfall than from human failure: wells are sunk without surveys, pumps break down for want of trained technicians, and the powerful seize the few that work. Women and girls pay the price in long daily treks; where reliable wells exist near homes, girls' school attendance rises markedly.''

\textbf{Word count:} 60 words. \textbf{Checks:} no copied phrases beyond unavoidable technical terms; no illustrative detail (the kilometre figures are cut); no opinion added; every main point present; the text reads coherently on its own.

\textbf{Common errors avoided:} keeping the ``four, six, ten kilometres'' detail (an illustration, not a point); copying ``a common gift'' (the writer's metaphor, not yours); writing ``the passage says\ldots'' (wasted words).
\end{exoresolu}

\courssub{Skill 4 — Grammar Mastery: Subordinate Clauses}

\begin{definition}[Clause Types]
A \cle{clause} is a group of words containing a subject and a finite verb. A \cle{main (independent) clause} can stand alone as a sentence; a \cle{subordinate (dependent) clause} cannot, and is introduced by a subordinator (\emph{that, which, who, because, although, when, if, unless\ldots}). According to its function in the sentence, a subordinate clause is a \cle{noun clause}, an \cle{adjective (relative) clause} or an \cle{adverb clause}.
\end{definition}

\begin{methode}[Identifying and Using Noun, Adjective and Adverb Clauses]
\begin{enumerate}
\item \textbf{Find the clause.} Confirm it has a subject and a finite verb, and identify its subordinator.
\item \textbf{Ask what job it does} in the main sentence:
\begin{itemize}
\item \textbf{Noun clause} — does the job of a noun: subject, object or complement. Test: replace it with \emph{something} or \emph{it}. \emph{What he said} surprised us. (\emph{It} surprised us.)
\item \textbf{Adjective (relative) clause} — qualifies a noun. Test: it answers ``which one?''. The girl \emph{who won the prize} is my cousin.
\item \textbf{Adverb clause} — modifies the verb: time, reason, condition, concession, purpose, result or manner. Test: it answers ``when? why? under what condition?''. \emph{Although it rained}, the match continued.
\end{itemize}
\item \textbf{Punctuation reflex:} a \emph{defining} relative clause (essential to identify the noun) takes \textbf{no commas}; a \emph{non-defining} one (extra information) is enclosed in commas. An adverb clause placed \emph{before} the main clause is followed by a comma; placed after it, it usually takes none.
\item \textbf{Reduce when elegant:} \emph{When he was asked}, he smiled becomes \emph{When asked}, he smiled — but only if the implied subjects match.
\end{enumerate}
\textbf{Exam reflex:} in transformation exercises, changing the clause type changes the label, not the meaning: \emph{Because he was ill, he stayed home} (adverb clause) = \emph{His illness kept him home} (noun phrase as subject).
\end{methode}

\begin{exoresolu}[Clause Analysis and Transformation]
\textbf{(a)} Identify the type and function of each underlined clause:
(i) The chief announced \emph{that the festival would be postponed}.
(ii) Students \emph{who revise regularly} rarely panic in examinations.
(iii) \emph{Unless the rains come early}, the harvest will be poor.

\textbf{(b)} Rewrite using the clause type in brackets without changing the meaning: ``The minister admitted that the project had failed.'' (adverb clause of concession)
\tcblower
\textbf{Detailed solution.}

\textbf{(a)(i)} \emph{that the festival would be postponed} — \textbf{noun clause}, functioning as the \emph{direct object} of ``announced''. \emph{Test:} ``The chief announced \textbf{something}.'' It can be replaced by a pronoun, so it is nominal. Note that ``that'' here is a pure subordinator and could be omitted: \emph{The chief announced the festival would be postponed}.

\textbf{(a)(ii)} \emph{who revise regularly} — \textbf{adjective (relative) clause}, qualifying the noun ``students''. It answers ``which students?'' It is \emph{defining}: without it, ``students rarely panic'' would make a false general claim; hence it is correctly written with \textbf{no commas}.

\textbf{(a)(iii)} \emph{Unless the rains come early} — \textbf{adverb clause of condition} (a negative condition: ``if\ldots not''), modifying ``will be''. It answers ``under what condition?'' Because it precedes the main clause, the comma after it is obligatory.

\textbf{(b)} The original contains a noun clause (object of ``admitted''). To express the same content with a \emph{concession}, we oppose the admission to an expectation:

\emph{Model answer:} ``\textbf{Although the project had been expected to succeed}, the minister admitted that it had failed.''

An alternative (noun clause removed entirely) — ``\textbf{Although it had failed}, the minister still defended the project publicly'' — is acceptable only if a slight shift of meaning is permitted; the first version is safer because it preserves every element of the original.

\textbf{Marking reflex:} for each item state (1) the clause type, (2) its grammatical function, (3) a replacement or question test. Three elements earn full marks.
\end{exoresolu}

\courssub{Skill 5 — Grammar Mastery: Dangling and Misplaced Modifiers}

\begin{definition}[Faulty Modifiers]
A \cle{modifier} is a word, phrase or clause that describes or limits another element. A \cle{dangling modifier} is an introductory phrase whose implied doer is not the subject of the main clause (\emph{Walking to school, the rain started}). A \cle{misplaced modifier} is positioned so that it appears to describe the wrong word (\emph{She served fufu to the guests on plastic plates}).
\end{definition}

\begin{methode}[Detecting and Correcting Faulty Modifiers]
\begin{enumerate}
\item \textbf{The Golden Rule:} a modifier must stand \emph{next to} the word it modifies, and an introductory phrase must describe the \emph{subject of the main clause}.
\item \textbf{Test for dangling modifiers:} when a sentence opens with a participle phrase (``Walking down the road,\ldots''), ask \emph{who is doing the walking?} If the subject after the comma cannot logically do it, the modifier \cle{dangles}.
\item \textbf{Test for misplaced modifiers:} can the phrase be read as describing the wrong noun? If yes, move it.
\item \textbf{Correct a dangling modifier} in one of two ways: (a) name the real doer as subject — \emph{Walking down the road, \textbf{I} saw an accident}; or (b) expand the phrase into a full clause — \emph{\textbf{While I was walking} down the road, I saw an accident}.
\item \textbf{Correct a misplaced modifier} by moving it beside its true noun: \emph{She served the guests fufu on plastic plates}.
\item \textbf{Beware ``only'', ``almost'', ``nearly'':} their position changes the meaning. \emph{I only lent him 500 francs} (I did not give) is not \emph{I lent him only 500 francs} (not more than 500).
\end{enumerate}
\end{methode}

\begin{exoresolu}[Correcting Modifier Errors]
Correct the following sentences and explain each correction.
(i) Flying over Mount Cameroon, the plantations stretched endlessly below us.
(ii) The teacher praised the boy with a smile who had won the scholarship.
(iii) Having finished the examination, the papers were collected by the invigilator.
\tcblower
\textbf{Detailed solution.}

\textbf{(i)} \emph{Fault:} dangling modifier. The opener ``Flying over Mount Cameroon'' must describe the subject ``the plantations'' — but plantations cannot fly. \emph{Correction (a), rename the doer:} ``Flying over Mount Cameroon, \textbf{we saw} the plantations stretch endlessly below us.'' \emph{Correction (b), full clause:} ``\textbf{As we flew} over Mount Cameroon, the plantations stretched endlessly below us.'' Both are acceptable; (b) keeps the original main clause intact.

\textbf{(ii)} \emph{Fault:} misplaced modifier. The relative clause ``who had won the scholarship'' sits next to ``a smile'', absurdly suggesting the smile won the scholarship. \emph{Correction:} move the clause beside its true antecedent: ``The teacher praised \textbf{the boy who had won the scholarship} with a smile.'' Now ``who\ldots'' clearly qualifies ``the boy'', and ``with a smile'' modifies ``praised''.

\textbf{(iii)} \emph{Fault:} dangling modifier (a classic GCE trap, with a passive main clause). ``Having finished the examination'' describes whoever finished it — the candidates — but the subject given is ``the papers''. \emph{Correction (a):} ``Having finished the examination, \textbf{the candidates} handed in their papers to the invigilator.'' \emph{Correction (b):} ``\textbf{When the candidates had finished} the examination, the papers were collected by the invigilator.''

\textbf{Reflex to memorise:} \emph{opener + comma + subject: the subject must perform the opener's action.} Apply this test mechanically to every sentence you write in Paper 2 essays; examiners penalise dangling modifiers as serious structural errors.
\end{exoresolu}

\courssub{Skill 6 — Phonetics: Transcription and Connected Speech}

\begin{definition}[Key Phonetic Terms]
A \cle{phoneme} is the smallest distinctive sound unit of a language; \cle{phonemic transcription} records one IPA symbol per sound between slant brackets / /. \cle{Elision} is the disappearance of a sound in connected speech; \cle{assimilation} is the modification of a sound under the influence of a neighbouring sound; \cle{intrusion} (or linking) is the insertion of an extra sound between words; \cle{catenation} is the smooth linking of a final consonant to a following initial vowel.
\end{definition}

\begin{methode}[Transcribing Words and Analysing Connected Speech]
\begin{enumerate}
\item \textbf{Transcribe sounds, not letters.} Say the word aloud in Standard British (Received Pronunciation) or educated Cameroonian English as specified; write one IPA symbol per \emph{sound}. \emph{Ship} = /ʃɪp/ (3 sounds, 4 letters).
\item \textbf{Master the trap symbols:} /ʃ/ (\emph{shoe}), /ʒ/ (\emph{vision}), /tʃ/ (\emph{church}), /dʒ/ (\emph{judge}), /θ/ (\emph{think}), /ð/ (\emph{this}), /ŋ/ (\emph{sing} — never /ŋg/ in RP), /j/ (\emph{yes}). Vowels: /iː/ versus /ɪ/, /uː/ versus /ʊ/, /ɜː/ (\emph{bird}), /ə/ (the schwa — the commonest English vowel), and the diphthongs /eɪ, aɪ, ɔɪ, əʊ, aʊ/.
\item \textbf{Mark the stress} with /ˈ/ before the stressed syllable: \emph{photograph} /ˈfəʊtəɡrɑːf/ but \emph{photography} /fəˈtɒɡrəfi/.
\item \textbf{Connected speech — four processes:}
\begin{itemize}
\item \cle{Elision}: a sound disappears — \emph{next day} /neks deɪ/ (/t/ elided).
\item \cle{Assimilation}: a sound changes under a neighbour's influence — \emph{good boy} /ɡʊb bɔɪ/ (/d/ becomes /b/ before /b/).
\item \cle{Intrusion (linking)}: a sound is added — \emph{law(r) and order} pronounced with linking /r/; \emph{go(w) out} with intrusive /w/.
\item \cle{Catenation}: a final consonant links to the next initial vowel — \emph{an apple} /əˈnæpəl/.
\end{itemize}
\item \textbf{Intonation:} a falling tone for statements, commands and wh-questions; a rising tone for yes/no questions, unfinished lists, politeness and doubt.
\end{enumerate}
\end{methode}

\begin{exoresolu}[Transcription and Connected Speech]
\textbf{(a)} Give the phonemic transcription of: \emph{judge, think, measure, singer, comfortable}.

\textbf{(b)} Name and illustrate the connected-speech process in: \emph{``Would you come in?''} and \emph{``last night''}.
\tcblower
\textbf{Detailed solution.}

\textbf{(a) Transcriptions (RP):}
\begin{itemize}
\item \emph{judge} — /dʒʌdʒ/. Both \emph{j} and \emph{-dge} give the same affricate /dʒ/; the vowel is short /ʌ/.
\item \emph{think} — /θɪŋk/. \emph{th} is voiceless /θ/ (not /ð/, which is voiced as in \emph{this}); \emph{n} before /k/ is realised as /ŋ/.
\item \emph{measure} — /ˈmeʒə/. The \emph{s} here is the voiced fricative /ʒ/, not /s/ or /z/; final \emph{-ure} reduces to schwa /ə/.
\item \emph{singer} — /ˈsɪŋə/. Trap: no /g/ is pronounced — /ŋ/ alone. Compare \emph{finger} /ˈfɪŋɡə/, where /g/ \emph{is} sounded.
\item \emph{comfortable} — /ˈkʌmftəbəl/ in natural speech. Note the compression: four written syllables are typically reduced to three spoken syllables, the second vowel being elided.
\end{itemize}

\textbf{(b) Connected speech.}

\emph{``Would you come in?''} — natural pronunciation /wʊdʒu kʌm ɪn/. Two processes: \textbf{assimilation} — /d/ + /j/ coalesce into /dʒ/ (\emph{would you} becomes /ˈwʊdʒu/), a case of coalescent assimilation; and \textbf{catenation} in \emph{come in} — the final /m/ links directly to the following vowel /ɪ/, so the two words flow as one unit.

\emph{``last night''} — /lɑːs naɪt/. \textbf{Elision}: the /t/ of \emph{last} is dropped between two consonants (the /st n/ cluster is simplified). The spelling keeps the \emph{t}, but fluent speech removes it; transcribing /lɑːst naɪt/ would mark an unnaturally careful pronunciation.

\textbf{Exam reflex:} always (1) name the process, (2) quote the sounds involved, (3) show the result in transcription. That three-part answer earns full credit.
\end{exoresolu}

\courssub{Skill 7 — Writing the Balanced Argumentative Essay}

\begin{methode}[Structuring a Balanced Argumentative Essay (Paper 2)]
\begin{enumerate}
\item \textbf{Decode the topic.} Identify the two positions in debate (``Boarding schools are better than day schools. Discuss.''). Your thesis must take a clear stand \emph{after} weighing both sides.
\item \textbf{Plan before writing} (about five minutes): list three arguments for and three against; select the strongest; decide your final position.
\item \textbf{Structure (five paragraphs minimum):}
\begin{itemize}
\item \emph{Introduction:} hook + background + \cle{thesis statement} announcing your stand.
\item \emph{Body 1--2:} the side you \emph{reject}, presented fairly (steelman it) — one idea per paragraph, each with explanation and a concrete (Cameroonian) example.
\item \emph{Body 3--4:} your own side, stronger and more developed; rebut the opposing points (``Admittedly\ldots however\ldots'').
\item \emph{Conclusion:} restate your position in fresh words; end with a recommendation or forward-looking sentence — never introduce a new argument.
\end{itemize}
\item \textbf{Signpost with discourse markers:} staging (\emph{To begin with, Furthermore, Finally}), concession (\emph{Admittedly, Granted that}), rebuttal (\emph{Nevertheless, Yet the evidence shows}), conclusion (\emph{In the final analysis}).
\item \textbf{One paragraph = one idea + explanation + example + link.} A paragraph without an example is an assertion, not an argument.
\item \textbf{Proofread} for the three killers: subject--verb agreement, tense consistency, and dangling modifiers.
\end{enumerate}
\textbf{Length reflex:} GCE expects roughly 450--500 words; quality of argument and accuracy of language outweigh length.
\end{methode}

\begin{exoresolu}[Argumentative Essay — Guided Model]
\textbf{Topic:} ``Mobile phones should be banned in all secondary schools in Cameroon.'' Write a balanced argumentative essay (the introduction and one body paragraph are modelled in full; the plan shows the rest).
\tcblower
\textbf{Detailed solution.}

\textbf{Step 1 — Plan.} \emph{For the ban:} distraction in class; cheating in tests; cyber-bullying. \emph{Against the ban:} emergency contact with parents; access to digital learning tools; preparing youths for a digital economy. \emph{My stand:} regulate rather than ban.

\textbf{Step 2 — Model introduction.}

\emph{``In the courtyard of almost every secondary school in Yaound\'e or Bamenda, the glow of small screens competes with the blackboard for the attention of learners. Alarmed by falling concentration and examination malpractice, some educators demand that mobile phones be banned outright from school premises. While the concern behind this demand is legitimate, a total prohibition would punish responsible students and deprive schools of a powerful learning tool. This essay argues that regulation, not prohibition, is the wiser policy.''}

\emph{Why it works:} a concrete Cameroonian hook; both sides named fairly; an explicit thesis (``regulation, not prohibition'') in the final sentence — exactly what examiners reward.

\textbf{Step 3 — Model body paragraph (opposing side, steelmanned).}

\emph{``Admittedly, the case for a ban is not without force. A phone vibrating in a pocket during a mathematics lesson fragments the attention not only of its owner but of the whole class, and cases of candidates photographing question papers have shown how easily the device becomes an instrument of fraud. Schools, one may grant, have a duty to protect the integrity of learning.''}

\emph{Why it works:} a concession marker (``Admittedly''), two genuine arguments, a fair tone — setting up the rebuttal to come.

\textbf{Step 4 — Remaining skeleton.} \emph{Paragraph 3 (rebuttal):} ``Nevertheless, a ban treats the symptom rather than the disease\ldots'' — argue that disciplined use (phones deposited at the staff room, used only for supervised research) solves the problem without sacrificing benefits; example: a biology class using phones to identify plants on a nature walk. \emph{Paragraph 4:} the digital-skills argument — Cameroonian employers increasingly demand them. \emph{Conclusion:} ``In the final analysis, schools should teach the responsible use of the tool rather than pretend it does not exist\ldots''

\textbf{Marking checklist applied:} clear thesis, balanced treatment, discourse markers, one idea per paragraph, Cameroonian examples, no dangling modifiers, conclusion without new arguments.
\end{exoresolu}

\courssub{Skill 8 — Critical Reasoning: Facts, Opinions and the Evaluation of Arguments}

\begin{definition}[Fact, Opinion, Fallacy]
A \cle{fact} is a statement that can be verified or falsified by evidence; an \cle{opinion} expresses a judgement, feeling or prediction. An \cle{argument} consists of a \cle{conclusion} (what the writer wants you to accept) supported by \cle{premises} (the reasons offered). A \cle{fallacy} is a flaw in reasoning that makes an argument invalid or unconvincing, however persuasively it is worded.
\end{definition}

\begin{methode}[Evaluating Arguments and Evidence in Persuasive Texts]
\begin{enumerate}
\item \textbf{Separate fact from opinion.} Signal words of opinion include \emph{should, must, undoubtedly, the best, I believe}. Check whether each claim could in principle be proved or disproved.
\item \textbf{Beware the disguised opinion:} ``Everyone knows that\ldots'' and ``It is obvious that\ldots'' often introduce the \emph{least} proven claims.
\item \textbf{Identify the argument's skeleton:} conclusion + premises. Ask: does the conclusion genuinely \emph{follow} from the premises?
\item \textbf{Test the evidence:} Is it relevant? Sufficient? Recent? From a credible, unbiased source? One anecdote (``my uncle smoked and lived to ninety'') does not establish a general truth.
\item \textbf{Spot the classic fallacies:} \emph{ad hominem} (attacking the person, not the argument); \emph{false dilemma} (``either we ban phones or education collapses''); \emph{hasty generalisation} (a conclusion from too few cases); \emph{slippery slope} (one step must lead to disaster); \emph{appeal to emotion} replacing evidence.
\item \textbf{Write your evaluation} in three moves: state the writer's claim, assess the evidence offered, then pronounce the argument strong, weak or partly convincing — with reasons.
\end{enumerate}
\end{methode}

\begin{exoresolu}[Evaluating a Persuasive Argument]
Read the extract, then (a) distinguish the facts from the opinions, and (b) evaluate the strength of the argument.

\emph{``Every parent in this country knows that private schools are simply better than public schools. My neighbour's daughter attended a private college in Douala and scored four A grades at the GCE Advanced Level; she is now studying medicine. It is obvious, therefore, that any government serious about education should close the public schools and hand the entire system over to private proprietors, who alone understand quality.''}
\tcblower
\textbf{Detailed solution.}

\textbf{(a) Facts versus opinions.} The only verifiable \emph{facts} are the anecdotal details: the neighbour's daughter attended a private college, obtained four A grades, and is studying medicine — assuming they are true, they can in principle be checked. Everything else is \emph{opinion}: ``every parent knows'' (an unverified universal claim), ``private schools are simply better'' (a judgement presented without data), ``it is obvious'' (a signal phrase masking an unproven claim), and the recommendation to close public schools (a value judgement about policy).

\textbf{(b) Evaluation of the argument.}

\emph{Skeleton:} Premise — one private-school pupil succeeded brilliantly. Conclusion — the entire public system should be abolished and replaced by private schools.

\emph{Weakness 1 — hasty generalisation.} A single success story, however genuine, cannot establish the superiority of a whole sector; many public-school candidates also earn top grades every year, and private schools also produce failures.

\emph{Weakness 2 — false dilemma.} The writer presents only two options (close public schools or be unserious about education), ignoring middle paths: reforming public schools, regulating both sectors, or partnership between them.

\emph{Weakness 3 — appeal to popular opinion.} ``Every parent knows'' substitutes supposed consensus for evidence; what ``everyone knows'' has often turned out to be false.

\emph{Weakness 4 — suppressed evidence.} No comparison of costs, access for poor families, or examination results across both sectors is offered — precisely the data needed to judge the claim.

\emph{Verdict (model sentence):} ``The argument is therefore weak: its conclusion is sweeping while its evidence is a single anecdote, and it rests on a hasty generalisation and a false dilemma rather than on comparative data.''

\textbf{Exam reflex:} name each fallacy, quote the words that commit it, and explain \emph{why} it fails — three elements per weakness identified.
\end{exoresolu}

\newpage

\coursec{Exercises}

\courssub{I check my knowledge}

\begin{exercice}[MCQ: Reading and responding to texts]
For each question, choose the correct option (A, B, C or D).
\begin{enumerate}
\item The \cle{main idea} of a passage is:
\begin{itemize}
\item[\textbf{A.}] the first sentence of the passage;
\item[\textbf{B.}] the central point the writer wants the reader to retain;
\item[\textbf{C.}] any interesting detail found in the middle paragraph;
\item[\textbf{D.}] the writer's final word.
\end{itemize}
\item \cle{Implicit meaning} is meaning that is:
\begin{itemize}
\item[\textbf{A.}] stated directly in the text;
\item[\textbf{B.}] suggested but not directly expressed;
\item[\textbf{C.}] written in brackets;
\item[\textbf{D.}] found only in the title.
\end{itemize}
\item A writer's \cle{tone} is revealed mainly through:
\begin{itemize}
\item[\textbf{A.}] the length of the paragraphs;
\item[\textbf{B.}] the choice of words and sentence patterns;
\item[\textbf{C.}] the number of characters;
\item[\textbf{D.}] the font used.
\end{itemize}
\item In \emph{The Old Man and the Sea}, the sea is best described as:
\begin{itemize}
\item[\textbf{A.}] a mere background;
\item[\textbf{B.}] a symbol of life, struggle and dignity;
\item[\textbf{C.}] an enemy to be destroyed;
\item[\textbf{D.}] a comic device.
\end{itemize}
\end{enumerate}
\end{exercice}

\begin{exercice}[True or False — justify your answer]
Say whether each statement is \textbf{true} or \textbf{false}, and justify your answer in one or two sentences.
\begin{enumerate}
\item A summary of a passage should include the writer's examples and minor details.
\item A dangling modifier is one whose implied subject is not clearly stated in the sentence.
\item In connected speech, assimilation occurs when a sound changes to become more like a neighbouring sound.
\item A fact can be verified; an opinion expresses a personal belief or judgement.
\item The second conditional is used for real, likely future situations.
\end{enumerate}
\end{exercice}

\begin{exercice}[Definitions]
Define each of the following terms in one clear sentence, and give one example of each.
\begin{enumerate}
\item Noun clause
\item Irony
\item Synecdoche
\item Discourse marker
\item Denotation and connotation
\item Intonation
\end{enumerate}
\end{exercice}

\begin{exercice}[Quick identification drill]
\begin{enumerate}
\item Underline the subordinate clause in each sentence and state its type (noun, adjective or adverb clause):
\begin{enumerate}
\item What the minister said surprised everyone.
\item The novel that we studied last term is set in Sierra Leone.
\item Although the rains were late, the farmers of the North West Region planted their maize.
\end{enumerate}
\item Transcribe the following words into IPA: \emph{think, measure, ship, judge, though}.
\end{enumerate}
\end{exercice}

\courssub{I practise}

\begin{exercice}[Fact or opinion?]
Classify each of the following statements, taken from a persuasive article on social media use, as a \textbf{fact} or an \textbf{opinion}. Give a brief reason for each classification.
\begin{enumerate}
\item Many Cameroonian students own smartphones.
\item Social media is the greatest enemy of serious study ever invented.
\item Some universities in Cameroon use online platforms to publish results.
\item Any parent who allows a child unlimited internet access is irresponsible.
\item WhatsApp groups are sometimes used by teachers to share past GCE questions.
\end{enumerate}
\end{exercice}

\begin{exercice}[Correcting dangling and misplaced modifiers]
Each sentence below contains a dangling or misplaced modifier. Rewrite it correctly.
\begin{enumerate}
\item Walking to school in the rain, my books got soaked.
\item The principal praised the students in his office who had passed the examination.
\item Having finished the examination, the papers were collected by the invigilator.
\item She almost sold all her vegetables at the Mokolo market.
\item Flying over Yaoundé, the hills looked beautiful.
\end{enumerate}
\end{exercice}

\begin{exercice}[Grammar transformation]
Rewrite each sentence according to the instruction in brackets, without changing the meaning.
\begin{enumerate}
\item The boy was tired. He had been reading all night. (Join using a perfect continuous tense.)
\item I did not study hard, so I failed. (Rewrite as a third conditional beginning \emph{If I...}.)
\item The teacher asked me a question. I could not answer it. (Join using an adjective clause.)
\item ``Work hard, and you will succeed,'' the master said. (Report using a first conditional.)
\item People believe that the new library will open soon. (Rewrite beginning \emph{It is believed...} using a noun clause.)
\end{enumerate}
\end{exercice}

\begin{exercice}[Figures of speech]
Identify the figure of speech used in each sentence (irony, paradox, hyperbole, synecdoche, metonymy, understatement, sarcasm):
\begin{enumerate}
\item ``What a lovely day!'' she said, as the thunderstorm destroyed her stall.
\item I have told you a million times to close the door.
\item The child is father of the man.
\item All hands were on deck during the flood in Douala.
\item The Ministry announced new measures to support farmers. (\emph{The Ministry} standing for its officials.)
\item Losing one mark in the mock examination, he declared it was the end of the world; his friend replied that it was ``a slight inconvenience''.
\end{enumerate}
\end{exercice}

\courssub{I prepare for the GCE}

\begin{exercice}[GCE Paper 1 style — Comprehension and Summary (25 marks)]
Read the following passage carefully and answer the questions that follow.

\textbf{The Weight of Waiting}

In the towns and cities of Cameroon, a familiar scene repeats itself every morning. Young men and women, certificates carefully folded in plastic folders, gather at cybercafés to print yet another application letter, or sit on verandas scanning newspapers whose job pages grow thinner each year. They are the graduates of a system that promised them that education was the key, only to hand them a key that fits no visible door.

The figures commonly discussed in public debate are alarming, but behind every statistic stands a face. Take Amina, a twenty-six-year-old from Bamenda, who graduated with a degree in management four years ago. Since then, she has sold second-hand clothes at the Commercial Avenue, sent over a hundred applications, and attended three interviews that ended with the same polite silence. ``I am not lazy,'' she says, ``I am simply waiting for a door that was promised to me.'' Her words carry the quiet dignity of a generation that refuses to beg but is tired of knocking.

Some observers blame the graduates themselves. They argue that too many young people chase white-collar jobs while despising agriculture and the trades that built their parents' lives. There is some truth in this: the son of a prosperous cocoa farmer in the Centre Region may prefer unemployment in Yaoundé to prosperity in the village. Yet this argument, pushed too far, becomes a convenient excuse. It shifts the whole burden onto the shoulders of the young, while ignoring an economy that produces too few jobs and an education system that trains minds for examinations rather than for enterprise.

Others point to the informal sector, which already absorbs the majority of young workers. The \emph{bendskin} rider, the call-box operator, the tailor in a roadside workshop — these are not idle people. They work long hours for uncertain rewards, without contracts, insurance or pensions. To celebrate their ``entrepreneurship'' without asking why formality remains out of reach is to mistake survival for success.

What, then, is to be done? The solutions proposed are as numerous as the commentators: vocational training, easier credit for small businesses, honest public recruitment, and a change of mentality among the youth. None of these is wrong; none is sufficient alone. Perhaps the first step is simply to listen — really to listen — to the Aminas of this country, whose patience, contrary to what many believe, is not unlimited.

\begin{enumerate}
\item What does the writer mean by ``a key that fits no visible door'' (paragraph 1)? \textbf{(2 marks)}
\item Give two details from the passage that show that Amina's situation is not due to laziness. \textbf{(2 marks)}
\item In your own words, explain the argument of those who ``blame the graduates themselves'' and the writer's objection to it. \textbf{(4 marks)}
\item What is implied, but not directly stated, in the sentence ``their patience... is not unlimited'' (paragraph 5)? \textbf{(2 marks)}
\item Identify the writer's tone in paragraph 4 and justify your answer with two expressions from the paragraph. \textbf{(3 marks)}
\item Is the statement ``the informal sector already absorbs the majority of young workers'' presented as fact or opinion? Justify. \textbf{(2 marks)}
\item Summarise the passage in about \textbf{80 words}, using your own words as far as possible. \textbf{(10 marks)}
\end{enumerate}
\end{exercice}

\begin{exercice}[GCE Paper 2 style — Literature Essay (20 marks)]
Answer \textbf{one} of the following questions in an essay of 450--500 words. Your essay must have a clear introduction, developed paragraphs with close reference to the text, and a conclusion.
\begin{enumerate}
\item \emph{``The sea is more than a setting in The Old Man and the Sea.''} Discuss this statement with close reference to Hemingway's use of symbolism and style.
\item \emph{``How does Soyinka use the character of the Kadiye to expose corruption in The Swamp Dwellers?''} Answer with reference to language and dramatic technique.
\item Compare the treatment of dignity and resistance in Claude McKay's \emph{``If We Must Die''} and Roland Tombekai Dempster's \emph{``Africa's Plea''}.
\end{enumerate}
\end{exercice}

\begin{exercice}[GCE Paper 2 style — Guided Writing (20 marks)]
Your school recorded a sharp decline in English Language results at the last GCE Advanced Level. As Senior Prefect, you carried out a survey among 120 Upper Sixth students. Your findings are summarised below:
\begin{itemize}
\item 68\% of respondents say they rarely read English books outside class;
\item 55\% complain that there are not enough past questions and marked scripts available for practice;
\item 47\% admit they speak Pidgin or French during English lessons;
\item 30\% mention overcrowded classrooms and shortage of prescribed texts;
\item 72\% request a debate club and regular essay clinics.
\end{itemize}
Write a \textbf{formal report} to your Principal, presenting these findings and making at least three recommendations. Use appropriate report structure (title, terms of reference, procedure, findings, conclusions, recommendations) and impersonal, formal language. \textbf{(20 marks)}
\end{exercice}

\newpage

\coursec{Solutions to the exercises}

\begin{corrige}[Exercise 1 — MCQ: Reading and responding to texts]
\begin{enumerate}
\item \textbf{B.} The main idea is the central point the writer wants the reader to retain. The first sentence may contain it, but it can appear anywhere in the passage or be implied; supporting details and final words are not necessarily the main idea.
\item \textbf{B.} Implicit meaning is suggested but not directly expressed; the reader must infer it from textual clues such as word choice, context, examples, and rhetorical structure.
\item \textbf{B.} Tone is revealed mainly through the choice of words (diction) and sentence patterns (syntax). Paragraph length, character names, and font style do not determine tone.
\item \textbf{B.} In \emph{The Old Man and the Sea}, the sea (\emph{la mar}) is a symbol of life, struggle and dignity: it is the arena of Santiago's test, the source of his livelihood, and a feminine presence he loves and respects — not a mere background, an enemy, or a comic device.
\end{enumerate}
\end{corrige}

\begin{corrige}[Exercise 2 — True or False]
\begin{enumerate}
\item \textbf{False.} A summary retains only the essential ideas; examples, illustrations and minor details must be left out, otherwise the summary becomes a copy rather than a condensation.
\item \textbf{True.} A dangling modifier is a phrase (often a participle phrase) whose implied subject does not appear in the sentence, e.g. \emph{Walking to school, the rain soaked me} — the rain was not walking.
\item \textbf{True.} Assimilation is precisely the process by which a sound becomes more like a neighbouring sound, e.g. \emph{good boy} pronounced /gʊb bɔɪ/, where /d/ becomes /b/ under the influence of the following /b/.
\item \textbf{True.} A fact can be checked or verified against evidence (dates, figures, observable events); an opinion expresses a belief, judgement or feeling and cannot be proved true or false.
\item \textbf{False.} The second conditional (\emph{If + past simple, ... would + infinitive}) is used for unreal, imaginary or unlikely present/future situations. Real, likely future situations take the first conditional (\emph{If + present simple, ... will...}).
\end{enumerate}
\end{corrige}

\begin{corrige}[Exercise 3 — Definitions]
\begin{enumerate}
\item \textbf{Noun clause:} a subordinate clause that functions as a noun in the sentence (as subject, object or complement). \emph{Example: What she said surprised me} (subject).
\item \textbf{Irony:} a figure of speech in which the intended meaning is the opposite of, or strikingly different from, the literal meaning of the words used. \emph{Example: saying ``What lovely weather!'' during a violent storm.}
\item \textbf{Synecdoche:} a figure of speech in which a part stands for the whole, or the whole for a part. \emph{Example: ``All hands on deck,'' where \emph{hands} stands for the sailors.}
\item \textbf{Discourse marker:} a word or phrase that organises speech or writing by signalling staging, sequencing, addition, contrast or conclusion. \emph{Example: \emph{firstly, however, in addition, to conclude}.}
\item \textbf{Denotation and connotation:} denotation is the literal, dictionary meaning of a word, while connotation is the set of associations and emotional overtones it carries. \emph{Example: \emph{home} denotes a place of residence but connotes warmth, family and security.}
\item \textbf{Intonation:} the rise and fall of the voice in speech, which expresses grammatical function and the speaker's attitude. \emph{Example: a rising tone in ``You're coming?'' signals a question or surprise.}
\end{enumerate}
\end{corrige}

\begin{corrige}[Exercise 4 — Quick identification drill]
\textbf{1. Subordinate clauses}
\begin{enumerate}
\item \underline{What the minister said} surprised everyone. — \textbf{Noun clause} (subject of the verb \emph{surprised}).
\item The novel \underline{that we studied last term} is set in Sierra Leone. — \textbf{Adjective (relative) clause}, qualifying the noun \emph{novel}.
\item \underline{Although the rains were late}, the farmers of the North West Region planted their maize. — \textbf{Adverb clause} of concession, modifying the main verb \emph{planted}.
\end{enumerate}
\textbf{2. IPA transcription (Received Pronunciation / Standard British model used in GCE)}
\begin{center}
\begin{tabular}{|l|l|}
\hline
\rowcolor{popblueL} \textbf{Word} & \textbf{Transcription} \\
\hline
think & /θɪŋk/ \\
\hline
measure & /ˈmeʒə/ \\
\hline
ship & /ʃɪp/ \\
\hline
judge & /dʒʌdʒ/ \\
\hline
though & /ðəʊ/ \\
\hline
\end{tabular}
\end{center}
Note the contrast between the voiceless dental fricative /θ/ in \emph{think} and the voiced /ð/ in \emph{though}, and the short vowel /ɪ/ in \emph{ship} (not the long /iː/ of \emph{sheep}). The schwa /ə/ appears in the weak syllable of \emph{measure} and the final position of \emph{though}.
\end{corrige}

\begin{corrige}[Exercise 5 — Fact or opinion?]
\begin{enumerate}
\item \textbf{Fact.} Ownership of smartphones by many students can be verified by surveys or observation.
\item \textbf{Opinion.} ``The greatest enemy of serious study ever invented'' is an exaggerated personal judgement that cannot be verified.
\item \textbf{Fact.} Whether universities publish results on online platforms can be checked directly on the platforms concerned.
\item \textbf{Opinion.} Calling such parents ``irresponsible'' is a value judgement, not a verifiable statement.
\item \textbf{Fact.} The use of WhatsApp groups by teachers to share past GCE questions can be verified by examining the groups themselves.
\end{enumerate}
\textbf{Examiner's tip:} words such as \emph{greatest, best, irresponsible, should, must} are typical markers of opinion; verifiable data, observable practices and checkable events signal facts.
\end{corrige}

\begin{corrige}[Exercise 6 — Correcting dangling and misplaced modifiers]
\begin{enumerate}
\item \emph{Walking to school in the rain, I got my books soaked.} (Or: \emph{As I was walking to school in the rain, my books got soaked.}) — The implied subject of \emph{walking} must be \emph{I}, not \emph{my books}.
\item \emph{In his office, the principal praised the students who had passed the examination.} — The misplaced phrase \emph{in his office} wrongly suggested the students were in his office.
\item \emph{Having finished the examination, the candidates handed in their papers, which were collected by the invigilator.} (Or: \emph{When the examination was finished, the papers were collected by the invigilator.}) — The papers did not finish the examination; the candidates did.
\item \emph{She sold almost all her vegetables at the Mokolo market.} — \emph{Almost} must modify \emph{all}, not \emph{sold}; she did not ``almost sell''.
\item \emph{As we flew over Yaoundé, the hills looked beautiful.} (Or: \emph{Flying over Yaoundé, we saw that the hills looked beautiful.}) — The hills were not flying; the observers were.
\end{enumerate}
\textbf{Rule applied:} an introductory modifying phrase must logically refer to the subject of the main clause; a modifier must be placed as close as possible to the word it modifies.
\end{corrige}

\begin{corrige}[Exercise 7 — Grammar transformation]
\begin{enumerate}
\item \emph{The boy was tired because he had been reading all night.} — The past perfect continuous (\emph{had been reading}) shows the ongoing activity that preceded and explained the tiredness.
\item \emph{If I had studied hard, I would not have failed.} — Third conditional: \emph{If + past perfect, ... would have + past participle}, expressing an unreal past situation and its imagined result.
\item \emph{The teacher asked me a question which/that I could not answer.} — The adjective (relative) clause \emph{which I could not answer} modifies \emph{a question}.
\item \emph{The master said that if we worked hard, we would succeed.} (Accept also direct first conditional: \emph{``If you work hard, you will succeed,'' the master said.}) — In reported speech the first conditional shifts back: \emph{work} $\rightarrow$ \emph{worked}, \emph{will} $\rightarrow$ \emph{would}.
\item \emph{It is believed that the new library will open soon.} — The noun clause \emph{that the new library will open soon} is the extraposed subject after the impersonal passive \emph{It is believed}.
\end{enumerate}
\end{corrige}

\begin{corrige}[Exercise 8 — Figures of speech]
\begin{enumerate}
\item \textbf{Irony} (verbal irony): she says the day is ``lovely'' while meaning the exact opposite, as the storm destroys her stall.
\item \textbf{Hyperbole}: ``a million times'' is a deliberate exaggeration for emphasis; no one has literally counted a million warnings.
\item \textbf{Paradox}: ``The child is father of the man'' seems self-contradictory, yet it expresses the truth that childhood shapes the adult.
\item \textbf{Synecdoche}: \emph{hands} (a part) stands for the whole persons — the rescuers or crew — during the flood.
\item \textbf{Metonymy}: \emph{The Ministry} (an institution/associated name) stands for the officials who actually made the announcement.
\item \textbf{Sarcasm and understatement}: the boy's ``the end of the world'' is hyperbolic despair, and the friend's ``a slight inconvenience'' is an understatement used sarcastically to mock the overreaction.
\end{enumerate}
\end{corrige}

\begin{corrige}[Exercise 9 — GCE Paper 1 style: Comprehension and Summary (25 marks)]
\textbf{1. ``A key that fits no visible door'' (2 marks).} The metaphor means that education was presented as the key that would open the door of employment, but graduates find no jobs to unlock: their qualifications lead nowhere, because suitable employment simply does not exist for them. \emph{(1 mark for identifying the metaphor education = key / jobs = door; 1 mark for explaining the failure to find employment.)}

\textbf{2. Two details showing Amina is not lazy (2 marks).} Any two of: she sells second-hand clothes at the Commercial Avenue (she works); she has sent over a hundred applications; she has attended three interviews. \emph{(1 mark each.)}

\textbf{3. The argument and the writer's objection (4 marks).} The critics argue that graduates are responsible for their own unemployment because they chase white-collar jobs and despise agriculture and manual trades, preferring idleness in town to prosperity in the village \emph{(2 marks)}. The writer concedes a grain of truth but objects that, pushed too far, this becomes a convenient excuse: it puts all the blame on the young while ignoring an economy that creates too few jobs and an education system that trains students for examinations rather than for enterprise \emph{(2 marks)}.

\textbf{4. Implied meaning of ``their patience... is not unlimited'' (2 marks).} It implies that if nothing changes, the young people's patience may run out — frustration could turn into anger, protest or social unrest; the warning is left unspoken to make it more powerful. \emph{(2 marks for the idea of possible explosion/revolt; 1 mark for a vague ``they will be tired''.)}

\textbf{5. Tone in paragraph 4 (3 marks).} The tone is critical/indignant (also sympathetic towards informal workers). Justification with any two expressions: ``these are not idle people'' (defensive, respectful); ``to mistake survival for success'' (biting criticism of complacent commentators); ``without contracts, insurance or pensions'' (indignant enumeration of injustice). \emph{(1 mark for the tone; 1 mark per justified expression.)}

\textbf{6. Fact or opinion? (2 marks).} It is presented as a \textbf{fact}: the writer asserts it as an objective, verifiable statement about the labour market (``which already absorbs the majority of young workers''), not as a personal judgement; in principle it could be checked against employment data. \emph{(1 mark for ``fact''; 1 mark for justification.)}

\textbf{7. Summary in about 80 words (10 marks).} \emph{Model answer:}

Many Cameroonian graduates remain unemployed despite their qualifications, like Amina, who works and applies tirelessly without success. Some blame the youths for rejecting agriculture and trades, but this ignores an economy that creates too few jobs and an exam-oriented education system. The informal sector keeps most young people surviving rather than prospering. Proposed solutions — vocational training, credit, fair recruitment, changed mentalities — are each insufficient alone; the first step is genuinely listening to the patient but increasingly frustrated youth. \emph{(78 words.)}

\textbf{Indicative mark scheme for the summary:} content points (5 marks): graduate unemployment illustrated by Amina (1); rejection of the ``blame the youth'' argument with its two counter-arguments (2); informal sector = survival, not success (1); multiple partial solutions + need to listen (1). Expression (5 marks): own words, no lifting (2); grammatical accuracy (2); respect of the word limit and fluency (1).
\end{corrige}

\begin{corrige}[Exercise 10 — GCE Paper 2 style: Literature Essay (20 marks)]
\textbf{General mark scheme (all three questions):} introduction with a clear thesis (3 marks); development — three or four organised paragraphs with close, accurate textual reference (10 marks); quality of analysis (symbolism, style, technique — not mere narration) (4 marks); conclusion and language accuracy (3 marks).

\textbf{Question 1 — Indicative content (\emph{The Old Man and the Sea}).}
\begin{itemize}
\item \textbf{Introduction:} the sea is a living presence and symbol, not a backdrop.
\item \textbf{Symbolism:} the sea as life itself — beautiful and cruel, giving and taking; Santiago calls it \emph{la mar} (feminine, loved), unlike the younger fishermen's hostile \emph{el mar}; the marlin as noble brother and worthy opponent; the sharks as destructive forces; the struggle on the sea mirrors the human condition — dignity in defeat (``A man can be destroyed but not defeated'').
\item \textbf{Style:} Hemingway's simple, spare prose, short declarative sentences and repetition (the ``iceberg principle'') give the sea a biblical, parable-like resonance; interior monologue and dialogue with the fish and the sea personalise it.
\item \textbf{Conclusion:} the sea is the true stage of Santiago's spiritual test — setting, character and symbol at once.
\end{itemize}

\textbf{Question 2 — Indicative content (\emph{The Swamp Dwellers}).}
\begin{itemize}
\item \textbf{Introduction:} the Kadiye, the fat priest of the Serpent, embodies religious and economic corruption.
\item \textbf{Exposure through character:} his physical grossness (Soyinka's stage directions stress his bulk) contrasts with the emaciated villagers he exploits; he controls the villagers through fear and superstition, demanding sacrifices and tribute while the swamp ruins their harvests.
\item \textbf{Language:} his pompous, evasive rhetoric versus the plain speech of Makuri and Alu; the confrontation with Igwezu, whose blunt accusations strip the priest of mystique.
\item \textbf{Dramatic technique:} the tense confrontation scene, dramatic irony (the audience sees what the villagers refuse to see), and the Kadiye's humiliation symbolise the collapse of corrupt traditional authority allied with a corrupt modern city (the beggar's revelations).
\item \textbf{Conclusion:} through the Kadiye, Soyinka dramatises how religion can be weaponised to exploit a suffering community.
\end{itemize}

\textbf{Question 3 — Indicative content (comparative poetry).}
\begin{itemize}
\item \textbf{Introduction:} both poems confront oppression, but with different strategies of dignity and resistance.
\item \textbf{McKay, ``If We Must Die'':} the Shakespearean sonnet form dignifies the oppressed; the extended metaphor of the ``mad and hungry dogs'' dehumanises the oppressors, not the victims; the call to die ``like men'' — fighting back so that ``even the monsters we defy / Shall be constrained to honor us though dead'' — turns certain death into honourable resistance; the collective ``we'' and the martial tone.
\item \textbf{Dempster, ``Africa's Plea'':} resistance takes the form of a plea — dignified appeal rather than armed defiance; personified Africa speaks with pathos and pride, asserting identity and demanding justice; the tone is supplicatory yet proud, refusing self-abasement.
\item \textbf{Comparison:} McKay's defiance is militant and confrontational, Dempster's is moral and petitionary; both reject humiliation and affirm the humanity and honour of the oppressed.
\item \textbf{Conclusion:} two complementary modes of resistance — the sword and the voice — both rooted in dignity.
\end{itemize}
\emph{Examiner's expectation:} candidates who merely retell plots score in the lower band; top-band answers quote or closely echo the texts, analyse technique, and (for Q3) sustain a genuine comparison rather than treating the two poems separately.
\end{corrige}

\begin{corrige}[Exercise 11 — GCE Paper 2 style: Guided Writing — Formal Report (20 marks)]
\textbf{Mark scheme:} format and structure — title, terms of reference, procedure, findings, conclusions, recommendations, signature block (5 marks); accurate presentation of all the survey data (5 marks); at least three relevant, clearly formulated recommendations (4 marks); formal, impersonal register, correct grammar and cohesion (6 marks).

\textbf{Model answer:}

\begin{center}
\textbf{REPORT ON THE CAUSES OF THE DECLINE IN GCE ADVANCED LEVEL ENGLISH LANGUAGE RESULTS AND PROPOSED REMEDIES}
\end{center}

\textbf{To:} The Principal\par
\textbf{From:} The Senior Prefect\par
\textbf{Date:} 15 November 2024

\textbf{1. Terms of Reference}\par
Following the sharp decline in English Language results at the last GCE Advanced Level, I was mandated to investigate the causes of the decline among Upper Sixth students and to propose remedies.

\textbf{2. Procedure}\par
A structured questionnaire was administered to 120 Upper Sixth students in October. The responses were collated and analysed, and the main findings are presented below.

\textbf{3. Findings}
\begin{itemize}
\item 68\% of respondents stated that they rarely read English books outside class;
\item 55\% complained that past questions and marked scripts are not sufficiently available for practice;
\item 47\% admitted speaking Pidgin or French during English lessons;
\item 30\% mentioned overcrowded classrooms and a shortage of prescribed texts;
\item 72\% requested the creation of a debate club and regular essay clinics.
\end{itemize}

\textbf{4. Conclusions}\par
The decline appears to result from a combination of poor reading habits, inadequate practice materials, insufficient use of English during lessons, and unfavourable learning conditions. Encouragingly, the students themselves show a strong willingness to participate in remedial activities.

\textbf{5. Recommendations}
\begin{enumerate}
\item The school library should be stocked with the prescribed texts and past GCE questions, and marked scripts should be made available for consultation.
\item An English-only policy should be enforced during English lessons.
\item A debate club and weekly essay clinics should be established, in response to the demand expressed by 72\% of respondents.
\item The administration should examine ways of decongesting overcrowded classrooms.
\end{enumerate}

It is hoped that these measures, if implemented promptly, will reverse the decline before the next examination session.\par
\emph{Signature}\par
Senior Prefect

\emph{Examiner's expectations:} impersonal constructions (``it was found that...'', ``respondents stated that...''), no slang or contractions, logical section headings, and recommendations that flow naturally from the findings. Reports written as letters or essays lose the format marks.
\end{corrige}

\newpage

\coursec{Summary sheet}

\begin{popfigure}
\begin{center}
\begin{tikzpicture}[mindmap, grow cyclic, every node/.style={concept, text=white, font=\small\bfseries},
root concept/.append style={concept color=popink, font=\small\bfseries},
level 1/.append style={level distance=3.4cm, sibling angle=51.4},
level 2/.append style={level distance=2.2cm, font=\scriptsize, text width=2.2cm}]
\node[root concept, text width=3cm, align=center]{Module One: Listen, Read and Respond}
child[concept color=popblue]{node[text width=2.6cm, align=center]{Diagnostic Assessment \& Syllabus Overview}
child{node{Exam structure}}
child{node{Placement test}}}
child[concept color=poporange]{node[text width=2.6cm, align=center]{Advanced Comprehension}
child{node{Main ideas \& implicit meaning}}
child{node{Tone, attitude, point of view}}
child{node{Facts vs opinions}}}
child[concept color=popgreen]{node[text width=2.6cm, align=center]{Summary, Synthesis \& Reports}
child{node{Selecting essentials}}
child{node{Multiple sources}}
child{node{Formal reports}}}
child[concept color=poppurple]{node[text width=2.6cm, align=center]{Grammar Mastery}
child{node{Clauses \& conditionals}}
child{node{Tenses \& modifiers}}
child{node{Idioms \& nuance}}}
child[concept color=popgold]{node[text width=2.6cm, align=center]{Phonetics \& Speech}
child{node{IPA transcription}}
child{node{Connected speech}}
child{node{Intonation}}}
child[concept color=poppink]{node[text width=2.6cm, align=center]{Essay \& Creative Writing}
child{node{Argumentative essays}}
child{node{Discourse markers}}
child{node{Short plays}}}
child[concept color=popdark]{node[text width=2.6cm, align=center]{Prescribed Texts \& Poetry}
child{node{Novels \& drama}}
child{node{Cosmic Anthology}}};
\end{tikzpicture}
\end{center}
\legende{Mind map of Module One: the seven skill strands of the Upper Sixth English Language course.}
\end{popfigure}

\begin{aretenir}[Module One --- Essentials at a Glance]
\textbf{1. Purpose of the module.} This module trains you to \cle{listen}, \cle{read} and \cle{respond} orally and in writing; to \cle{select}, \cle{order}, \cle{evaluate} and \cle{present} facts, ideas and opinions from varied texts; and to develop a \cle{positive attitude} towards listening and reading. Every GCE paper (1, 2 and 3) draws on these skills.

\textbf{2. Key definitions.}
\begin{itemize}
\item \cle{Comprehension}: extracting main ideas, supporting details and \cle{implicit meaning} (what is suggested, not stated) from passages of 900--1000 words.
\item \cle{Tone}: the writer's attitude (ironic, critical, nostalgic, sarcastic\ldots); \cle{point of view}: the angle from which events are presented.
\item \cle{Fact} vs \cle{opinion}: a fact is verifiable; an opinion expresses judgement, often signalled by modals (\emph{should, must}) and evaluative adjectives.
\item \cle{Summary}: a faithful, condensed restatement in your own words; \cle{synthesis}: integrating ideas from \emph{several} sources into one coherent text.
\item \cle{Cohesion} (grammatical and lexical links) and \cle{coherence} (logical organisation of ideas).
\item \cle{Denotation} (dictionary meaning) vs \cle{connotation} (suggested associations).
\end{itemize}

\textbf{3. Grammar to master.}
\begin{itemize}
\item \cle{Subordinate clauses}: noun clauses (\emph{What he said surprised us}), adjective (relative) clauses, adverb clauses of time, reason, condition, concession.
\item \cle{Conditionals}: zero (general truths), first (real future), second (unreal present), third (unreal past), and mixed (\emph{If I had studied, I would be a doctor now}).
\item Perfect and perfect continuous tenses; dangling and misplaced modifiers (\emph{Walking to school, the rain started} $\rightarrow$ \emph{Walking to school, I felt the rain start}).
\end{itemize}

\textbf{4. Phonetics.} Full IPA chart (vowels, diphthongs, consonants); connected speech: \cle{elision}, \cle{assimilation}, \cle{intrusion}, \cle{catenation}; intonation for attitude and grammar; self-monitoring strategies in speech.

\textbf{5. Methods.}
\begin{itemize}
\item \textbf{Summary}: read twice $\rightarrow$ identify topic sentences $\rightarrow$ paraphrase $\rightarrow$ count words $\rightarrow$ check no examples, figures of speech or quotations remain.
\item \textbf{Argumentative essay}: introduction with thesis $\rightarrow$ balanced body (for/against) $\rightarrow$ your justified position $\rightarrow$ conclusion; use discourse markers (\emph{firstly, moreover, admittedly, in conclusion}).
\item \textbf{Comprehension}: read questions first, locate line references, answer in your own words.
\end{itemize}

\textbf{6. Prescribed texts.} \emph{Twilight of Misty Foliage} (Ben Jama), \emph{Unanswered Cries} (Osman Conteh), \emph{The Old Man and the Sea} (Hemingway), \emph{The Swamp Dwellers} (Soyinka); poems from \emph{Cosmic Anthology}, including \emph{If We Must Die} (McKay) and \emph{Africa's Plea} (Dempster). Know context, characters, themes, symbols, style.

\textbf{7. Common mistakes.}
\begin{itemize}
\item Copying sentences in summaries instead of paraphrasing.
\item Confusing tone with mood; stating opinion as fact.
\item Mixed conditional errors: *\emph{If I would have known}\ldots{} is wrong.
\item Dangling modifiers after introductory phrases.
\item Quoting the text instead of analysing it in literature essays.
\end{itemize}

\textbf{8. GCE tips.} Manage time per mark; plan essays for 5 minutes; leave 5 minutes to proofread; in Paper 2, answer the question asked, not the one you hoped for.
\end{aretenir}

\findecours

\end{document}
